\documentclass[11pt]{article}
\usepackage[T1]{fontenc}
\usepackage[utf8]{inputenc}
\usepackage{lmodern,amsmath,amssymb,amsthm,mathtools}
\usepackage[a4paper,margin=27mm]{geometry}
\usepackage{booktabs,longtable,array,enumitem,microtype,xurl}
\usepackage[colorlinks=true,linkcolor=blue,citecolor=blue,urlcolor=blue,hypertexnames=false]{hyperref}
\usepackage{fancyhdr}
\pagestyle{fancy}\fancyhf{}\fancyhead[L]{THETA3: three signatures}\fancyhead[R]{21 November 2026}\fancyfoot[C]{\thepage}
\setlength{\headheight}{14pt}
\newcommand{\Q}{\mathbb Q}\newcommand{\C}{\mathbb C}\newcommand{\Z}{\mathbb Z}\newcommand{\R}{\mathbb R}\newcommand{\A}{\mathbb A}\newcommand{\OO}{\mathcal O}
\newcommand{\F}{\mathbb F}
\providecommand{\mathscr}[1]{\mathcal{#1}}
\newcommand{\Hdg}{\operatorname{Hdg}}\newcommand{\CH}{\operatorname{CH}}\newcommand{\Hom}{\operatorname{Hom}}\newcommand{\Tr}{\operatorname{Tr}}\newcommand{\SU}{\operatorname{SU}}\newcommand{\U}{\operatorname{U}}
\newcommand{\im}{\operatorname{im}}\newcommand{\vol}{\operatorname{vol}}\newcommand{\Res}{\operatorname{Res}}\newcommand{\Nm}{\operatorname{N}}\newcommand{\cl}{\operatorname{cl}}\newcommand{\Span}{\operatorname{span}}
\newcommand{\framework}[1]{\par\smallskip\noindent\textit{Framework and scope.} #1\par}
\newtheorem{theorem}{Theorem}[section]\newtheorem{proposition}[theorem]{Proposition}\newtheorem{lemma}[theorem]{Lemma}\newtheorem{corollary}[theorem]{Corollary}
\theoremstyle{definition}\newtheorem{definition}[theorem]{Definition}\newtheorem{remark}[theorem]{Remark}
\title{THETA3\\Theta kernels, definite periods and integral models\\\large Exact comparisons and quantitative certificates}
\author{Run THETA3-20261121\\GPT-6 Astra Pro, with independent agent proof checks}
\date{21 November 2026}
\begin{document}\maketitle
\begin{abstract}
We compare the CM Hodge construction with the theta and period mechanisms in Alp\"oge's archive, retaining the Hermitian spaces, splitting characters, coefficient types, measures and integral structures. Exact zero-orbit and restriction-seesaw identities give finite quantitative certificates and a positive finite-frame formula for the ball pairing. A direct construction on an $n$-ball realizes balanced $2n$-factor CM tensors by codimension-$n$ cycles; Liu's componentwise theorem supplies their precise CM sources in all ranks. A finite-dimensional kernel argument realizes the wedge as a genuine single theta lift; its true adjoint supplies a common double-lift and last-step Bessel identity. The completed Taylor tensor agrees with the matched canonical harmonic jet. For non-CM sources, finite \'etale covers of one arithmetic ball surface provide infinitely many pairwise nonisogenous elliptic Albanese factors, explicit new $j$-invariants, a simple genus-two Jacobian with endomorphism ring $\Z$, and positive integral periods. An algebraic Weil projector and a curve-product calculation isolate the remaining exceptional-Hodge-class problem. The integral part separates derived signs, horizontal cycles, intrinsic sharp height-two derivative bounds and actual reduced theta frames. Checked implications and cited inputs are distinguished from the unresolved instantiation of the complete character-independent iterated functional. No Masullo claim is used, and no general non-CM Hodge-conjecture conclusion is asserted.
\end{abstract}
\noindent\textbf{Working status.} This is the evolving research paper for the requested run. Checked mathematical statements below have the scope stated in their hypotheses; an unchecked global identification is never used as a theorem. The final status and current gates are recorded in \texttt{QUESTIONS.tsv} and \texttt{LEDGER.md}.

\tableofcontents\clearpage
\part{Scope, sources and the geometric dictionary}
\section{What is being compared}
The requested sketch is a map of proposed connections, not a collection of background theorems. In particular, sketch~1.7 does not specify the test vectors, splittings, inner forms or normalizing constants of its ``one seesaw pairing.'' We therefore distinguish exact identities, exact nonvanishing implications, and a still-uninstantiated identification with the character-independent iterated functional in Alp\"oge's step~12. This is a mathematical distinction: a local Bessel integral, a global period on a definite Shimura set, and a cup-product integral on a surface have different domains.

Throughout, $E/F$ is a CM extension with $F$ totally real. A skew-Hermitian multiplicity space is denoted $W$; writing $W=\delta P$ for a Hermitian space $P$ requires a fixed $\delta\in E^\times$ with $\bar\delta=-\delta$. The associated symplectic form is the trace of the tensor-product form. A change of $\delta$, an additive character, a Weil splitting or a determinant-line trivialization can change a numerical period. These choices are not suppressed.

In the surface construction, $V$ is Hermitian of dimension~3, with signature $(2,1)$ at one real place $\mu$ and $(3,0)$ at the others. There is at least one definite place. Thus $V$ is anisotropic over $F$: an isotropic rational vector would remain isotropic at a definite real place. A sufficiently small level gives a compact projective ball surface. We distinguish $\SU(V)$, used in the original geometric construction, from $\U(V)$, used for adelic theta identities. Passing between their appropriate components and finite covers retains determinant characters and component measures.

\subsection{Background accepted and what was checked}
The supplied 032 theorem for CM abelian varieties, the supplied human proof of van Hoften--Sempliner, and the stated Alp\"oge RH and Dirichlet GRH theorems are carried as cited background, as requested. This run does not purport to referee those global theorems. It reads and checks the particular hypotheses and formulas used. No Masullo claim enters any deduction.

The supplied ZIP contains two Lean entry files and a source map, not the full tree. Named modules were subsequently obtained from a public mirror at commit
\begin{center}\small\nolinkurl{7cbf7a0c32fca1badf05931879a37c3f8a2b03b4}.\end{center}
The mirror's \texttt{Main.lean} and \texttt{Statements.lean} match the supplied files byte for byte. This does not authenticate every mirror module against the original multi-part archive. No Lean compilation was performed. A Lean \texttt{def ... : Prop} is a definition of a proposition; it is not itself a proof of that proposition. Source readings and elementary consequences of those definitions are recorded accordingly.

A mathematical finding here means a deduction formulated and proved for this comparison; it does not mean a claim of research priority. Standard consequences included to prevent an incorrect inference are marked as such. The independent checks are whole-proof readings or reconstructions, not independent formal certifications. The model was GPT-6 Astra Pro; no numerical reasoning-effort setting was exposed. Actual check events, including background uses, appear in the TSV register.

\subsection{Realizations and correspondences}
A nonzero mixed period is a nonzero functional on a Betti realization. It produces an algebraic cycle only after an algebraic morphism or correspondence has been constructed and an extraction criterion has been verified. In the 032 setting rational maps from CM Hodge structures into $H^1(X,\Q)$ come from maps to abelian varieties after multiplying by an integer. Pushing forward $[X]$ then gives the surface cycle. An analytic inverse, an evaluation of a differential at a CM point, or a positive height is not by itself such an algebraic correspondence.

We use motives modulo homological equivalence when discussing projectors whose idempotence is proved on cohomology. We do not silently assert idempotence in the Chow ring. A homologically trivial cycle is numerically trivial, although it can carry nonzero Abel--Jacobi or height data. These distinctions are essential in comparing the ordinary and arithmetic theta correspondences.

\subsection{Grothendieck anchors}
The archival anchors in this paper locate the framework of a finding, rather than claim that its theta formula appears in Grothendieck's text. The original scans of cote~44, scans~7--8, were read: the letter of 12 August 1964 discusses the intermediate Jacobians $J^i$, their duality and the correspondence pairing used to reduce positivity to a polarization question. Cote~16, scans~43--46, the letter to Coates dated 4 January 1967, discusses the standard conjectures, equality of dimensions of complementary algebraic classes, and algebraic inverses associated with hyperplane sections. In particular the date and the first dimension statement are on scan~43, not scan~44. The published \emph{Standard conjectures on algebraic cycles}, pp.~196--198, supplies the complementary realization/polarization language. Later sections give verified SGA/EGA locators for specialization, cycle classes and spreading out.

No separate ``text of record'' transcription of these archival pages was supplied. Consequently a typographical problem in extracted modern-paper text is not entered as a correction to Grothendieck's page. The transcription register records only actual comparisons for which both readings exist.


\section{Oscillator and arithmetic source conventions}

% THETA3 contribution by /root/theta_032, 21 November 2026.
% GPT-6 Astra Pro; inherited reasoning-effort label not separately exposed.
% External theorems are cited inputs. A1--A5 deductions were checked in full.

\subsection{The exact theta dictionary and an applicable Rallis formula}

The attached 032 paper fixes a Galois CM field $E$, its real subfield
$F=E^+$ with $[F:\Q]\geq2$, and a Hermitian $E$-space $V$ of rank three,
of signature $(2,1)$ at the distinguished embedding $\mu$ and $(3,0)$
elsewhere. Its group is $\SU(V)$, and its domain $D$ parametrizes negative
lines $L\subset V_\mu$. Put $Q=V_\mu/L$; thus
\[
 T_D=\Hom(L,Q),\qquad \Omega_D^1=L\otimes Q^*.
\]
For a multiplicity space $P$ of rank $k$, the active positive Lagrangian
has pieces $L\otimes P_\mu$ and $Q^*\otimes P_\mu^*$. The half-determinant
line is $L^k$ times a constant line. These are the precise identifications
in 032, Section 4.1 and Lemma 4.1, printed pp.13--16. The character signs
are $d(\tau)=v(\tau^{-1})$, not $v(\tau)$. The finite Schwartz input, the
rational symplectic space, the embedding of $\SU(V)$, and $V$ itself remain
fixed in Proposition 4.6. Only multiplicity subspaces at compact places
move. This operation does not change the signature of $V$.

The first Taylor tensor projects to a holomorphic one-form. Proposition
4.5, pp.21--22, proves that finitely many actual finite inputs span any
specified cotangent fibre; the common level is chosen after selecting
those inputs. For $P=P_1\perp P_2$, double alternation of the completed
second Taylor tensor gives a two-form. Proposition 4.6, pp.22--24, proves
the product identity. In its divided-power convention the degree-two
coefficient of a product is
\[
 T_2^{(1)}T_0^{(2)}+T_1^{(1)}T_1^{(2)}+T_0^{(1)}T_2^{(2)}.
\]
Double alternation kills the two unmixed terms and sends the mixed term
to the wedge of the two first tensors with coefficient one. The general
completed second tensor need not be holomorphic. Intermediate compact
parameters are asserted only to give smooth forms; the rational
decomposition endpoints are holomorphic. Theorem 5.1, p.26, proved in
Section 6.3, and Corollary 6.7, p.45, then extract the rational CM Hodge
maps on a common principal cover. These arithmetic extraction results
are distinct from the analytic product identity.

The companion paper of van Hoften--Sempliner makes the global character
data explicit. In Sections 3.1 and 4.1, the $\mu_i$ are unitary
conjugate-symplectic Hecke characters of weight one, with restriction
$\epsilon_{E/F}$; their CM types satisfy
$\Phi_{\mu_i}=\Phi_i^{-1}$. The pair consists of $\U(V)$ and a
skew-Hermitian line $W_i$, with splitting characters $(\mu_i,\nu_5)$.
The character $\pi_i$ on $[\U(W_i)]$ is a source automorphic character,
and is not the splitting character $\mu_i$. Section 4.2 imposes the
global equality $\mu_1\mu_2=\mu_3\mu_4$. Lemma 4.2.1 supplies
$W_1\perp W_2\simeq W_3\perp W_4$. Doubling gives a split rank-four
skew-Hermitian space and the model $\mathcal S(V(\A_E)^2)$.
Remark 4.2.5 explains the cancellation of the diagonal $\U(V)$ character.
With that cancellation its action is $\Phi(x)\mapsto\Phi(g^{-1}x)$.
The doubling map has normalization
\[
 \delta(f\otimes\bar f)(0)=\|f\|_2^2.
\]
Proposition 4.2.6 and Section 4.2.8 express the fourfold cup period as a
specified geometric scalar times a doubled theta integral. Lemmas
4.2.7, 4.2.9, and 4.2.10 provide, respectively, nonvanishing, the
archimedean zero evaluation, and finite torus-character extraction.

For the rank-one to rank-three lift, an applicable explicit Rallis
formula is Haan, \emph{The Bessel period of $\U(3)$ and $\U(2)$ involving
a non-tempered representation}, arXiv:1403.5061v7, Section 3.3,
Theorem 3.4, printed p.8. Its proof cites J.-S. Li, \emph{Non-vanishing
theorems for the cohomology of certain arithmetic quotients}, Theorem
2.1, and the doubling factorization. This is a convergent-range
formula; no appeal to a second-term-range theorem with incompatible
dimensions is needed.

Here is its normalization. Let $V_1$ be Hermitian of rank one and $W_3$
skew-Hermitian of rank three, and put $T=\U(V_1)$ and $G=\U(W_3)$.
Haan uses the symplectic form
$\frac12\Tr_{E/F}(h_{V_1}\otimes h_{W_3})$, a fixed nontrivial additive
character $\psi$, and a unitary character $\gamma$ of
$E^\times\backslash\A_E^\times$ with
$\gamma|_{\A_F^\times}=\epsilon_{E/F}$. The splitting pair is
$(\gamma,\gamma^3)$. The theta lift is
\[
 \theta(\bar f,\varphi)(g)
 =\int_{[T]}\theta(t,g;\varphi)\overline{f(t)}\,dt,
\]
with Tamagawa measure. If $\tau$ is a unitary automorphic character of
$T$, $f_i=\bigotimes_v f_{i,v}\in\tau$, and
$\varphi_i=\bigotimes_v\varphi_{i,v}$, then
\begin{equation}\label{A-Rallis}
 \left\langle\theta(\bar f_1,\varphi_1),
                    \theta(\bar f_2,\varphi_2)\right\rangle_G
 =\frac{L_E(3/2,\operatorname{BC}(\tau)\otimes\gamma^3)}
        {L_F(3,\epsilon_{E/F})}
       \prod_v Z_v^{\#}(1,f_{1,v},f_{2,v},\Phi_{1,v}).
\end{equation}
The target Petersson measure is also Tamagawa. Local pairings are
chosen so that their product is the global Petersson pairing on
$\tau$, and local Haar measures multiply to the relevant global
Tamagawa measure. The section $\Phi_{s,v}$ is the holomorphic
Siegel--Weil section obtained from
$\rho_{1,3}(\varphi_{1,v}\otimes\bar\varphi_{2,v})$. With these
conventions,
\[
\begin{split}
 Z_v(1)&=\int_{T(F_v)}\Phi_{1,v}(\iota(t,1))
               \langle\tau_v(t)f_{1,v},f_{2,v}\rangle_v\,dt,\\
 Z_v^{\#}(1)&=
 \frac{L_v(3,\epsilon_{E_v/F_v})}
 {L_{E_v}(3/2,\operatorname{BC}(\tau_v)\otimes\gamma_v^3)}Z_v(1),\\
 \Phi_{1,v}(\iota(t,1))&=
       \langle\omega_v(t)\varphi_{1,v},\varphi_{2,v}\rangle_v.
\end{split}
\]
The normalized local factors equal one for normalized unramified data
outside a finite set. These conventions, including the last matrix
coefficient identity, are in Haan Sections 2.1, 2.4, and 3.1--3.2,
pp.3--8. The $L$-values at $3/2$ and $3$ are nonzero for the unitary
characters occurring here. Thus choosing nonzero local factors gives a
nonzero global lift. For the anisotropic rank-three group in the present
application, square integrability and cuspidality cause no additional
global convergence issue. Local occurrence and local test vectors still
must be checked; the formula does not assert that every prescribed
Schwartz vector has nonzero norm.

This displayed formula fixes one splitting convention completely.
To transport it to arbitrary splitting pairs, first compute the actual
unitary characters $\lambda(g)$ and $\kappa(t)$ in an identity
$\omega'(g,t)=\lambda(g)\kappa(t)\omega(g,t)$. Then the new lift with
input $\bar\tau$ is the old lift with input
$\overline{\tau\kappa^{-1}}$, multiplied by $\lambda$; its norm is
unchanged by the latter unitary twist. This gives the correct character
substitution without silently choosing a conjugation convention.
The factor $1/2$ in Haan's symplectic form must likewise be matched by
the form or additive-character normalization. Replacing both quotient
measures by probability measures multiplies the displayed target norm
by $\operatorname{vol}([G])^{-1}\operatorname{vol}([T])^{-2}$ when the
underlying input function is kept fixed. One must also account for
passage from $\U(V)$ and all adelic components to a connected
$\SU(V)$ quotient. No such factor is implicitly set equal to one here.



\subsection{The arithmetic dictionary in 032, Sections 5--6}

The geometric target of the degree-one theta tensor is the Albanese
realization. The arithmetic extraction uses numerical characters of a
commuting algebra of algebraic Hecke endomorphisms; these are not the
global splitting characters $\mu_i$ or the torus inputs $\pi_i$.
Here are the exact formulas, with the conventions that determine them.

At a suitable split good prime with residue cardinality $q$, let
$\Phi$ denote the $q$-power Frobenius endomorphism of the reduced
Albanese. In 032, Proposition~5.5, p.~32, the integral endomorphisms
$D_1,D_2,D_3$ satisfy
\[
 \Phi^3-D_1\Phi^2+qD_2\Phi-q^3D_3=0.
\]
The Frobenius on crystalline degree-one cohomology is its pullback
$\varphi=\Phi^*$, with the same polynomial in $D_i^*$.
Replacing this by reciprocal Frobenius would invert the slopes and
would not be this formula. The valuation is normalized by
$\nu_\iota(q)=1$ in Section~5.2, p.~30.

Fix the finite theta input away from the active split place $w$.
Lemma~6.2, pp.~37--38, constructs $r\in E^\times$, $M>0$, and
$\omega\in\overline\Q^\times$ such that
\[
 r\bar r=1,\quad \operatorname{val}_w(r)=M,
 \quad\operatorname{val}_{\bar w}(r)=-M,
 \quad D_q^M=\omega,
 \qquad
 \omega^2=\zeta\prod_{\tau\in\Sigma}\tau(r)^{-3d(\tau)}.
\]
Here $D_q$ is the normalized radial dilation, the relation holds on
the image in theta classes, and $\zeta$ is a root of unity. The
square records the metaplectic ambiguity. The exponent $-3$ is
the active rank, not the dimension of the multiplicity line.

Proposition~6.3, pp.~39--40, gives the half-modulus projective
spherical parameter
\[
 \{q^{1/2},q^{-1/2},e\},\qquad
 e=e_0^{-1},\quad e_0^M=\omega,
 \qquad
 \nu_\iota(\sigma e)=-\frac32v(\sigma|_E).
\]
Every complex conjugate of $e$ has modulus one. A scalar-conjugate
active character has third parameter $e'=\pm\sigma e$.
For a full numerical character $\chi$, write its full parameter as
$b(q^{1/2},q^{-1/2},e')$. The central relation in Lemma~6.4
implies that $b^3e'$ is a root of unity. In equation~(6.9), p.~42,
\[
 \chi(D_1)=q e_1(\alpha),\qquad
 \chi(D_2)=q e_2(\alpha),\qquad
 \chi(D_3)=e_3(\alpha),
 \quad\alpha=b(q^{1/2},q^{-1/2},e').
\]
The $e_i$ are elementary symmetric polynomials. These are numerical
eigenvalues; the Hecke operators on their generalized spaces can
still have nilpotent parts.

Proposition~6.6, pp.~43--44, combines this polynomial, Weil weight,
filtered comparison, and weak admissibility. It gives the pure
Hodge index
\[
 t_\eta=\frac12-\frac13\nu_\iota(e')\in\{0,1\}.
\]
The generalized space is holomorphic for $t_\eta=1$ and
antiholomorphic for $t_\eta=0$. The proof uses polynomial
idempotents in extending algebraic Hecke endomorphisms, so the
projectors preserve the Hodge filtration. For each scalar
conjugate the filtration ranks are calculated over the actual
compositum of the definition field and the projector coefficient
field. Thus it does not assume independence of those fields or
semisimplicity of all ambient Hecke operators.

Theorem~5.1, p.~26, proved in Section~6.3, pp.~44--45, then applies
scalar descent to the whole selected generalized space. Formula
(6.13) reads
\[
 t_{\sigma\eta}=\frac{1+v(\sigma|_E)}2.
\]
It supplies the exact rational-Hodge maps from $U(v)$ into
$H^1(X,\Q)$; Corollary~6.7, p.~45, places a finite collection
on a common principal cover. Only then does the Albanese convert
the maps into algebraic maps to the prescribed CM abelian varieties.
This completes the original rank-three dictionary. A6 uses Liu's
all-rank component theorem to obtain the corresponding identification
in higher ball dimensions, without claiming that the displayed
cubic Frobenius calculation itself has already been generalized.

The relation to the definite machine is therefore at two levels:
the oscillator kernels and splitting rules are shared, whereas the
selected numerical Hecke eigenspaces, their integral projectors,
and the definite class-function normalization have to be matched.
Q9 makes the projector denominators explicit when their integral
matrices and annihilating polynomials have been supplied.


\part{Checked findings and proved conditional implications}

\section{Common kernels, ball tensors and exact pairing recovery}

\subsection{THETA3-A1: the anisotropic Siegel zero-orbit identity}\label{A1}

\paragraph{Statement.}
Let $E/F$ be CM and let $V$ be an anisotropic Hermitian $E$-space.
Let $r\geq1$, and let $W^{\square}$ be split skew-Hermitian of rank
$2r$, with chosen complementary maximal isotropic subspaces. Use its
Schr\"odinger model $\mathcal S(V(\A_E)^r)$ and compatible splittings
for which
\[
 (\omega(g)\Phi)(x)=\Phi(g^{-1}x),\qquad
 (\omega(n(b))\Phi)(x)=\psi(\langle b,Q_V(x)\rangle)\Phi(x).
\]
Here $N$ is the Siegel unipotent, identified with the additive
$F$-space of Hermitian $r$ by $r$ matrices, $Q_V(x)$ is the Gram
matrix, and the pairing is nondegenerate. A conventional factor
$1/2$ may be absorbed in this identification. Give $[G]$ and $[N]$
probability Haar measures, where $G=\U(V)$, and define
\[
 T_\Phi(h)=\int_{[G]}\sum_{x\in V(E)^r}(\omega(g,h)\Phi)(x)\,dg.
\]
Then
\begin{equation}\label{A1-eq}
                    \int_{[N]}T_\Phi(n)\,dn=\Phi(0).
\end{equation}

\paragraph{Proof.}
Anisotropy makes $[G]$ compact; the centre is an anisotropic norm-one
torus as well. The quotient $[N]$ of an adelic vector group is compact.
On compact sets of $g$, the finite support condition confines rational
points to finitely many lattice cosets, and uniform archimedean
Schwartz decay dominates their polynomial counting function. Thus the
theta series converges absolutely and uniformly there. The $N$-action
adds only phases of modulus one. Summation and both integrals may
therefore be interchanged.

Since $Q_V(g^{-1}x)=Q_V(x)$, the inner unipotent integral is
\[
 \int_{[N]}\psi(\langle b,Q_V(x)\rangle)\,db
 =\begin{cases}1&Q_V(x)=0,\\0&Q_V(x)\ne0.\end{cases}
\]
For a nontrivial character the vanishing follows by translation by an
element on which its value differs from one. Nondegeneracy of the
pairing ensures that the character is trivial precisely for zero Gram
matrix. If the Gram matrix vanishes, every $x_i$ has zero norm, so
anisotropy forces $x=0$. The surviving term is fixed by $G$ and is
$\Phi(0)$, proving the formula. $\square$

The identity applies to both the ball Hermitian space and a totally
definite Hermitian space. It identifies the normalized zero-orbit
functional, not their nonzero Fourier coefficients, unaveraged periods,
or motives. A residual diagonal character changes its right-hand side
to $\Phi(0)\int_{[G]}\chi(g)\,dg$; isotropic $V$ instead admits
nonzero zero-Gram tuples. These hypotheses cannot be omitted.

\subsection{THETA3-A2: finite translation and quantitative nonvanishing}\label{A2}

\paragraph{Statement.}
In A1 write $\Phi=\Phi_\infty\otimes\Phi_f$, assume
$a=\Phi_\infty(0)\ne0$ and $\Phi_f(0)=1$, and keep the infinite input
fixed. For every $0<\varepsilon<1$ there is $h_f\in N(\A_{F,f})$
such that
\[
 \left|T_{\Phi_\infty\otimes\omega_f(h_f)\Phi_f}(1)\right|
                   \geq(1-\varepsilon)|a|.
\]
The finite Schwartz norm is unchanged. Every compact subgroup of
$G(\A_{F,f})$ fixing the total doubled input still fixes its translate.

\paragraph{Proof.}
The normalized mean in A1 is $a$. Compactness and continuity give
$n=(n_\infty,n_f)$ with $|T_\Phi(n)|\geq|a|$. Additive weak
approximation gives $\gamma_j\in N(F)$ with
$\gamma_{j,\infty}\longrightarrow n_\infty$. Automorphy and
continuity imply
\[
 T_\Phi(\gamma_{j,\infty},n_f)
 =T_\Phi(1,\gamma_{j,f}^{-1}n_f)
 \longrightarrow T_\Phi(n).
\]
Choose $j$ making the error at most $\varepsilon|a|$. Commutation of
the dual-pair actions identifies the value with the translated-input
integral. The finite Weil action is unitary. The same commutation shows
that every old $G$-stabilizer fixes the translated total input.
$\square$

This assertion changes finite theta data. Its finite tensor factors
may require a smaller common level after decomposition. It supplies
no uniform bound for their number or coefficients. The explicit
additive grid and Gaussian derivative certificate in THETA3-Q7
strengthen this existence argument under named effective data.

\subsection{THETA3-A3: the exact finite extraction cost}\label{A3}

\paragraph{Statement.}
Suppose an inverse-doubling decomposition and the finite torus
expansions of van Hoften--Sempliner Lemma 4.2.10 give
\[
 P=\sum_{a=1}^{R}c_a
       \sum_{\chi_i\in\mathcal X_{i,a},\ 1\leq i\leq4}
                       P_{a,\chi_1,\chi_2,\chi_3,\chi_4}.
\]
The inner terms are actual periods of integrated character lifts, with
conjugation on the last two factors and one fixed geometric
normalization. If $P\ne0$, at least one satisfies
\[
 |P_{a,\chi_1,\chi_2,\chi_3,\chi_4}|
 \geq \frac{|P|}{\sum_a|c_a|\prod_{i=1}^4|\mathcal X_{i,a}|}.
\]

\paragraph{Proof.}
Let $M$ be the maximum modulus of the finitely many periods in the
sum. The triangle inequality gives
$|P|\leq M\sum_a|c_a|\prod_i|\mathcal X_{i,a}|$; divide by the
positive denominator. $\square$

For an effective certificate use finite permissible character sets or
computable upper bounds for their sizes; determining the exact nonzero
support may itself require zero testing. Further rational Hodge-map or
connected-component extraction has its own finite coefficient cost.
The estimate is the same finite-sum principle developed in THETA3-Q5.

\subsection{THETA3-A4: higher-rank ball tensors and a surface section}\label{A4}

\paragraph{Statement.}
Replace the rank-three $V$ in the analytic part of 032 by a Hermitian
space of rank $m\geq3$, of signature $(m-1,1)$ at one real place and
positive definite at the others. Write $d=m-1$. For a multiplicity line
the first projected Taylor tensor again takes values in $\Omega_D^1$;
for a multiplicity plane, double alternation of the second tensor takes
values in $\Omega_D^2$. Finite actual inputs span every cotangent
fibre, and decomposable inputs satisfy the same wedge identity as in
032. Balanced compact multiplicity signs therefore give a nonzero
mixed pairing
\begin{equation}\label{A4-pair}
 \int_X\xi\wedge\bar\eta\wedge c_1(A)^{d-2}\ne0
\end{equation}
on a sufficiently deep compact principal ball quotient $X$, where
$A$ is ample and $\xi,\eta$ are the rational-decomposition endpoint
two-forms. A smooth complete-intersection surface $S\subset X$ of
divisors in $|A^{n_j}|$, $1\leq j\leq d-2$, satisfies
\[
 \int_S\xi|_S\wedge\overline{\eta|_S}
 =\left(\prod_jn_j\right)
      \int_X\xi\wedge\bar\eta\wedge c_1(A)^{d-2}\ne0.
\]
For $d=2$ take $S=X$ and the empty product one.

\paragraph{Proof.}
Let $L$ be the negative tautological line and $Q=V_\mu/L$, of rank
$m-1$. For multiplicity rank $k$ the active Lagrangian is
$L\otimes P\oplus Q^*\otimes P^*$. The determinant identity is
\[
 \det\mathcal F_P
 =L^{2k}\otimes(\det V_\mu)^{-k}\otimes(\det P)^{2-m}.
\]
The last two factors, and all compact-place factors, are constant on
$D$. On $\SU(V)$ the first of those constant representations is
trivial. After choosing the same constant-line trivializations as in
032, $(\det\mathcal F_P)^{1/2}=L^k$ times a constant line. Thus
\[
 (\det\mathcal F_{P_1})^{1/2}\otimes Q^*\otimes P_1^*
       \simeq L\otimes Q^*=\Omega_D^1
\]
for a line, and
\[
 (\det\mathcal F_P)^{1/2}\otimes
             \bigwedge^2Q^*\otimes\bigwedge^2P^*
       \simeq L^2\otimes\bigwedge^2Q^*=\Omega_D^2
\]
for a plane. The latter is the canonical bundle only when $m=3$.

The Gaussian-isolation argument of 032 Proposition 4.5 uses a position
space of dimension $m$ instead of three. In a fixed graph frame,
indicators of $x_0+k\widehat{\OO}_F^{\,m}$ give first tensors tending
to a nonzero Gaussian scalar times the Minkowski vector $x_0$.
Uniform Gaussian decay proves the convergence. These lattice vectors
span the complexified position space; the image of the span of actual
inputs is a finite-dimensional, hence closed, subspace. Projection
onto $L\otimes Q^*$ is surjective. Consequently actual theta values
span every cotangent fibre, and a finite selected list has a common
principal level. The covariance argument uses the same unitary Weil
splittings and triviality of determinant twists on $\SU(V)$, and is
independent of the value three.

For a multiplicity decomposition the completed theta series multiply.
The degree-two coefficient and double-alternation calculation stated
above is valid for any rank of $Q$. Select two first tensors with
independent values at one point; their wedge $\xi$ is nonzero.
Compactness of $X$ follows from a definite real place, and projectivity
from the arithmetic ball-quotient construction. Choose a positive
curvature representative $\omega$ of $c_1(A)$. Pointwise
$\xi\wedge\bar\xi\wedge\omega^{d-2}$ is a positive constant
times $\|\xi\|^2\omega^d$, so its integral is positive.

Keep the rational total multiplicity plane and its finite input fixed.
On a compact parameter neighborhood, the same Gaussian domination
with polynomial derivatives proves $C^\infty$ continuity of the
projected second tensor, exactly as in 032 Proposition 4.6(iii).
Therefore its mixed pairing with $\xi$, using $\omega^{d-2}$, remains
nonzero near the original decomposition. Balanced signs permit a
rational orthogonal decomposition with the requested endpoint signs
arbitrarily close to this compact-place datum, by the same weak
approximation argument. Only multiplicity subspaces move; $V$ is
unchanged. At that rational endpoint the tensor is a finite sum of
wedges of holomorphic first tensors, so is holomorphic and closed.
Pass to a common principal cover for the selected finite families.
The integral scales by the positive covering degree. This proves
\eqref{A4-pair}.

For sufficiently large $n_j$, Bertini gives a smooth projective
complete-intersection surface. Its cohomology class is
$(\prod_jn_j)c_1(A)^{d-2}$. Both endpoint forms are closed because
holomorphic forms on a projective manifold are closed. The
cohomological restriction formula gives the displayed surface
identity. This also shows directly that every such smooth complete
intersection retains a nonzero holomorphic two-form when $\eta=\xi$.
$\square$

This extends the analytic and geometric wedge mechanism. It does not
by itself extend 032's precise $U(v)$ arithmetic extraction theorem
with its rank-three Frobenius calculation. The exact all-rank source
identification is supplied separately in A6 using Liu's componentwise
Theorems 4.15 and 4.18 and the reflex-type calculation. Liu's
Corollary 4.20 also gives the CM Albanese conclusion in
all ranks $m\geq3$, so simply raising this ball rank does not create
non-CM $H^1$. Nor does adding positive directions turn this family
into a hyperbolic Witt tower or a chain of totally definite theta
lifts. A theorem imposing signatures at every real place still
requires those signatures to match: one active ball and positive
definite compact places do not satisfy such a hypothesis merely by
increasing $m$.

\subsection{THETA3-A5: an exact finite sum of definite Bessel periods}\label{A5}

\paragraph{Setup.}
Return to rank three and let $G=\U(V)$, with the ball signature at one
real place and positive definite signature elsewhere. Choose
$z\in V(E)$ negative at the active place, using weak approximation;
then $H=\U(z^\perp)$ is totally definite of rank two. View $H$ as a
subgroup of $G$ fixing $z$. Let $K_f\subset G(\A_{F,f})$ be compact
open and let $K_\infty$ be the stabilizer of the negative line $Ez$ at
the active place, times the compact groups at the other real places.
Fix a unitary scalar character $\kappa$ of $K_\infty$.

Let $\mathcal E$ be a nonzero finite-dimensional space of smooth
automorphic functions on $[G]$, right $K_f$-invariant, satisfying
$f(gk)=\kappa(k)^{-1}f(g)$ for $k\in K_\infty$. Give it the
Petersson inner product inherited from probability measure on $[G]$.
The application is a fixed-level span of holomorphic canonical theta
forms in their equivariant adelic realization. On a ball surface the
canonical fibre is one-dimensional, which gives precisely a scalar
$K_\infty$ type. Fix the splitting and central type and use the full
adelic realization, as in van Hoften--Sempliner, before extracting
connected components. This is not an assertion that an arbitrary
unrealized form on a connected $\SU(V)$ quotient has already been
extended to $G$.

For $a\in G(\A_{F,f})$ put
\[
 R_af(h)=f(ha),\qquad
 K_H(a)=H(\A_{F,f})\cap aK_fa^{-1}.
\]
Give $[H]$ probability measure, and let $\mathcal M_a$ be the space of
automorphic functions on $H$ right $K_H(a)$-invariant and of right
$H(F\otimes_\Q\R)$-type $\kappa^{-1}|_H$.

\paragraph{Statement.}
Writing $D=\dim_\C\mathcal E$, there are $s\leq D$ finite adelic
elements $a_1,\ldots,a_s$ for which
\[
 \mathcal Q=\sum_{j=1}^sR_{a_j}^*R_{a_j}
\]
is positive definite on $\mathcal E$. For all $u,v\in\mathcal E$,
\begin{equation}\label{A5-frame}
 \langle u,v\rangle_G
  =\sum_{j=1}^s
          \langle R_{a_j}\mathcal Q^{-1}u,R_{a_j}v\rangle_H.
\end{equation}
Choose an orthonormal basis $e_{j,\alpha}$ of $\mathcal M_{a_j}$ and
put
\[
 \mathcal P_{j,\alpha}(f)=
    \int_{[H]}f(ha_j)\overline{e_{j,\alpha}(h)}\,dh.
\]
With inner products conjugate-linear in the first variable, the
identity becomes the finite sum
\begin{equation}\label{A5-periods}
 \langle u,v\rangle_G
  =\sum_{j,\alpha}
       \overline{\mathcal P_{j,\alpha}(\mathcal Q^{-1}u)}
                   \mathcal P_{j,\alpha}(v).
\end{equation}
These are genuine global adjacent-unitary-group periods of $\U(3)$
against totally definite $\U(2)$. Moreover
\begin{equation}\label{A5-bounds}
 \lambda_{\min}(\mathcal Q)\|u\|_G^2
 \leq\sum_{j,\alpha}|\mathcal P_{j,\alpha}(u)|^2
 \leq\lambda_{\max}(\mathcal Q)\|u\|_G^2.
\end{equation}

\paragraph{Proof: well-defined finite restriction spaces.}
The left $H(F)$-invariance is inherited from $G(F)$-invariance. If
$k\in K_H(a)$, then $a^{-1}ka\in K_f$, so
$f(hka)=f(ha)$. At infinity, finite $a$ commutes with $H(F\otimes\R)$,
which is contained in $K_\infty$; this proves the asserted right type.
Thus $R_a$ maps into $\mathcal M_a$.

Because $H$ is totally definite,
$H(F)\backslash H(\A_{F,f})/K_H(a)$ is finite. After choosing finite
representatives, a scalar-type function is determined by its values
at those representatives. A finite stabilizer can force a value to
vanish when its character is nontrivial, but introduces no infinite
dimensional space. Thus $\mathcal M_a$ is finite-dimensional. Its
norm is a finite weighted sum of squared values, with positive masses
proportional to the reciprocal orders of the finite stabilizers. In
particular, a smooth nonzero restriction has positive norm.

\paragraph{Proof: separation and at most $D$ translates.}
The group $G$ satisfies weak approximation at the archimedean places
needed here. A direct proof uses the Cayley transform
\[
 X\longmapsto(1-X)(1+X)^{-1},\qquad X^*=-X.
\]
The skew-adjoint matrices form an $F$-vector space; its $F$-points
are dense in the product of its real completions. The Cayley transform
is an isomorphism of open sets, with inverse
$(1-g)(1+g)^{-1}$. Its real image, excluding the proper locus with
eigenvalue $-1$, is dense. Hence $G(F)$ is dense in
$G(F\otimes_\Q\R)$. This uses $\U(V)$, not an unproved
approximation assertion for a chosen $\SU(V)$ presentation.

If $f\in\mathcal E$ is nonzero, choose $(g_\infty,g_f)$ where its
value is nonzero. Continuity and weak approximation provide
$\gamma\in G(F)$ such that $f(\gamma_\infty,g_f)\ne0$.
Automorphy gives
\[
 f(\gamma_\infty,g_f)=f(1,\gamma_f^{-1}g_f).
\]
For $a=\gamma_f^{-1}g_f$, therefore $R_af(1)\ne0$, and
$R_af$ has positive norm. The family of restrictions separates
$\mathcal E$.

Start with its full kernel space $\mathcal E_0=\mathcal E$.
If $\mathcal E_i\ne0$, choose nonzero $f_i\in\mathcal E_i$ and
one restriction not annihilating it, and intersect $\mathcal E_i$
with that restriction's kernel. Dimension drops strictly. After at
most $D$ steps the common kernel is zero. Thus
$\langle u,\mathcal Q u\rangle_G=\sum_j\|R_{a_j}u\|_H^2$
is positive for every nonzero $u$.

\paragraph{Proof: identity, Parseval, and bounds.}
The operator $\mathcal Q$ is positive and self-adjoint, hence
invertible. By its definition,
\[
 \sum_j\langle R_{a_j}\mathcal Q^{-1}u,R_{a_j}v\rangle_H
 =\langle\mathcal Q^{-1}u,\mathcal Q v\rangle_G
 =\langle u,v\rangle_G.
\]
Finite-dimensional Parseval in each $\mathcal M_{a_j}$ gives
\eqref{A5-periods}. The spectral bounds for a positive self-adjoint
operator give \eqref{A5-bounds}. $\square$

\paragraph{Exact matrices and conditional quantitative bounds.}
In a basis of $\mathcal E$, let $G_{\rm ball}$ be its Petersson Gram
matrix. Let $R_j$ be the finite evaluation matrix and $D_j$ the
positive diagonal mass matrix for $\mathcal M_{a_j}$. Then
\[
 H_{\rm fin}=\sum_jR_j^*D_jR_j,\qquad
                  \mathcal Q=G_{\rm ball}^{-1}H_{\rm fin}.
\]
This formula keeps the analytic Gram matrix separate from definite
finite masses. If $D>1$, $G_{\rm ball}\leq M I$, and a square
$D$ by $D$ minor $B$ of the stacked matrix with blocks
$D_j^{1/2}R_j$ has determinant of modulus at least $\delta>0$, then
Cauchy--Binet and the eigenvalue product give the valid lower bound
\[
 \lambda_{\min}(\mathcal Q)
 \geq\frac{\delta^2}{M\bigl(\operatorname{tr}H_{\rm fin}\bigr)^{D-1}}.
\]
The case $D=1$ is immediate. A height or denominator argument can
supply $\delta$ only after algebraicity of appropriately normalized
entries has actually been proved. No such algebraicity, and no
rationality of $\mathcal Q^{-1}$, follows from finite dimensionality.

\paragraph{Meaning of the comparison.}
Formula \eqref{A5-periods} is an exact finite-sum refinement of the
desired comparison with definite periods, with an explicitly defined
frame correction. A fixed restriction alone computes its own norm,
not the ball norm. The evaluation of canonical forms at these CM
points is a fibrewise operation, not a pullback of a degree-two
cohomology class to a zero-dimensional variety. The analytic operator
$\mathcal Q^{-1}$ is not automatically an algebraic correspondence.
An algebraic, rational, or uniform-integral realization is an
additional theorem required for a motivic transfer.

\subsection{Direct sums, genuine seesaws, and iterated lifts}

For compatible splittings and $W=W_1\perp W_2$, the common oscillator
representation gives
\[
 \theta_W(g,h_1\oplus h_2;\varphi_1\otimes\varphi_2)
 =\theta_{W_1}(g,h_1;\varphi_1)
             \theta_{W_2}(g,h_2;\varphi_2).
\]
The differential version is the double-alternating wedge identity.
It belongs to the actual seesaw whose pairs are
$\U(V)\times\U(V)$ against $\U(W_1)\times\U(W_2)$, and diagonal
$\U(V)$ against $\U(W)$. Integrating against characters or test
automorphic forms gives a seesaw identity by Fubini on the compact
quotients in this application. A second orthogonal decomposition of
the same $W$ supplies two other tori in the same group. It does not
change $V$ to a definite Hermitian space.

This construction is bilinear in two independent inputs and combines
two kernels by a direct sum. An iterated lift is instead a composition
of linear integral operators, with the output automorphic
representation of one step serving as the input to the next. One
cannot identify the two merely from their number of theta factors.
The literal step 8 of Alp\"oge's note, p.4, describes a chain of
totally definite groups and changes of automorphic parameters. Its
recovered mirror contains actual successive operators \texttt{up}
and \texttt{dn}, rather than the above wedge operation. A comparison
therefore needs intertwining maps and matching representations,
splittings, and vectors. The finite-frame identity A5 and the
restriction seesaw for $V=z^\perp\perp Ez$ are concrete alternative
comparisons; neither asserts that the two independent degree-one
lifts already form that iterated chain.


% THETA3-A6 contribution by /root/theta_032, 21 November 2026.
% GPT-6 Astra Pro; inherited reasoning-effort label not separately exposed.
% Liu/BMM/CM classification are cited inputs, not independently re-proved.

\subsection{THETA3-A6: a direct construction for balanced tensors of arbitrary even length}\label{A6}

\paragraph{Statement and scope.}
Let $n\geq2$, let $E$ be a Galois CM field with real subfield $F$
and $[F:\Q]\geq2$, and let $c\in\operatorname{Gal}(E/\Q)$ be complex
conjugation. Suppose that the odd sign functions
$v_1,\ldots,v_{2n}:\operatorname{Gal}(E/\Q)\to\{1,-1\}$ satisfy
\begin{equation}\label{A6-balance}
 v_i(1)=1,\qquad \sum_{i=1}^n v_i=\sum_{i=n+1}^{2n}v_i.
\end{equation}
Put $\lambda=(1,\ldots,1,c,\ldots,c)$, with $n$ entries of each kind.
The theta construction on a compact ball quotient of dimension $n$
gives a direct cycle construction for the line
\begin{equation}\label{A6-line}
 \bigotimes_{i=1}^{2n}U(v_i)_{\lambda_i}
 \ \subset\ H^{2n}\!\left(\prod_{i=1}^{2n}B(v_i),\overline\Q\right).
\end{equation}
The arithmetic input is Liu's exact oscillator-component theorem in
arbitrary rank, together with the reflex-type argument of
van Hoften--Sempliner. The proof below does not infer the prescribed
$U(v_i)$ sources from CM-ness of the Albanese alone. It does not use
the reduction to four-factor relations.

All auxiliary abelian varieties in this construction remain CM. The algebraicity
conclusion is already covered by the supplied CM Hodge-conjecture
results; the contribution here is the direct $2n$-factor construction,
its exact multiplicity and differential degree, and the comparison
with the stated theta inputs. No claim of research priority is made.

\paragraph{The exact arithmetic source in rank $m\geq3$.}
Let $V$ be Hermitian of rank $m$, with signature $(m-1,1)$ at the
distinguished real place and $(m,0)$ elsewhere. It is anisotropic
over $E$, because there is a definite real place. Replacing its active
real component by a positive definite component while retaining the
finite components gives a totally definite incoherent Hermitian
space. The original $V$ is its active nearby coherent space.
Consequently Liu's setup applies to this precise Shimura tower:
Definitions C.3--C.4 and Proposition C.5, printed pp.108--109, and
Section 4.2, pp.45--47. The compact case holds because $[F:\Q]>1$.
There is no all-place indefinite-signature hypothesis in this use
of Liu's theorem.

For an odd $v$ with $v(1)=1$, write
$\Phi(v)=\{g:v(g)=1\}$ and $d(\tau)=v(\tau^{-1})$. Choose a
unitary conjugate-symplectic Hecke character $\mu$ of weight one with
\[
 \mu|_{\A_F^\times}=\epsilon_{E/F},\qquad
 \Phi_\mu=\Phi(v)^{-1}.
\]
Existence is van Hoften--Sempliner, Lemma 3.1.5, pp.9--10.
Liu's admissibility condition is the existence of a skew element
$e\in E^\times$, $e^c=-e$, such that
$\operatorname{Im}\tau(e)<0$ for $\tau\in\Phi_\mu$, with the
specified finite norm classes; see Definition 4.12, p.47. This is
exactly the 032 sign prescription: for $e=\delta h_P$, it says
\[
 \operatorname{sgn}(h_{P,\tau})
   =-d(\tau)\operatorname{sgn}(\operatorname{Im}\tau(\delta)).
\]
The skew-Hermitian source line therefore has active signature
$(0,1)$.

Liu, Definition 4.3, p.41, puts
$\mu_{\rm alg}=\mu|\cdot|_E^{-1/2}$ and defines a coefficient
field $M_\mu$ containing the reflex field $M'_\mu$ of
$(E,\Phi_\mu)$, with reflex type $\Psi_\mu$.
Definition 4.5, p.42, specifies a CM abelian variety $A_\mu$ and its
distinguished $M_\mu$-inclusion eigenline, whose Galois character is
$\mu_{\rm alg}$; Proposition 4.6, pp.42--43, supplies such data.
Theorem 4.15, pp.50--51, identifies the Galois character of the
particular admissible oscillator component with this exact character
for every rank $m\geq3$. Theorem 4.18, pp.52--54, then gives
\begin{equation}\label{A6-Liu}
 \Hom_E(A_\infty,A_\mu)_\Q\otimes_{M_\mu}\C
       \simeq\bigoplus_{\varepsilon,\chi}
                         \omega(\mu,\varepsilon,\chi).
\end{equation}
The sum uses $\mu$-admissible $\varepsilon$. The map in its proof
is pullback of a basis of the distinguished inclusion eigenline;
its finite-level assertion identifies $K$-invariants with
$\Hom_E(A_K,A_\mu)_\Q$. Here $\chi$ is the finite central
automorphic character of Liu, Definition 4.11. It is neither the
splitting character $\mu$ nor the source torus character.
Corollary 4.20, pp.54--55, describes the whole CM Albanese, but the
stronger componentwise statement \eqref{A6-Liu} is what is needed.

Let $(E_0,\Phi_0)$ be the primitive CM type inducing
$(E,\Phi(v))$. Then $B(v)\sim B_0^{[E:E_0]}$. Since $E$ is
Galois and $\Phi_\mu=\Phi(v)^{-1}$, the reflex pair
$(M'_\mu,\Psi_\mu)$ is $(E_0,\Phi_0)$, and
$A_{\mu,\C}\sim B_0^{[M_\mu:E_0]}$. These are precisely the
CM-type computations in van Hoften--Sempliner, Corollary 3.1.13,
p.12, and Proposition 4.1.8, pp.15--16. Their proof with the
primitive factor and the selected inclusion line gives
\begin{equation}\label{A6-extraction}
 \omega(\mu,\varepsilon,\chi)
 \subseteq
 \operatorname{span}_{\C}\{h(U(v)_1):
       h\in\Hom_{\Hdg}(U(v),H^1(X,\Q))\},
\end{equation}
on a sufficiently deep common cover. This argument is independent
of $m$. The inverse in $\Phi(v)^{-1}$ is essential.

\paragraph{The first theta vector and the scalar conventions.}
The local comparison needed to use \eqref{A6-extraction} is explicit.
In a graph frame at a base point, 032's first Taylor tensor is the
ordinary gradient: a nonzero scalar times a linear position function
multiplied by the positive Gaussian. The Bargmann transform sends
linear position times the vacuum to a nonzero scalar times the
corresponding linear Fock coordinate. After projection to the active
cotangent summand, this is the contraction
\[
 \psi_{1,0}=\sum_{\alpha=1}^{m-1}
                 z''_\alpha\otimes\xi'_{\alpha,m}.
\]
It is the $q=1$ case of BMM, published Lemma 5.5, pp.43--44;
equation (5.11), p.54, defines its Schwartz realization by the
inverse Bargmann transform. One may check the comparison on this
fibre and transport it by $\SU(V)$, since both sides have the same
covariance. The Gaussian creation-operator formulas in BMM,
Appendix C, Lemmas C.2--C.3, pp.115--116, also verify directly that
the coefficient is nonzero.

The numerical scalar depends on the Bargmann, Gaussian and
constant-line normalizations. It need not be algebraic: only a
complex span and nonvanishing are used in this argument. It must be
retained if a numerical period identity is later asserted.
BMM, Lemma 5.18, p.49, fixes the twist
$\det^{(a-b)/2}$ required for the special cochain on $\U(V)$.
For the holomorphic degree-one cochain $a=0,b=1$, it is
$\det^{-1/2}$ in their oscillator convention. The reciprocal
half-determinant in a covariant Taylor-value convention is a dual
representation convention, not a change that may be silently made
on one side of an identity. On $\SU(V)$ the determinant twists
are trivial. As in 032, Lemma 4.4, pp.20--21, the circle and double
lifts agree there: their ratio is a continuous character of a
connected semisimple real group and hence is trivial.

At the compact places the appropriately split Gaussian is the
trivial representation; the exact possibilities are Liu,
Lemma D.2(1), pp.127--128. At the active place the holomorphic
component is $\omega_{m-1,1}^{-1,-,0}$ by Lemma D.2(2).
Fix a source splitting $\nu_m$ with
$\nu_m|_{\A_F^\times}=\epsilon_{E/F}^m$. Changing $\nu_m$
twists the source $\U(1)$ character, denoted by $\pi$; it does not
change the indicated $\mu$-component after that twist. In the
van Hoften--Sempliner convention the lift integrates
$\theta(g,t)\pi(t)$. Writing it instead as
$\theta(g,t)\eta(t)^{-1}$ requires $\eta=\pi^{-1}$.

At fixed finite level and prescribed infinity type the source
torus quotient has only finitely many character components.
Thus the raw first theta gradient is a finite sum of the corresponding
integrated lifts, all of whose nonzero components satisfy
\eqref{A6-extraction}. The degree-one argument of
van Hoften--Sempliner, Proposition 4.1.8, applies in every rank:
BMM, Section 9.1 and Theorem 9.1, pp.72--75, require
$3(a+b)+|a-b|<2m$, which here is $4<2m$. This theorem is being
used in degree one, not for the subsequent top-degree wedge.

For a finite list of forms, construct the full $\U(V)$ lifts with
these splittings and restrict to a connected component. A sufficiently
deep principal $\SU(V)$ level is contained in their common finite
stabilizer. Equivalently, at a lattice-preserving unitary level the
determinants of rational arithmetic elements are norm-one algebraic
units, hence roots of unity; a deep principal level kills them.
This gives the required component cover and its rational Hodge
pullbacks. No extension of an arbitrary unrealized special-unitary
form is being assumed.

There is also a split-local check on the rank parameter. Liu,
Lemma D.1, proof p.126, gives, for
$\mu=\nu\boxtimes\nu^{-1}$,
\[
 \omega(\mu,\varepsilon,\chi)
 =\operatorname{Ind}_{\mathrm{GL}_{m-1}\times\mathrm{GL}_1}
                       ^{\mathrm{GL}_m}
       ((\nu\circ\det)\boxtimes\chi\nu^{1-m}),
\]
with normalized induction. Its projective unramified Satake tuple is
\[
 \{q^{(m-2)/2},q^{(m-4)/2},\ldots,q^{-(m-2)/2},
                   \chi(\varpi)\nu(\varpi)^{-m}\}.
\]
Thus the exponent is $-m$, agreeing with $-3$ in 032. An inverse
geometric Hecke convention inverts the tuple. This is a consistency
check; Theorems 4.15 and 4.18 supply the actual CM source.

\paragraph{Degree $n$ alternation on the $n$-ball.}
Now set $m=n+1$. Choose multiplicity lines $P_1,\ldots,P_n$ with
the first $n$ prescribed signs, and put $P=\perp_{i=1}^nP_i$.
At the active place all line signs agree, and
\[
 \mathcal F_P=(L\otimes P)\oplus(Q^*\otimes P^*),
 \qquad \operatorname{rank}Q=n.
\]
The determinant computation in A4 gives
\[
 \det\mathcal F_P=L^{2n}(\det V_\mu)^{-n}(\det P)^{2-m}
\]
times compact constant factors. Its square root is $L^n$ times a
constant line. Define
\[
 \operatorname{Alt}_n:\operatorname{Sym}^n(Q^*\otimes P^*)
       \longrightarrow\bigwedge^nQ^*\otimes\bigwedge^nP^*
\]
by the product of the two exterior wedges. Swapping two factors
changes both wedge signs, so it factors through the symmetric power.
After the fixed constant-line trivializations, its target with the
half-determinant is
$L^n\otimes\bigwedge^nQ^*=\Omega_D^n$.

For a decomposable finite input the completed theta series is the
product of its $n$ line theta series. Its degree-$n$ coefficient,
in the divided-power convention, is the sum indexed by
$k_1+\cdots+k_n=n$. A summand with some $k_i\geq2$ contains
two factors in the same multiplicity line and is killed by
$\bigwedge^nP^*$. The sole surviving index is $(1,\ldots,1)$.
Consequently the projected degree-$n$ tensor is exactly the wedge
of the $n$ first tensors, with coefficient one.

For an arbitrary finite input the same assertion holds as a finite
sum of such wedges. Indeed a compactly supported locally constant
function on a finite-adelic direct sum is a finite sum of rectangle
indicators after choosing a sufficiently fine open lattice, so its
Schwartz space is the algebraic tensor product of the factor spaces.
A change to a rational adapted polarization uses the rational
metaplectic covariance of 032, Lemma 4.2 and Remark 4.3, in the
corresponding dimension. Its nonzero lift scalar is retained in the
constant-line normalization; it is not discarded outside $\SU(V)$.

Here is the spanning argument in the new dimension. Fix a point
$x_*\in D$ and, for one prescribed multiplicity line, choose rational
Schr\"odinger positions $F^m$. Their Minkowski image is a real vector
space of dimension $D_0=m[F:\Q]$, with full lattice
$\Lambda=\OO_F^m$. In the graph frame of its period point
$Z=(Z_\tau)_\tau$, put
\[
 P_i:\C^{D_0}\longrightarrow
       (L\otimes Q^*)_{x_*}.
\]
This is the complex-linear projection of the first Taylor fibre
onto its $c\mu$ summand, followed by the fixed nonzero
half-determinant and multiplicity-line trivializations. It is
surjective. The complex span of $\Lambda$ is $\C^{D_0}$, so choose
$x_{i,1},\ldots,x_{i,n}\in\Lambda$ such that
$P_i x_{i,1},\ldots,P_i x_{i,n}$ is a basis of the cotangent fibre.
For each $x_0=x_{i,j}$ use the actual finite input
$\phi_{x_0,k}=1_{x_0+k\widehat{\OO}_F^{\,m}}$. Its rational
support is exactly $x_0+k\OO_F^m$, and its projected first tensor is
\[
 2\pi i\sum_{x\in x_0+k\Lambda}
       P_i x\,
       \exp\left(\pi i\sum_\tau x_\tau^tZ_\tau x_\tau\right).
\]
Choose $\lambda_0>0$ below the least eigenvalue of every
$\operatorname{Im}Z_\tau$, and let $\rho>0$ be a lower bound for
the norm of a nonzero lattice vector. For
$k\geq2\|x_0\|/\rho$ and $0\ne y\in\Lambda$,
$\|x_0+ky\|\geq k\|y\|/2$. The tail is therefore bounded by
\[
 2\|P_i\|\sqrt{\frac{\pi}{e\lambda_0}}
   \sum_{r\geq1}(2r+3)^{D_0}
      \exp\left(-\frac{\pi\lambda_0 k^2\rho^2 r^2}{8}\right).
\]
Indeed $u e^{-\pi\lambda_0u^2/2}\leq
(\pi e\lambda_0)^{-1/2}$ absorbs the linear factor, and disjoint
balls of radius $\rho/2$ give at most $(2r+3)^{D_0}$ lattice
points in the shell $r\rho\leq\|y\|<(r+1)\rho$.
The displayed series tends to zero by dominated convergence;
THETA3-Q2 also gives a finite certified tail bound. Thus the
projected tensors tend to
\[
 2\pi i\,P_i x_{i,j}\,
 \exp\left(\pi i\sum_\tau
              x_{i,j,\tau}^tZ_\tau x_{i,j,\tau}\right).
\]
Each exponential is nonzero. The determinant of these $n$ limit
columns is nonzero, so continuity of the determinant gives $n$
actual finite inputs whose values form a basis. This construction
works separately for every prescribed type, with its own positive
Gaussian and projection $P_i$.

Now choose one input from each of the first $n$ prescribed families
in order. The already chosen values span a proper subspace until
the last step, while the next family's selected finite basis spans
the full fibre. One of its values lies outside that subspace.
The resulting $n$ values are independent. Their first theta forms
are holomorphic because the first completed tensor is the ordinary
holomorphic gradient and the active embedding projection is
holomorphic. Their wedge $\xi$ is nonzero. Only after this finite
selection do we choose a common principal level, using smoothness
of the finite Weil action. On the resulting smooth projective
compact ball quotient $X$,
\begin{equation}\label{A6-positive}
 i^{n^2}\int_X\xi\wedge\bar\xi>0.
\end{equation}
The phase is necessary when $n$ is odd: in complex coordinates,
$i^{n^2}dz_1\wedge\cdots\wedge dz_n\wedge
d\bar z_1\wedge\cdots\wedge d\bar z_n=2^n\,d\operatorname{vol}$.

\paragraph{A nearby rational orthogonal decomposition.}
Keep $V$, $P$, its rational symplectic realization, the finite total
input, and the total special-unitary level fixed. At a compact real
place, \eqref{A6-balance} says that the first and last lists contain
the same number of each sign. Relabel an initial real orthogonal
basis to have the desired last-list labels. This leaves its aggregate
positive and negative multiplicity subspaces unchanged. At the
active place every line has the same sign, so the total active
positive Lagrangian is unchanged as well.

Approximate these ordered orthogonal bases at all real places by
vectors in $P(E)$, using weak approximation. Apply Hermitian
Gram--Schmidt. The denominators are nonzero norms; being nonzero
with the required signs is an open condition. The resulting
rational orthogonal basis can therefore be chosen arbitrarily close
and with all the prescribed signs. It gives
$P=P_{n+1}\perp\cdots\perp P_{2n}$ over $E$.

Positive Gaussian majorants, with polynomial factors for derivatives,
give $C^\infty$ continuity of the projected total tensor on a compact
parameter neighborhood, exactly as in 032, Proposition 4.6(iii),
pp.22--24. Choose its constant-line trivializations continuously,
agreeing with the original ones at the reference point. The mixed
pairing with $\xi$ remains nonzero near that point by
\eqref{A6-positive}. Intermediate parameters need only give smooth
forms. At the chosen rational endpoint, finite tensor decomposition
gives a finite sum of wedges of holomorphic first tensors of the last
$n$ types. Select a summand whose mixed period is nonzero. After
a common principal cover for the selected inputs, one has
\begin{equation}\label{A6-period}
 \int_X\theta_1\wedge\cdots\wedge\theta_n
       \wedge\bar\theta_{n+1}\wedge\cdots\wedge\bar\theta_{2n}
        \ne0.
\end{equation}
The cover multiplies the integral by its positive degree.

The full unitary splitting characters may also be matched globally.
Choose $\mu_1,\ldots,\mu_{2n-1}$ with their prescribed types and set
\[
 \mu_{2n}=\frac{\prod_{i=1}^n\mu_i}
                    {\prod_{i=n+1}^{2n-1}\mu_i}.
\]
Its restriction to $\A_F^\times$ is
$\epsilon_{E/F}^{n-(n-1)}=\epsilon_{E/F}$, and
\eqref{A6-balance} gives its required weight-one infinity type.
Thus $\prod_{i\leq n}\mu_i=\prod_{i>n}\mu_i$ globally, with the
same $\nu_m$ retained on the source side. The special-unitary
argument has already removed these determinant twists, but they
are available for a compatible full-unitary tensor rule.

\paragraph{The direct cycle extraction, including its sign.}
By \eqref{A6-extraction}, expand the finitely many classes in
\eqref{A6-period} in complex spans of rational Hodge images.
Conjugate the last $n$ expansions. Nonvanishing selects rational
Hodge maps
$h_i:U(v_i)\to H^1(X,\Q)$ and nonzero
$\alpha_i\in U(v_i)_{\lambda_i}$ with algebraic Betti
coefficients, whose integrated product is nonzero. The arbitrary
analytic scalars are coefficients of the expansion; each selected
eigenvector itself is chosen over $\overline\Q$.

By the Albanese and the rational Riemann theorem, realize suitable
positive integer multiples of the $h_i$ by morphisms
$f_i:X\to B(v_i)$, as in 032, Lemma 2.2, p.5. Put
$A=\prod_iB(v_i)$, $g=\dim A$, $f=(f_i)$, and
$Z=\operatorname{cl}(f_*[X])\in H^{2g-2n}(A,\Q)$.
Writing $t_\alpha=\bigotimes_i\alpha_i$ in the distinct
$H^1$ slots, the projection formula gives
\[
 \int_A Z\smile t_\alpha=k\in\overline\Q^\times.
\]
For each $i$ take a copied factor $B(v_i)'$ and a nonzero
algebraic-Betti vector $\beta_i$ in its line labelled by
$c\lambda_i$. The pure class
$D_i=\alpha_i\otimes\beta_i\in H^2(B_i\times B_i',\overline\Q)$
is algebraic. Indeed every scalar conjugate puts its two factors
in opposite Hodge types, by centrality of $c$; scalar descent and
the Lefschetz $(1,1)$ theorem give its exact class as an algebraic
linear combination of divisor classes. No projection of an
unspecified divisor is involved.

Let $A'=\prod_i B_i'$, let $p,p'$ be the two projections from
$A\times A'$, and let $q_i$ select $B_i\times B_i'$. With the old
factors ordered before the new ones,
\[
 K:=\prod_iq_i^*D_i=(-1)^n t_\alpha\otimes t_\beta,
 \qquad t_\beta=\bigotimes_i\beta_i.
\]
The number of odd-degree transpositions is
$(2n)(2n-1)/2=n(2n-1)$, of parity $n$. Hence
\begin{equation}\label{A6-cycle}
 p'_*(p^*Z\smile K)=(-1)^n k\,t_\beta.
\end{equation}
Its codimension is $(g-n)+2n-g=n$. All operations on the left
preserve algebraicity, and $k$ is nonzero algebraic. Divide by $k$
and apply scalar complex conjugation, which preserves the
$\overline\Q$ span of rational cycle classes and changes the labels
back from $(c^n,1^n)$ to $(1^n,c^n)$. Identifying the copied factors
proves algebraicity of \eqref{A6-line}. This is a cycle on the line,
with the codimension and coefficient given by \eqref{A6-cycle}.
\hfill$\square$

\paragraph{The Grothendieck framework.}
This is an explicit composition of divisor correspondences,
pullback and proper pushforward in homological motives. The
relevant framework locators in the supplied Grothendieck scans
are \emph{cote 44}, scans 7--8, for intermediate Jacobians and
polarizations through correspondences, and \emph{cote 16},
scans 43--46, for the dimensions of Lefschetz operators and
algebraic inverse correspondences. Here the actual correspondence
$K$ is displayed and its output is computed in
\eqref{A6-cycle}. No general algebraic inverse Lefschetz
operator or Chow projector is assumed. Its action on the
specified tensor line follows directly from the divisor classes
and the projection formula.

\paragraph{An explicit family not generated by divisor products.}
For every $n\geq2$, choose a Galois CM field with at least $n+1$
real pairs of embeddings; a sufficiently large cyclotomic field
does so. Choose compact-place representatives
$\tau_1,\ldots,\tau_n$. Let $v_0$ be $+1$ on a chosen set of
representatives, including the distinguished one, and extend it
oddly. For $1\leq i\leq n$, let $v_i$ differ from $v_0$ only at
$\tau_i$ and its conjugate. Let $v_{n+1}$ differ at all those
$n$ pairs and put $v_{n+2}=\cdots=v_{2n}=v_0$. Then
\[
 \sum_{i=1}^n v_i=v_{n+1}+(n-1)v_0.
\]
No first-list function equals a second-list function.

The selected line \eqref{A6-line} is not in the subalgebra generated
by divisor classes. To prove this using the full Galois criterion,
project a product of $n$ divisor classes to that pure
one-$H^1$-per-factor line. A potentially nonzero contribution must
pair the $2n$ factors. At the distinguished embedding, each pair
must join a first-list factor to a last-list factor. For its selected
two eigenlines to have type $(1,1)$ under every scalar conjugation,
the two complete type functions must be equal. There are no such
pairs. Every divisor product therefore has zero projection to the
line. This argument uses the fixed eigenline labels, not a claim
that the underlying abelian varieties are pairwise nonisogenous.
Galois translates or isogenies involving other labels do not change
the conclusion. Formula \eqref{A6-cycle} nevertheless constructs
the line by a codimension-$n$ algebraic cycle, for all $n$, including
$n\geq4$.

\paragraph{Transfer to actual Weil pieces.}
Let a CM abelian variety $A_0$ of dimension $2n$ have multiplication
by an imaginary quadratic field $K$ with signature $(n,n)$.
Each of its two Weil lines over $\overline\Q$ is a wedge of
$2n$ CM embedding lines and has exactly $n$ of each Hodge type
under every scalar conjugation. Choose a Galois CM field containing
the CM fields involved, and let $s_i(g)$ be the complete type
functions of those lines. Reorder by their type at the distinguished
embedding and put $v_i=s_i$ for the first $n$, $v_i=-s_i$ for the
last $n$. Then \eqref{A6-balance} holds, and the selected original
line has the same complete type as $U(v_i)_{\lambda_i}$.

Extend a map between these matching lines by zero on the other
source embedding lines. Every scalar conjugate is a Hodge map,
so scalar descent expresses it as a $\overline\Q$ linear
combination of rational Hodge maps; this is 032, Lemma 2.4, p.8.
Lemma 2.2 realizes those maps by pullbacks of abelian morphisms.
Apply their tensor product to \eqref{A6-line} and take cup product.
After expansion this is a linear combination of pullbacks along
product morphisms from $A_0$ to the auxiliary product, and hence
preserves algebraicity. It proves the prescribed Weil line
algebraic, without an assumed algebraic projector.

For an explicit family in every $n$, choose a totally real Galois
field $F$ of degree $2n$, put $E=F(i)$, and select a CM type
containing $n$ embeddings above each embedding of $\Q(i)$.
The CM abelian variety has dimension $2n$ and the required Weil
signature. One field construction takes a prime
$\ell\equiv1\pmod{4n}$, the unique degree-$2n$ subfield of
$\Q(\zeta_\ell)$, which is real, and then adjoins $i$.
This does not assert simplicity or that every such Weil class is
outside the divisor algebra.

The same argument applies to the embedding-line monomials of any
rational CM Hodge class: rationality makes every occurring monomial
balanced under all scalar conjugations. Thus the direct
$2n$-factor construction can replace the four-factor combinatorial
reduction, with Liu's all-rank theorem retained as a deep cited input.

The remaining requirements for the proposed iterated Bessel bridge
are different. They include an actual tuple satisfying the recovered
\texttt{LawsX} and link conditions, identification of these
cohomological vectors and their normalizers, and the relevant
integral theta frames. Increasing $m$ does not satisfy the
every-place signature hypotheses of a different tower theorem,
and an independent $n$-fold product is not thereby an iteration
of theta operators.


% THETA3-A7 contribution by /root/theta_032, 21 November 2026.
% GPT-6 Astra Pro; inherited reasoning-effort label not separately exposed.

\subsection{THETA3-A7: an independent wedge is a genuine single theta lift}\label{A7}

\paragraph{Setup and statement.}
Let $V$ have rank $n+1$, signature $(n,1)$ at one real place and
positive definite signature elsewhere, with $[F:\Q]\geq2$.
Let $P=P_1\perp\cdots\perp P_n$ be the skew-Hermitian multiplicity
space of the first $n$ types in A6. Its active component is definite.
Put
\[
 G=\U(V),\qquad S=\U(P),\qquad T=\prod_{i=1}^n\U(P_i).
\]
The quotient $[S]=S(F)\backslash S(\A_F)$ is compact: a nonzero
global isotropic vector would remain isotropic in the definite
active component. This assertion does not require $S$ to be compact
at every real place.

Fix compatible theta splittings, a finite Schwartz input, and the
degree-$n$ holomorphic Fock cocycle. At the active place this is
the determinant coefficient of the wedge of $n$ degree-one
cocycles; at the other places take the specified Gaussian
coefficients. Let
\[
 \mathcal K_h=\theta_V(\,\cdot\,,h;\Phi)\in\mathcal E,\qquad
 \mathcal E=H^0(X,\Omega_X^n),
\]
in the equivariant adelic realization on a fixed sufficiently small
compact Shimura level, including its finitely many components.
The notation includes the indicated differential-form projection.
Fix a product of unitary source characters
$\pi(t)=\prod_i\pi_i(t_i)$ on $[T]$, and use the
van Hoften--Sempliner convention
\[
 F=\int_{[T]}\mathcal K_t\,\pi(t)\,dt.
\]
For a decomposable input the tensor rule and degree-$n$ alternation
identify $F$ with the wedge of the corresponding $n$ integrated
degree-one lifts. For a finite sum of inputs it is the corresponding
finite sum of wedges.

There is a smooth automorphic function $v_\pi$ on $[S]$, of finite
level and finite maximal-compact type, such that the exact identity
\begin{equation}\label{A7-lift}
 F=\int_{[S]}\theta_V(\,\cdot\,,h;\Phi)\,
                                  \overline{v_\pi(h)}\,dh
\end{equation}
holds with the displayed measures. Thus the independent wedge is
an actual single theta lift from $\U(P)$. The representative
$v_\pi$ is allowed to depend on the source characters and on the
kernel. No character independence is asserted.

\paragraph{Why the kernel has finite rank.}
First we verify that $\mathcal K_h$ belongs to the same
finite-dimensional holomorphic space for every $h\in S(\A_F)$.
At infinity the source Weil action commutes with the target
$G$-action, its maximal compact action, and its relative Lie
algebra differential. Applying the source action to a closed
cochain of Hodge type $(n,0)$ therefore gives another closed
cochain of that type. The corresponding differential form is
$\bar\partial$-closed and hence holomorphic. The original cochain
is the wedge of the degree-one holomorphic cocycles; its
closedness follows from the derivation rule for the cochain
differential. The explicit source is BMM, Lemma 5.4 and
Proposition 5.17(1),(4), pp.43 and 48. No stable-range assertion
about generation in degree $n$ is used.

At finite places the source and target Weil actions commute.
Every element of a compact open target group fixing the original
finite input consequently fixes its source translate. The target
level is thus independent of $h$. The chosen target splitting
and canonical compact type are retained. Compactness of $X$
implies that $\mathcal E$ is finite-dimensional.

The theta series and its derivatives converge uniformly on compact
sets of representatives, by Schwartz decay and the positive
Gaussian estimates. This gives a smooth $\mathcal E$-valued
function of $h$. It is left $S(F)$-invariant by rational theta
automorphy. A sufficiently small finite source subgroup fixes
$\Phi_f$, giving a common finite level. At infinity a polynomial
times a Gaussian is finite under a source maximal compact
subgroup, so the resulting coefficient functions have finite
maximal-compact type. These statements use compact representative
sets for the anisotropic adelic quotients, not a finite-dimensional
assertion about all functions with one compact type on a
noncompact real group.

Choose a Petersson orthonormal basis
$f_1,\ldots,f_D$ of $\mathcal E$, with inner products conjugate
linear in the first variable. Then
\begin{equation}\label{A7-kernel}
 \mathcal K_h=\sum_{j=1}^D f_j\,e_j(h),
 \qquad e_j(h)=\langle f_j,\mathcal K_h\rangle_G .
\end{equation}
The functions $e_j$ are smooth automorphic functions on $[S]$
with the level and compact-type properties just proved.
Let $\mathcal C=\operatorname{span}_{\C}\{e_1,\ldots,e_D\}$.
It is finite-dimensional. Its $L^2([S])$ inner product is positive
definite: a nonzero smooth function is nonzero on an open set
of positive Haar measure. In particular the passage from smooth
functions to their $L^2$ classes is injective here.

\paragraph{The exact Riesz representative.}
Define a complex-linear functional on $\mathcal C$ by
\[
 \Lambda_\pi(e)=\int_{[T]}e(t)\pi(t)\,dt.
\]
The inclusion $T\hookrightarrow S$ gives a well-defined continuous
map $[T]\to[S]$. The torus quotient is compact, so this integral
exists for every $e\in\mathcal C$. Its boundedness as a functional
on $\mathcal C$ follows from finite dimensionality; no assertion
about boundedness of pointwise restriction on an unrestricted
$L^2$ space is needed.

By finite-dimensional Riesz representation there is a unique
$v_\pi\in\mathcal C$ such that
\[
 \Lambda_\pi(e)=\langle v_\pi,e\rangle_S
               =\int_{[S]}\overline{v_\pi(h)}e(h)\,dh
 \quad(e\in\mathcal C).
\]
Substitute \eqref{A7-kernel} and interchange the finite sum and
the integrals:
\[
 F=\sum_jf_j\Lambda_\pi(e_j)
   =\int_{[S]}\left(\sum_jf_je_j(h)\right)
                                  \overline{v_\pi(h)}\,dh.
\]
This is \eqref{A7-lift}. If $F\ne0$, then
$\Lambda_\pi$ is nonzero on $\mathcal C$, so $v_\pi\ne0$.
For an orthonormal basis $\eta_1,\ldots,\eta_r$ of $\mathcal C$,
the representative and its norm are explicitly
\[
 v_\pi=\sum_{a=1}^r
              \overline{\Lambda_\pi(\eta_a)}\,\eta_a,
 \qquad
 \|v_\pi\|_S^2=\sum_{a=1}^r|\Lambda_\pi(\eta_a)|^2.
\]
These are analytic finite-dimensional quantities; their values
or normalizations have not been proved algebraic.
\hfill$\square$

The assertion is an equality of theta integrals with a smooth
automorphic input. It does not select a prescribed irreducible
source representation. When $P$ is totally definite, the entire
archimedean source group is compact; the stated finite source
level and finite archimedean type then place $v_\pi$ in a
finite-dimensional automorphic module, which has a finite
spectral decomposition. For the more general anisotropic source,
finite rank of the particular coefficient space $\mathcal C$
is what proves the identity.

\paragraph{The signature condition for a totally definite middle source.}
At a compact real place $\tau$, the sign of the Hermitian form on
$P_i$ is
$-d_i(\tau)\operatorname{sgn}(\operatorname{Im}\tau(\delta))$.
Consequently $P$ is definite there if and only if all
$d_i(\tau)$ agree. It is already definite at the active place.
Thus the actual multiplicity space is totally definite if and
only if
\[
 v_1=v_2=\cdots=v_n
\]
as complete type functions. For $n=2$, a place where $v_1$
and $v_2$ differ gives signature $(1,1)$ on $P$. Multiplying
the form by a scalar can reverse the signs but cannot make that
component definite.

There are always nonzero wedges from one repeated type: by A6's
fibre-spanning proof, one theta family supplies $n$ forms with
independent values. Such a wedge has a totally definite middle
source in \eqref{A7-lift}. But for the particular balanced tensor
of A6, equality of the first $n$ types forces equality of all
$2n$ types. Indeed at each $g$, the first sum is either $n$ or
$-n$; a sum of $n$ numbers in $\{1,-1\}$ attains that value
only when each summand has the same sign. Hence every last-list
$v_i$ equals the common first-list function. The selected
balanced tensor line is then a product of divisor lines.

This separates an arbitrary nonzero canonical wedge from the
particular wedge needed to reach a selected non-Lefschetz
Weil component. Changing to isogenous CM factors does not
remove this condition: the signs follow the complete types of
the selected eigenlines, including their embedding labels.

In the direct $m=2$ instantiation of the recovered
\texttt{MainX}, the middle group is required to be totally
definite. Formula \eqref{A7-lift} therefore provides the actual
middle source only in the just-described definite case.
For a general exceptional balanced input, it supplies a genuine
single lift from an anisotropic group with compact quotient,
but not the required totally definite middle space. It does
not rule out another source realization of the same canonical
automorphic representation. A second reverse lift, its chosen
vector, and the complete \texttt{LawsX} and arithmetic link data
remain separate requirements for that iterated comparison.


\subsection{THETA3-A8: the completed Taylor tensor is the canonical harmonic jet}\label{A8}

This comparison concerns a fixed polarized abelian variety and its
rigidified theta line bundle.  When the polarized torus and its CM
action in 032 are realized by an archive datum, it identifies the
actual analytic section and its derivatives.  It does not identify
an arbitrary representation in a theta packet with that section.
The distinction between a divided Taylor coefficient and a
differential operator is retained throughout.

\paragraph{The completed theta function is the Appell--Humbert reading.}
First use principal symplectic coordinates.  Write $Z=X+iY$ with
$Y>0$, put $\Lambda=\Z^g+Z\Z^g$, and set
\[
 H(z,w)=z^tY^{-1}\overline w,\qquad
 \widehat\vartheta(Z,z)
       =\exp\left(\frac\pi2z^tY^{-1}z\right)\vartheta(Z,z).
 \tag{A8a}
\]
For $\lambda=a+Zb$ with $a,b\in\Z^g$, the usual theta multiplier
and a direct calculation give
\[
 \widehat\vartheta(Z,z+\lambda)
  =(-1)^{a^tb}
    \exp\left(\pi H(z,\lambda)+\frac\pi2H(\lambda,\lambda)\right)
       \widehat\vartheta(Z,z).
 \tag{A8b}
\]
Indeed, the ordinary multiplier is
$\exp(-\pi i b^tZb-2\pi i b^tz)$.  The new linear term is
$\pi z^t(Y^{-1}\lambda-2ib)=\pi z^tY^{-1}\overline\lambda$.
For the constant term, use
$\lambda-\overline\lambda=2iYb$ to obtain
\[
 \frac\pi2\lambda^tY^{-1}\lambda-\pi i b^tZb
  =\frac\pi2\lambda^tY^{-1}\overline\lambda+\pi i a^tb.
\]
This proves (A8b), including its sign.

For a finite Schwartz input, choose a common lattice level $N$ so
that the positions in its support lie in $N^{-1}\Z^g$ and its
coefficients are invariant under $N\Z^g$.  The same reindexing
argument gives (A8b) on $N\Lambda$.  Equivalently,
$w\mapsto\widehat\vartheta_\varphi(Z,Nw)$ is a section of the
$N^2$-th power of the line bundle on $\C^g/\Lambda$ with the
displayed multiplier.  A nonprincipal polarization is handled in
an integral symplectic lattice with its actual Riemann form; the
same calculation uses its corresponding Hermitian form and
semicharacter.  In the restriction-of-scalars coordinates of 032,
the exponent and Hermitian form are the sums over real embeddings.

The archive's \texttt{AnFr.IsTh}, in
\path{Definitions/ThetaValues/ThConst.lean}, is exactly the law
\[
 \theta(z+\lambda)
   =\chi(\lambda)^k
       \exp\left(k\pi H(z,\lambda)
                +\frac{k\pi}{2}H(\lambda,\lambda)\right)\theta(z).
\]
Thus (A8a), with the stated level scaling, is its analytic reading.
Here the polarized lattice, semicharacter, and rigidification must
be matched, and the two readings must represent the same actual
section.  For that section, two nowhere-vanishing trivialization
identifications with the same multiplier have a holomorphic
periodic frame ratio.  This ratio descends to a compact complex
torus, hence is constant, and the rigidification fixes that
constant.  This does not claim that the multiplier determines an
arbitrary section of a higher tensor power.  The $E$-adjointness of
$H$ is the Rosati adjointness of
the polarization and supplies \texttt{AnPt.adj}.  The tangent
frame supplies \texttt{AnPt.der}.

There is a conjugation in the embedding labels that must be kept.
The torus with period $Z$ is naturally
$A=\mathcal F_P^\vee/\iota(\Lambda_{\mathrm{rat}})$, with
$\iota(u)(v)=\psi_P(v,u)$ for $v\in\mathcal F_P$.
In the graph frame this is $\iota(p,q)=Zq-p$, which gives the
displayed period lattice.  Rosati adjointness gives
$\iota(au)(v)=\iota(u)(\overline a v)$.
Consequently the tangent variable dual to the tensor block
$(\mathcal F_P)_\tau$ has the archive tangent label $c\tau$.
The $c\mu$ Taylor-tensor block of 032 therefore corresponds to
the $\mu$ tangent variables in this convention.

\paragraph{What the archive's canonical derivative means.}
The exact definitions are as follows.  In
\path{Formalism/LocalToGlobal/ProjectiveLattices/IsThRule.lean},
\texttt{NabV} assigns to a framed point $y$, a section $F$ of
$L^k$, and a monomial $\alpha$ in the $E$-eigenvectors of the
tangent space a scalar $\operatorname{nab}_y^\alpha(F)$.
\texttt{ThGerm.lean}, \texttt{ChartAt} and \texttt{IsNab}, require
that, for $k\geq0$, these scalars be computed from a formal-group
chart, a compatible biextension frame, a normalized primitive,
and a canonical gauge.  They are not arbitrary ordinary
derivatives in an unspecified local generator.

The precise correction is in
\texttt{ExpTiePf.lean}, theorem \texttt{tie\_holds}: if $f$ is the
section's formal coefficient in such a chart, $\phi$ is the formal
exponential with the prescribed tangent frame, and $h_L$ is the
chart's canonical potential, then the corrected exponential germ
is
\[
       g=\bigl(f\exp(kh_L)\bigr)\circ\phi,
 \qquad
       \operatorname{nab}_y^\alpha(F)=\alpha!\,[z^\alpha]g.
 \tag{A8c}
\]
This is the actual sign and factorial convention.  The definition
\texttt{expVal} in \texttt{IsExpChart.lean} includes $\alpha!$.
The analytic chart constructed in
\nolinkurl{ExistsAnalyticBiextensionChartComplexPoint.lean} has
additive formal group and $h=0$, and its germ is the Taylor series
of the Appell--Humbert reading.  The two proof files
\nolinkurl{HarmonicValueCanonicalConnectionEqPolyDeriv.lean} and
\nolinkurl{ThetaValueAnalyticFrameEqPolyDerivFock.lean} use exactly
these data, rather than identifying a raw theta series with an
integral local coefficient.

For completeness, the residual ambiguity has a precise and much
smaller form.  The normalized exponential-frame comparison
\texttt{GaugeUniqSt}, proved in \texttt{GaugeUniqPf.lean}, gives
\[
       g'_b=g_b\exp(-kq),\qquad
 q(z)=\sum_{\tau,i,j}q_{\tau ij}z_{\tau i}z_{c\tau,j},
 \tag{A8d}
\]
where one representative of each conjugate pair is used.
In particular $q$ is homogeneous of degree two; there are no
linear or higher terms.  These restrictions come from the
rigidification and the CM adjointness cocycle in
\texttt{IsExpGauge}, not from a choice to disregard other terms.
The formal exponential is unique for its tangent frame.
For the chart used by \texttt{IsNab}, compatibility under jet
truncation follows from existence of germs for every section of
every nonnegative tensor power: this is
\nolinkurl{GlobalTheory.NablaPinCompare.compat_of_germs} in
\nolinkurl{ExistsIsNabNotIsNablaPin.lean}, using properness and
ampleness.  It is therefore legitimate to apply the exponential
and gauge comparison to the actual \texttt{IsNab} data.  All the
archive statements in this paragraph are source inputs read at
the pinned mirror commit; no Lean build or transitive dependency
certificate is claimed here.

\paragraph{A stronger harmonic comparison, with proof.}
Let $z_{\tau i b}$ denote the coordinate in tangent embedding
$\tau$, row $i$, and section column $b$.  Let
\[
 \Delta_{\tau ij}
     =\sum_b
       \frac{\partial^2}
       {\partial X_{\tau i b}\partial X_{c\tau,jb}}.
\]
For any polynomial $P$ satisfying
$\Delta_{\tau ij}P=0$ for every conjugate pair and all rows, and
sections $F_b$ of the same nonnegative power $L^k$, one has
\[
 \sum_\alpha P_\alpha
      \prod_b\operatorname{nab}_y^{\alpha_b}(F_b)
   =P(\partial)\left(\prod_b\theta_b(z_b)\right)\bigg|_{z=0},
 \tag{A8e}
\]
where $\theta_b$ are their Appell--Humbert readings in the matched
tangent frame.  The scalar coefficients can be complex; the
arithmetic version uses the indicated coefficient field.

Here is a direct proof, which also explains why highest-weight
conditions are unnecessary for this particular equality.  Pair a
polynomial and a formal series by
$[P,F]=P(\partial)F(0)$.  On monomials,
\[
       [P,z_aF]=[\partial_{X_a}P,F].
\]
Indeed both sides are zero unless the relevant exponents agree,
and in the remaining case both are the same product of
factorials.  Iteration proves
$[P,Q(z)F]=[Q(\partial_X)P,F]$ for every polynomial $Q$.
Only finitely many terms of an exponential contribute to this
pairing.  Therefore, writing $Q(z)=\sum_b q(z_b)$, (A8d) gives
\[
 [P,\exp(-kQ)F]
      =[\exp(-kQ(\partial_X))P,F]=[P,F],
\]
because $Q(\partial_X)$ is a linear combination of the displayed
$\Delta_{\tau ij}$ and annihilates $P$.  Apply (A8c) in each
chart, use uniqueness of the exponential coordinates, and compare
with the analytic chart.  This proves (A8e).  In particular every
polynomial supported in just one embedding is covered, regardless
of its row weights.

The archive's \texttt{HHa} is more restrictive: by
\texttt{RowDeg.lean}, it is the joint harmonic space together with
lower-row-unipotent invariance and specified degree in every row.
Those extra conditions select highest-weight vectors for the
automorphic construction.  They are not needed for the
coefficient identity just proved.  Equation (A8e) thus extends the
archive's stated harmonic comparison at this precise point; it
does not assert an extension of its automorphic generation or
integrality theorems to arbitrary weights.

\paragraph{The determinant for A6 and all normalization factors.}
If $n$ theta sections live on the same torus and the chosen
embedding has at least $n$ tangent rows, then
\[
       P=\det(X_{\sigma i b})_{1\leq i,b\leq n}
 \tag{A8f}
\]
belongs to \texttt{HHa}.  It is harmonic because it uses no
$c\sigma$ variables.  Its row degrees are one in the first $n$
rows of $\sigma$ and zero elsewhere.  Replacing a selected lower
row by that row plus a multiple of an earlier row preserves its
determinant; operations on an unselected row leave it unchanged.
This proves the required row invariance, including all rows beyond
$n$.  Its coefficients are integers, so it also satisfies the
coefficient condition in \texttt{harmA}.

In the A6 construction the prescribed line types can differ at
compact places.  Their honest common torus is then the product of
the individual polarized tori.  At the active embedding $c\mu$,
the tensor block has dual tangent label $\mu$ by the preceding
calculation.  Write those tangent coordinates as $z_{ij}$, where $i$ is a
cotangent direction and $j$ is the multiplicity-line factor.  The
product theta section is
$\widehat\Theta(z)=\prod_{j=1}^n\widehat\theta_j(z_j)$, and the
required operator is
\[
 \mathcal D=\det(\partial_{z_{ij}}),\qquad
 \mathcal D\widehat\Theta(0)
       =\det\bigl(\partial_{z_{ij}}
                         \widehat\theta_j(0)\bigr)_{i,j}.
 \tag{A8g}
\]
Expand the determinant.  In each summand precisely one derivative
acts on each separate factor, so (A8g) follows from the product
rule with coefficient one.  This polynomial need not satisfy the
row-degree or row-invariance conditions for the full
$n^2$-dimensional embedding block of the product torus.  It does
use only $\mu$ tangent coordinates, so the proved extension (A8e)
applies.  This closes the analytic comparison for the actual
different-type product used in A6, whenever the indicated
polarized product is placed in the archive datum.

To check the factorial separately, write
$\widehat\Theta(z)=\sum_d\langle T_d,z^{\otimes d}\rangle$
as in 032, (4.10).  For homogeneous $P$ of degree $d$, with the
natural symmetric-tensor pairing,
\[
          P(\partial)\widehat\Theta(0)
                 =d!\,\langle P,T_d\rangle.
 \tag{A8h}
\]
For a multiindex $\alpha$ of degree $d$, both sides equal
$\alpha!$ times the coefficient of $z^\alpha$; this proves the
formula.  The double alternation in A6 maps a symmetric product
of $n$ degree-one tensors to the two exterior products with
coefficient one.  Its determinant functional is $n!$ times the
natural symmetric-tensor pairing in (A8h).  Consequently (A8g),
the A6 wedge, and the archive's ordinary differential convention
have precisely the same determinant coefficient.  There is no
extra $n!$ in the wedge itself.

For the archive's translated theta values, its second comparison
theorem retains the additional scalar $M^{-\deg P}$ and uses
\[
 \mathcal F_{q,a}\theta(z)
   =\exp\left(-q\pi H(z,a)-\frac{q\pi}{2}H(a,a)\right)
       \theta(z+a),\qquad q=M^2.
 \tag{A8i}
\]
The translation is $a=M^{-2}\operatorname{mk}(x_2)$ and the
untranslated function is $w\mapsto\theta(Mw)$.  The oddness,
ideal-divisibility and representative conditions on $M,x,x_2$ in
that theorem remain part of its hypothesis.  They are not
consequences of an arbitrary finite Schwartz input.  Likewise,
the final automorphic convention must retain the archive's
complex conjugations and its negative standard additive
character; the analytic equality (A8e) does not silently change
032's positive additive character.

\paragraph{Specialization to the CM elliptic frame in C18.}
For a CM elliptic curve with the invariant coordinate $z$ and
Weierstrass sigma function, the reduced theta function has the
form
\[
 \theta_{\mathrm{red}}(z)
       =\exp(-e_2^*z^2/2)\sigma(z).
 \tag{A8j}
\]
This is the standard reduced-theta normalization used in
Bannai--Kobayashi, Example 1.9, pp.~10--11 of
\url{https://arxiv.org/abs/math/0610163v4}.
If $H(z,a)=h z\overline a$ with $h>0$ and
$\eta(a)=e_2^*a+\pi h\overline a$, direct substitution in (A8i)
gives
\[
 \exp\left(-\pi H(z,a)-\frac\pi2H(a,a)\right)
          \theta_{\mathrm{red}}(z+a)
       =\exp(-e_2^*z^2/2)\,
             e^{-\eta(a)(z+a/2)}\sigma(z+a).
 \tag{A8k}
\]
Thus the common factor for the canonical analytic reading of
C18's tame translated sections is
\[
      q_{\mathrm{can}}(t)
       =\exp\bigl(-e_2^*\lambda(t)^2/2\bigr)
                         \frac{\sigma(\lambda(t))}{t}.
 \tag{A8l}
\]
Equal-degree ratios cancel this entire factor.  For tensor degree
$d$ the exact normalized operator on its integral coefficient is
$q_{\mathrm{can}}^{-d}D(q_{\mathrm{can}}^d\,\cdot)$.
On an elliptic curve with a CM field acting through a single
tangent embedding, the opposite tangent multiplicity is zero.
The quadratic ambiguity (A8d) is therefore zero, and (A8c)
identifies every nonnegative-order canonical jet with the actual
Taylor derivative of the reduced-theta reading.  Harmonicity
does not discard the nonquadratic factor in (A8l): it is the
specified conversion from the integral local generator to the
canonical analytic generator.

\paragraph{Exact archive source locations.}
All these locators refer to the recovered mirror at commit
\nolinkurl{7cbf7a0c32fca1badf05931879a37c3f8a2b03b4}.
The common prefix for \texttt{IsThRule.lean},
\texttt{ThGerm.lean}, \texttt{IsExpChart.lean}, and
\texttt{RowDeg.lean} is
\path{Formalism/LocalToGlobal/ProjectiveLattices/}.
The prefix for \texttt{ExpTiePf.lean},
\texttt{GaugeUniq.lean}, and \texttt{GaugeUniqPf.lean} is
\path{Formalism/SharpRestriction/Divisibility/GeneratorPowers/}.
The prefix for the three files beginning
\nolinkurl{ExistsAnalyticBiextensionChart},
\nolinkurl{HarmonicValueCanonicalConnection}, and
\nolinkurl{ThetaValueAnalyticFrame}, whose full basenames are given
above, is
\path{Development/AbelianVarietyTheory/Proofs/CuspFormEvaluations/}.
The file \nolinkurl{ExistsIsNabNotIsNablaPin.lean} is under
\path{Development/AbelianVarietyTheory/Proofs/UnitaryExistence/}.
Finally \texttt{ThConst.lean} is under
\path{Development/AbelianVarietyTheory/Definitions/ThetaValues/}.

The arithmetic assertion requires that the chosen imaginary
quadratic CM order, its action, the rigidified polarization and
the invariant differential be the ones in the archive datum.
Neither a general CM field nor an unspecified analytic torus
provides those data automatically.  For the weak tensor-degree
estimate of C18, passage from its raw sigma normalization to
(A8l) additionally requires control of $e_2^*$; equal-degree
ratios do not require it.  If $e_2^*$ is integral, the same Bell
proof extends by allowing the quadratic term: its new weight
inequality is $1/(p-1)\leq2p/(p^2-1)$.  This last observation is
an elementary implication, not an unproved integrality assertion
about $e_2^*$.

\paragraph{The resulting inert-prime specialization.}
Under C19's hypotheses, the required integrality of $e_2^*$ is
proved.  Suppose in addition that the algebraic differential,
rigidification and tame torsion lifts are matched as above, and
that any change of the integral tangent frame has unit determinant.
For a local tensor expression of degree $d$, its canonical
analytic reading is $Q^dv$ with $Q=q_{\mathrm{can}}$ and
$v\in R[[t]]$.  Its order-$r$ canonical local jet is therefore
$D^r(Q^dv)(0)$.  C19 gives the explicit bound
\[
 v_p\bigl(D^r(Q^dv)(0)\bigr)
 \geq\max\left\{0,\frac{r}{p^2-1}
                         -\frac{s_p(r)}{p-1}\right\}.
\]
For equal-degree quotients the common factor cancels and the
bound improves to
\[
 v_p\bigl(D^rv(0)\bigr)
 \geq\max\left\{0,
          \left\lfloor\frac{rp}{p^2-1}\right\rfloor-1\right\}.
\]
For actual archive sections of nonnegative power these are the
jets in (A8c).  For local ratios and negative tensor degree, they
are the induced local connection expressions, obtained by the
product and quotient rules.  They make no assertion about the
unconstrained negative-power values of an arbitrary
\texttt{NabV} satisfying only \texttt{IsNab}.

The level normalization can also be checked exactly.  For the
degree-one reduced theta function and integer $M$, linearity of
$H$ gives
\[
 \mathcal F_{M^2,a}\bigl(z\mapsto\theta_{\mathrm{red}}(Mz)\bigr)(z)
      =\bigl(\mathcal F_{1,Ma}\theta_{\mathrm{red}}\bigr)(Mz).
\]
An order-$r$ derivative on the right contributes $M^r$, which
cancels the archive's $M^{-r}$ in (A8i).  Thus, when $M$ and the
relevant torsion order are prime to $p$, the displayed C19 bounds
apply to this normalized value without an additional level
denominator.  This is an elliptic specialization with specified
normalizations, not an assertion that the higher-dimensional
032 torus has such an elliptic realization.

\framework{The comparison fixes a rigidified line bundle, its
formal completion, and the correspondence between an algebraic
section and its analytic reading.  It is compatible with the
correspondence and polarization viewpoint in Grothendieck's
\emph{cote44}, scans 7--8, and with the dimension and Lefschetz
discussion in \emph{cote16}, scans 43--46.  No algebraic inverse
Lefschetz operator or Chow projector is inferred from equality of
analytic jets.}


\section{Definite theta periods and the exact iterated seesaw}

There are now two precise global comparisons. First, an orthogonal decomposition of the rank-three Hermitian space gives a totally definite rank-two subgroup. Restriction to that subgroup is a genuine global unitary Bessel period, and the theta tensor rule gives an exact seesaw identity; this is proved in THETA3-B11--B12. Second, the recovered Alp\"oge source supplies a common normalized functional for an entire family of iterated lifts. Its exact statement is \texttt{MainX}, and its source witness and elementary scalar consequence are recorded in THETA3-B13. Identifying the particular cohomological inputs with that family, including its finite lattices and normalization, remains an instantiation problem.

The inner-form calculations below describe a different possible route. Their obstructions do not obstruct the definite-subgroup construction. All new proofs have been read in full for the targeted checks recorded in the accompanying event files. Source statements used as background are distinguished from these proofs. No Lean compilation, complete dependency audit, or validation of the cited RH theorem is claimed or attempted.

\subsection{THETA3-B0: what the supplied and recovered sources establish}\label{B0}

The attached four-page Alp\"oge note gives a timeline. Step 4 mentions determinant positivity, step 7 GGP periods and interlacing, step 8 iterated theta lifts on totally definite unitary groups, and step 12 a seesaw producing a Bessel period independent of the Hecke characters. The note does not specify the integral, test vectors, measures, or comparison with differential forms. The attached source map is explicitly a text-parsing report, with no Lean build. Its paths are useful search locators; they are not substitutes for the definitions at those paths. The attached \texttt{Main.lean} and \texttt{Statements.lean} contain the final analytic declarations, not the theta-period definitions. This is an evidence finding, not a refutation of the intended comparison.

During this run, the named definitions and several proof bodies were recovered from the secondary public mirror \path{CocoSgt/riemann-hypothesis-lean}, pinned to commit
\begin{center}\small\nolinkurl{7cbf7a0c32fca1badf05931879a37c3f8a2b03b4}.\end{center}
The commit date is 20 November 2026. The recovered \texttt{Main.lean} and \texttt{Statements.lean} are byte-for-byte equal to the attachments. A manifest records the pinned URL, SHA256, Git blob hash, size, and retrieval time of each fetched file. This is partial identity evidence; it does not assert a complete primary-archive hash match. Consequently the initial missing-definition gate has been reduced to the exact mathematical matching problem in B13.

For reference, the 032 setting is a Galois CM extension $E/F$, with $[F:\Q]\geq2$, and an $E$-Hermitian space $V$ of dimension three. At one real place $\mu$ it has signature $(2,1)$, and elsewhere signature $(3,0)$. The compact ball surface is attached to $\SU(V)$ at sufficiently small principal level. For multiplicity space $P$ of dimension $k=1$ or $2$, the alternating form is
\[
 \psi_P=\Tr_{E/F}\bigl(\delta h_V\otimes h_P\bigr),
 \qquad \overline\delta=-\delta.
\]
At the active place the positive Lagrangian is
\[
 \mathcal F_P=(L\otimes P_\mu)\oplus(Q^*\otimes P_\mu^*),
 \qquad Q=V_\mu/L.
\]
The determinant square root and projected Taylor tensors give $\Omega^1$ and $K_X$; see 032, Lemma 4.1 and equations (4.5)--(4.8), pp.14--16. Proposition 4.6 and Corollary 4.7 vary the multiplicity-space data at compact places while keeping $V$ fixed. This operation does not change the signature of $V$.

\subsection{THETA3-B1: coherent rank-three definite neighbors}\label{B1}

\textbf{Statement.} There is a global Hermitian space $V^-$ with signature $(0,3)$ at $\mu$ and $(3,0)$ elsewhere, locally isometric to $V$ at every finite place. There is no everywhere positive definite $V^+$ with all those finite local isometry classes. For any finite nonsplit place $v_0$, there is an everywhere positive definite $V^{(v_0)}$ locally isometric to $V$ outside $\{\mu,v_0\}$, with the opposite determinant norm class at $v_0$.

\textbf{Proof.} Let $\epsilon_v$ be the local quadratic character of the determinant. At a real place it is $(-1)^{q_v}$, where $q_v$ is the negative index. A coherent collection has $\prod_v\epsilon_v=1$. Local dimension and determinant classify Hermitian spaces at finite places; together with the specified real signatures, the product condition gives global existence by Landherr's theorem. An exact source is B.H. Gross, \emph{Incoherent definite spaces and Shimura varieties}, 25 February 2018, \S3, Theorem 3.1, pp.6--7.

Changing the negative index at $\mu$ from one to three preserves $\epsilon_\mu$, so all finite data can be retained. Changing it from one to zero reverses $\epsilon_\mu$, contradicting the product condition if all other local data are retained. Reversing the determinant class at one nonsplit finite place reverses one more sign and restores the product condition. Landherr's theorem gives the asserted global spaces. \hfill$\square$

This is a statement about Hermitian isometry data and oscillator embeddings. Multiplying an odd-dimensional Hermitian form by a scalar can leave its abstract unitary group unchanged. Such a group identification does not identify the dual-pair embedding, splitting, or test vector.

\subsection{THETA3-B2: the scalar type distinguishes the definite neighbors}\label{B2}

\textbf{Statement.} Keep the active-place $\delta$ and multiplicity-space sign used in 032, and put $r=t^2$ on the double cover of the scalar circle of $\U(P_\mu)$. In the covariant Taylor model of 032, the degree-$k$ cohomological tensor has weight $t^{-3k}$. For the negative definite replacement in B1 the polynomial model
\[
 \bigoplus_{j\geq0}(\det\mathcal F^-)^{1/2}\otimes
                    \operatorname{Sym}^j\mathcal F^-
\]
has only weights $t^{3k+2j}$. The required weight is absent. For the positive definite replacement, its vacuum has weight $t^{-3k}$.

\textbf{Proof.} The scalar $r$ acts by $r$ on the $k$-dimensional summand $L\otimes P_\mu$ and by $r^{-1}$ on the $2k$-dimensional summand $Q^*\otimes P_\mu^*$. Hence $\det\mathcal F_P$ has weight $r^{-k}$, and its chosen square root has weight $t^{-k}$. The projected degree-$k$ tensor contributes $r^{-k}=t^{-2k}$, giving $t^{-3k}$.

Write $\delta_\mu=ia$, with $a\in\R^*$. The prescribed sign of $h_P$ is $-\operatorname{sgn}(a)$, as in 032, Lemma 4.1. When $h_V$ is negative definite, $a h_Vh_P$ is positive definite. The positive Lagrangian is then the nonconjugate summand of dimension $3k$, on which the scalar acts by $r$. Its half determinant contributes $t^{3k}$ and polynomial degree $j$ contributes $t^{2j}$. None gives $t^{-3k}$. When $h_V$ is positive definite, the positive Lagrangian is the conjugate summand, with scalar $r^{-1}$; its vacuum therefore has weight $t^{-3k}$. \hfill$\square$

The actual infinite Schr\"odinger coefficient vectors carry the dual action in 032, equation (4.13). Inverting both displayed spectra preserves both conclusions. A common honest character twist also preserves the mismatch. The same names for two splitting characters do not establish the same correction from the bare metaplectic action. This computation does not assert that an unspecified global theta packet is empty. Positive definite vacuum compatibility is necessary for this particular tensor comparison, not a sufficient identification of periods.

\subsection{THETA3-B3: an explicit local epsilon obstruction}\label{B3}

Fix a nonsplit quadratic local extension $E/F$, an additive character, $\delta$, a skew-Hermitian line $W$, a splitting character $\chi_W$, and a character $\alpha$ of the norm-one torus. If $H_+$ and $H_-$ are the two Hermitian lines of opposite determinant classes, then exactly one of the corresponding lifts of $\alpha$ to $W$ is nonzero.

Indeed, Borade--Franzel--Girsch--Yao--Yu--Zelingher, \emph{On tori periods of Weil representations of unitary groups}, Selecta Math. 31 (2025), article 49, Theorem 3.5, gives the condition
\[
 \epsilon_{E/F}(\chi_W^{-1}\alpha_E,\psi,\delta)
       =\epsilon(H)\epsilon_\delta(W),
 \qquad \alpha_E(x)=\alpha(x/\overline x).
\]
The left side is fixed and the two right sides are opposite signs, so exactly one equality holds. This proves the assertion as an application of that theorem. The source's \S\S3.2.1 and 4.2.1 state the splitting compatibilities before the seesaw identities. The example shows why local signs cannot be suppressed merely because a global parameter or a branching dimension agrees. It does not identify every Alp\"oge Bessel period with this particular torus period.

\subsection{THETA3-B4: normalization is indispensable for a numerical bound}\label{B4}

\textbf{Statement and proof.} Nonvanishing alone gives no positive lower bound independent of scalar normalization. If a nonzero theta input is replaced by $\varepsilon$ times itself, with $\varepsilon>0$, the resulting sesquilinear norm or simultaneous mixed pairing is multiplied by $\varepsilon^2$. It remains nonzero and tends to zero. Likewise, a nonzero constant determinant-line trivialization can be rescaled. Thus every numerical lower bound must specify the input norms and period conventions. \hfill$\square$

\subsection{THETA3-B5: a quantitative compact projection certificate}\label{B5}

\textbf{Statement.} Let $M$ be an $N$-dimensional unitary representation of a compact group $H$, with probability Haar measure. Let
\[
 P=\int_H\rho(h)\,dh
\]
have rank $r>0$. Some $r\times r$ principal minor of $P$ is at least $\binom Nr^{-1}$. In particular, when $r=1$, a normalized invariant functional has squared value at least $1/N$ on some vector of any orthonormal basis.

\textbf{Proof.} Haar invariance gives $P^2=P$; unitarity and invariance under inversion give $P^*=P$. Thus $P$ is the orthogonal projection onto $M^H$. Choose an $N\times r$ matrix $U$ whose columns form an orthonormal basis of that subspace. Then $P=UU^*$ and $U^*U=I_r$. Cauchy--Binet gives
\[
 1=\det(U^*U)=\sum_{|I|=r}|\det U_I|^2.
\]
There are $\binom Nr$ terms. For some $I$, its term is at least the reciprocal of this number, and that term is $\det P_{I,I}$. For $r=1$ these are exactly the squared coefficients of a normalized invariant vector, or equivalently of its normalized invariant functional. \hfill$\square$

For compact branching, take $M=\pi\otimes\sigma^\vee$. A multiplicity-one statement then gives a selected normalized period of squared size at least $1/(\dim\pi\dim\sigma)$. This does not set an unnormalized global automorphic integral equal to one.

\subsection{THETA3-B6: iteration requires an overlap estimate}\label{B6}

\textbf{Statement.} Positivity and nonzero eigenvalues of individual stages do not bound their composite. If instead
\[
 \|T_jw\|\geq b_j\|w\|,\qquad b_j>0,
\]
holds on the image of $T_{j-1}\cdots T_1$, then
\[
 \|T_m\cdots T_1v\|\geq\left(\prod_{j=1}^m b_j\right)\|v\|.
\]

\textbf{Proof.} In $\C^2$, let $P$ project onto $e_1$ and $Q_\theta$ onto $(\cos\theta,\sin\theta)$. Both have rank one and nonzero eigenvalue one, but $\|Q_\theta P\|=|\cos\theta|$, which can be arbitrarily small or zero. Under the stated additional hypothesis, apply the bound for $T_1$, then that for $T_2$ to $T_1v$, and continue. Each argument lies in the specified image, and multiplying the inequalities proves the result. \hfill$\square$

Thus multiplicity gained at each theta-lift stage does not, by itself, give a lower bound for a fixed composite period.

\subsection{THETA3-B7: an arithmetic Gram-matrix lower bound}\label{B7}

\textbf{Statement.} Let $K$ be CM with real subfield $K_0$ of degree $d$. Let $G\in M_r(K)$ be Hermitian and positive definite at every embedding $\sigma:K_0\hookrightarrow\R$. Suppose $DG$ has algebraic-integer entries for a positive integer $D$, and
\[
 \operatorname{tr}(G_\sigma)\leq T_\sigma,\qquad T_\sigma>0.
\]
For a chosen $\sigma_0$ and $r\geq2$,
\[
 \lambda_{\min}(G_{\sigma_0})\geq
 D^{-rd}\prod_{\sigma\ne\sigma_0}\left(\frac r{T_\sigma}\right)^r
       \left(\frac{r-1}{T_{\sigma_0}}\right)^{r-1}.
 \tag{B7}
\]
For $r=1$, omit the last factor. The right side also bounds
$a^*G_{\sigma_0}a/\|a\|_2^2$ from below for nonzero $a$.

\textbf{Proof.} The determinant $\Delta=\det G$ lies in $K_0$. The number $D^r\Delta$ is a totally positive algebraic integer, so its norm to $\Q$ is a positive integer and at least one. Consequently
\[
 \prod_\sigma\det G_\sigma\geq D^{-rd}.
\]
For $\sigma\ne\sigma_0$, the arithmetic--geometric mean inequality gives
$\det G_\sigma\leq(T_\sigma/r)^r$. Divide to obtain the corresponding lower bound for $\det G_{\sigma_0}$. The product of the other $r-1$ eigenvalues at $\sigma_0$ is at most $(T_{\sigma_0}/(r-1))^{r-1}$, so division gives (B7). When $r=1$ the determinant is the single eigenvalue. The final assertion follows from the spectral theorem. \hfill$\square$

The bounds at the other embeddings are essential: the totally positive unit $(3-2\sqrt2)^m$ tends to zero at one real embedding while its conjugate tends to infinity. On a finite definite adelic module, algebraic-valued vectors and rational mass weights can supply such a matrix once the field, denominator, and conjugate bounds are specified. For cohomological periods one must likewise identify an algebraic normalization and the relevant positivity at all embeddings; arbitrary complex theta coordinates do not have these properties automatically.

\subsection{THETA3-B8: a conditional two-place transfer inequality}\label{B8}

Suppose compatible normalized positive period formulas have been established:
\[
 I_{\mathrm{ball}}=C_{\mathrm{ball}}L\prod_{v\in S}Z_{\mathrm{ball},v},
 \qquad
 I_{\mathrm{def}}=C_{\mathrm{def}}L\prod_{v\in S}Z_{\mathrm{def},v},
\]
with the same nonzero $L$-factor and identical local factors outside
$\{\mu,v_0\}$. If
\[
 I_{\mathrm{def}}\geq b,\qquad
 \frac{C_{\mathrm{ball}}}{C_{\mathrm{def}}}\geq c>0,
 \qquad \frac{Z_{\mathrm{ball},v}}{Z_{\mathrm{def},v}}\geq a_v>0
 \quad(v=\mu,v_0),
\]
then
\[
 I_{\mathrm{ball}}\geq bc\,a_\mu a_{v_0}.
\]
The proof is division of the two identities followed by cancellation and multiplication of the displayed lower bounds. This proves an implication; it does not supply the two identities for the particular 032 and Alp\"oge inputs.

An applicable Rallis statement must be checked in its actual direction and range. Gan--Qiu--Takeda, arXiv:1207.4709v3, \S11.3, p.52, displays the doubling seesaw and its determinant-character twist. Section 11.4 assumes $d(n)<m\leq2d(n)$ before Proposition 11.1, including its convergent clause, and Theorem 11.4 repeats that range. For unitary groups $d(n)=n$, so it does not directly cover the naive $\U(1)\to\U(3)$ substitution. Dividing a local zeta integral by its standard local $L$-factor can still leave the normalizing denominator $d_v(s,\chi_V)$; one must not silently declare all such spherical factors to be one.

An applicable positive locator for the $\U(1)\to\U(3)$ direction is J. Haan, arXiv:1403.5061v7, \S3.3, Theorem 3.4, p.8. With Hermitian source dimension one, skew-Hermitian target dimension three, splitting $(\gamma,\gamma^3)$, $\gamma|_{\A_F^*}=\chi_{E/F}$, and the paper's Tamagawa and symplectic conventions, it states
\[
 \langle\theta(\overline f_1,\phi_1),\theta(\overline f_2,\phi_2)\rangle
 =\frac{L_E(3/2,\operatorname{BC}(\tau)\otimes\gamma^3)}
        {L_F(3,\chi_{E/F})}
   \prod_v Z_v^{\#}(1,f_{1,v},f_{2,v},\Phi_{1,v}).
\]
Here $Z_v^{\#}=L_v(3,\chi_{E_v/F_v})Z_v/L_{E_v}(3/2,\operatorname{BC}(\tau_v)\gamma_v^3)$ and the doubled section comes from $\rho_{1,3}(\phi_1\otimes\overline\phi_2)$. The source proves it by the convergent formula of J.-S. Li; it is not an application of the restricted GQT theorem just audited. Transport of the splitting, the conjugated source input, and the measures is still required. The formula controls an individual torus theta lift; it does not automatically give the norm of an arbitrary wedge of two such lifts. Haan's main Bessel theorem has additional special-pair hypotheses.

The recovered archive also retains a quotient-volume factor explicitly. In
\path{Development/AutomorphicArithmetic/ThetaCorrespondence/RallisInnerProduct/GraphRelationExtendsFromFockVectorsToKFinite.lean}, lines 791--802, the proof of \nolinkurl{rallisInnerProductConvergent_of} passes through the identity
\[
 \frac{\langle\Theta(\Phi_1,f_1),\Theta(\Phi_2,f_2)\rangle_B}
      {\operatorname{vol}([B])}
 =\operatorname{dblInt}_{A,\chi_B^{-1}}
       (\operatorname{ThetaMean}_B(\delta(\Phi_1,\Phi_2));f_1,f_2).
\]
Lines 776--789 explicitly convert the integral over $[B]$ to its volume times the normalized theta mean. The proof uses the conjugated source functions in its Fubini step; its lift convention is not interchangeable with \texttt{liftS}. Indeed, \nolinkurl{GlobalTheory.RallisIP.thetaLift}, defined in \path{PublishedTheorems/TwistedTraceFormula/LowerLiftsZeroOfWittIdxEqZero.lean}, lines 22--26, is $h\mapsto\int_{[A]}\theta(g,h;\Phi)\overline{f(g)}\,dg$; lines 41--59 state the resulting range and volume-normalized zeta identity. Its input \nolinkurl{Published.SiegelWeilUnitaryConvergent}, defined in \path{PublishedTheorems/SiegelWeil$_2$/SiegelWeilUnitaryConvergent.lean}, lines 18--35, requires
\[
 \operatorname{rank}A<\operatorname{rank}B\leq2\operatorname{rank}A,
 \qquad \operatorname{rank}A+\operatorname{wittIdx}B<\operatorname{rank}B.
\]
Thus this archive formula does include the anisotropic adjacent-rank alternatives $1\to2$ and $2\to3$ of B15.

A \emph{different} law, \nolinkurl{SphericalVolumeLift.Convergent}, in \path{BlockTheory/TwistedDiscreteSeries/SphericalVolumeLift.lean}, lines 75--91, imposes $\operatorname{rank}B=\operatorname{rank}A+t$ and $\operatorname{rank}A<t$. Its range is $\operatorname{rank}B>2\operatorname{rank}A$; it is this larger-rank law that the recovered definite-positivity argument uses. Neither of these two ranges includes equal rank. There is also a separate \nolinkurl{Published.SiegelWeilAnisotropicUnitary} statement in \path{PublishedTheorems/TwistedTraceFormula/SiegelWeilAnisotropicUnitary.lean}, lines 109--119: with anisotropic $B$ and positive ranks, its Eisenstein value is $(1\text{ if }\operatorname{rank}A<\operatorname{rank}B,\,2\text{ otherwise})$ times the normalized theta mean. This supplies an analytic identity at equal rank, not automatic nonvanishing of a given theta lift. The volume, character, and rank hypotheses must be kept separately for each formula.

\subsection{THETA3-B9: the canonical local Bessel integral equals one}\label{B9}

The recovered \path{Formalism/PeriodIntegrals/BesselPeriodGGP.lean} defines a real local nonvanishing proposition. Its ambient matrix \texttt{Hpk} has signature $(k+r,r+1)$; the compact subgroup \texttt{UWset} fixes the hyperbolic coordinates and a negative vector. Its integral has the form
\[
 \int_{\U(W_z)\times N}
 \langle\sigma(x_0)w,\sigma(hn)\sigma(x)w\rangle
       \overline{\psi_N(n)}\,\pi(h)^{-1}\,dh\,dn.
 \tag{B9a}
\]
It is operator-valued in the finite-dimensional $\pi$-space. The proposition asks for a discrete-series $\sigma$, a lowest-$K$-type unit vector, compact $\pi$, interlacing and half-integral differences, and asserts existence of $x,x_0$ making the integral nonzero. It does not say that arbitrary periods are zero or one. \texttt{GapBig} is in a separate numerical setting, not in the actual proposition's assumptions.

\textbf{Statement.} Put $k=2$, $r=0$, and $z=e_3$, and use probability Haar measures. Suppose a unitary representation of $\U(2,1)$ contains a unit vector $w$ with
\[
 \sigma(\operatorname{diag}(A,d))w=\det(A)d^{-2}w,
 \qquad A\in\U(2),\ d\in\U(1).
\]
Take $\pi(\operatorname{diag}(A,1))=\det A$ on $\C$. Then (B9a), at $x=x_0=1$, equals one.

\textbf{Proof.} The matrix equations defining \texttt{Nset} at $r=0$ force $N=\{1\}$. A unitary matrix fixing $e_3$ preserves its positive orthogonal complement, so \texttt{UWset} is precisely $\operatorname{diag}(\U(2),1)$. In the convention linear in the second inner-product variable,
\[
 \langle w,\sigma(\operatorname{diag}(A,1))w\rangle=\det A.
\]
Multiplication by $\pi(h)^{-1}$ cancels this factor. The integrand is identically one; compactness gives absolute convergence, and probability Haar gives the stated value. \hfill$\square$

This is the correct type for the right-regular coefficient of a canonical two-form. At a negative line $o$, the tangent space is $\Hom(L,Q)$. For $k=\operatorname{diag}(A,d)$, its derivative is $A/d$. If
\[
 f_\xi(g)=(g^*\xi)_o(e_1,e_2),
\]
then
\[
 f_\xi(gk)=\det(A)d^{-2}f_\xi(g).
\]
Left lattice invariance follows from that of $\xi$. On a compact quotient this coefficient has finite nonzero $L^2$ norm when $\xi\ne0$, and can be normalized. The displayed character is trivial on the center of $\U(2,1)$, so it is compatible with the projective special-unitary action. It is a coefficient convention; the action on the canonical-bundle fibre is dual.

The exact numerical choice in the archive is
\[
 \lambda=(1,0,-1),\quad v=(1),\quad b=(0,1),\quad
 u=(3/2,1/2),\quad\epsilon=(+,+,-).
\]
The definitions in \texttt{UpqRep.lean} give
\[
 \rho_c=(1/2,-1/2,0),\quad\rho_n=(1/2,1/2,-1),\quad
 \lambda+\rho_n-\rho_c=(1,1;-2).
\]
For compact $\U(2)$, $u-\rho_c=(1,1)$, so $\pi=\det$. The interlacing is
$3/2>1>1/2>0$; the Harish--Chandra integrality and half-integral differences follow directly from these tuples. The direct evaluation does not require the general discrete-series theorem. Invoking that theorem instead would additionally require the relevant \texttt{IsDSHpk} representation identification. No global period has been set equal to one.

The same direct calculation works for a pluricanonical coefficient of $K_X^q$, $q\geq1$. Its compact character is $(\det A)^qd^{-2q}$, and taking $\pi=\det^q$ again makes the local period one. The corresponding parameters are
\[
 \lambda=(q,q-1,1-2q),\quad v=q,\quad
 b=(1-q,2q-1),\quad u=(q+1/2,q-1/2).
\]
The same half-sum calculation gives lowest weight $(q,q;-2q)$, and the interlacing follows from
$q+1/2>q>q-1/2>q-1$. The node gap is $3q-2$. Thus powers can enlarge the numerical gap, whereas a determinant twist shifts both nodes equally and cannot do so. A section $\xi^q$ is pluricanonical; it is not the original ordinary cohomology class in $H^2(X)$. Passing a quantitative period bound back to that class needs an additional argument.

\subsection{THETA3-B10: a separate high-rank matching theorem}\label{B10}

The recovered \path{Development/AbelianVarietyTheory/BlockTheory/DivisorBase/WeilTensorModels.lean} defines one global matching theorem. It is distinct from the iterated Bessel family in B13. Its two spaces have dimensions $a,a'$, and its common target has dimension $a+a'$. The exact equality \texttt{IsConst}, line 491, is
\[
 I_V(\Phi)=c\,I_{V'}(\Phi')
\]
for matched data. \texttt{OneConst}, line 503, requires uniqueness and $c\ne0$. \texttt{StepAlpha} and \texttt{StepBeta}, lines 571 and 559, contain recursive constants involving partial Hecke $L$-values and local epsilon/Gauss factors. \texttt{LocalMatching} and \texttt{ThetaIntegralsAgree} occur at lines 581 and 585.

Its \texttt{Geo.IsGood}, lines 360--370, requires
\[
 a'<a,\quad a+a'\equiv0\pmod2,\quad
 2(p+q)+1<\left\lfloor\frac{a-a'}2\right\rfloor,\quad
 2(p+q)+1<a',
\]
and signatures
\[
 V:(a-p,p),\quad V':(q,a'-q),\quad
 W:(a+a'-p-q,p+q).
\]
The definitions \texttt{Spc.IsHerm} and \texttt{Spc.IsSkew}, lines 182--200, impose these signatures at every real place. The base also has CM/Galois and integral unramified-extension restrictions.

For $p=1,q=0$, the inequalities force $a'\geq4$ and the even difference $a-a'\geq8$; the smallest pair is $(12,4)$, with common target dimension sixteen. Thus assigning $a=3$ does not instantiate this theorem. Increasing the dimension of a one-active-place ball does not by itself repair its all-place signature or source-dimension requirements. This calculation audits this one matching theorem; it does not exclude the different definite-subgroup route or the abstract family theorem below.

\subsection{THETA3-B11: the exact restriction theta seesaw}\label{B11}

\textbf{Statement.} Let $V=V_H\perp L$ be Hermitian, with $L$ a line. Put $G=\U(V)$ and $H=\U(V_H)$, embedded by fixing $L$ pointwise. Let $W$ be a skew-Hermitian source and $T=\U(W)$. Assume $[H]$ and $[T]$ are compact, or assume absolute convergence of all sums and integrals below. Fix compatible oscillator splittings for the orthogonal tensor rule. In particular the splitting character for $V$ is the product of those for $V_H$ and $L$, and the source splittings agree.

For a finite sum of Schwartz tensors
$\Phi=\sum_a\Phi_{H,a}\otimes\Phi_{L,a}$ and a unitary automorphic character $\eta$ of $T$, set
\[
 F_\eta(g)=\int_{[T]}\theta_V(g,t;\Phi)\eta(t)^{-1}\,dt,
 \qquad
 B_e(F)=\int_{[H]}F(h)\overline{e(h)}\,dh.
\]
For a continuous automorphic test function $e$ on $[H]$, define
\[
 \Theta_{H\to T}(e;\Phi_H)(t)=
 \int_{[H]}\theta_{V_H}(h,t;\Phi_H)\overline{e(h)}\,dh.
\]
Then
\[
 B_e(F_\eta)=\sum_a\int_{[T]}
 \Theta_{H\to T}(e;\Phi_{H,a})(t)
 \theta_L(1,t;\Phi_{L,a})\eta(t)^{-1}\,dt.
 \tag{B11}
\]
All measures are the displayed measures, so no unspecified scalar occurs.

\textbf{Proof.} A polarization compatible with the orthogonal sum identifies the rational indexing lattice with the direct sum of the two indexing lattices. The tensor identity for the oscillator action identifies each summand for a pure tensor with a product. Absolute convergence permits separation of the sums, giving
\[
 \theta_V((h,1),t;\Phi_H\otimes\Phi_L)
       =\theta_{V_H}(h,t;\Phi_H)\theta_L(1,t;\Phi_L).
\]
Sum over $a$, multiply by $\overline{e(h)}\eta(t)^{-1}$, and integrate. Fubini applies by the convergence hypothesis. The line factor and $\eta$ are independent of $h$, and the inner $h$-integral is the displayed reverse theta lift. This proves (B11). In the compact case, uniform convergence on compact representative sets and boundedness of $e$ supply absolute integrability. \hfill$\square$

For a theta differential form, the statement applies to its polynomial Fock coefficient once that coefficient and its orthogonal tensor decomposition have been identified. This is a required identification of vectors, not merely of form degrees.

The actual archive declaration is
\nolinkurl{GlobalTheory.T.thetaTensorRule_orthogonalSum_source}, in
\path{Proofs/DecoratedLatticeProofs/ThetaTensorRuleOrthogonalSumSourceStatement.lean}, with the proof in the matching file without \texttt{Statement}. The property \texttt{TensorSource}, in \path{BlockTheory/TwistedDiscreteSeries/ConjugateAdjointDagger.lean}, lines 27--53, explicitly imposes the splitting product and tensor action. The adjacent target tensor rule is used by \nolinkurl{LastSeeSawThetaIdentity.lean}.

The precise global Bessel orientation is also recovered. In
\path{Definitions/TwistedInnerProducts/WittSmallGroupEmbeddingMemPointwiseStabilizerIso.lean}, lines 176--181, its integrand is
\[
 \phi(h)f(n\,\operatorname{emb}_A(h))\overline{\psi_{\rm Bes}(n)}.
\]
Its \texttt{liftS\_apply}, line 163, integrates $\theta(g,h)G(h)$ without conjugating $G$. Thus (B11) matches by taking $\phi=\overline e$ and $G=\eta^{-1}$. At Bessel rank zero the unipotent is trivial.

The source character $\eta$ remains visible in (B11). The identity therefore supplies a genuine global definite-theta/Bessel comparison, but does not on its own prove the cancellation claimed for the special iterated family. Nor does it equate the full ball norm with its restriction norm.

\subsection{THETA3-B12: every nonzero canonical form has a definite Bessel coefficient}\label{B12}

\textbf{Statement.} Let $X$ be the compact congruence ball surface of B0 and $0\ne\xi\in H^0(X,K_X)$. There is an $E$-rational negative line $L$ such that $H=\U(L^\perp)$ is totally definite of rank two and a finite-level automorphic test function $e$ with $B_e(f_\xi)\ne0$. This evaluates a canonical-bundle coefficient, not the cohomological pullback of $H^2(X)$ to a point.

More precisely, use probability Haar on $[H]$, a compact open finite level $K_H$ fixing $f_\xi$, and the finite-dimensional space $R$ with this level and the same one-dimensional infinity type. Put $N_H=\dim R$, let $m_0>0$ be the mass of the identity double coset, and let $b=|f_\xi(1)|$. For any orthonormal basis $(e_j)$ of $R$,
\[
 \sum_{j=1}^{N_H}|B_{e_j}(f_\xi)|^2
      =\|f_\xi|_{[H]}\|_2^2\geq m_0b^2,
 \qquad
 \max_j|B_{e_j}(f_\xi)|^2\geq\frac{m_0b^2}{N_H}.
 \tag{B12}
\]
In particular, $\xi\ne0$, $\int_X\xi\wedge\overline\xi>0$, and existence of some such nonzero Bessel coefficient are equivalent.

\textbf{Proof.} First make the adelic realization precise. Fix the $\mathcal O_E$-lattice defining the principal congruence quotient. An element of $\U(V)(F)$ preserving this lattice has determinant in $\mathcal O_E^\times$ with $\alpha\overline\alpha=1$. Every complex conjugate of $\alpha$ has absolute value one, so $\alpha$ is a root of unity. There are finitely many such roots in $E$. Choose a principal congruence level at which no nonidentity root is congruent to one. Then
\[
 \Gamma'=\U(V)(F)\cap K_f\subset\SU(V)(F).
\]
Make the level deeper so that $\Gamma'\subset\Gamma$, where $X=\Gamma\backslash D$; this inclusion has finite index. Thus $X'=\Gamma'\backslash D$ is both a finite principal cover of $X$ and a connected component of the adelic $\U(V)$ quotient. Pull $\xi$ to $X'$ and extend it by zero on the other components. This gives the automorphic canonical-bundle coefficient $f_\xi$ used below. Its central archimedean character is trivial, since $\det(tI_2)t^{-2}=1$. If the covering degree is $d$, then
\[
 \int_{X'}p^*\xi\wedge\overline{p^*\xi}
       =d\int_X\xi\wedge\overline\xi,
\]
so nonvanishing and positivity are unchanged. The masses and finite level in (B12) refer to this chosen realization. We retain the notation $\xi$ for its pullback.

The nonzero locus of $\xi$ is a nonempty open subset of the ball. Rational vectors are dense at the active embedding: the image of $E$ in $\C$ contains $\Q+\Q\alpha$ for a nonreal $\alpha$, hence is dense. In coordinates, $V(E)$ is therefore dense in $V_\mu$. Choose $z\in V(E)$ whose negative line lies in this open locus. At every other real place the form on $V$ is positive definite, so $h_\tau(z,z)>0$. Thus $L^\perp$ has signature $(2,0)$ everywhere and $H$ is totally definite.

Base the canonical coefficient at $L$, choosing a nonzero tangent two-frame there. The calculation of B9 gives the right compact character $\det(A)d^{-2}$, and on $H_\infty$, which fixes $L$ pointwise, it is $\det A$ at the active place. The other compact factors have the prescribed canonical coefficient type, trivial for the trivial coefficient system. Its absolute value is therefore constant along each right $H_\infty$ orbit. By construction $f_\xi(1)\ne0$. The adelic realization just constructed gives an open compact $K_H$ fixing the coefficient.

The set $H(F)\backslash H(\A_F)/(H_\infty K_H)$ is finite. On each double coset a function of the specified scalar infinity type is determined by one complex number, subject to its finite stabilizer condition. Thus $R$ is finite-dimensional. Let $m_i>0$ be the masses of these cosets. The identity coset contributes exactly $m_0b^2$ to the norm square, and all other terms are nonnegative. Consistent Haar choices give masses proportional to reciprocal finite stabilizer orders, normalized by $\sum_i m_i=1$. At neat level the stabilizers are trivial and the masses are equal.

Parseval in $R$ proves the equality in (B12); one term is at least their average. The orthonormal basis may be refined within the automorphic spectral summands of the compact quotient, so these are periods against definite automorphic test vectors. This assertion does not require their values or their Hilbert normalizations to be algebraic.

Conversely, a nonzero Bessel coefficient forces $f_\xi$, and hence $\xi$, to be nonzero. In local holomorphic coordinates, if $\xi=a\,dz_1\wedge dz_2$, then
\[
 \xi\wedge\overline\xi
   =4|a|^2\,dx_1\,dy_1\,dx_2\,dy_2.
\]
Its integral in the complex orientation is positive exactly when $\xi\ne0$. This proves the equivalence. \hfill$\square$

There is also a certified point-selection version. If a unit-frame coefficient has size $b_0>0$ at a point, has Lipschitz constant $L_\xi$ nearby, and the Hermitian form has a known negative margin, choose a rational point within the margin-preserving radius and within $b_0/(2L_\xi)$ when $L_\xi>0$. Then $b\geq b_0/2$ and the lower bound is $m_0b_0^2/(4N_H)$. THETA3-Q3 gives an explicit denominator from the coordinate matrix, margin, and radius. The resulting level, mass, and dimension still require bounds or computation.

For two restrictions in the same $R$, Parseval gives the mixed identity
\[
 \int_{[H]}f\overline{f'}=
 \sum_j B_{e_j}(f)\overline{B_{e_j}(f')}.
\]
THETA3-A5 passes from finitely many restrictions to the full ball inner product by a positive inverse frame operator. A naive quotient disintegration does not do this: it unfolds to $H(F)\backslash G(\A)$, rather than $G(F)\backslash G(\A)$. The rational or algebraic nature of the frame correction is a further issue.

Finally, a central lift of the Shimura cocharacter at $L$ may factor through $\U(L^\perp)\times\U(L)$ rather than $H$ alone. The proof uses the finite definite adelic module and canonical-fibre evaluation. It does not assert an unspecified morphism of Shimura data $H\to G$.

\subsection{THETA3-B13: the recovered character-independent functional}\label{B13}

\paragraph{The exact family.}
The source \path{Definitions/LatticeFrames/ValLastX.lean}, under
\path{Development/AbelianVarietyTheory}, contains the relevant definitions. The following locators separate data from conclusions:
\begin{itemize}
\item \texttt{Comm}, lines 28--57, fixes the indefinite spaces $V_j$, discrete-series models, definite spaces $X_j$, Bessel/Jacobi decompositions, measures, final $T$, and $X_m=T\perp L_z$.
\item \texttt{Tup}, lines 79--113, supplies varying Weil representations, finite translation combinations, row functions, and constants. \texttt{RealEq}, lines 182--184, requires their real Weil actions to agree.
\item \texttt{perF}, lines 119--120, is a finite sum of genuine Bessel periods of finite right translates. \texttt{per}, lines 122--123, inserts the row function and automorphic realization. \texttt{perT}, lines 137--138, is the final definite restriction period.
\item \texttt{up}, lines 166--167, iterates two definite theta operators with the intervening $S$ operation; \texttt{dn}, lines 169--170, composes the reverse automorphic lift maps.
\end{itemize}
The normalizing products, lines 172--174, are
\[
 c_{\rm Top}(n)=
 \left(\prod_{i=0}^{n-1}c_V(2i)c_K(2i)c_K(2i+1)\right)c_P,
 \qquad
 c_{\rm Bot}=\prod_{j=0}^{m-1}c_I(j).
\]
\texttt{LawsX}, lines 212--221, includes common-data validity, tuple validity, finite support, common real action, the double-step and last-step integral laws, definite kernel reduction, final period factorization, and row covariance.

The actual source witness is
\nolinkurl{GlobalTheory.SeeSaw.mainX_of_lawsX}, in
\path{Proofs/DecoratedLatticeProofs/SeesawTransportRelationsOfFamilyLaws.lean}, lines 25--37. Its inputs are \texttt{CuspRapidDecay}, \texttt{HarishChandraModules}, and \texttt{LawsX}; its conclusion is \texttt{MainX}. The matching \texttt{Statement.lean}, line 12, exposes the wrapper
\nolinkurl{GlobalTheory.T.seesawTransportRelations_of_familyLaws}.
These declarations and the displayed proof body were read, but their complete transitive dependencies were not audited or compiled.

\paragraph{Mathematical content.}
Fix $m=2n+2$. For an allowed tuple, input column, and terminal automorphic vector, indexed collectively by $s$, put
\begin{align*}
 A_s&=c_{{\rm Top},s}(n+1)
      \operatorname{tail}_s(\operatorname{up}_s(n+1)\psi_s,F_s),\\
 B_s(y,v)&=c_{{\rm Bot},s}
      \operatorname{per}_s\bigl(\psi_s,
       \operatorname{dn}_s(2n+1)\operatorname{ioT}_s(F_s);y,v\bigr).
\end{align*}
\texttt{MainX}, lines 244--251, says that every finite scalar relation
$\sum_s\alpha_s A_s=0$ implies
$\sum_s\alpha_s B_s(y,v)=0$ for every fixed $y,v$.
Consequently there is a single scalar $a(y,v)$ such that
\[
 B_s(y,v)=a(y,v)A_s
 \quad\text{for every allowed tuple, column, and terminal vector.}
 \tag{B13}
\]
The scalar is unique if some $A_s\ne0$. This is the precise normalized independence supplied by the recovered theorem. It is not independence under arbitrary changes of characters and Schwartz data.

\textbf{Proof of the consequence.} If all $A_s$ vanish, one-term relations force all $B_s(y,v)=0$, and $a=0$ works. Otherwise choose $s_0$ with $A_{s_0}\ne0$ and put $a=B_{s_0}(y,v)/A_{s_0}$. For each $s$, the two-term relation
$A_{s_0}A_s-A_sA_{s_0}=0$ gives
$A_{s_0}B_s-A_sB_{s_0}=0$. Division proves (B13), and also uniqueness. This is a complete elementary re-proof of the logical passage from relation transfer to the common scalar. It does not re-prove the analytic inputs of \texttt{MainX}. \hfill$\square$

\paragraph{The exact instantiation and arithmetic gates.}
The abstract family is more flexible than the particular tower used in the construction. In \path{Definitions/LatticeFrames/JointSmooth.lean}, lines 56--62, \texttt{DefSp.Is} requires a compact quotient and a Cartan involution equal to $+1$ or $-1$ at each real place, hence definiteness, without requiring the same sign everywhere. Its \texttt{BesAt}, lines 107--115, consists of a line orthogonal decomposition and a hyperbolic decomposition. Its rank-zero local geometry therefore includes the active $\U(2,1)/\U(2)$ pair of B9.

By contrast, the concrete \texttt{PiLocalization$_2$.Is} in
\path{Definitions/LatticeFrames/ChainParSucc.lean}, lines 56--60, requires the same signature \texttt{sigYD D.R j} at every embedding. In \texttt{ExistsDatOfTower.lean}, lines 151--171, the Bessel step is built with a globally negative line, of signature $(0,1)$ everywhere. Hence that constructed tower does not immediately equal the one-active-place 032 datum. The theorem's abstract shape does not itself supply a witness satisfying all these hypotheses, nor the discrete-series span and gap requirements used to construct the tower.

The gap obstruction can be stated numerically. In \path{BlockTheory/TwistedDiscreteSeries/ThetaSourceAction.lean}, lines 209--211, \texttt{NodeGap} requires a natural number $N\geq\operatorname{rank}X'+1$ such that
\[
 t+2N^3+1\leq b_{i+1}-b_i
\]
for every adjacent pair and every infinite place. Consequently its right-hand gap is at least three. The canonical choice $b=(0,1)$ in B9 cannot satisfy this condition when identified with the two nodes. This requirement is used to construct the tower's spanning data, for example in \texttt{StepVanOfStepX.lean}, lines 71--84; it is not an additional assumption of the direct local integral evaluated in B9, nor of the elementary consequence (B13) once \texttt{LawsX} has already been supplied.

The bridge to the arithmetic columns is explicitly named. In
\path{Definitions/LatticeFrames/RatioOKOfRootOfUnity.lean}, the property
\texttt{FamAll}, lines 60--62, combines \texttt{LawsX} with
\texttt{LinkZero}, \texttt{LinkTail}, \texttt{LinkDown},
\texttt{LinkOnto}, \texttt{LinkComm}, and \texttt{RatioOK}.
The last condition compares $c_{\rm Top}/c_{\rm Bot}$ across tuples.
\texttt{UnitRatio}, in \path{Definitions/LatticeFrames/UnitRatio.lean}, lines 22--23, means
\[
 y\ne0,\quad \exists n>0,\ r\in\Q^*:
 (x/y)^n=r,
 \qquad v_p(r)\ne0\Longrightarrow p\in S.
\]
Thus the source controls radical-of-rational ratios with specified prime support; it does not automatically identify their absolute values with one. To derive a quantitative period bound, one must instantiate these links for the cohomological vectors, retain the normalizers, and bound the relevant common scalar or normalized Gram matrix. B9 and B12 give explicit local and definite-subgroup nonvanishing inputs, while THETA3-A5 supplies a separate exact finite-frame comparison of the global mixed pairing.

\subsection{THETA3-B14: large weights, tensor contractions, and the return to cohomology}\label{B14}

The large-weight device has an exact scope. For $q\geq1$, the canonical-power coefficient has compact type $(\det A)^qd^{-2q}$ and the parameters displayed after B9. Its local Bessel integral is one and its node gap is $3q-2$. Thus, for a fixed numerical threshold $M$, the condition $3q-2\geq M$ is attained by $q\geq\lceil(M+2)/3\rceil$. This does not alone construct the remaining tower data, and changes the input from a two-form to a pluricanonical section.

\paragraph{A cohomological distinction.}
For $q>1$, the coherent multiplication map
\[
 \operatorname{Sym}^q H^0(X,K_X)\longrightarrow H^0(X,K_X^q)
\]
is not the exterior-product map followed by diagonal pullback. If
$0\ne\xi\in H^{2,0}(X)$, its external product lies in
$H^{2q,0}(X^q)$, whereas diagonal pullback lies in
$H^{2q,0}(X)=0$, since $2q>\dim_\C X=2$.
The coherent section $\xi^q$ is nonzero: in a local canonical frame it is the $q$th power of its coefficient, tensor the $q$th power of the frame. These are different operations.

There is no nonzero global holomorphic bundle map $K_X^q\to K_X$ on a compact ball surface. Indeed, $K_X$ is ample, and such a map would be a section of $K_X^{1-q}$. A nonzero section would give an effective divisor $D$ with
\[
 D\cdot K_X=(1-q)K_X^2<0,
\]
contrary to ampleness; the nowhere-zero case would make $K_X^{1-q}$ trivial, also contradicting $K_X^2>0$. This rules out a global fibrewise holomorphic contraction of that specified kind. It does not rule out other correspondences between auxiliary varieties or maps requiring extra data.

Division by $\xi^{q-1}$ is an inverse on the particular image
$\xi^{q-1}H^0(X,K_X)\subset H^0(X,K_X^q)$, because cancellation gives a holomorphic result on that image. On an arbitrary section it can have poles along the zero divisor of $\xi$. Thus using this inverse requires proof that the high-weight input belongs to this image. A nonzero vector produced somewhere in the full high-weight space does not supply that proof.

\paragraph{The precise tensor contraction and its nonvanishing.}
Let $V=H^2(X,\Q)$ and
\[
 Q_X(u,v)=\int_Xu\smile v:\ V\otimes V\longrightarrow\Q(-2).
\]
Put $A=Q_X(\xi,\overline\xi)>0$. Every contraction of two factors in the pure tensor $\xi^{\otimes q}$ by $Q_X$ vanishes, because $Q_X(\xi,\xi)=0$ by Hodge type. By contrast the contraction
\[
 C_q:V^{\otimes q}\otimes V^{\otimes(q-1)}
       \longrightarrow V(-2(q-1)),
\]
\[
 C_q(u_1\otimes\cdots\otimes u_q\otimes v_2\otimes\cdots\otimes v_q)
       =u_1\prod_{i=2}^q Q_X(u_i,v_i)
\]
satisfies
\[
 C_q(\xi^{\otimes q}\otimes\overline\xi^{\otimes(q-1)})
           =A^{q-1}\xi\ne0.
 \tag{B14a}
\]
There is no new nonvanishing hypothesis in this calculation beyond $\xi\ne0$.

The map $C_q$ is induced by algebraic diagonals and push-forward. On $X^{2q-1}$, pair the $i$th and $(q+i-1)$st factors, for $2\leq i\leq q$, by the pullback of the diagonal class of $X\times X$, and push to the first factor. For an external product, the projection formula gives one factor $\int_Xu_i\smile v_i$ from each pair. All factors have even degree, so no permutation sign occurs. The resulting Tate twist is exactly $-2(q-1)$. This proves both the formula and its algebraic-correspondence realization.

This construction needs an actual external-product cohomological input, including the conjugate factors. A coherent section of $K_X^q$ alone is not such an input. Nor does the fact that $C_q$ is algebraic turn a non-Hodge vector such as $\xi^{\otimes q}\otimes\overline\xi^{\otimes(q-1)}$ into an algebraic cycle. In a proposed cycle construction, the ambient Hodge types and the required projected component must first be identified; after that, algebraic contraction preserves algebraicity and (B14a) supplies the appropriate nonzero test value.

A parallel representation-theoretic check applies to a specified rational Weil component. Suppose a polarizing algebraic torus acts on
$W_\C=\ell_\chi\oplus\ell_{\chi^{-1}}$, where $\chi$ has infinite order, and its invariant pairing pairs the two lines. Any torus-equivariant map from $W^{\otimes q}$ to $W$ sends the pure weight $\chi^q$ to zero for $q>1$, since neither $\chi$ nor $\chi^{-1}$ equals $\chi^q$. Pairwise polarization contractions are equivariant and therefore have this obstruction. With $q-1$ opposite-character factors the total weight is $\chi$, and the contraction is nonzero whenever their pairings are nonzero. This is a statement about that explicitly specified torus and equivariant contractions; it is not an impossibility theorem for every algebraic correspondence. It pinpoints which conjugate component a tensor-power argument must retain.

\paragraph{Analytic return with an explicit norm loss.}
Fix a Hermitian metric on $K_X$ and a positive measure on $X$. Write
\[
 A_q=\int_X|\xi|^{2q},\qquad M=\|\xi\|_\infty>0.
\]
Pointwise $|\xi|^{2q}\leq M^{2q-2}|\xi|^2$, hence
\[
 A_1\geq \frac{A_q}{M^{2q-2}}.
 \tag{B14b}
\]
Therefore a certified lower bound for the global norm of the particular pure power $\xi^q$, together with a certified upper bound for $M$, gives a lower bound for the original norm. A local Bessel integral normalized to one is not by itself that global bound.

There is also a nonzero analytic contraction. Use the metric to contract
$\xi^q\otimes\overline{\xi}^{\otimes(q-1)}$ to the smooth canonical section
$|\xi|^{2q-2}\xi$, and take its orthogonal projection to $H^0(X,K_X)$. Its inner product with $\xi$ is $A_q>0$, so the projection is nonzero. This projection depends on the metric and the original form and is not thereby an algebraic correspondence. Equations (B14a) and (B14b) give two valid return mechanisms with their required input data; neither identifies an arbitrary high-weight output with the original Weil component automatically.

\subsection{THETA3-B15: a small-rank local alternative and its exact global gates}\label{B15}

There is an alternative to the large-weight construction: the canonical discrete series itself has a nonzero local theta correspondence with a compact rank-two group. We give the characters and parameters explicitly. This removes a local occurrence obstruction; it does not identify a particular global theta vector with the iterated lift of B13.

\paragraph{The literal one-active-place datum forces the rank-zero step.}
Suppose that $V_0$ in \texttt{Comm} is the rank-three space $V$ of B0 and that $m=2n+2$. Then $n=0$, hence $m=2$. Indeed, \texttt{Comm.bes}, at \texttt{ValLastX.lean}, line 46, requires a \texttt{BesAt} datum of rank $n$. By \texttt{JointSmooth.lean}, lines 107--115, this contains a global hyperbolic decomposition of $V$ with $n$ hyperbolic planes. A nonzero isotropic vector over $E$ remains a nonzero isotropic vector at every completion. Since $V$ is positive definite at another real place, it has no such vector. Its global Witt index is zero, and therefore $n=0$.

At $m=2$ all the intermediate conditions indexed by $j+2(\sigma+1)+2=m$ are empty. The rank-zero \texttt{JacAt} identifies $V_1$ with the definite space $X_1$; write this middle space as $S$. Thus the remaining geometric pattern is
\[
 V=A\perp L,\qquad S=X_1\text{ totally definite},\qquad
 X_2=T\perp L_z\text{ totally definite}.
 \tag{B15a}
\]
Here $A$ is the definite rank-two complement in B12. The abstract \texttt{Comm} does not impose the concrete ranks \texttt{rankY} on its spaces $V_j$. Its field \texttt{par}, line 33, may vary with the real place. The numerical \texttt{Ranks} conditions do allow $m=2$: the choices $k=6,t=13,\operatorname{ord}=1$ satisfy the inequalities in \path{Formalism/LocalRingCardinality/Params.lean}, lines 6--22. This numerical observation does not construct the other fields of \texttt{D7} or a full \texttt{Comm.Is} witness.

The last seesaw imposes a precise line relation, not an arbitrary real signature change. In \texttt{PlaceLaws.lean}, lines 165--166, \texttt{WeilOK.tau} is
\[
 \operatorname{rank}L_z=1,\qquad
 \tau^*G_L\tau=-N_{E/F}(\xi_0)G_{L_z}.
 \tag{B15b}
\]
The norm is positive at every real place. Hence $L_z$ has the opposite sign to $L$: positive at $\mu$, negative elsewhere. This has a coherent global realization: set $L_z=-N_{E/F}(\xi_0)^{-1}L$ and $\tau=1$. Taking $T=L_z$ gives a rank-one top and a definite rank-two $X_2=T\perp L_z$, with positive sign at $\mu$ and negative sign elsewhere. These are actual global forms; no unverified collection of finite determinant classes is being prescribed. Choosing them does not yet prove the needed finite local theta occurrences.

\paragraph{Convention and local theorem.}
Use Ichino's convention in \emph{Theta lifting for tempered representations of real unitary groups}, arXiv:2008.06174v2, \S4.1, p.10. The input space $W$ is skew-Hermitian of dimension $n$ and signature $(p,q)$, the target space is Hermitian of dimension $r+s=m$, and
\begin{equation}
\begin{gathered}
 \chi_V(z)=(z/|z|)^{m_0},\qquad
 \chi_W(z)=(z/|z|)^{n_0},\\
 m_0\equiv m\pmod2,\qquad n_0\equiv n\pmod2,
 \qquad \psi(x)=e^{-2\pi i x}.
\end{gathered}
 \tag{B15c}
\end{equation}
The local lift is characterized by
$\Hom_{\U(W)\times\U(V)}(\omega,\pi\boxtimes\theta(\pi))\ne0$.
Let $\pi_\mu$ have Harish--Chandra parameter $(1,0;-1)$ on $\U(2,1)$, and let $\pi_\tau=1$ on $\U(3)$ for $\tau\ne\mu$. For the pair of dimensions $(n,m)=(3,2)$, take $m_0=-2,n_0=-1$. Then
\[
 \theta_{2,0}(\pi_\mu)=\det,\qquad
 \theta_{2,0}(\pi_\tau)=\det\quad(\tau\ne\mu).
 \tag{B15d}
\]
Here the character $\det$ of compact $\U(2)$ has Harish--Chandra parameter $(3/2,1/2)$.
The representation $\pi_\mu$ is the holomorphic discrete series with lowest compact type $(1,1;-2)$, namely $(A,d)\mapsto\det(A)d^{-2}$. This is the canonical right-regular coefficient convention fixed in B9: $\lambda+\rho_n-\rho_c=(1,1;-2)$. Replacing the right-action convention by its dual would require dualizing this parameter and the subsequent pairings as well.

For the next pair, let $\sigma=\det^{-1}$ on $\U(2)$, with Harish--Chandra parameter $(-1/2,-3/2)$. The following local lifts are nonzero and trivial:
\[
\begin{array}{c|c|c|c|c}
\text{place}&(n,m)&(m_0,n_0)&\text{target signature}&\theta(\sigma)\\ \hline
\mu&(2,1)&(-3,-2)&(1,0)&1\\
\tau\ne\mu&(2,1)&(-1,2)&(0,1)&1.
\end{array}
 \tag{B15e}
\]
Thus there are compatible middle representations for a two-step composition at every real place: the first oscillator has a quotient $\pi_\tau\boxtimes\det$, and the second has a quotient $\det^{-1}\boxtimes1$.

\textbf{Proof.}
In all three calculations the rank drops by one. After subtracting $m_0/2$ from the input Harish--Chandra parameter, the nonvanishing criterion has a single zero chain, of length $k=1$. For (B15d) the shifted parameters are
\[
 (2,1;0)\quad\text{and}\quad(2,1,0),
\]
respectively. In Atobe's notation both have $(k_\lambda,r_\lambda,s_\lambda)=(1,2,0)$. Atobe, \emph{On the non-vanishing of theta liftings of tempered representations of $\U(p,q)$}, Adv. Math. 363 (2020), Definition 1.6 and Theorem 1.7(2), pp.4--5, gives nonvanishing at the base point $(2,0)$. The theorem for discrete series is unconditional, as stated immediately after it. Ichino, \S6.2, explicitly identifies Atobe's additive-character convention with (B15c).

The parameter rule for adjacent ranks deletes the zero, transfers positive entries in the first block and negative entries in the second block to the positive target block, and then adds $n_0/2$. This is the adjacent-rank case of Ichino's Theorem 4.1(ii), pp.11--12, already covered by Paul's adjacent-rank correspondence cited there. In (B15d) it gives $(2,1)-(1/2,1/2)=(3/2,1/2)$. For the first row of (B15e), the shifted parameter is $(1,0)$, giving base point $(1,0)$ and output $1-1=0$. For the second row it is $(0,-1)$, giving base point $(0,1)$ and output $-1+1=0$. The same nonvanishing theorem applies in both cases. Every parity condition in (B15c) holds. Compact rank-one Harish--Chandra parameter zero is the trivial representation. This proves the claims. \hfill$\square$

The abstract groups in this calculation must be identified with the Hermitian/skew-Hermitian spaces of 032 and the archive, including the sign of the imaginary scalar used to pass between the forms. Equations (B15c)--(B15e) specify the convention rather than silently making that identification. Likewise, the archive's \texttt{liftS} integrates the kernel against the input function bilinearly. A quotient $\pi\boxtimes\det$ therefore pairs with an automorphic input of type $\det^{-1}$ to produce $\pi$. This is why the middle input in (B15e) is the dual of the B9 compact test representation. The two pairs are allowed different splitting characters; \texttt{WeilOK.last}, lines 160--164, separately names the second pair's character $\chi'_S$.

\paragraph{The archimedean characters can be made global.}
Let $E/F$ be CM, let $\epsilon$ be its quadratic idele-class character, and fix integers $a_\tau$ all congruent to $d$ modulo two. There is a unitary Hecke character $\chi$ of $E$ satisfying
\[
 \chi|_{\A_F^\times}=\epsilon^d,
 \qquad \chi_\tau(z)=(z/|z|)^{a_\tau}.
 \tag{B15f}
\]
Finite ramification is allowed; prescribed finite local values are not part of this assertion.

\textbf{Proof.}
The natural map $C_F\to C_E$ is a closed embedding. Indeed it is injective, and the idele norm of $C_E$ restricts to the square of the norm on $C_F$; the norm-one idele-class group of $F$ is compact, so the embedding is proper. The character-extension theorem for closed subgroups of locally compact abelian groups supplies a unitary extension $\chi_0$ of $\epsilon^d$. This is precisely the argument in 032, Lemma 4.4, using Folland, \emph{A Course in Abstract Harmonic Analysis}, second edition, \S4.3, Corollary 4.42. Since $\chi_0$ is trivial on each positive real scalar, its archimedean characters are $(z/|z|)^{a^0_\tau}$, with $a^0_\tau\equiv d\pmod2$.

Put $T_1=\operatorname{Res}^{1}_{E/F}\mathbf G_m$. Its real group $T_1(F\otimes\R)$ embeds as a compact subgroup of $T_1(F)\backslash T_1(\A_F)$ by placing one at the finite components: the intersection with the diagonal rational group is trivial. Character extension on this locally compact abelian quotient gives an automorphic character $\rho$ with archimedean weights $(a_\tau-a^0_\tau)/2$. Then
\[
 \chi(z)=\chi_0(z)\rho(z/\overline z)
\]
is trivial on $E^\times$, has the same restriction to $\A_F^\times$ as $\chi_0$, and has the prescribed weights. Its finite components have open kernels on compact unit groups, so it has finite conductor. \hfill$\square$

This proves that the integer choices in (B15d)--(B15e) themselves have no global parity obstruction. It does not force the finite components of those characters to be those in a previously fixed theta construction.

\paragraph{What changing a splitting transports.}
For a fixed dual pair, changing the splitting multiplies the oscillator action by determinant characters of the same two groups. Locally, Atobe's Proposition 3.3 and Proposition 3.5, pp.20--22, give
\begin{equation}
\begin{aligned}
 \omega_{m_0+2a}|_{\U(p,q)}
       &\simeq\omega_{m_0}|_{\U(p,q)}\otimes\det^a,\\
 \theta^{m_0+2a}(\pi)
       &\simeq\theta^{m_0}(\pi\otimes\det^{-a}),
\end{aligned}
 \tag{B15g}
\end{equation}
when the other splitting character is kept fixed. The latter formula includes zero lifts. A simultaneous change of the other character tensors the output by its corresponding determinant character.

The archive exposes its orientation in
\path{Formalism/LocalToGlobal/IndefiniteReal/LocalNewformFunction.lean}, line 51:
\[
 \texttt{splitFacO}(\chi_S,\chi_T;x,y)=\chi_T(x)^{-1}\chi_S(y).
 \tag{B15h}
\]
If $\kappa_S(z)=\rho_S(z/\bar z)$ and $\kappa_T(z)=\rho_T(z/\bar z)$ are character ratios, its split-place scalar changes by
$\rho_T(\det g)^{-1}\rho_S(\det h)$.
The determinants here follow from \texttt{TiedWeil.lean}, lines 71--78, and the actual \nolinkurl{IsWeilDatumO.om_split} condition at lines 72--75 of \nolinkurl{LocalNewformFunction.lean}. Thus an inverse in one convention cannot simply be removed when comparing it with (B15g).

All these determinant factors restrict trivially to the appropriate special unitary group. The canonical $\SU(V)$ geometry can consequently survive a determinant twist of its $\U(V)$ realization. However, the underlying source space is unchanged: a mixed-sign rank-two $P$ in A7 remains that same $P$, and a local group $\U(1,1)$ does not become $\U(2)$ by a character twist. Compensating the target twist as in (B15g) preserves the original local correspondence. To replace the source by a different definite inner form while keeping a fixed target vector, one must establish occurrence in the new dual pair. The original A7 realization does not alone provide that occurrence at every finite place, or an automorphic representation on the new source.

\paragraph{A precise global nonvanishing test for the alternative.}
There is a usable criterion with no unspecified real-place correction. In the conventions of Gan--Qiu--Takeda, arXiv:1207.4709v3, Theorem 11.7, pp.56--57, let $U_n$ be the source of dimension $n=1$ or $2$, and let $V$ be an anisotropic target of dimension $m=n+1$. For an irreducible cuspidal automorphic source representation $\Sigma$, its lift to this particular $V$ is nonzero if and only if
\begin{equation}
\begin{gathered}
 \theta_{V_v}(\Sigma_v)\ne0\quad\text{for every place }v,
 \\ L(1,\Sigma\times\chi_V)\ne0
       \quad\text{with a holomorphic value}.
\end{gathered}
 \tag{B15i}
\end{equation}
Here $L$ is the source's completed standard doubling $L$-function, with the character as named in that paper. The hypotheses are $d(n)=n<m=n+1\leq2n$, target Witt index $r=0$, and hence no lower step in its global Witt tower. The equality $m=d(n)+1$ is exactly the final case in Theorem 11.7(ii); Proposition 11.6 says that it removes the otherwise possible real-sign replacement. No temperedness hypothesis is needed for this criterion.

With Tamagawa measures and the paper's normalized doubling factors, Theorem 11.4(ii), p.54, specializes to
\begin{equation}
\begin{gathered}
 \langle\theta(\phi_1,f_1),\theta(\phi_2,f_2)\rangle
 \\
 =2L(1,\Sigma\times\chi_V)
       Z^*(1/2,\phi_1\otimes\overline{\phi_2},f_1,f_2),
 \\
 Z_v^*(s)=\frac{Z_v(s)}{L(s+1/2,\Sigma_v\times\chi_{V,v})}.
\end{gathered}
 \tag{B15j}
\end{equation}
The global product includes the unramified denominator factors of that normalization; they must not be replaced by one. The archive's bilinear input and Ichino's character names must be converted before substituting a representation in this formula. For comparison, Ichino's \S8.2, p.37, uses an incomplete $L$-function twisted by $\chi_V^{-1}$ and, for dimensions $2\to3$ and tempered source, divides it by $L^S(2,1)L^S(3,\epsilon)$ while retaining the local matrix-coefficient integrals at the excluded places. These are different displayed normalizations of the standard method, not an asserted equality of the named characters.

Thus both arrows in a proposed rank-one to rank-two to rank-three global construction have a specific value-at-one criterion, together with finite local occurrence. This is not a central-value criterion. It still must be shown that the resulting rank-three representation contains the prescribed canonical theta vector, rather than merely some canonical vector. For a fixed rank-three representation, the reverse first-occurrence question is the pole side of the Rallis theory; an upward criterion for an independently chosen $\Sigma$ does not establish that reverse occurrence automatically.

The definite side of the small choice in (B15a) has a different rank pattern: $A_2\to S_2\to X_2$, with $\operatorname{rank}X_2=2$. Its two arrows have equal rank. The adjacent-rank criterion (B15i) does not cover them; their standard Rallis parameter would be the central point, and their nonvanishing requires its own argument. In particular, the strict range of the archive's \nolinkurl{SphericalVolumeLift.Convergent} just recorded in B8 is not satisfied. No claim that every definite-side lift is nonzero is part of the local calculation.

\paragraph{The universal identity and the arithmetic equation left to prove.}
The last-step analytic identity is already a universal source theorem, rather than a large-gap instance. In \path{Proofs/DecoratedLatticeProofs/DoubleLiftPeriodRankZeroFubiniCompactQuotient.lean}, line 136, the source proves \texttt{lastLaw : \ensuremath{\forall} D, LastLaw D}; in \nolinkurl{LastSeeSawThetaIdentity.lean}, line 22, it proves \texttt{seeLast : \ensuremath{\forall} D, SeeLast D}. The proof bodies use compact quotient integration and the exact tensor identity with (B15b). With the requisite Weil data and
\[
 \mu_{V_1}=(\operatorname{emb}_0)_*\mu_{X_1},
 \tag{B15k}
\]
which is \texttt{Comm.VolLast} in \texttt{HzUnit.lean}, lines 23--24, the theorem \nolinkurl{GlobalTheory.LastRed.main}, in \nolinkurl{LastTwoStepsDoubleLiftPeriodEqIntegral.lean}, line 293, gives
\begin{equation}
\begin{gathered}
 B_A\bigl(\varphi,
   \Theta_{V\leftarrow S}(\Theta_{S\leftarrow T}(F;\Phi_2);\Phi_1)\bigr)
 \\
  =\int_{[T]} F(t)\,
       \operatorname{biLift}_{A\to S\to X_2}
       (\varphi;C_{\Phi_1,\Phi_2})(\iota(t,1))\,dt.
\end{gathered}
 \tag{B15l}
\end{equation}
This is the untranslated specialization; the source additionally gives every finite/right translate explicitly. The contracted kernel is its actual \texttt{lastBi}, defined in \texttt{CombAtSingleZero.lean}, lines 29--34: on the decomposition $X_2=T\perp L_z$ it is
\[
 C_{\Phi_1,\Phi_2}(u;v_T,v_z)
       =\Phi_2(v_T)\Phi_1(u,\operatorname{ofAdj}_{L,S}(\tau v_z)).
\]
All characters and measures remain in this identity. There is no arbitrary comparison scalar at this stage. The required global input is the concrete double theta lift appearing on its left.

The additional equation needed to identify these analytic kernels with the arithmetic columns is explicit. In \path{Definitions/LatticeFrames/OnFramesRdApply.lean}, lines 108--113, the condition \texttt{ClsK.Core} is
\begin{equation}
\begin{gathered}
 \operatorname{volLev}(X_j,U_{j,\alpha})\,
 \theta_{K,j}(h_{j,\alpha,a},h_{j+1,\gamma,a'};\Phi_{j,P})
 \\
 =c_{K,j}\,\mathcal A_{j,\alpha,\gamma,a,a'}
       \bigl(\operatorname{mvZ}(1,s_{j,\tau}I)(\operatorname{castP}P)\bigr).
\end{gathered}
 \tag{B15m}
\end{equation}
Here $\mathcal A$ is the specified \texttt{abelian\ensuremath{_2}} functional, not an arbitrary linear functional defined to make the equality true. Its class sums carry the reciprocal stabilizer orders in lines 115--121. The terminal condition \texttt{ClsT.Is.tail}, lines 147--150, similarly includes the top-level covolume and stabilizer masses. Their conjunction with the real branching and level conditions is \texttt{ClsPinD}, lines 152--153, from which the source derives the corresponding \texttt{KerX} and \texttt{PerX} laws. The universal Fubini identity (B15l) does not prove (B15m), supply those arithmetic functionals, or establish that the given 032 vector lies in the indicated double-lift image. These are now concrete representation and kernel/covolume equations, rather than an unspecified appeal to global matching.


% THETA3-B16 contribution by /root/definite_bridge, 21 November 2026.
\subsection{THETA3-B16: a pole criterion that also supplies finite coherence}\label{B16}

The value-at-one criteria in B15 start with an automorphic representation on a chosen smaller group. There is also a useful converse route: a pole on the rank-three group can construct a coherent smaller group. Its character conventions must be retained.

\paragraph{A precise conditional descent statement.}
Use exactly the notation of Jiang--Wu, \emph{On $(\chi,b)$-factors of cuspidal automorphic representations of unitary groups I}, arXiv:1409.0767v3, \S3.1 and Theorem 3.3(1), pp.15--17. Let $X$ be a skew-Hermitian space of rank three, let $\chi$ be a unitary Hecke character of $E$ with $\chi|_{\A_F^\times}=1$, and let $\sigma$ be irreducible and cuspidal on $\U(X)(\A_F)$. Write
\[
 \rho=\sigma\otimes\chi^{-1}
 \tag{B16a}
\]
in the source's notation for its twist on the unitary group. Suppose that:
\begin{enumerate}
\item $L^S(s,\sigma\times\chi)$, with the source's convention, has a pole at $s=1$;
\item $\rho$ has a real component which is not one dimensional;
\item at each real place, the only nonzero rank-two local theta lift of the corresponding input, with this splitting and additive character, is on a definite signature.
\end{enumerate}
Then there is a global Hermitian space $S$ of rank two, definite at every real place, for which the global theta lift of $\rho$ to $\U(S)$ is nonzero and cuspidal. In particular its finite local isometry classes are coherent, and its local theta lifts are nonzero at all finite places. No finite local form or lattice has been prescribed in this statement.

\textbf{Proof.}
In Jiang--Wu's Theorem 3.3(1), take their $m=3$, $\epsilon_\chi=0$, and $j=1$. Its distinguished point is
\[
 \tfrac12(m+1-\epsilon_\chi)-j=1,
\]
and it gives $\operatorname{LO}_{\psi,\chi}(\rho)\leq2$. The lowest occurrence minimizes the first-occurrence dimension over global Hermitian Witt towers of the prescribed even parity; see \S3.2, p.16. Therefore the dimension realizing the minimum is zero or two. At dimension zero the theta series is the determinant character $\nu_\chi\circ\det$, as the source states immediately before \S3.2. A nonzero global pairing with it would force every local component of the irreducible input to be one dimensional. Hypothesis 2 excludes this. Thus the lowest occurrence has dimension two, and by its definition its global theta lift is nonzero and cuspidal.

Nonzero global theta occurrence implies the corresponding local occurrences. Hypothesis 3 forces the real signatures of the constructed space to be definite. The theorem constructed a global space, so its finite local forms already satisfy the global coherence relation. This proves the assertion. \hfill$\square$

This result replaces separate finite-place guesses by one global pole test, once its representation and character dictionary has been established. It does not assert that the pole is known for the 032 wedge. In particular, the twist in (B16a) must be resolved against the bilinear integral in Jiang--Wu's \S3.1 and the contragredient in the local Howe quotient before substituting the canonical representation in (B15d). The local computations of B15 give the required unique definite occurrence after that convention matching; they do not by themselves prove hypothesis 1.

\paragraph{The invariant under a compensating splitting change.}
This remaining test is not removed by the independence of the special-unitary action from splitting characters. Let $\tau$ be a global character of $\operatorname{Res}^{1}_{E/F}\mathbf G_m$ and put
\[
 \tau_E(z)=\tau(z/\overline z),\qquad
 \Pi'=\Pi\otimes(\tau\circ\det),\qquad
 \chi'=\chi\tau_E.
\]
In the standard local-parameter convention of (B15c), changing the splitting by $\tau_E$ multiplies the input action by $\tau\circ\det$. Hence the compensating simultaneous twist transports the same local Howe correspondence on the same underlying Hermitian spaces. This is the global version of (B15g). Jiang--Wu, equation (3.1), p.15, gives the same determinant-character rule, with $\nu_\kappa(z/\bar z)=\kappa(z)$.

There is an equally explicit invariant of unramified standard parameters. Whenever standard base change is used to write them, one has
\[
 \operatorname{BC}(\Pi')\otimes(\chi')^{-1}
   =\operatorname{BC}(\Pi)\otimes\chi^{-1}.
 \tag{B16b}
\]
Indeed, twisting a representation by $\tau\circ\det$ tensors its standard base-change parameter by $\tau_E$; the two additional factors then cancel. This also proves equality of the corresponding incomplete standard $L$-functions term by term at every unramified place. Formula (B16b) is stated in this parameter convention; it is not an unexplained replacement of Jiang--Wu's named $L$-function by a differently twisted one.

To change the relative splitting further by a character $\eta$ changes the right side of (B16b) by $\eta^{-1}$. That is a different standard parameter and a different pole question. The original nonzero A7 lift supplies information for its original source $P$ and characters. The compensating transport preserves $P$, including any real signature $(1,1)$; an additional relative change has not inherited that nonvanishing. This identifies the missing representation-theoretic calculation without asserting that another definite realization is impossible.

Finally, a successful application to each of several rank-three representations separately need not produce the same finite local source forms. The \texttt{Comm} in B13 is common to the whole family. Its use still requires common spaces and real action across tuples, the specified finite lattice data, and the kernel/covolume equation (B15m). The pole construction therefore closes a single-representation coherence problem under its hypotheses; it does not supply the entire arithmetic family by itself.


% THETA3-B17 contribution by /root/definite_bridge, 21 November 2026.
\subsection{THETA3-B17: the actual wedge admits an adjoint double lift and last-step seesaw}\label{B17}

The compact-quotient condition in the analytic last-step theorem is
strictly weaker than the totally definite condition in the arithmetic
\texttt{MainX} family. This distinction permits a complete analytic
specialization for the actual multiplicity plane of A7. The second
lift can be its adjoint, with terminal space equal to the original
rank-three space.

\paragraph{Fixed data and the finite-rank operator.}
Use the rank-three case of A7, with its exact scalar automorphic
realization of the canonical coefficient. Thus $V$ has signature
$(2,1)$ at $\mu$ and $(3,0)$ elsewhere, $[F:\Q]\geq2$, and the
multiplicity plane $P=P_1\perp P_2$ is definite at $\mu$.
Write
\[
 G=\U(V),\qquad S=\U(P).
\]
Both adelic quotients $[G]$ and $[S]$ are compact. For $V$, use
the definite component at a place other than $\mu$; for $P$, use
its definite component at $\mu$. A rational isotropic vector would
remain a nonzero isotropic vector in that definite completion.
The quotient compactness is the anisotropic case of the usual
arithmetic compactness theorem already used in A7.

For the nonvanishing assertion about a particular wedge, first
choose its rational negative line $L$ as in B12 and rebase its
canonical coefficient frame at $L$. Make this choice before
fixing the kernel and the operators below. The original
coefficient at an unrelated base point need not be nonzero,
and its maximal compact subgroup need not contain $\U(L^\perp)$.
After this rebasing the following data remain common to all
source characters under consideration.

Fix the additive character, the two splitting characters
$(\chi_V,\chi_P)$, the theta input $\Phi$, finite levels, the
canonical coefficient frame, and positive Haar quotient measures
$dg,dh$. Let
\[
 K(g,h)=\theta_{V,P}(g,h;\Phi).
\]
The coefficient identification here is the one in A7: after
evaluation on the fixed canonical frame, this is an actual scalar
theta kernel, not a separately chosen function having the same
cohomology class. More explicitly, let $\Psi$ be the holomorphic
top-degree relative Lie algebra cochain and choose a unit vector
$\kappa$ in its one-dimensional canonical tangent-frame line.
Its scalar Schwartz coefficient is $\Phi=\Psi(\kappa)$; evaluation
of the theta cochain on the frame translated by $g$ is
$\theta(\omega(g,h)\Phi)$. This fixes the archimedean coefficient
pairing before any adjoint is taken. A fixed conversion between
the Hermitian/skew-Hermitian presentation and the archive's
Gram-matrix presentation is part of the oscillator identification;
the subsequent dagger is applied to this very coefficient.
By A7, $K(\,\cdot\,,h)$ belongs to one
finite-dimensional space $\mathcal E$ of holomorphic canonical
coefficients for every $h$. Its smooth automorphy and the
compactness of the two quotients imply that $K$ is bounded on
$[G]\times[S]$.

Equip $\mathcal E$ with the Petersson inner product inherited from
$L^2([G],dg)$, conjugate linear in the first variable. Define
\begin{equation}
\begin{aligned}
 \mathcal T:L^2([S],dh)&\longrightarrow\mathcal E,\\
 (\mathcal T u)(g)&=\int_{[S]}K(g,h)u(h)\,dh,\\
 \mathcal E_0&=\operatorname{im}\mathcal T.
\end{aligned}
\tag{B17a}
\end{equation}
This is the archive's bilinear \texttt{liftS} convention. Boundedness
follows from Cauchy--Schwarz on the product of the two finite-measure
quotients. Its range is finite dimensional and therefore closed.
Moreover $K(\,\cdot\,,h)\in\mathcal E_0$ for every $h$:
if $f\in\mathcal E$ is orthogonal to the range, then its
continuous coefficient $\langle f,K(\,\cdot\,,h)\rangle_G$
vanishes almost everywhere by testing against all $u$, hence
everywhere. Taking the orthogonal complement proves the claim.
All the A7 character integrals with these fixed data lie in
$\mathcal E_0$: A7 writes the integral for a character $\pi$ as
$f_\pi=\mathcal T(\overline{v_\pi})$.

\paragraph{The adjoint is another genuine theta lift.}
The Hilbert-space adjoint is
\[
 (\mathcal T^*f)(h)
   =\int_{[G]}\overline{K(g,h)}f(g)\,dg.
 \tag{B17b}
\]
Indeed, expanding the inner products and using Fubini gives
$\langle\mathcal T u,f\rangle_G
 =\langle u,\mathcal T^*f\rangle_S$.
In particular $\mathcal T^*f$ is a smooth automorphic function
when $f\in\mathcal E$.

This adjoint kernel is a theta kernel for the exchanged pair with
the same underlying spaces. Since the canonical compact type is
one dimensional, conjugating the chosen scalar coefficient is
exactly the adjoint operation on that coefficient line; there is
no unrecorded sum over compact-type coordinates. Any volume of
the compact fibre is already included in $dg$ and in the
inherited Petersson norm. In the recovered source let
\[
 d_{V,P}(x)=G_V^{-1}x^*G_P,\qquad
 \Phi^\dagger(w)=\overline{\Phi(d_{P,V}(w))},
 \qquad \omega^\dagger=\operatorname{conjAdjoint}(\omega).
\]
Here $G_V,G_P$ are the Gram matrices and $x^*$ is conjugate
transpose. The exact source rule is
\begin{equation}
\begin{gathered}
 \omega\text{ has data }(V,P;\chi_V,\chi_P)
 \quad\Longrightarrow\quad
 \omega^\dagger\text{ has data }(P,V;\chi_P,\chi_V),\\
 \theta_{P,V}^{\dagger}(h,g;\Phi^\dagger)
       =\overline{\theta_{V,P}(g,h;\Phi)}.
\end{gathered}
\tag{B17c}
\end{equation}
There is no inverse on the displayed exchanged splitting
characters in this archive convention.

For precision,
\path{Development/AbelianVarietyTheory/BlockTheory/TwistedDiscreteSeries/ConjugateAdjointDagger.lean}, lines
112--124, defines the rational involution and the conjugated
exchanged representation. Lines 144--156 give its
antisymplectic, theta-function, and Weil-data clauses.
The following proof files have the full directory prefix
\path{Development/AbelianVarietyTheory/Proofs/DecoratedLatticeProofs/}.
The actual proof \nolinkurl{ConjWeilDatumViaAdjointExchangedPair.lean},
theorem \nolinkurl{isWeilDatumO_conjAdjoint}, proves the exchanged
Weil data, including the split-place normalization.
\nolinkurl{HermitianAdjointSymplecticActionThetaHeisenberg.lean}
proves the rational, antisymplectic, and theta-function
transport. The complete rule is proved by
\nolinkurl{WeilConjugationAdjointRuleStandardChar.lean},
theorem \texttt{conjAdjoint\_holds}, and exposed by
\nolinkurl{GlobalTheory.T.weilConjugationAdjointRule_standard_char}
in its statement file.

One can also check the kernel equality directly. The involution
$d$ carries the rational lattice bijectively onto the exchanged
rational lattice, reverses the symplectic form, and exchanges
the two group actions. Complex conjugation reverses the additive
character at the same time, so the standard additive character
is retained on the exchanged pair. The defining formula
$\omega^\dagger(h,g)\Phi^\dagger
 =\overline{\omega(g,h)\Phi}\circ d_{P,V}$
then gives (B17c) at zero, equivalently after summing the
rational Schwartz model. Thus (B17b) is the bilinear theta lift
from $V$ to $P$ with input $f$ and kernel $\Phi^\dagger$.

\paragraph{An exact double lift, common to all source characters.}
On $\mathcal E_0$ put
\[
 Q=\mathcal T\mathcal T^*.
\]
This is positive definite. In fact
$\langle f,Qf\rangle_G=\|\mathcal T^*f\|_S^2$;
if it vanishes, then $f$ is orthogonal to
$\operatorname{im}\mathcal T=\mathcal E_0$, and hence $f=0$.
Consequently, for every $f\in\mathcal E_0$,
\begin{equation}
\begin{gathered}
 F_f=Q^{-1}f,\qquad
 f=\mathcal T\mathcal T^*F_f\\
   =\Theta_{V\leftarrow P}
          \bigl(\Theta_{P\leftarrow V}(F_f;\Phi^\dagger);
                \Phi\bigr).
\end{gathered}
\tag{B17d}
\end{equation}
All functions in this identity are continuous automorphic
functions on compact adelic quotients. It requires neither a new
global nonvanishing theorem nor a rank-one reverse descent:
the second source is $T=V$, of rank three.

The same $Q$ works for every A7 source character with the fixed
spaces, kernel, splittings, levels, coefficient frame, and
measures. Thus (B17d) is a common linear operator identity on the
whole space $\mathcal E_0$, rather than a normalization chosen
separately for each vector or character. Changing the complete
types may change $P$; changing the Schwartz input may change
$Q$. Uniformity across those changes is not asserted.

The exact mixed-pairing identities are
\begin{equation}
\begin{aligned}
 \langle f_1,f_2\rangle_G
    &=\langle\mathcal T^*Q^{-1}f_1,
                        \mathcal T^*f_2\rangle_S\\
    &=\langle\mathcal T^*Q^{-1/2}f_1,
                        \mathcal T^*Q^{-1/2}f_2\rangle_S .
\end{aligned}
\tag{B17e}
\end{equation}
They follow by adjunction and self-adjointness of $Q$.
The positive square root is defined by the finite-dimensional
spectral theorem. In particular, the second identity is an
isometry from $\mathcal E_0$ to a space of actual reverse theta
lifts on the original anisotropic $P$. The conversion from this
Petersson pairing to the geometric canonical cup-product
pairing uses the fixed coefficient frame and quotient measure;
it is the same conversion for every $f_\pi$.

For completeness, when $\mathcal E_0\ne0$, fix an orthonormal basis $f_1,\ldots,f_r$ of
$\mathcal E_0$ and write $K(g,h)=\sum_i f_i(g)e_i(h)$. The matrix
of $Q$ is the positive Gram matrix
\[
 Q_{ij}=\int_{[S]}e_i(h)\overline{e_j(h)}\,dh .
\]
With $t=\operatorname{Tr}Q$ and $\Delta=\det Q>0$, for $r\geq2$
its least eigenvalue satisfies
\[
 \lambda_{\min}(Q)\geq
       \Delta\left(\frac{r-1}{t}\right)^{r-1},
 \qquad
 \|Q^{-1}\|\leq
       \frac{(t/(r-1))^{r-1}}{\Delta}.
 \tag{B17f}
\]
Indeed the product of the other $r-1$ eigenvalues is at most
$(t/(r-1))^{r-1}$ by the arithmetic--geometric mean inequality.
For $r=1$, $\|Q^{-1}\|=\Delta^{-1}$.
Also
$\|\mathcal T^*Q^{-1}f\|_S^2
 =\langle Q^{-1}f,f\rangle_G
 \leq\lambda_{\min}(Q)^{-1}\|f\|_G^2$.
These give an explicit analytic matrix certificate. B7's
arithmetic determinant bound applies only if its arithmetic
normalization hypotheses for this Gram matrix are separately
verified.

\paragraph{Specialization of the actual rank-zero Bessel identity.}
Choose a rational negative line $L\subset V$ at $\mu$ and put
$A=L^\perp$, $H=\U(A)$. Then $A$ is rank two and totally
positive definite. Put
\[
 L_z=-(N_{E/F}\xi_0)^{-1}L,\qquad
 T=V,\qquad X=T\perp L_z .
\]
The identity map between the underlying one-dimensional vector
spaces supplies the exact line relation
$G_L=-N_{E/F}(\xi_0)G_{L_z}$ required by
\texttt{SeeLast}. Use identity rank-zero Witt decompositions
for $V$ and $P$ and the same measures on the identified copies.
The zero-dimensional unipotent measure has mass one.

Choose a splitting $\chi_z$ for $L$ and put
$\chi_A=\chi_V\chi_z^{-1}$. The orthogonal source decomposition
$V=A\perp L$ then has character identity
$\chi_V=\chi_A\chi_z$. The exchanged second pair in (B17c)
has characters $(\chi'_P,\chi_T)=(\chi_P,\chi_V)$, so the
upper pair $(P,X)$ has characters
\[
 (\chi'_P,\chi_T\chi_z^{-1})=(\chi_P,\chi_A),
 \qquad (P,L_z)\text{ has }(\chi_P,\chi_z^{-1}).
\]
These are exactly the character equations in
\path{Development/AbelianVarietyTheory/Definitions/LatticeFrames/IsQuasiFam.lean},
lines 114--117. As in B15f,
such global splitting characters exist; their selection is
part of the fixed input, with no assertion about prescribed
unramified finite values.

Let $e$ be a continuous automorphic test function on $[H]$, put
$\varphi=\overline e$, and use the rank-zero Bessel period
\[
 B_e(f)=\int_{[H]}f(\iota_A(a))\overline{e(a)}\,da .
\]
In the notation of (B15l), let
$C=C_{\Phi,\Phi^\dagger}=\texttt{lastBi}(\Phi,\Phi^\dagger)$
be the exact contracted two-variable theta input. Its formula
on $X=V\perp L_z$ is
\[
 C(u;v_T,v_z)
     =\Phi^\dagger(v_T)\,
       \Phi\bigl((o.\operatorname{res}_S)^{-1}
                       (u,\operatorname{ofAdj}_{L,P}(v_z))\bigr).
\]
The line map here is the identity chosen above; the source
definition incorporates its Gram matrices and $\xi_0$ in
$\operatorname{ofAdj}$. Define $\mathcal J_e(x)$ to be the
archive's upper double lift
$\operatorname{biLift}_{A\to P\to X}(\overline e;C)(x)$,
with the same $da,dh$ as above. Then the exact identity is
\begin{equation}
 B_e(f)=\int_{[G]}(Q^{-1}f)(t)\,
                 \mathcal J_e(\iota_T(t,1))\,dt
       \qquad(f\in\mathcal E_0).
\tag{B17g}
\end{equation}

\textbf{Proof.}
Set
$\beta_e(h)=\int_{[H]}\overline{e(a)}K(\iota_A(a),h)\,da$.
Substitute (B17d) in the left side and use Fubini on
$[H]\times[S]\times[G]$. The result is
\[
 \int_{[G]}F_f(t)\int_{[S]}
       \overline{K(t,h)}\,\beta_e(h)\,dh\,dt.
\]
All factors are continuous on the respective compact quotients,
so their product is absolutely integrable. Formula (B17c)
identifies the first kernel in the inner integral with
$\theta^\dagger_{P,V}(h,t;\Phi^\dagger)$.
The orthogonal tensor identity and the line relation now give
\[
 \mathcal J_e(\iota_T(t,1))
       =\int_{[S]}\theta^\dagger_{P,V}(h,t;\Phi^\dagger)
                         \beta_e(h)\,dh .
\]
This is precisely the integrand equality of
\texttt{SeeLast}, after integration over $[H]$ and $[S]$.
It proves (B17g).

These are the actual hypotheses used by the recovered analytic
theorems. Both definition files in this paragraph have directory
prefix
\path{Development/AbelianVarietyTheory/Definitions/LatticeFrames/}.
\texttt{IsThetaFam.lean}, line 17, defines
\texttt{CoCpt} by a compact set of adelic representatives;
its structures \texttt{BesLev}, \texttt{JacLev}, and
\texttt{Last} ask for this condition, not for definiteness at
every real place. \texttt{IsQuasiFam.lean}, lines 71--75,
states \texttt{LastLaw}, and lines 110--123 state
\texttt{SeeLast}. The complete proofs, in
\path{Development/AbelianVarietyTheory/Proofs/DecoratedLatticeProofs/}, are
\nolinkurl{DoubleLiftPeriodRankZeroFubiniCompactQuotient.lean},
lines 136--159, and \nolinkurl{LastSeeSawThetaIdentity.lean},
lines 22--54. The proof above instantiates their analytic
content with $S=P$ and $T=V$. It does not require the
additional \texttt{Comm.Is} hypotheses used to package
\nolinkurl{GlobalTheory.LastRed.main}. \hfill$\square$

\paragraph{What is now obtained, and what remains arithmetic.}
For a nonzero actual A7 wedge, make the B12 choice of rational
$H$ and rebase the coefficient frame before fixing $K,Q$, as
specified above. B12 then supplies an $e$ for which $B_e(f)\ne0$.
Equation (B17g) is therefore a
nonzero actual last-step Bessel seesaw involving that wedge.
It is uniform in the source characters for the fixed data.
The identity uses a common positive operator $Q^{-1}$, not
an assertion that this operator is a scalar or is algebraic.

Equation (B17g) is an identity for one Bessel coefficient.
The entire geometric mixed period is recovered from the
Petersson pairing in (B17e), with the fixed geometric measure
conversion. If a description solely through definite Bessel
coefficients is wanted, apply the finite Bessel frame of A5 to
$\mathcal E_0$, and then apply (B17g) to each of its coefficients.
For an adelic translate in that frame, translate the first
kernel by the same group element; \texttt{LastLaw} permits the
two kernels to be different. This combines actual last-step
seesaws to recover the whole pairing. A single restriction
period is not being identified with the whole ball integral.

The upper space is completely determined here. Since
$N_{E/F}\xi_0$ is a norm, $L\perp L_z$ is a hyperbolic plane:
if $e$ and $e_z$ have Gram entries $a$ and
$-a/N_{E/F}\xi_0$, the vector $e+\xi_0e_z$ is isotropic.
Thus
\[
 X\simeq A\perp\mathbb H,\qquad
 \operatorname{rank}X=4,\qquad
 \operatorname{sig}_\tau X=(3,1)
        \text{ for every real }\tau .
\]
Its Witt index is one, because $A$ is definite at every real
place. The quotient $[X]$ is not being integrated over in
(B17g); evaluation on the embedded compact $[G]$ is sufficient.
This feature causes no convergence gap in the displayed
identity.

The arithmetic \texttt{MainX} family, by contrast, asks for
totally definite spaces, specified arithmetic functionals,
and the common kernel/covolume and terminal mass equations
(B15m). Neither the mixed $P$ nor this rank-four $X$ is replaced
by a definite inner form in the argument. The result closes
the analytic double-lift and last-step comparison for the
actual A7 kernel while keeping those arithmetic matching
conditions separate. In particular it proves no algebraicity
of $Q^{-1}$ and no new projected algebraic cycle.


% THETA3-B18 contribution by /root/definite_bridge, 21 November 2026.
\subsection{THETA3-B18: the remaining definite alternatives have a finite parameter test}\label{B18}

The analytic identity of B17 uses the actual multiplicity plane and
is unconditional on a new global descent. A different question is
whether a prescribed irreducible rank-three constituent also comes
from a totally definite rank-two group. The following statements
make that question more explicit. The global base-change hypotheses
below are inputs, not calculations of the actual A7 constituent.

\paragraph{The complete real deletion table in the canonical case.}
Use B15's real convention, including $\psi(x)=e^{-2\pi i x}$.
Let the rank-three input have Harish--Chandra parameter
$(1,0;-1)$ at $\mu$, and $(1,0,-1;\varnothing)$ at a compact
place. Write the target rank-two splitting character as
$\chi_2(z)=(z/|z|)^{m_0}$, with $m_0$ even; the other character
has any fixed odd exponent $n_0$. A rank-two lift is nonzero
only when $m_0\in\{-2,0,2\}$. Its unique possible signature is
then given by
\[
\begin{array}{c|c|c|c|c}
 m_0&\lambda_\mu-m_0/2&(r_\mu,s_\mu)
       &\lambda_\tau-m_0/2&(r_\tau,s_\tau)\\ \hline
 -2&(2,1;0)&(2,0)&(2,1,0;\varnothing)&(2,0)\\
 0 &(1,0;-1)&(2,0)&(1,0,-1;\varnothing)&(1,1)\\
 2 &(0,-1;-2)&(1,1)&(0,-1,-2;\varnothing)&(0,2)
\end{array}
\tag{B18a}
\]
All six listed lifts are nonzero. In particular total definiteness
requires $m_{0,\mu}\in\{-2,0\}$ and
$m_{0,\tau}\in\{-2,2\}$ elsewhere. The other splitting character
changes the output by a determinant twist and does not change this
table.

\textbf{Proof.}
Ichino's Theorem 4.1(ii), pp.11--12, with $n=3,m=2$, requires a
zero in the shifted parameter: its condition $p_1+q_1>0$ is
necessary. This proves the asserted restriction on $m_0$.
Delete that zero. Each remaining positive entry in the first
block and negative entry in the second block contributes to the
positive target block; the other entries contribute to the
negative target block. The same theorem gives exactly the six
signatures in (B18a), so no other rank-two signature is possible.
For nonvanishing, apply Atobe's Definition 1.6 and Theorem 1.7(2),
pp.4--5. In five cases the zero chain has length one and the
listed signature is its base point $(r_\lambda,s_\lambda)$.
For the compact input with $m_0=0$, the maximal chain is instead
$(1,0,-1)$ of length three; its base point is $(0,0)$, and the
listed $(1,1)$ is obtained by the theorem's equal increment
$l=1$. This checks that case without mistaking its chain length
for one. The output Harish--Chandra parameter is obtained from
the retained entries by adding $n_0/2$, as in B15.
\hfill$\square$

The conversion between directions is explicit. Start with a
Hermitian rank-three $V$ and skew-Hermitian $W$ of rank two.
For a trace-zero scalar $\delta$, switch to the Hermitian space
$(W,\delta^{-1}h_W)$ and the skew-Hermitian space
$(V,\delta h_V)$. The tensor symplectic form is unchanged.
Ichino, \S2.2, p.5, states that this switches the two factors of
the same oscillator. At a real place take $\delta=i$, or a
positive multiple of it in the fixed complex embedding, so the
signatures in (B18a) are precisely those of the Hermitian
multiplicity space. The splitting characters switch their names;
none is inverted in this factor-switch rule. These identifications
must be made before applying the table to the archive convention.

\paragraph{The exact finite-place parameter and inner-form rule.}
Now label the original pair by $V_3$ Hermitian and $W_2$
skew-Hermitian, with characters $\chi_3|_{\A_F^\times}=\epsilon$
and $\chi_2|_{\A_F^\times}=1$. In the local quotient convention
\[
 \Hom_{\U(V_3)\times\U(W_2)}
       (\omega,\Pi\boxtimes\Sigma)\ne0,
\]
Gan--Ichino, \emph{The Gross--Prasad conjecture and local theta
correspondence}, arXiv:1409.6824v2, Theorem 4.4, pp.14--15,
gives at every nonsplit finite place
\[
 \phi_{\Pi_v}
    =(\phi_{\Sigma_v}\otimes\chi_{3,v}^{-1}\chi_{2,v})
                                      \oplus\chi_{2,v}.
 \tag{B18b}
\]
This theorem covers all local $L$-parameters, not just tempered
ones. For fixed $\Pi_v$ and fixed splitting characters, the
smaller parameter is obtained by deleting one occurrence of
$\chi_{2,v}$. The theorem's packet bijections then determine a
unique smaller representation, including its Hermitian inner
form. If the smaller parameter does not contain $\chi_{3,v}$,
use Theorem 4.4(i)(b) for the fixed target sign; if it does,
use Theorem 4.4(ii)(c). The enhanced packet character is the
restriction prescribed in \S4.6, proved by Theorem 1.3. This
last assertion is a finite-place statement; (B18a) supplies the
real-place assertion separately.

The same parameter formula holds at split finite places, where
there is a unique space of each dimension. One can check the
convention rather than silently omit a contragredient.
M\'inguez, \emph{Correspondance de Howe explicite: paires duales
de type II}, Ann. Sci. ENS 41 (2008), Corollary 6.3, p.734,
gives $\phi^\vee\oplus1$ for his untwisted matrix model in
ranks $2\to3$. His equation (1.4) and footnote on p.721 state
that replacing the first action by transpose inverse replaces
the input by its contragredient. Thus the switched model has
parameter $\phi\oplus1$. Twisting its first action by
$\chi_3\circ\det$ and its second by $\chi_2\circ\det$ gives
(B18b). This is the split normalization in Ichino's \S8.1.3,
Lemma 8.4, p.36; M\'inguez's statement also covers non-tempered
inputs. In particular (B18b) is available at every unramified
place, without a Ramanujan assumption.

\paragraph{A conditional global uniqueness test.}
Suppose that a fixed irreducible automorphic $\Pi$ on $\U(V_3)$
has the real components specified above and occurs in a global
theta pair with an irreducible automorphic $\Sigma$ on $\U(W_2)$.
Here $\Sigma$ labels the local oscillator quotient in the
preceding display. In the archive bilinear integral of B17, the
integrating automorphic input has type $\Sigma^\vee$, or its
chosen unitary conjugate realization. Assume that the standard
weak base changes $\mathcal B_3$ and $\mathcal B_2$ are specified isobaric
sums of unitary cuspidal automorphic representations of total
ranks three and two over $E$, with the local parameters in (B18b)
at almost all places. Then
\[
 \mathcal B_3=
   (\mathcal B_2\otimes\chi_3^{-1}\chi_2)\boxplus\chi_2.
 \tag{B18c}
\]
Indeed (B18b) identifies the unramified parameters at almost all
places, and strong multiplicity one for unitary isobaric
representations gives (B18c). The version being used is the
usual finite-exception case of Rajan, \emph{On strong
multiplicity one for automorphic representations},
arXiv:math/0210235, Theorem 1, p.2, and its isobaric extension
stated after the proof on p.6.

If $\mathcal B_2$ is cuspidal on $\mathrm{GL}_2$, then
$\chi_2$ is the unique one-dimensional constituent of the
isobaric cuspidal support in (B18c). Consequently any other
rank-two global realization of this same $\Pi$, in the same
standard-parameter convention and with base changes of the
same specified isobaric kind, must have the same rank-two character
$\chi'_2=\chi_2$. To prove this, apply (B18c) to the second
realization and compare the unique isobaric cuspidal supports.
Its displayed $\mathrm{GL}_1$ summand must be the unique one
in the first support.

Changing $\chi_3$ while fixing $\chi_2$ only twists the smaller
representation by the corresponding determinant character.
The oscillator action on the smaller group is multiplied by
that character, whereas its action on $\U(V_3)$ is unchanged.
Thus the local space is unchanged. More explicitly, write
$\chi'_3/\chi_3=\eta$, and use Hilbert 90 to define
$\nu_\eta(z/\bar z)=\eta(z)$. The source representation is
transported by $\nu_\eta\circ\det$; its parameter in (B18b)
is multiplied by $\eta$, exactly cancelling $\chi_3'^{-1}$.
The finite-place packet bijection and the real table then
force the original smaller local isometry classes.
Therefore, under these base-change and cuspidality hypotheses,
a mixed real place of the original smaller space obstructs a
totally definite rank-two realization of the same $\Pi$.
This is a conditional obstruction for the specified constituent;
no cuspidal $\mathrm{GL}_2$ transfer has been computed here for
an actual A7 wedge.

A compensating determinant twist of the rank-three realization
does not remove this obstruction. It multiplies every summand
of $\mathcal B_3$ and the distinguished $\chi_2$ by the same
character. At a real place it adds $a$ to every entry of
$\lambda$ and adds $2a$ to $m_0$, leaving
$\lambda-m_0/2$ and hence (B18a) unchanged. This recovers the
invariance proved in B16 at the level of the explicit table.

If instead $\mathcal B_2=\eta_1\boxplus\eta_2$ is a sum of
characters, (B18c) has at most three distinct character
constituents. Any alternate $\chi'_2$ must be one of them,
must have the allowed restriction to $\A_F^\times$, and must
pass (B18a) at every real place. For each candidate, deletion
in (B18b) and its enhanced packet rule determine the finite
local smaller representation. Coherence and global occurrence
remain to be proved for that candidate, for example through
the pole criterion of B16. This is a finite list of specified
character tests, not a claim that every deletion occurs globally.
Nontrivial Arthur $\mathrm{SL}_2$ blocks or base changes outside
the unitary isobaric hypothesis require their own analysis;
they are not covered by the preceding two cases.

\paragraph{An equal-rank central-value implication.}
There is also a precise sufficient condition for constructing a
new definite rank-two group from rank-two input. Use exactly
Jiang--Wu's notation in \S3.1--3.2 and Theorem 3.3(1), p.17.
Let $X$ be skew-Hermitian of rank two, let $\sigma$ be irreducible
cuspidal on $\U(X)$, and let $\chi$ be unitary with
$\chi|_{\A_F^\times}=1$.
Put $\rho=\sigma\otimes\chi^{-1}$ in their unitary-group
character convention. Suppose
\[
 L^S(1/2,\sigma\times\chi)\ne0,
 \tag{B18d}
\]
with a finite value, suppose $\rho$ has at least one local
component which is not one dimensional, and suppose that
at each real place every nonzero rank-two local theta lift
of the indicated input is definite. Then some global definite
Hermitian $Y$ of rank two has a nonzero cuspidal theta lift of
$\rho$. Its finite local forms are coherent automatically.

\textbf{Proof.}
In that theorem take $m=2$, $\epsilon_\chi=0$ and $j=1$.
Then $m-\epsilon_\chi$ is even and
$(m+1-\epsilon_\chi)/2-j=1/2$. Hence
$\operatorname{LO}_{\psi,\chi}(\rho)\leq2$.
The dimension is zero or two. The rank-zero theta series is
a determinant character, as in Jiang--Wu's \S3.1; the local
hypothesis excludes its occurrence. The first occurrence is
therefore two and is nonzero cuspidal by the definition of
lowest occurrence. The assumed real local condition makes its
global space definite. The global construction supplies
coherence at the finite places. \hfill$\square$

The real condition in this last statement can itself be made
explicit when the fixed input $\rho$ is discrete series at every
real place. For a mixed input $\U(1,1)$ with parameter $(a;b)$,
$a,b\in\Z+1/2$, $a\ne b$, choose the even exponent $m_0$ of
the target splitting character so that $m_0/2$ lies strictly
between $a$ and $b$. The equal-rank case $k=0$ of Ichino's
Theorem 4.1(ii) then gives the unique signature $(2,0)$ if
$a>b$, and $(0,2)$ if $a<b$. Atobe's base-point criterion
proves nonvanishing. For a compact input, choose $m_0/2$
outside the interval spanned by its two half-integral parameters;
the same rules again give a unique definite target and a
nonzero lift. Suitable integers exist in all these cases.
B15's character-extension lemma globalizes these even infinity
exponents, allowing finite ramification. If the desired input
is fixed as $\rho$, set $\sigma=\rho\otimes\chi$ so that
$\rho=\sigma\otimes\chi^{-1}$ holds in Jiang--Wu's convention.
Condition (B18d) must then be tested for this resulting pair
$(\sigma,\chi)$; the freedom to choose infinity characters
does not prove the central value is nonzero.

At a nonsplit finite place the equal-rank local test is also
explicit. In Gan--Ichino's conventions, their Theorem 4.1(ii),
p.13, says
\[
 \begin{gathered}
 \theta_{Y_v}(\rho_v)\ne0
 \quad\Longleftrightarrow\quad
 \epsilon\bigl(1/2,\phi_{\rho_v}\chi_{Y,v}^{-1},
                      \psi_{2,v}^E\bigr)
               =\epsilon(Y_v)\epsilon(X_v),\\
 \psi_{2,v}^E(x)=\psi_v(\operatorname{Tr}_{E_v/F_v}(\delta x)).
 \end{gathered}
 \tag{B18e}
\]
Its two sides determine one of the two finite inner forms.
The additive character in this epsilon factor has no factor
$1/2$; the $\psi^E$ used for the Whittaker normalization in
that paper does have that factor. These characters must not be
confused. When (B18d) applies, the global lowest-occurrence
theorem constructs a coherent collection of these finite forms;
otherwise guessing their signs separately does not prove
coherence or a nonzero global lift.

This central-value condition is different from B15's upward
value-at-one test and B16's rank-three pole-at-one test. It can
produce an equal-rank definite transfer when its hypotheses are
verified. It does not say that the resulting smaller
representation lifts back to the prescribed $\Pi$, or that it
satisfies the arithmetic kernel/covolume equation of B15.
In particular the conditional uniqueness obstruction above
must still be respected when trying to use such a transfer in
a rank-two-to-rank-three realization of the original wedge.


\section{Explicit quantitative certificates}
\label{sec:quantitative}
The estimates in this section have completely named normalizations and inputs. They do not claim a scalar-invariant positive lower bound, or a uniform degree bound for the final algebraic cycles. The latter would require arithmetic bounds for the rational source maps and extraction coefficients.

\subsection{THETA3-Q1: fixed-input perturbation of the geometric period}\label{Q1}
Fix the compact surface $X$, the total finite theta input and the determinant-line trivializations of 032 Proposition~4.6. On smooth $(2,0)$-forms use the Hermitian pairing
\[
 (\xi,\eta)=\int_X\xi\wedge\overline\eta,
 \qquad \|\xi\|^2=(\xi,\xi).
\]
For a nonzero $\xi$ put $A=\|\xi\|^2$. If
$\|\xi'-\xi\|\leq\eta\|\xi\|$, $0\leq\eta<1$, then
\begin{equation}\label{eq:perturb}
 \left|\int_X\xi\wedge\overline{\xi'}\right|
 \geq (1-\eta)A.
\end{equation}
If $\xi'=\Xi(t)$, $\xi=\Xi(t_0)$, on a convex coordinate neighborhood with $\|D\Xi\|\leq L$, it suffices that
$\|t-t_0\|\leq\eta\sqrt A/L$. For $L=0$ the assertion is immediate. A certified lower bound $A\geq A_0>0$ can replace $A$ in this sufficient radius.

On a finite etale cover $p:Y\to X$ of degree $d$, suppose
\[
 p^*\xi=\sum_r a_r\alpha_r\wedge\beta_r,
 \qquad p^*\xi'=\sum_s b_s\gamma_s\wedge\delta_s.
\]
Set $S=(\sum_r|a_r|)(\sum_s|b_s|)$. At least one of the resulting four-form periods has absolute value at least
$d(1-\eta)A/S$.
\begin{proof}
Cauchy--Schwarz gives
$|(\xi,\xi'-\xi)|\leq\|\xi\|\|\xi'-\xi\|\leq\eta A$.
The reverse triangle inequality proves \eqref{eq:perturb}. Integrating $D\Xi$ on the segment from $t_0$ to $t$ proves the sufficient radius. The pairing of $(2,0)$-forms remains positive along a smooth intermediate path; holomorphy is needed at the endpoints for the cohomological interpretation. Pullback by a degree-$d$ cover multiplies integration by $d$. Expand the two finite sums and apply the triangle inequality. Both forms are nonzero because $\eta<1$, so $S>0$.
\end{proof}
\framework{This quantifies the continuity step of 032 Corollary~4.7 on a fixed Betti realization. The cycle is still the image of a surface under the rational Hodge-source maps. The relevant polarization/correspondence framework is cote~44, scans~7--8; it does not assert that this analytic constant is a degree of a cycle.}

\subsection{THETA3-Q2: a certified Gaussian isolation bound}\label{Q2}
Let $\Lambda\subset\R^d$ be a lattice with minimum nonzero norm at least $\rho>0$. Assume $\operatorname{Im}Z\geq\lambda I$, $\lambda>0$, for the period matrix in the first theta tensor. Fix $0\ne x_0\in\Lambda$ and an integer $k\geq2\|x_0\|/\rho$. The contribution of $x_0+k\Lambda$ differs from its $x_0$ term by a vector of norm at most
\begin{equation}\label{eq:gaussian-tail}
 E(k)=2\sqrt{\frac{\pi}{e\lambda}}
 \sum_{j\geq1}(2j+3)^d e^{-b j^2},
 \qquad b=\frac{\pi\lambda k^2\rho^2}{8}.
\end{equation}
For $f_j=(2j+3)^d e^{-bj^2}$ and
\[
 q_J=\left(\frac{2J+5}{2J+3}\right)^d e^{-b(2J+1)}<1,
\]
the tail beginning at $J$ is at most $f_J/(1-q_J)$.
If projected leading gradients $u_0,v_0$ satisfy
$\|u_0\wedge v_0\|=\Delta>0$, and the two projected errors are at most $e_1,e_2$, then
\begin{equation}\label{eq:wedge-error}
 \|u\wedge v-u_0\wedge v_0\|
 \leq e_1\|v_0\|+e_2\|u_0\|+e_1e_2.
\end{equation}
An upper bound strictly smaller than $\Delta$ certifies nonvanishing.
\begin{proof}
For $m\ne0$ in $\Lambda$, $\|x_0+km\|\geq k\|m\|/2$.
The first theta tensor summand has norm at most $2\pi r e^{-\pi\lambda r^2}$, where $r=\|x_0+km\|$. The inequality
$r e^{-\pi\lambda r^2/2}\leq1/\sqrt{\pi e\lambda}$ bounds this by
$2\sqrt{\pi/(e\lambda)}e^{-\pi\lambda k^2\|m\|^2/8}$.
Balls of radius $\rho/2$ around lattice points are disjoint. Points with
$j\rho\leq\|m\|<(j+1)\rho$ number at most $(2j+3)^d$, by comparing the volume of these balls with the ball of radius $(j+3/2)\rho$. Summing proves \eqref{eq:gaussian-tail}. The successive ratios $f_{j+1}/f_j$ decrease with $j$, which gives the geometric tail bound. A projection multiplies the error by its operator norm. Expanding
$(u_0+(u-u_0))\wedge(v_0+(v-v_0))$ gives \eqref{eq:wedge-error}.
\end{proof}
\paragraph{Numerical certificate in a specified Gaussian model.}
For $d=6$, $\rho^2=2$, $\lambda=1$, $k=6$, and orthogonal projected leading terms of amplitude $c=2\pi e^{-2\pi}$, the interval-arithmetic computation supplied with the package gives
\[
 c>0.01173348,\qquad E(6)<1.765593\cdot10^{-8},
 \qquad c^2-2cE(6)-E(6)^2>\frac{137}{10^6}.
\]
The six-dimensional lattice can be taken as the Minkowski image of
$\OO_{\Q(\sqrt5)}^3$: for one coordinate its squared norm is
$2a^2+2ab+3b^2$, with minimum~2. The computation certifies the displayed Gaussian estimate. Identification with a particular global theta convention still requires its symplectic coordinates, projection and determinant-line normalization; the numerical value is not asserted to be an unnormalized global period.

\framework{This is a quantitative form of 032 Proposition~4.5, at the level of sections of the cotangent bundle. It certifies a nonzero realization before the algebraic surface correspondence is extracted; cote~44, scans~7--8, locates that later correspondence/polarization step.}

\subsection{THETA3-Q3: rational compact data with a denominator}\label{Q3}
Choose a $\Q$-basis of $E$ and coordinates on a two-dimensional $E$-space $W$. Its diagonal real embedding has dimension $D=2[E:\Q]$. Let $B$ be the invertible real embedding matrix. For any archimedean target vector $u$ and $N\geq1$, rounding $NB^{-1}u$ coordinatewise to $a\in\Z^D$ gives $w=Ba/N\in W(E)$ with
\begin{equation}\label{eq:round}
 \|w-u\|\leq\frac{\|B\|\sqrt D}{2N}.
\end{equation}
At a place let the Hermitian matrix be $H$, with $\|H\|=h$, $\|u\|\leq U$, and $|u^*Hu|\geq\eta>0$. If $\epsilon=\|w-u\|\leq1$ and
$h(2U+1)\epsilon\leq\eta/2$, the signs of $w^*Hw$ and $u^*Hu$ agree. Their Hermitian line projectors satisfy
\begin{equation}\label{eq:projector}
 \left\|\frac{w(w^*H)}{w^*Hw}-\frac{u(u^*H)}{u^*Hu}\right\|
 \leq \frac{2h(2U+1)}{\eta}
       \left(1+\frac{hU^2}{\eta}\right)\epsilon.
\end{equation}
\begin{proof}
Each rounded coordinate has error at most $1/2$, proving \eqref{eq:round}. Both the denominator difference and the numerator difference in \eqref{eq:projector} are bounded by
$h(2U\epsilon+\epsilon^2)$. Hence $|w^*Hw|\geq\eta/2$. Writing the difference of the fractions as the numerator difference divided by $w^*Hw$, plus
$u(u^*H)(u^*Hu-w^*Hw)/((w^*Hw)(u^*Hu))$, and using
$\|u(u^*H)\|\leq hU^2$, proves the formula. Apply it at every real place. The complementary projector is $1-P$, so it has the same error bound.
\end{proof}
\framework{This makes the weak-approximation endpoint in 032 Corollary~4.7 explicit relative to a field basis and prescribed analytic neighborhood. It changes multiplicity subspaces, not the signature of $V$. The resulting rational datum can be spread over a finite type base; it is not a bound for its final cycle degree.}

\subsection{THETA3-Q4 and Q5: zero-orbit selection and extraction}\label{Q4}\label{Q5}
These are the quantitative versions of A2 and A3. Let the anisotropic theta integral in A1 use probability measures, with
$\Phi=\Phi_\infty\otimes\Phi_f$, $\Phi_f(0)=1$, and
$a=\Phi_\infty(0)\ne0$. For every $0<\epsilon<1$ there is
$h_f\in N(\A_f)$ such that
\begin{equation}\label{eq:finite-selection}
 |I_{\Phi_\infty}(\omega_f(h_f)\Phi_f)|\geq(1-\epsilon)|a|,
 \qquad\|\omega_f(h_f)\Phi_f\|_2=\|\Phi_f\|_2.
\end{equation}
Suppose inverse doubling gives the finite tensor expansion
$\sum_a c_a\bigotimes_{i=1}^4\phi_{i,a}$ and the $i$th torus expansion has at most $N_{i,a}$ permitted characters. Then an individual integrated fourfold period has absolute value at least the total absolute period divided by
\begin{equation}\label{eq:loss}
 S=\sum_a |c_a|\prod_{i=1}^4N_{i,a}.
\end{equation}
Zero character coefficients may be retained; exact zero testing is unnecessary.
\begin{proof}
By A1 the continuous function $T_\Phi$ has mean $a$ on the compact quotient $[N]$. Its maximum absolute value is at least $|a|$. At a maximizing representative $(n_\infty,n_f)$ choose rational elements $\gamma_j\in N(F)$ with $\gamma_{j,\infty}\to n_\infty$. Automorphy gives
$T_\Phi(\gamma_{j,\infty},n_f)=T_\Phi(1,\gamma_{j,f}^{-1}n_f)$.
Continuity proves the inequality for some $j$. The finite Weil action is unitary. Expand the finite products and apply
$|\sum_a c_a\sum_{\mathbf j}P_{a,\mathbf j}|
\leq S\max_{a,\mathbf j}|P_{a,\mathbf j}|$.
The relevant finite torus quotient bounds the permitted character counts by its order. For actual geometric integrals the explicit measure conversion in van Hoften--Sempliner \S4.2.8 must be included, together with connected-component coefficients. More generally, if the exact extracted identity is $I=\sum_\lambda b_\lambda J_\lambda$, use $S=\sum_\lambda|b_\lambda|$.
\end{proof}
\framework{The doubled oscillator pairing is a realization-level construction. Finite rational Hodge-source extraction and pushforward of the surface supply the cycle. These are different operations, and the displayed coefficient loss records their interaction.}

\subsection{THETA3-Q6: the local positive determinant}\label{Q6}
On $\C^2$ use $d\gamma(z)=\pi^{-2}e^{-\|z\|^2}dz$. Put
$\ell_u(z)=\sum_i u_i z_i$ and
$P_{u,v}(z,w)=\ell_u(z)\ell_v(w)-\ell_v(z)\ell_u(w)$.
Then
\begin{equation}\label{eq:fock-det}
 \|P_{u,v}\|^2_{L^2(d\gamma\otimes d\gamma)}
 =2\det\begin{pmatrix}\|u\|^2&\langle u,v\rangle\\
 \langle v,u\rangle&\|v\|^2\end{pmatrix}.
\end{equation}
For orthonormal $u,v$, $P_{u,v}/\sqrt2$ has norm~1. Under a unitary Bargmann transform and the compatible doubling intertwiner,
$\delta(f\otimes\bar f)(0)=\|f\|^2$, so this normalized local doubled value is~1.
\begin{proof}
The covariance identity $\int z_i\bar z_j\,d\gamma=\delta_{ij}$ shows that the two diagonal terms of $|P|^2$ each integrate to $\|u\|^2\|v\|^2$, and the two cross terms each integrate to $|\langle u,v\rangle|^2$ with negative sign. This proves \eqref{eq:fock-det}. The remaining assertions follow from the specified unitary transforms and the doubling identity.
\end{proof}
The determinant polynomial is the shape of the $q=1,p=2$ Fock wedge cocycle in the published Bergeron--Millson--M\oe glin, Proposition~5.17 and Proposition~5.19(1), pp.~48--49. The value~1 here is a chosen convention, not an assertion about a paper's unmodified numerical constants. With the raw convention retain
$a_0=\Phi_{\infty,\mathrm{split}}(0)\ne0$ from van Hoften--Sempliner Lemma~4.2.9. Its matrix-coefficient continuity and rational isometry approximation give
$|a|\geq(1-\eta)|a_0|$, and Q4--Q5 give
$(1-\epsilon)(1-\eta)|a_0|/S$.
\framework{This is local Hermitian positivity in the automorphic realization. It is compatible with the polarization language of cote~44, scans~7--8, but is not itself an algebraic intersection form or the global Alp\"oge determinant.}

\subsection{THETA3-Q7: a finite search with a certified grid size}\label{Q7}
View the additive Siegel unipotent $N$ over $\Q$, of dimension $d$. Let $K_N\subset N(\A_f)$ be a compact open additive subgroup fixing $\Phi_f$. Set
$\Lambda=N(\Q)\cap K_N$ and
$F(b)=T_\Phi(n(b),1_f)$. Let $B$ have as columns a $\Z$-basis of $\Lambda$. If $F$ is $L$-Lipschitz on $N(\R)/\Lambda$, then among the $N_0^d$ rational grid elements
\[
 \gamma=\sum_{i=1}^d\frac{j_i}{N_0}\lambda_i,
 \qquad 0\leq j_i<N_0,
\]
one has
\begin{equation}\label{eq:grid}
 |T_\Phi(1,n_f(-\gamma))|
 \geq |a|-\frac{L\|B\|\sqrt d}{2N_0}.
\end{equation}
Thus $N_0\geq\max(1,\lceil L\|B\|\sqrt d/(2\epsilon|a|)\rceil)$ suffices for the Q4 bound. The finite Schwartz norm is unchanged; any compact subgroup of $\U(V)(\A_f)$ fixing the total doubled input still fixes its translate.
\begin{proof}
Additive strong approximation gives $N(\Q)+K_N=N(\A_f)$ and identifies
the compact quotient as
\[
 N(\Q)\backslash N(\A)/K_N \simeq N(\R)/\Lambda.
\]
Its probability mean is $a$ by A1. Choose a maximum point of $|F|$ and round its coordinates modulo $\Z^d$ to the nearest grid point. The torus distance is at most $\|B\|\sqrt d/(2N_0)$. The Lipschitz estimate and reverse triangle inequality prove the bound for $F(\gamma)$. Rational automorphy gives $F(\gamma)=T_\Phi(1,n_f(-\gamma))$. Unitarity proves the norm assertion, and commuting dual-pair actions prove the compact-level assertion. Extracting separate four-factor tensors can still require a smaller common level.
\end{proof}
Here is an explicit sufficient value of $L$. On a chosen compact fundamental set for the anisotropic quotient suppose all rational supports lie in a lattice $L_0\subset\R^D$ of minimum at least $\rho$, the integrand vanishes outside $L_0$, and
\[
 |\omega(g)\Phi(x)|\leq C(1+\|x\|)^m e^{-\lambda\|x\|^2}
\]
uniformly, with known $C,\lambda>0$ and $m\geq0$. Write the unipotent phase as $e^{2\pi i\langle b,Q(x)\rangle}$, with
$\|Q(x)\|\leq c_Q\|x\|^2$. One may take
\begin{equation}\label{eq:L-bound}
 L=2\pi c_Q C\sum_{j\geq1}(2j+3)^D((j+1)\rho)^2
       (1+(j+1)\rho)^m e^{-\lambda j^2\rho^2}.
\end{equation}
Indeed termwise differentiation in a unit $b$-direction introduces at most $2\pi c_Q\|x\|^2$; probability integration cannot enlarge the resulting uniform bound. The shell packing argument of Q2 gives \eqref{eq:L-bound}. The differentiated series converges uniformly. Its terms have polynomial growth times a Gaussian, and their successive ratios eventually decrease below~1, giving an explicit geometric tail certificate. Compactness guarantees the existence of an envelope, but a numerical application must supply its constants or a fundamental set; no effective reduction-theory bound is being silently invoked.

With the Q6 convention $a_0=1$, choose the rational isometry so $|a|\geq1/2$. Taking
$N_0\geq\max(1,\lceil2L\|B\|\sqrt d\rceil)$ gives $|I|\geq1/4$. For an exact finite extraction $I=\sum_\lambda b_\lambda J_\lambda$, some geometric mixed period therefore satisfies
\begin{equation}\label{eq:quarter}
 |J_\lambda|\geq\frac1{4\sum_\lambda|b_\lambda|}.
\end{equation}
\framework{This gives a finite search in the adelic realization with a named denominator. It preserves the total doubled level, but does not bound the complexity of every subsequent rational Hodge map. The distinction between a realization and an algebraic correspondence is the one in cote~16, scans~43--45.}

\subsection{THETA3-Q8: a pointwise lower bound for the surface norm}\label{Q8}
Let an injective holomorphic chart in $X$ contain the Euclidean ball $B_R\subset\C^2$, and write
$\xi=f(z)\,dz_1\wedge dz_2$. Suppose $|f(0)|\geq b>0$ and the real differential of $f$ has norm at most $M$ on $B_R$. For $0<r<R$ and $Mr\leq b/2$,
\begin{equation}\label{eq:local-norm}
 \int_X\xi\wedge\bar\xi\geq\frac{\pi^2}{2}b^2r^4.
\end{equation}
\begin{proof}
Integrating the differential along a segment gives $|f(z)|\geq b/2$ on $B_r$. In the complex orientation,
$dz_1\wedge dz_2\wedge d\bar z_1\wedge d\bar z_2=4\,d\vol_{\R^4}$.
Since $\vol(B_r)=\pi^2r^4/2$, integration over $B_r$ gives the bound. The integrand is nonnegative on the rest of $X$. If $M=0$, any $r<R$ works.
\end{proof}
\framework{Together with Q2 and bounds for differentiated theta series, this turns a certified nonzero coefficient into the $A_0$ needed by Q1. The injective chart is required data. This is a Hodge norm estimate, not an integral intersection multiplicity.}

\subsection{THETA3-Q9: denominators in finite Hecke extraction}\label{Q9}
Let $\Lambda=H^1(X,\Z)/\mathrm{torsion}$, and let $T_j$ be commuting integral Hodge endomorphisms induced by rational algebraic correspondences. Suppose monic $f_j,g_j\in\Z[t]$ are coprime over $\Q$ and
$f_j(T_j)g_j(T_j)=0$. Set $R_j=\Res(f_j,g_j)\ne0$. An integral Bezout identity
$u_jf_j+v_jg_j=R_j$ gives
\[
 e_j=\frac{v_j(T_j)g_j(T_j)}{R_j},\qquad e=\prod_j e_j.
\]
Then $e$ is a rational Hodge projector induced by a rational algebraic correspondence, with
$e\Lambda\subset R^{-1}\Lambda$, $R=\prod_j|R_j|$. Its image is the intersection of the $f_j$-primary factors.

Let $K$ have degree $n$ over $\Q$. For $i=1,\ldots,4$, take such a projector $e_i$ with denominator $R_i$ and $u_i\in\Lambda\otimes K$ satisfying
$d_i u_i\in\Lambda\otimes\OO_K$. Put
\[
 P=\int_X e_1u_1\cup e_2u_2\cup e_3u_3\cup e_4u_4,
 \qquad D=\prod_{i=1}^4R_i d_i.
\]
Then $DP\in\OO_K$. If $P\ne0$ and
$|\sigma(P)|\leq M_\sigma>0$ for every $\sigma\ne\sigma_0$, then
\begin{equation}\label{eq:norm-bound}
 |\sigma_0(P)|\geq D^{-n}\prod_{\sigma\ne\sigma_0}M_\sigma^{-1}.
\end{equation}
\begin{proof}
The adjugate of the Sylvester matrix gives the integral Bezout identity. The Chinese remainder theorem gives $e_j=1$ on $\ker f_j(T_j)$ and $e_j=0$ on $\ker g_j(T_j)$; equivalently $e_j^2-e_j$ is divisible by $f_jg_j$ as a rational polynomial. Commutativity gives the product projector and its image. Integral polynomial numerators give the denominators. Polynomial operations preserve Hodge endomorphisms and are realized by addition and composition of algebraic correspondences.

Cup product evaluated on the integral fundamental class is a $\Z$-valued multilinear form on $\Lambda$. Multiplication of its four arguments by $R_i d_i$ therefore gives $DP\in\OO_K$. The nonzero algebraic integer $DP$ has nonzero integral norm, whose absolute value is at least~1. Consequently
$\prod_\sigma|\sigma(P)|\geq D^{-n}$. Divide by the other bounds to get \eqref{eq:norm-bound}.
\end{proof}
In the CM application, the cited Hecke/Frobenius and weak-admissibility arguments of 032 first identify the relevant eigenspaces as pure CM Hodge sources. Their eigenvectors can be chosen over a number field in the Betti realization. Harmonic representatives have the same cup product. Arbitrarily rescaled analytic theta forms need not be those algebraic Betti vectors; their comparison scalars must be retained. The matrices $T_j$, resultants, eigenvectors and conjugate bounds are required data. The supplied papers do not give universal numerical values for them.
\framework{This is a projector in the homological realization, not an assertion of an idempotent modulo rational equivalence. It implements the algebraic-correspondence requirement of cote~16, scans~44--45, with an explicit denominator. B7 supplies a positive Gram-matrix variant when total positivity under all embeddings is established.}


% THETA3-Q10 contribution by /root/theta_032, 21 November 2026.
% GPT-6 Astra Pro; inherited reasoning-effort label not separately exposed.

\subsection{THETA3-Q10: determinant normalization and a finite certificate in every degree}\label{Q10}

\paragraph{The exact Fock determinant norm.}
Let $u_1,\ldots,u_n\in\C^n$, put
$\ell_u(z)=\sum_a u_a z_a$, and let $z_1,\ldots,z_n$ be independent
standard complex Gaussian vectors, each with probability density
$\pi^{-n}e^{-\|z\|^2}$ relative to real Lebesgue measure. Define
\[
 P_u(z_1,\ldots,z_n)=\det(\ell_{u_i}(z_j))_{1\leq i,j\leq n},
 \qquad G_u=(\langle u_i,u_j\rangle)_{i,j},
\]
with the Hermitian inner product conjugate linear in its first
variable. Then
\begin{equation}\label{Q10-Fock}
 \|P_u\|^2=n!\det G_u,\qquad
 \langle P_u,P_v\rangle
       =n!\det(\langle u_i,v_j\rangle)_{i,j}.
\end{equation}

\paragraph{Proof.}
The Gaussian covariance is
$\mathbb E[\overline{z_a}z_b]=\delta_{ab}$, so
$\mathbb E[\overline{\ell_u(z)}\ell_v(z)]
=\langle u,v\rangle$.
Expand both determinants. Independence of the $n$ columns gives
\[
 \langle P_u,P_v\rangle
 =\sum_{\sigma,\tau\in\mathfrak S_n}
   \operatorname{sgn}(\sigma)\operatorname{sgn}(\tau)
       \prod_{j=1}^n\langle u_{\sigma(j)},v_{\tau(j)}\rangle.
\]
For fixed $\sigma$, relabel $i=\sigma(j)$ and put
$\rho=\tau\circ\sigma^{-1}$. The sign product is
$\operatorname{sgn}(\rho)$ and the inner sum is the determinant
of the displayed mixed Gram matrix. There are $n!$ choices of
$\sigma$. Taking $v=u$ proves the norm identity.
\hfill$\square$

For independent $u_i$, divide by $\sqrt{n!\det G_u}$ to obtain
a unit vector. In particular the coordinate determinant divided
by $\sqrt{n!}$ has norm one. BMM, Proposition 5.19(1), p.49,
with $q=1$ and $b=n$, identifies the coefficient of the
holomorphic wedge cocycle with precisely this determinant shape.
The other active Fock coordinates and all remaining compact
Gaussians may be normalized to norm one. The unitary Bargmann
transform preserves the resulting norm. Under the standard
doubling map,
\[
 \delta(f_\infty\otimes\bar f_\infty)(0)=\|f_\infty\|^2=1.
\]
This is a chosen complete normalization, not an assertion that
an unmodified source convention already has the number one.
For the actual mixed endpoint vector retain its value
$a=\Phi_\infty(0)$; a certified continuity estimate
$|a-1|\leq\eta<1$ gives $|a|\geq1-\eta$.

\paragraph{Dimensions in the doubled model.}
For the ball of dimension $n$, $\dim_E V=n+1$. Doubling the
rank-$n$ skew-Hermitian source gives a split source of rank $2n$.
The Siegel unipotent is the additive group
$N\simeq\operatorname{Herm}_n(E)$ over $F$. A Hermitian matrix has
$n$ diagonal $F$-entries and $n(n-1)/2$ off-diagonal $E$-entries;
therefore
\[
 d_N:=\dim_\Q N=n^2[F:\Q].
\]
The Schr\"odinger position space is $V^n$ over $E$, of dimension
$n(n+1)$, and hence its real dimension at all archimedean places is
\[
 D=2n(n+1)[F:\Q].
\]
These are different dimensions: $d_N$ controls the unipotent grid,
whereas $D$ controls the position-lattice Gaussian tail.
In the grid statement below, $N$ is viewed by restriction of
scalars as a $\Q$-vector group; thus $N(\Q)$ denotes its
$F$-rational Hermitian matrices and $N(\A_f)$ their finite
$F$-adelic counterparts.

\paragraph{The explicit unipotent grid.}
Use probability measures on the compact quotients, fix
$\Phi_f(0)=1$, and fix the splitting for which the $G=\U(V)$
action in the doubled model is untwisted pullback. Let $K_N$ be
a compact open additive subgroup of $N(\A_f)$ fixing $\Phi_f$.
Set
\[
 \Lambda_N=N(\Q)\cap K_N,
 \qquad B:\R^{d_N}\longrightarrow N(\R),
\]
where the columns of $B$ are a $\Z$-basis of this full lattice.
Suppose the function
$F(b)=T_\Phi(n(b),1_f)$ is $L$-Lipschitz on
$N(\R)/\Lambda_N$. For $N_0\geq1$ define the explicit grid
\[
 \Gamma_{N_0}=
 \left\{\sum_{j=1}^{d_N}\frac{a_j}{N_0}\lambda_j:
                         0\leq a_j<N_0\right\}.
\]
Then one of its $N_0^{d_N}$ elements satisfies
\begin{equation}\label{Q10-grid}
 |T_\Phi(1,n_f(-\gamma))|
    \geq |a|-\frac{L\|B\|\sqrt{d_N}}{2N_0}.
\end{equation}

\paragraph{Proof.}
Additive strong approximation gives
$N(\A_f)=N(\Q)+K_N$ and identifies the quotient by $K_N$
with $N(\R)/\Lambda_N$. The quotient measure is its probability
Haar measure. The general rank-$r$ statement A1, applied with
$r=n$, gives its mean $a=\Phi_\infty(0)$: Siegel-unipotent
averaging annihilates every nonzero Gram matrix, and anisotropy
of $V$ leaves only the zero tuple. Hence
$\max |F|\geq|a|$.
Round the coordinates of a maximum point modulo $\Z^{d_N}$ to
the nearest grid point. Their Euclidean error is at most
$\sqrt{d_N}/(2N_0)$, so their images differ on the torus by at
most $\|B\|\sqrt{d_N}/(2N_0)$. The Lipschitz and reverse-triangle
inequalities give the bound for $|F(\gamma)|$.
Rational theta automorphy gives
$F(\gamma)=T_\Phi(1,n_f(-\gamma))$, proving
\eqref{Q10-grid}. The finite Weil action is unitary and commutes
with $G(\A_f)$, so the translate preserves both the finite
Schwartz norm and every compact target level that fixed the total
input. A factorization into individual inputs can require a
smaller common level. \hfill$\square$

If $|a|\geq1/2$, the explicit choice
\[
 N_0\geq\max\{1,\lceil2L\|B\|\sqrt{d_N}\rceil\}
\]
makes the error at most $1/4$, so the selected total integral
has absolute value at least $1/4$. Suppose the exact finite
tensor, source-character and connected-component expansion is
\[
 I=\sum_\lambda b_\lambda J_\lambda,\qquad
 S_{\rm ext}=\sum_\lambda|b_\lambda|>0,
\]
where the geometric normalization is included in $b_\lambda$.
Then at least one actual $2n$-factor integrated period satisfies
\begin{equation}\label{Q10-extract}
 |J_\lambda|\geq\frac{1}{4S_{\rm ext}}.
\end{equation}
This is the triangle inequality. If the tensor decomposition has
coefficients $c_a$ and its $2n$ torus factors have at most
$N_{i,a}$ permitted characters, a valid cost before the geometric
conversion is
$\sum_a|c_a|\prod_{i=1}^{2n}N_{i,a}$.
Zero character components may be retained in the count. The
certificate permits the displayed finite input translation; it
is not a lower bound for every preassigned unaveraged period.

\paragraph{A fully specified Lipschitz bound.}
Here is the all-rank version of the sufficient envelope in Q7.
Fix a compact representative set for the anisotropic target
quotient and a position lattice of minimum norm at least $\rho$,
containing the rational support for all its finite components.
Suppose, in the chosen Euclidean coordinates,
\[
 |\omega(g)\Phi(x)|
      \leq C(1+\|x\|)^M e^{-\lambda\|x\|^2}
\]
on that set, with $C,\lambda>0$. Write the unipotent phase as
$\exp(2\pi i\langle b,Q(x)\rangle)$ and suppose
$\|Q(x)\|\leq c_Q\|x\|^2$. Then one may take
\begin{equation}\label{Q10-L}
 L=2\pi c_QC\sum_{j\geq1}(2j+3)^D
       ((j+1)\rho)^2(1+(j+1)\rho)^M
                   e^{-\lambda j^2\rho^2}.
\end{equation}
To check this, differentiate the phase in a unit direction.
Its derivative has absolute value at most
$2\pi c_Q\|x\|^2$. The zero term contributes zero.
The shell $j\rho\leq\|x\|<(j+1)\rho$ contains at most
$(2j+3)^D$ points by the disjoint-ball argument. The displayed
bound follows by summing absolute values and integrating against
probability measure. Gaussian decay makes the differentiated
series uniformly convergent and justifies the operations.
The positive terms have an explicit polynomial factor times
$e^{-\lambda\rho^2j^2}$; after any index for which their
successive ratios are bounded by a number $q<1$, their tail is
at most its first term divided by $1-q$. This supplies a finite
certified evaluation.

For the determinant polynomial and its doubled vector one may
use $M=2n$ after the Gaussian model and its linear transformations
have been fixed: partial Fourier transform and metaplectic
changes of a polynomial Gaussian preserve the bound on its
polynomial degree, though they change the Gaussian and its
coefficients. In coordinates $Q_a(x)=x^tH_ax$ one can take
$c_Q=(\sum_a\|H_a\|^2)^{1/2}$. The lattice, envelope constants,
and geometric conversion must actually be supplied for numerical
use; no unrecorded effective reduction-theory bound is assumed.

\paragraph{A pointwise certificate for the top-form norm.}
If an injective holomorphic chart contains $B_R\subset\C^n$ and
$\xi=f(z)\,dz_1\wedge\cdots\wedge dz_n$ there, suppose
$|f(0)|\geq b>0$ and the real derivative of $f$ has norm at most
$M_1$ on the ball. For $0<r<R$ with $M_1r\leq b/2$,
\begin{equation}\label{Q10-topnorm}
 i^{n^2}\int_X\xi\wedge\bar\xi
       \geq \frac{2^{n-2}\pi^n}{n!}\,b^2r^{2n}.
\end{equation}
Indeed the segment estimate gives $|f|\geq b/2$ on $B_r$.
The positive top-form density is $2^n|f|^2$ times real
Lebesgue measure, and
$\operatorname{vol}(B_r\subset\R^{2n})=\pi^nr^{2n}/n!$.
Integration proves the formula. A determinant of first-tensor
values can be certified from individual errors $\varepsilon_i$
using the multilinear estimate
\[
 \left\|\bigwedge_i u_i-\bigwedge_i u_i^0\right\|
 \leq\prod_i(\|u_i^0\|+\varepsilon_i)-\prod_i\|u_i^0\|.
\]
Expand the wedge by subsets and apply the norm bound for each
term to prove this inequality. Together with A6's lattice tail,
it gives explicit input data for \eqref{Q10-topnorm}.

\paragraph{The definite subgroup and its local constant.}
Choose an $E$-rational line $L$ negative at the active place.
Then $H=\U(L^\perp)$ is totally definite of rank $n$. For a
canonical $n$-form, its right-regular coefficient at that line
has compact character
\[
 \kappa(\operatorname{diag}(A,d))=\det(A)d^{-n}.
\]
This follows because the tangent representation is
$\Hom(L,Q)$, whose derivative is $A/d$ and whose determinant
is $\det(A)d^{-n}$. It is trivial on the scalar center.
Suppose a unit vector $w$ in a unitary representation has that
compact type. With $\pi(A)=\det A$ on $H_\infty=\U(n)$ and
probability Haar,
\[
 \int_{\U(n)}
   \langle w,\sigma(\operatorname{diag}(A,1))w\rangle
                      \pi(A)^{-1}\,dA=1.
\]
The matrix coefficient is $\det A$ in the convention linear
in the second variable, so the integrand is one. The rank-zero
unipotent is trivial. This direct local computation is valid
for every $n$; the particular archive proposition imposing
even $k=n$ is invoked only when its additional hypotheses hold.
No global norm is identified with this local number.

\paragraph{Increasing the rank does not enlarge the canonical gap.}
For the compact sign pattern $(+,\ldots,+,-)$, the half sums are
\[
 \rho_c=\left(\frac{n-1}{2},\frac{n-3}{2},\ldots,
                         \frac{1-n}{2};0\right),\qquad
 \rho_n=\left(\frac12,\ldots,\frac12;-\frac n2\right).
\]
The right-regular canonical lowest compact weight is
$(1,\ldots,1;-n)$. Solving the archive's lowest-weight relation
$\lambda+\rho_n-\rho_c=(1,\ldots,1;-n)$ gives
\[
 \lambda=(1,\ldots,1;-n)-\rho_n+\rho_c
       =\left(\frac n2,\frac n2-1,\ldots,
                          1-\frac n2;-\frac n2\right).
\]
The compact determinant representation has tuple
$u_i=(n+3-2i)/2$, $1\leq i\leq n$.
Taking $v$ to be the first $n-1$ entries of $\lambda$ leaves
the archive's two nodes
\[
 b=\left(\frac n2-1,\frac n2\right).
\]
Their gap is exactly one in every rank.
The interlacing follows from
$u_i=\lambda_i+1/2>\lambda_i>u_{i+1}$ and the final inequality
$u_n>-\!b_0$; all cross differences are half-integral.
Thus increasing the dimension does not repair the canonical
\texttt{NodeGap} condition, whose required gap is at least
three. Invoking a general discrete-series theorem also requires
the representation identification demanded by that theorem;
the displayed direct compact integral does not require it.

For a pluricanonical power $K_X^q$, $q\geq1$, the compact weight
is $(q,\ldots,q;-nq)$. The same computation gives
\[
 \lambda_i=q+\frac n2-i\quad(1\leq i\leq n),\qquad
 \lambda_{n+1}=-nq+\frac n2,
 \qquad
 b=\left(\frac n2-q,nq-\frac n2\right).
\]
Its gap is $(n+1)q-n$. For $n=2$ this is $3q-2$, as in B9.
It enlarges the gap by changing the line bundle and its vector;
it does not identify the resulting section with the original
ordinary degree-$n$ cohomology class.


% THETA3-Q11 contribution by /root/theta_032, 21 November 2026.
% GPT-6 Astra Pro; inherited effort label not separately exposed.

\subsection{THETA3-Q11: definite periods give an explicit lower bound for the ball norm}\label{Q11}

This bound replaces the unknown upper bound for the ball Gram
matrix in A5 by local holomorphic evaluation constants.  It uses
the actual finite definite restrictions and a fixed smooth
quotient throughout.  No derivative bound for the local
holomorphic coefficient is needed.

\paragraph{Fixed geometric and period normalizations.}
Let $X$ be a compact smooth ball quotient of complex dimension
$n$, or the finite disjoint union of components of a fixed neat
Shimura level.  Let $\mathcal E\subset H^0(X,\Omega_X^n)$ be a
finite-dimensional space in its specified scalar adelic
realization.  Use the geometric Hermitian norm
\[
 \langle\xi,\zeta\rangle_X
       =i^{n^2}\int_X\zeta\wedge\overline\xi,
 \qquad \|\xi\|_X^2=i^{n^2}\int_X\xi\wedge\overline\xi.
 \tag{Q11a}
\]
It is conjugate linear in the first variable.  If an adelic
Petersson measure is used instead, its specified conversion to
(Q11a) must be included; the argument does not identify these
measures without that conversion.

Take a nonempty finite collection of restrictions $R_{a_j}$ of A5, with finite adelic
$a_j$, totally definite $H=\U(n)$, and the fixed canonical
scalar type.  The same argument applies to any such finite
collection; separation is needed only when positive definiteness
on all of $\mathcal E$ is desired.  Give each $[H]$ probability
Haar measure, choose an orthonormal basis $e_{j,\alpha}$ of its
finite restriction module, and write
\[
 B_{j,\alpha}(\xi)
     =\int_{[H]}(R_{a_j}\xi)(h)\,
                              \overline{e_{j,\alpha}(h)}\,dh .
\]
For representatives of
$H(F)\backslash H(\A_{F,f})/K_H(a_j)$, let $r_{j,t}(\xi)$
be the corresponding scalar values.  Their exact masses are
\[
 m_{j,t}=
    \frac{|\Gamma_{j,t}|^{-1}}
         {\displaystyle\sum_s|\Gamma_{j,s}|^{-1}}>0,
 \qquad
 \Gamma_{j,t}=H(F)\cap h_{j,t}K_H(a_j)h_{j,t}^{-1}.
 \tag{Q11b}
\]
Here the finite stabilizers are viewed in the compact real
group.  The formula follows by decomposing the quotient into
$\Gamma_{j,t}\backslash H(F\otimes_\Q\R)$, first giving the
real group volume one and then normalizing total quotient mass.
A stabilizer character can force a coordinate to vanish; such
coordinates may be retained with value zero.  On a neat level
the relevant finite stabilizers are trivial.  In particular the
masses are concrete rational numbers, although the form values
have not thereby been proved algebraic.

Finite-dimensional Parseval gives the exact positive quantity
\[
 \mathcal P(\xi):=\sum_{j,\alpha}|B_{j,\alpha}(\xi)|^2
          =\sum_{j,t}m_{j,t}|r_{j,t}(\xi)|^2.
 \tag{Q11c}
\]
Flatten the pairs $(j,t)$ into a single index $i$.  Each scalar
value is a linear evaluation in the canonical fibre at a point
$x_i\in X$.  Choose an injective holomorphic chart
$\phi_i:B_{R_i}\subset\C^n\longrightarrow X$ with
$\phi_i(0)=x_i$ and $0<r_i<R_i$.  Write
\[
 \phi_i^*\xi=g_i(z)\,dz_1\wedge\cdots\wedge dz_n,
 \qquad r_i^{\rm ev}(\xi)=\gamma_i g_i(0),
 \tag{Q11d}
\]
where $r_i^{\rm ev}$ denotes the scalar evaluation formerly
indexed by $(j,t)$, so it is distinct from the radius $r_i$.
The nonzero constant $\gamma_i$ is fixed by the actual
coefficient frame: explicitly it is the value of that
evaluation functional on the coordinate canonical generator
at $x_i$.  Thus it includes every prescribed frame or
rigidification factor.  Neither $\gamma_i=1$ nor its
algebraicity is assumed.

Put
\[
 A_i=\frac{2^n\pi^n r_i^{2n}}{n!},\qquad
 c_i=\frac{m_i|\gamma_i|^2}{A_i},\qquad
 C_{\rm eval}=\sum_i c_i.
 \tag{Q11e}
\]
The radii, frame constants and masses are all data for the
chosen quotient; changing its level requires recording the
corresponding changed data.

\paragraph{The explicit bound.}
For every $\xi\in\mathcal E$ one has
\[
 \boxed{\quad
       \mathcal P(\xi)\leq C_{\rm eval}\|\xi\|_X^2,
       \qquad
       \|\xi\|_X^2\geq
                \frac{\mathcal P(\xi)}{C_{\rm eval}}.
       \quad}
 \tag{Q11f}
\]
In particular any explicit lower certificate for the sum of
definite period squares gives a lower certificate for the
geometric ball norm.  This does not use an upper bound for
$G_{\rm ball}$.

\textbf{Proof, including the local constant.}
For a holomorphic function $g$ on $B_R$ and $0<r<R$,
\[
       \int_{B_r}|g(z)|^2\,dV(z)
          \geq \frac{\pi^n r^{2n}}{n!}|g(0)|^2.
 \tag{Q11g}
\]
To see the mean identity directly, fix $w\in B_r$.  The
one-variable function $u\mapsto g(uw)$ is holomorphic on a
disc containing the closed unit disc, so its circle average is
$g(0)$.  Integrate this identity in $w$.  Rotation invariance
of $B_r$ and Fubini show that the average of $g$ on $B_r$ is
$g(0)$.  Cauchy--Schwarz then gives (Q11g), using
$\operatorname{vol}(B_r)=\pi^nr^{2n}/n!$.

In coordinates $z_j=x_j+iy_j$, the order of the exterior
factors gives exactly
\[
 i^{n^2}\,dz_1\wedge\cdots\wedge dz_n\wedge
        d\bar z_1\wedge\cdots\wedge d\bar z_n
        =2^n\,dx_1\,dy_1\cdots dx_n\,dy_n.
\]
Applying (Q11g) to $g_i$ therefore proves
\[
 \int_{\phi_i(B_{r_i})}i^{n^2}\xi\wedge\bar\xi
          \geq A_i |g_i(0)|^2.
 \tag{Q11h}
\]
The chart is injective, so this is an integral over a genuine
subset of $X$, with no covering multiplicity.  Multiplication
by $m_i|\gamma_i|^2/A_i$ gives
\[
 m_i|r_i^{\rm ev}(\xi)|^2
       \leq c_i\int_{\phi_i(B_{r_i})}i^{n^2}\xi\wedge\bar\xi
       \leq c_i\|\xi\|_X^2.
\]
Summing and using (Q11c) proves (Q11f).
\hfill$\square$

The same proof retains overlap information.  If
$U_i=\phi_i(B_{r_i})$, one may replace $C_{\rm eval}$ by
\[
 C_{\rm ov}=\sup_{x\in X}\sum_{i:x\in U_i}c_i
         \leq C_{\rm eval}.
 \tag{Q11i}
\]
Indeed the sum of the local integrals is the integral against
the displayed nonnegative weight function.  In particular,
if no point lies in more than $\kappa$ chosen charts, then
$C_{\rm ov}\leq\kappa\max_i c_i$.
The single-point consequence is also sharper than a derivative
envelope: if $|g_i(0)|\geq b$, then
$\|\xi\|_X^2\geq A_i b^2$ for every recorded injective radius.
It removes the Lipschitz hypothesis from the pointwise estimate
in Q10 and its factor-four loss.

\paragraph{A mixed-period and component-extraction certificate.}
Suppose $\xi,\xi'\in H^0(X,\Omega_X^n)$ are expressed on the
same quotient and in the same canonical line bundle, and
\[
       \|\xi'-\xi\|_X\leq\eta\|\xi\|_X,
       \qquad 0\leq\eta<1.
 \tag{Q11j}
\]
They may have been constructed from different theta types;
that origin does not replace the asserted norm comparison.
Cauchy--Schwarz and (Q11f) give
\[
 \begin{split}
 \operatorname{Re}\langle\xi,\xi'\rangle_X
      &\geq(1-\eta)\|\xi\|_X^2,\\
 |\langle\xi,\xi'\rangle_X|
      &\geq\frac{1-\eta}{C_{\rm eval}}
                    \sum_{j,\alpha}|B_{j,\alpha}(\xi)|^2.
 \end{split}
 \tag{Q11k}
\]
For completeness, subtract $\|\xi\|_X^2$ from the pairing:
the difference is $\langle\xi,\xi'-\xi\rangle_X$, whose
absolute value is at most $\eta\|\xi\|_X^2$.  This proves
the first inequality and hence the second.

If the constructed comparison form has a finite expansion
$\xi'=\sum_\nu b_\nu\xi'_\nu$, set
$S=\sum_\nu|b_\nu|$.  For $\xi\ne0$, (Q11j) implies $S>0$.
Linearity in the second variable and the triangle inequality
then prove that at least one actual component satisfies
\[
 |\langle\xi,\xi'_\nu\rangle_X|
       \geq\frac{1-\eta}{S C_{\rm eval}}
                    \sum_{j,\alpha}|B_{j,\alpha}(\xi)|^2.
 \tag{Q11l}
\]
If both forms are expanded, the same proof uses the product
of their two coefficient $\ell^1$ sums.  The finite tensor and
character-extraction costs in A3 or Q7 can be inserted into
$S$ once their exact expansions are fixed.  No uniform bound
for those costs is silently supplied here.

\paragraph{A determinant certificate without the ball Gram matrix.}
Choose any basis $\xi_1,\ldots,\xi_D$ of $\mathcal E$ and form
the weighted evaluation matrix
\[
 E_{ia}=\sqrt{m_i}\,r_i^{\rm ev}(\xi_a),\qquad
 H_{\rm fin}=E^*E.
\]
For $\xi=\sum_a u_a\xi_a$ one has
$\mathcal P(\xi)=u^*H_{\rm fin}u$.  Suppose a square
$D$-row minor has determinant of modulus at least $\delta>0$
and $\operatorname{tr}H_{\rm fin}\leq T$.  If $D>1$, then
\[
 \|\xi\|_X^2
    \geq\frac{\delta^2}{C_{\rm eval}}
               \left(\frac{D-1}{T}\right)^{D-1}
                         \sum_a|u_a|^2 .
 \tag{Q11m}
\]
To prove this, Cauchy--Binet gives
$\det H_{\rm fin}\geq\delta^2$.  The product of all but
the smallest eigenvalue is at most
$(T/(D-1))^{D-1}$ by the arithmetic--geometric mean inequality.
Thus its smallest eigenvalue is at least the numerator in
(Q11m) before division by $C_{\rm eval}$.  Apply (Q11f).
For $D=1$ the corresponding bound is
$\delta^2|u_1|^2/C_{\rm eval}$.

This is a finite matrix certificate in prescribed coordinates;
an absolute positive number also requires a lower bound for
$\sum_a|u_a|^2$ for the particular input.  The exact mass and
height/denominator arguments in B5 and B7 may provide the
definite data only under their stated arithmetic normalization
hypotheses.  Rationality of the masses alone proves neither
algebraicity of the entries $r_i^{\rm ev}(\xi_a)$ nor a height
bound for their determinants.

\paragraph{Meaning for the comparison problem.}
The new constant $C_{\rm eval}$ depends on local injective
charts, the exact coefficient frames and finite definite
masses.  These data replace an unknown bound for
$G_{\rm ball}$ in this lower-bound argument.  They still must
be supplied to obtain a numerical certificate, as must any
positive definite-period lower bound and the relative error
in (Q11j).  The argument does not compute those numbers from
signature alone, nor identify a single Bessel period with the
whole ball pairing.  It gives an explicit route from the
finite definite periods to the same geometric norm that
controls the mixed theta period.

\framework{This is a comparison of positive Hermitian forms
and specified fibre evaluations.  It fits the polarization
and correspondence discussion in Grothendieck's
\emph{cote44}, scans 7--8, while retaining the distinction
between fibrewise evaluation and the cohomological operations
discussed in \emph{cote16}, scans 43--46.  The analytic estimate
does not turn the evaluation matrix or the frame correction
of A5 into a rational algebraic correspondence.}


% Non-CM and integral-model part. Uses only standard LaTeX, amsmath/amssymb,
% and the common macros authorized in the main paper.
\section{Non-CM sources and algebraic extraction}

\subsection{THETA3-C01: the precise congruence $H^1$ obstruction}\label{C01}

\paragraph{Statement.}
Let $X$ be a smooth projective congruence unitary ball quotient in the
setup of Y.~Liu, \emph{Fourier--Jacobi cycles and arithmetic relative
trace formula}, with Hermitian rank $m\geq3$. If $B$ is a simple complex
abelian variety which is not of CM type, then
\[
 \Hom_{\mathrm{Hdg}}\bigl(H^1(B,\Q),H^1(X,\Q)\bigr)=0.
\]
The same conclusion holds for any character-generated $H^1$ subspace
which is a sum of images of CM Hodge structures. Sums, tensor products,
duals, and Hodge subquotients of such structures remain of CM type.

\paragraph{Proof.}
Liu's Corollary~4.20, pp.~54--55, gives the CM isogeny decomposition of
the whole Albanese in rank $m\geq3$. Consequently $H^1(X,\Q)$ has
commutative Mumford--Tate group. A subquotient of a CM Hodge structure
is CM, because its Hodge homomorphism factors through an image of the
same torus. A nonzero Hodge morphism from $H^1(B,\Q)$ is injective:
its kernel is a rational sub-Hodge structure, and simplicity follows
from the equivalence between abelian varieties up to isogeny and
polarizable rational Hodge structures of weight one. Such an
injection into a CM structure would force $B$ to be CM, a
contradiction. A finite collection of CM structures is controlled by
a product of tori; all its tensor constructions and subquotients are
therefore controlled by a torus as well. $\square$

This is a deduction from the cited CM decomposition, not a new proof
of that decomposition. Liu's Hermitian rank is $m$, so the ball has
dimension $m-1$. In rank two the cited statement concerns the
endoscopic Albanese part, not the entire Albanese. The attached
van Hoften--Sempliner paper gives the relevant rank-three result in
Proposition~3.1.7, Theorem~3.1.11, and \S3.1.12. Increasing the ball
dimension does not remove this $H^1$ obstruction. Neither statement
applies to arbitrary noncongruence subgroups or to all higher
cohomology of the ball quotient.

\paragraph{Framework anchor.}
The relevant realization is the Albanese, not an arbitrary higher-degree
Hodge summand.  Grothendieck's letter of 12 August 1964, cote~44,
scans~7--8, explicitly places $\operatorname{Pic}^0$ and
$\operatorname{Alb}^0$ in the correspondence and polarization framework.
The CM restriction above is the additional automorphic input, not a
statement extracted from that letter.

\subsection{THETA3-C02: a trace-zero selection rule on covers}\label{C02}

\paragraph{Statement.}
Let $f:X\to S$ be a finite \'etale cover of connected smooth
projective complex surfaces of degree $d$. Then
\[
 H^1(X,\Q)=f^*H^1(S,\Q)\oplus
       \ker\bigl(f_*:H^1(X,\Q)\to H^1(S,\Q)\bigr).
\]
If $u$ belongs to the second summand and $v_1,v_2,v_3\in H^1(S,\Q)$,
then
\[
 \int_Xu\smile f^*v_1\smile f^*v_2\smile f^*v_3=0.
\]
If $H^1(S,\Q)$ is CM, every simple non-CM sub-Hodge structure of
$H^1(X,\Q)$ belongs to this trace-zero summand.

\paragraph{Proof.}
The transfer satisfies $f_*f^*=d\,\mathrm{id}$. Therefore
$d^{-1}f^*f_*$ is a projector with image $f^*H^1(S,\Q)$ and kernel
$\ker f_*$. The projection formula gives
\[
 \int_Xu\smile f^*v_1\smile f^*v_2\smile f^*v_3
 =\int_S f_*u\smile v_1\smile v_2\smile v_3=0.
\]
The restriction of $f_*$ to a simple non-CM substructure vanishes by
the argument of C01. $\square$

Thus one new non-CM factor and three old theta factors cannot give
the desired nonzero four-form period. Two new factors, for example a
holomorphic form and its conjugate, are the first possible case.

\paragraph{Framework anchor.}
The maps here are the pullback and trace of an actual finite geometric
correspondence.  Their action on the Albanese realization fits cote~44,
scans~7--8.  The vanishing calculation is a projection-formula statement;
it is not an assertion that arbitrary projectors in a realization are
represented by algebraic cycles.

\subsection{THETA3-C03: a mixed non-CM and CM period}\label{C03}

\paragraph{Hypotheses and source.}
Let $S_0=\Gamma_0\backslash D$ be a compact principal congruence
ball surface in the exact setup of 032, \S4. Suppose a compact ball
surface $X_0$ commensurable with $S_0$ has a nonconstant morphism
$a_0:X_0\to E_0$ to a non-CM elliptic curve.
Schoen's \emph{An arithmetic ball quotient surface whose Albanese
variety is not of CM type}, pp.~133--134, supplies such a
commensurability class over $\Q(\zeta_5)/\Q(\sqrt5)$, with
$j(E_0)=2^{11}$. To use it, choose 032's Hermitian space to be the
particular rank-three space in Schoen's construction. Multiplication
of the form by an element of the totally real field with suitable
real signs gives signatures $(2,1)$ and $(3,0)$. Equality of CM
fields alone would not imply commensurability. The determinant image
of an arithmetic unitary subgroup is finite: it is commensurable
with an integral subgroup, and integral norm-one units in a CM field
form a finite group. A finite-index subgroup therefore lies in
$\SU(V)$. Keep the principal congruence surface as the source of the
theta construction.

\paragraph{Statement.}
Fix an allowed CM sign function $v$ of 032 with $v(1)=1$.
There exist a sufficiently deep principal cover $S$ of $S_0$, a
common finite \'etale cover $X$ of $X_0$ and $S$, and a theta
one-form $\beta$ on $S$ in the complex span of rational source maps
from $U(v)$ such that, for a nonzero invariant form $\alpha_0$ on
$E_0$,
\[
 \xi=a^*\alpha_0\wedge f^*\beta\ne0,
 \qquad \int_X\xi\wedge\overline\xi>0,
\]
where $a:X\to E_0$ and $f:X\to S$ are the induced maps. Moreover
there are a CM abelian variety $B(v)$, a morphism $b:X\to B(v)$,
and a form $\gamma\in H^0(B(v),\Omega^1)$ with
$a^*\alpha_0\wedge b^*\gamma\ne0$. The image of $(a,b)$ has
dimension two.

\paragraph{Proof.}
The pullback $\alpha$ of $a_0^*\alpha_0$ to $D$ is nonzero. Choose
$x\in D$ with $\alpha_x\ne0$. By 032 Proposition~4.5,
pp.~21--22, finitely many actual finite Schwartz inputs give theta
one-form values spanning $T_x^*D$. At least one resulting $\beta_x$
is independent of $\alpha_x$. These finitely many inputs descend to
a sufficiently deep principal cover $S$ of $S_0$. Refine it using
032 Theorem~5.1, p.~26, and Corollary~6.7, p.~45, so that $\beta$
has its rational CM-source expansion. Only now form $X$ from the
intersection of the torsion-free commensurable lattices defining
$X_0$ and $S$. The common finite cover preserves the nonzero wedge
at a lift of $x$.

On the projective surface $X$, a nonzero holomorphic two-form has
strictly positive integral of $\xi\wedge\overline\xi$. Indeed, in
local coordinates $\xi=g\,dz_1\wedge dz_2$, this is a nonnegative
real volume form and is positive wherever $g\ne0$.

The CM extraction writes $[\beta]$ as a finite complex sum of
$h_j(u_j)$, where $h_j:U(v)\to H^1(S,\Q)$ are rational Hodge maps
and $u_j$ lie in the selected $(1,0)$ line. Holomorphic one-forms on
a compact K\"ahler manifold are the unique holomorphic
representatives of their classes, so the same expansion holds for
the forms. At least one term satisfies
$\alpha\wedge f^*h_j(u_j)\ne0$. By the Albanese realization used
in 032 Lemma~2.2, a positive integer multiple of $h_j$ is induced
by an algebraic map $S\to B(v)$. Composing with $f$ gives $b$.
If $(a,b)$ had image dimension at most one, its differential would
have rank at most one on a dense open set, forcing this wedge to
vanish. Therefore its image has dimension two, and
\[
 (a,b)_*[X]\in
 \CH^{\dim B(v)-1}(E_0\times B(v))_{\Q}
\]
is an algebraic image-surface cycle with the stated nonzero pairing.
$\square$

This construction uses two new factors, $\alpha$ and
$\overline\alpha$, and two CM theta factors. It does not realize
the new forms as adelic character theta lifts on a noncongruence
subgroup. It also does not identify a previously unknown Hodge class
as algebraic; C05 explains the limited Hodge content of this target.

\paragraph{Framework anchor.}
The surface image is an actual algebraic correspondence, and the nonzero
fourfold period is its Betti pairing.  Cote~44, scans~7--8, is the
correspondence/polarization anchor for that passage.  The extra non-CM
Albanese factor and the transverse theta form are supplied by the two
modern source constructions, not by an archival universality statement.

\subsection{THETA3-C04: the exact projector and rank gate}\label{C04}

\paragraph{Statement.}
Let $A$ be smooth projective, and let
$T\subset H^{2p}(A,\Q)(p)$ be a $d$-dimensional space of rational
Hodge classes. Suppose an algebraic correspondence acts as a
projector $e$ onto $T$. For algebraic cycles $Z_1,\ldots,Z_r$ of
codimension $p$, a basis $t_1,\ldots,t_d$ of $T$, and a
nondegenerate rational bilinear form $q_T$ on $T$, put
\[
 M_{ij}=q_T\bigl(e\,\operatorname{cl}(Z_i),t_j\bigr).
\]
If $\operatorname{rank}M=d$, every element of $T$ is algebraic.
Alternatively, if an algebra $\mathcal D$ of algebraic
correspondences preserves $T$, it suffices that the
$\mathcal D$-orbit of $e\operatorname{cl}(Z)$ spans $T$ for one
cycle $Z$. If $T$ is a simple $\mathcal D$-module, a nonzero
$e\operatorname{cl}(Z)$ suffices.

\paragraph{Proof.}
Every $e\operatorname{cl}(Z_i)$ is algebraic. Nondegeneracy
identifies $T$ with its dual, so full row rank means that these
vectors span $T$. In the orbit version, the finite-dimensional
orbit span is generated by algebraic correspondences applied to an
algebraic cycle. A nonzero submodule of a simple module is the whole
module. $\square$

A single nonzero scalar pairing does not give this rank assertion.
For example, $\Q(1,0)\subset\Q^2$ pairs nontrivially with $(1,1)$
under the ordinary dot product, although $(1,1)$ is outside that
subspace. This is a linear-algebra counterexample to an inference,
not an example of a nonalgebraic Hodge class. The pure embedding-line
divisors of 032 Proposition~3.1 supply the necessary CM extraction
in that paper. For a general non-CM $H^1$, the polarization provides
a sum of rank-one tensors; it does not supply individual
holomorphic eigenline projectors as algebraic correspondences.

\paragraph{Framework anchor.}
The distinction between an inverse or projector on cohomology and one
induced by an algebraic correspondence is explicit in cote~16,
scan~44, conditions~(iv) and~(iv bis), with the discussion continuing on
scans~45--46.  This finding isolates the same existence requirement and
adds the finite-dimensional rank test needed for the specified target.

\subsection{THETA3-C05: one non-CM elliptic factor gives a known consequence}\label{C05}

\paragraph{Statement; no novelty claim.}
Let $E_0$ be a non-CM elliptic curve, and let $B$ be a CM abelian
variety for which the rational Hodge conjecture holds. Then every
Hodge class on $E_0\times B$ is algebraic. More precisely,
\[
 \Hdg^p(E_0\times B)
 =1\otimes\Hdg^p(B)
 \oplus[\mathrm{pt}]\otimes\Hdg^{p-1}(B),
\]
where $\Hdg^p(Y)$ denotes the rational $(0,0)$ classes in
$H^{2p}(Y,\Q)(p)$.

\paragraph{Proof.}
The classical Hodge-group facts are
$\operatorname{Hg}(E_0)=\operatorname{SL}(H^1(E_0,\Q))$ and
$\operatorname{Hg}(B)=T$ for a torus $T$. For the first statement,
see Moonen, \emph{Notes on Mumford--Tate groups}, \S2,
Example~2.6, p.~11. The connected Hodge group $H$ of the product
projects onto both factors. Its derived group projects onto
$\operatorname{SL}_2$ and trivially to $T$. As a subgroup of
$\operatorname{SL}_2\times1$ with full projection to
$\operatorname{SL}_2$, it equals that group. The remaining quotient
surjects onto $T$, whence $H=\operatorname{SL}_2\times T$.
The only $\operatorname{SL}_2$ invariants in
$H^*(E_0,\Q)=\Q\oplus H^1(E_0,\Q)\oplus\Q(-1)$ occur in degrees
zero and two. Taking invariants in the K\"unneth decomposition gives
the formula. The two kinds of class are represented by a pullback of
a cycle on $B$, or by a point on $E_0$ times such a cycle. $\square$

\paragraph{Framework anchor.}
This is a consequence inside the tensor category generated by the
weight-one realizations, with the elliptic degree-two factor represented
by a divisor.  It uses the Albanese/correspondence framework of cote~44,
scans~7--8; it does not supply the additional higher-degree algebraic
correspondences discussed in cote~16, scans~43--46.



\subsection{THETA3-C12: four non-CM factors and a period of 500}\label{C12}
Use Schoen's surface $Y$ and curves from pp.~133--134 of the cited paper. The two projections to the Fermat quintic $C_5$ are
$(y_{01}:y_{03}:y_{13})$ and $(y_{02}:y_{03}:y_{23})$.
Let $C_0$ be its degree-five quotient, with equation $y^2=1-x^5$, and let $C\to C_0$ be Schoen's etale double cover, obtained from the elliptic curve
\[
 F:\quad s^2=1+x+x^2+x^3+x^4.
\]
The source proves $j(F)=2048$ and that $F$ is non-CM. There is a degree-two map $C\to F$.

There is a connected etale degree-four cover $X\to Y$ and a generically finite map $(a_1,a_2):X\to F^2$ of degree~500. If
$\int_F i\alpha\wedge\bar\alpha=1$, then
\begin{equation}\label{eq:500}
 \int_X(a_1^*\alpha\wedge a_2^*\alpha)
 \wedge\overline{(a_1^*\alpha\wedge a_2^*\alpha)}=500.
\end{equation}
All four degree-one sources are non-CM. The map
$(a_1,a_2,a_1,a_2):X\to F^4$ pushes its fundamental cycle to
$500[\Delta_{13}\cap\Delta_{24}]$, a product of diagonal divisor classes.
\begin{proof}
On $y_{03}=1$ write $u=y_{13}$ and $v=y_{23}$. The function field of $C_5^2$ has equations
$y_{01}^5=1-u^5$, $y_{02}^5=1-v^5$. Adjoining $y_{12}$ with
$y_{12}^5=u^5-v^5$ gives the function field of $Y$. The radicand has valuation~1 at the generic point of the diagonal of $C_5^2$, since $u^5-v^5=(u-v)$ times a unit there. It is not a fifth power, so the Kummer extension has degree~5. Resolution leaves function fields unchanged. Hence the product of the two curve maps gives a morphism $Y\to C_0^2$ of generic degree $5\cdot5\cdot5=125$.

The connected Galois etale cover $C^2\to C_0^2$ has degree~4 and group $(\Z/2)^2$. Its function field is linearly disjoint from the degree-125 extension, since the intersection degree divides both~4 and~125. Therefore
$X=Y\times_{C_0^2}C^2$ is connected and etale of degree~4 over $Y$. It is smooth projective and remains an arithmetic ball quotient. Its generic degree over $C^2$ is~125, and the two degree-two maps to $F$ give generic degree $125\cdot2\cdot2=500$ over $F^2$.

For $\eta=\mathrm{pr}_1^*\alpha\wedge\mathrm{pr}_2^*\alpha$ the Koszul sign gives
$\int_{F^2}\eta\wedge\bar\eta=-(-i)^2=1$.
Pushforward of the fundamental class under a proper generically finite morphism multiplies it by its degree, proving \eqref{eq:500}. Postcomposing with $(p,q)\mapsto(p,q,p,q)$ proves the cycle formula. The image is the transverse intersection of the two diagonal divisors.
\end{proof}
\framework{This realizes a four-factor non-CM period by an algebraic surface correspondence, in the sense of cote~44, scans~7--8. Its class is already generated by divisors. The new forms arise on a noncongruence cover, rather than from the original adelic character lifts. Thus neither a new exceptional Hodge-class case nor an extension of the original congruence theta packet is asserted.}


\subsection{THETA3-C13: an explicit algebraic projector onto the Weil space}\label{C13}

\paragraph{Statement.}
Let $A$ be a complex abelian variety of dimension $2n$, with an action of
an imaginary quadratic field $E$ of Weil type: $H^1(A,\Q)$ has rank
$2n$ over $E$, and the two embeddings of $E$ each occur with multiplicity
$n$ on $H^{1,0}(A)$. Put
\[
 W_E(A)=\bigwedge_E^{2n}H^1(A,\Q)\ \subset\ H^{2n}(A,\Q).
\]
Here the inclusion denotes the usual rational two-dimensional Weil
subspace, whose complexification is the sum of the two extreme
eigenspaces. Then the projector onto $W_E(A)$ is represented by an
explicit rational algebraic correspondence on $A\times A$.
Consequently a single algebraic codimension-$n$ cycle $Z$ with
$e_W\cl(Z)\ne0$ suffices to prove that every Weil class of $A$ is
algebraic. More concretely, it suffices to find a smooth projective
$n$-fold $X$, a morphism $b:X\to A$, and $w\in W_E(A)_{\C}$ such that
\begin{equation}\label{C13-period}
 \int_X b^*w\ne0.
\end{equation}
This is an extraction criterion for a supplied cycle, not an assertion
that such a cycle has been constructed for every $A$.

\paragraph{Proof and explicit denominators.}
Let $\mathcal R=E\cap\operatorname{End}(A)$, an order in $E$.
Choose $0\ne b_0\in\mathcal R$ with $\bar b_0=-b_0$, and choose an
integer $M>2n|b_0|$. Set $a=M+b_0$. After choosing one embedding of
$E$ and replacing $b_0$ by its negative if necessary, write
$a=|a|e^{i\vartheta}$ with $0<\vartheta<1/(2n)$.
The $2n+1$ numbers
\[
 \lambda_j=a^j\bar a^{\,2n-j},\qquad 0\le j\le2n,
\]
are distinct: the ratio of two of them is
$e^{2i(j-k)\vartheta}$, whose argument is nonzero and has absolute
value less than $2<2\pi$. Moreover $a^{2n}\notin\R$.
The existence of $b_0$ follows by taking $c-\bar c$ for any nonreal
$c\in\mathcal R$; since $\bar c=\Tr_{E/\Q}(c)-c$ and the trace is an
integer, this difference still belongs to $\mathcal R$.

Let $V_\tau,V_{\bar\tau}$ be the two complex eigenspaces of $H^1(A)$.
The endomorphism $T=[a]^*$ acts on
\[
 H^{2n}(A,\C)=\bigoplus_{j=0}^{2n}
       \bigwedge^jV_\tau\otimes\bigwedge^{2n-j}V_{\bar\tau}
\]
by $\lambda_j$. The two extreme summands are
$W_E(A)_\C$. Define
\[
 f(t)=(t-a^{2n})(t-\bar a^{\,2n}),\qquad
 g(t)=\prod_{j=1}^{2n-1}(t-\lambda_j).
\]
Both are monic in $\Z[t]$: their roots are algebraic integers and
their root sets are stable under the two automorphisms of $E/\Q$.
They are coprime. Put $R=\operatorname{Res}(f,g)\ne0$.
The adjugate of the integral Sylvester matrix supplies
$u,v\in\Z[t]$ with $uf+vg=R$. Thus
\[
 e_W=R^{-1}v(T)g(T)
\]
acts as one on the two extreme summands and as zero on every other
summand of $H^{2n}(A)$. For a projector on the whole cohomology, compose
with the explicit middle K\"unneth projector
\[
 \pi_{2n}=
 \prod_{\substack{0\le k\le4n\\k\ne2n}}
   \frac{[N]^*-N^k\,\mathrm{id}}{N^{2n}-N^k},\qquad N\ge2.
\]
Indeed $[N]^*$ acts by $N^k$ on $H^k(A)$. Endomorphism pullbacks,
their compositions, and the identity are induced by algebraic
correspondences. The displayed polynomial is therefore a rational
algebraic correspondence. Its idempotence is asserted modulo
homological equivalence; no Chow-idempotence is needed here.
An explicit common denominator is $R$ times the product of the
nonzero integers $N^{2n}-N^k$ in the last display.

The two extreme summands have Hodge type $(n,n)$ by the Weil-type
condition. Their rational sum has dimension two, and $T$ on it has
the irreducible quadratic minimal polynomial $f$, because
$a^{2n}\notin\Q$. If $0\ne\alpha=e_W\cl(Z)$, then
$\alpha,T\alpha$ are linearly independent over $\Q$; a dependence
would give a rational eigenvalue of $T$ on this space. Both are
classes of rational algebraic cycles. They therefore span $W_E(A)$.

Finally the middle cup pairing pairs the $j$-summand only with the
$(2n-j)$-summand. In particular $W_E(A)$ is nondegenerate for this
pairing and is orthogonal to the kernel of $e_W$. The projection
formula gives
\[
 \int_X b^*w=\int_A b_*[X]\smile w
           =\int_A e_W\cl(b_*[X])\smile w.
\]
Thus \eqref{C13-period} makes this projected class nonzero, and the
preceding argument applies. $\square$

\framework{This is a concrete projector in the homological category
of abelian motives, followed by algebraic correspondences on its
two-dimensional realization. It gives the explicit form of the
projector-and-rank step C04 in this Weil-type case. The correspondence
framework is Grothendieck, cote~16, scans~43--46, and
\emph{Standard conjectures on algebraic cycles}, pp.~193--196.
It uses no Hodge conjecture to construct the projector. The remaining
condition is a cycle with nonzero projection. In dimension four the
supplied Markman result already establishes algebraicity; this
criterion is not presented as a new proof of that theorem.}

\subsection{THETA3-C14: which curve-product constructions miss the Weil space}\label{C14}

\paragraph{Statement.}
Let $P=C_1\times\cdots\times C_r$ be a product of smooth projective
connected complex curves and let $b:P\to A$ be a morphism to an abelian
variety. The cohomology class $b_*[P]$ belongs to the subalgebra
generated by rational divisor classes on $A$ (and is zero if the
image has dimension less than $r$). If $X$ is smooth projective,
$\pi:X\to P$ is generically finite of degree $d$, and $c=b\circ\pi$, then
$c_*[X]=d\,b_*[P]$ has the same property.

Suppose now $r=n$, $A$ is as in C13, its polarization $\theta$ has
Rosati involution restricting to complex conjugation on $E$, and
$\operatorname{NS}(A)_\Q=\Q\theta$. Then
\[
 e_W\cl(c_*[X])=0.
\]
Thus a construction factoring through a product of $n$ curves cannot
satisfy C13's nonzero-period criterion for this target.

\paragraph{Proof.}
If some $C_i$ has genus zero, every morphism from that factor to an
abelian variety is constant, and the asserted pushforward is zero.
Otherwise choose base points, let $J_i=J(C_i)$, and let
$\iota_i:C_i\to J_i$ be the Abel maps. The universal property of
the Albanese variety expresses $b$ as a translation followed by a
homomorphism
\[
 q:J_1\times\cdots\times J_r\longrightarrow A
\]
composed with $\prod_i\iota_i$. Translation acts trivially on
cohomology. If $g_i=\dim J_i$ and $\theta_i$ is its principal
polarization, the Abel-curve class is
\[
 (\iota_i)_*[C_i]=\frac{\theta_i^{g_i-1}}{(g_i-1)!}.
\]
For completeness, this formula follows by pairing both sides with
$H^2(J_i)$. In a symplectic integral basis $e_j,f_j$ of $H^1(C_i)$,
the Abel map identifies $H^1(J_i)$ with $H^1(C_i)$,
$\theta_i=\sum_j e_j\wedge f_j$, and
$\int_{J_i}\prod_j e_j\wedge f_j=1$. Pairing either displayed class
with $e_j\wedge f_k$ gives $\delta_{jk}$, and pairing with
$e_j\wedge e_k$ or $f_j\wedge f_k$ gives zero. These elements span
$H^2(J_i)$, so Poincar\'e duality proves the formula.

Consequently the class of $P$ in $\prod_iJ_i$ is a product of divisor
classes. Milne, \emph{Lefschetz classes on abelian varieties} (1999),
Corollary~5.5, printed p.~25, states that direct image under a regular
map of abelian varieties preserves such classes. Applying it to $q$
proves the first assertion. The degree formula
$\pi_*[X]=d[P]$ proves the second.

For the final assertion, all codimension-$n$ divisor polynomials on
$A$ are multiples of $\theta^n$. Write
$H^1(A,\C)=V_\tau\oplus V_{\bar\tau}$. The Rosati condition gives
$[a]^*\theta=\Nm_{E/\Q}(a)\theta$ for the element $a$ chosen in C13.
The three eigenvalues on $\bigwedge^2H^1(A)$ are
$a^2,a\bar a,\bar a^2$, which are distinct. It follows that
$\theta\in V_\tau\otimes V_{\bar\tau}\subset\bigwedge^2H^1(A,\C)$.
Thus $\theta^n$ belongs to the $j=n$ summand in C13, whereas the
Weil projector retains only $j=0,2n$. Since $n\ge1$, these summands
are different and $e_W\theta^n=0$. $\square$

\framework{These are Lefschetz classes and algebraic direct images,
in the setting of Milne's Corollaries~5.3 and~5.5, pp.~25--26;
homological and numerical equivalence agree on the Lefschetz
subalgebra there. The Grothendieck anchors are cote~16,
scans~43--46, and the correspondence and positivity framework of
cote~44, scans~7--8. The finding explains the scope of C12's explicit
degree-500 non-CM construction: its diagonal-product cycle is
algebraic, but it supplies no exceptional Weil class. It does not
assert that every map from a finite cover of a curve product factors
through that product. New Albanese factors or a genuinely different
correspondence remain possible.}


\subsection{THETA3-C15: transferring curve factors to ball covers, and the countability boundary}
\label{C15}

We work over $\overline{\Q}\subset\C$ unless otherwise indicated.  Write
$C\rightsquigarrow D$ when a connected finite \'{e}tale cover of the smooth
projective curve $C$ admits a nonconstant morphism to the smooth projective
curve $D$.  This relation is directed.  It is not permissible to reverse an
arrow merely because the two curves occur in a correspondence.

\paragraph{Curve-to-ball transfer.}
Let $Y$ be a smooth compact arithmetic ball surface with a surjective morphism
$f:Y\to C_0$.  If $C_0\rightsquigarrow D$, there is a connected finite
\'{e}tale cover $X\to Y$ and a surjective morphism $X\to D$.  Consequently
$\operatorname{Alb}(X)\to J(D)$ is surjective.  Every abelian quotient of
$J(D)$ is therefore a quotient of the Albanese of an arithmetic ball cover.

\emph{Proof.}
Choose a connected finite \'{e}tale cover $C'\to C_0$ and a nonconstant
map $C'\to D$.  Any connected component $X$ of $Y\times_{C_0}C'$ is finite
\'{e}tale over $Y$.  Its image in the connected variety $Y$ is nonempty,
open, and closed, so this projection is surjective.  The image of $X$ in
$C'$ is closed by properness.  It cannot be a point, since its image in
$C_0$ is all of $C_0$; thus it is $C'$.  Composition gives the asserted
surjection $h:X\to D$.  The universal covering of $X$ is again the ball,
and its lattice has finite index in that of $Y$, so arithmeticity persists.

The pullback of a nonzero holomorphic differential on $D$ is nonzero on
$X$: this can be checked at a point where $h$ is smooth.  Thus
$h^*:H^0(D,\Omega_D^1)\to H^0(X,\Omega_X^1)$ is injective.  The dual map
on the tangent spaces of Albanese varieties is surjective.  The resulting
homomorphism $\operatorname{Alb}(X)\to J(D)$ is therefore surjective.
This also proves the assertion about further quotients. \hfill$\square$

The relation $\rightsquigarrow$ is transitive.  Indeed, if
$C'\to C$ is finite \'{e}tale, $C'\to D$ is nonconstant, and
$D'\to D$ is finite \'{e}tale with $D'\to E$ nonconstant, take a connected
component of $C'\times_DD'$.  It is finite \'{e}tale over $C'$ and maps
surjectively to $D'$, hence to $E$.  This proves transitivity, including the
connectedness qualification, without identifying any of the maps as an
isomorphism.

\paragraph{The exact Bogomolov--Tschinkel direction.}
Let $C_0$ now be Schoen's curve $y^2=1-x^5$.  Schoen constructs the morphism
$Y\to C_0$ on pp.~133--134 of his 2014 paper.  Put
\[
 D_6:\quad w^6=z(1-z),
\]
with smooth projective normalization understood.  This curve has genus two.
Bogomolov--Tschinkel, \emph{Unramified correspondences}, Proposition~1.8,
p.~2, proves
\[
 C_0\rightsquigarrow D_6.
\]
The same statement is Proposition~1.5, p.~2, of their
\emph{On curve correspondences}.  Their statement applies to every
hyperelliptic curve of genus at least two over an algebraically closed field
of characteristic different from $2,3$.  It does not assert
$D_6\rightsquigarrow C_0$ or that $D_6$ is universal in characteristic zero.
Their universality theorem over $\overline{\mathbf F}_p$, $p\geq5$, is
Theorem~1.7 of the first paper and Theorem~1.4 of the second; that theorem
must not be transported to characteristic zero.

There is nevertheless a useful positive continuation.  In
\emph{Unramified correspondences}, Example~3.1, pp.~5--6, let
$E_*\to\mathbf P^1$ be the degree-three quotient of the elliptic curve
with $j=0$, branched over $0,1,\infty$.  The paper proves
$D_6\rightsquigarrow H$ for curves $H$ admitting a map to $\mathbf P^1$
with local ramification indices at most two and branch points contained in
the image of $E_*(\overline{\Q})_{\rm tors}$.  It also states density of
these families in the indicated Hurwitz spaces.  Transitivity and the
preceding transfer lemma give actual quotients $J(H)$ of Albaneses of
covers of the fixed arithmetic ball surface $Y$.  This is a composition
of arrows in their proved directions.  In C16 we give a direct version
starting from Schoen's non-CM elliptic curve, with a full ramification
proof, an explicit new factor, and infinitely many nonisogenous factors.

\paragraph{A precise conditional universality reduction.}
If $D_6$ is universal over $\overline{\Q}$ in the sense of
Bogomolov--Tschinkel, then every abelian variety over $\overline{\Q}$ is a
quotient of the Albanese of a connected finite \'{e}tale cover of $Y$.
To see this, let $A$ be an abelian variety.  There is a smooth projective
curve $H$ with a surjection $J(H)\to A$.  For $\dim A=1$ take $H=A$;
in larger dimension one may take a smooth sufficiently ample complete
intersection curve in $A$, whose map on $H^1$ is injective by weak
Lefschetz.  Universality gives $D_6\rightsquigarrow H$, and then
$C_0\rightsquigarrow H$.  Apply the transfer lemma and compose the two
surjective homomorphisms of abelian varieties.  The hypothesis here is a
characteristic-zero universality assertion, not a consequence of the
positive-characteristic theorem.  Even under that hypothesis, the conclusion
concerns $H^1$ source factors; it does not produce the higher-degree cycle
projection or spanning property required in C04 and C13.

\paragraph{Countability.}
For a fixed compact ball surface $Y$, only countably many simple complex
abelian isogeny classes occur in the Albaneses of its finite \'{e}tale
covers.  In particular, for each $g\geq1$, a very general simple complex
abelian variety of dimension $g$ is not such a factor.

\emph{Proof.}
The compact manifold $Y$ has finitely generated fundamental group.  For
each $n$, it has only finitely many subgroups of index $n$: each such
subgroup gives a transitive permutation action on $n$ letters, and the
images of a fixed finite generating set admit only finitely many choices
in the finite group $S_n$.  Thus the connected finite covers form a
countable set.  Each Albanese has finitely many simple isogeny factors,
by complete reducibility of abelian varieties.  A countable union of
finite sets remains countable.

For a fixed complex abelian variety, there are countably many lattices
commensurable with its period lattice.  Thus its isogeny class is
countable, and allowing polarizations still gives a countable set because
integral alternating forms on a fixed lattice are countable.  The
positive-dimensional moduli space of polarized abelian varieties of
dimension $g$ has uncountably many complex points; the countable set just
obtained is contained in a countable union of proper closed points.
The same exclusion holds on its simple locus.  Finally, a nonzero rational
Hodge map from the $H^1$ of a simple abelian variety to $H^1(X)$ would,
by the equivalence for polarizable Hodge structures of weight one, make
that variety an isogeny factor of $\operatorname{Alb}(X)$.  This gives the
claimed source obstruction. \hfill$\square$

The same proof applies to any countable collection of arithmetic towers.
In particular, the number-field Hermitian data and finite adelic levels
used in the present unitary constructions give such a collection: number
fields, finite matrices over them, and finite-level choices are countable.
This countability obstruction does not conflict with the conditional
statement over $\overline{\Q}$, where the set of algebraic varieties is
already countable.  Nor does countability preclude density; the explicit
families in the next subsection are countable and dense in the relevant
complex hyperelliptic moduli spaces.

\paragraph{Framework anchor.}
The transfer lemma uses actual curve maps and their Albanese realizations,
the weight-one case of the correspondence framework in Grothendieck's
letter of 12 August 1964, cote~44, scans~7--8.  It supplies source factors;
the existence of higher-degree algebraic projectors remains the separate
issue visible in cote~16, scans~44--46.

\subsection{THETA3-C16: infinitely many non-CM elliptic factors from torsion branch points}
\label{C16}

Schoen's construction supplies a connected curve $C$ with
\[
 C\longrightarrow C_0 \quad\hbox{finite \'{e}tale of degree }2,
 \qquad C\longrightarrow F \quad\hbox{of degree }2,
\]
where $C_0:y^2=1-x^5$ and
$F:s^2=x^4+x^3+x^2+x+1$; see his pp.~133--134.  The second double cover
is ramified.  Indeed, in the function field of $F$ it is obtained by
adjoining $\sqrt{1-x}$.  The two points over $x=1$ and the two points at
infinity have odd valuation of $1-x$, so they are branch points.
Choose one branch point $q_0$ as origin on $F$.  This does not change its
$j$-invariant, which Schoen computes as $2048$.

\paragraph{The torsion-branch transfer lemma.}
Let $E$ be any elliptic curve in characteristic zero, let $C\to E$ be a
connected double cover ramified at its origin, and put
$\pi:E\to E/\{\pm1\}\simeq\mathbf P^1$.  For an odd positive integer
$N$, define
\[
 S_N=\pi(E[N]),\qquad |S_N|=(N^2+1)/2.
\]
Suppose that $H$ is a smooth connected projective curve and
$\sigma:H\to\mathbf P^1$ is a finite morphism of degree $d$, all of whose
local ramification indices are at most two, with branch locus contained
in $S_N$.  Then a connected finite \'{e}tale cover $D\to C$ of degree
at most $dN^2$ admits a surjection $D\to H$ of degree at most four.

\emph{Proof.}
Form $C_N=C\times_{E,[N]}E$.  Its projection to $C$ is finite \'{e}tale
of degree $N^2$.  It is a smooth connected curve: as a double cover of
the second copy of $E$ it is still branched, at every point of $E[N]$,
and a split degree-two cover would have no branch points.  Its composite
map $a:C_N\to E\xrightarrow{\pi}\mathbf P^1$ has degree four.
Every local ramification index of $a$ over $S_N$ is even.  Namely, above
$P\in E[N]$ the map $C_N\to E$ has ramification index two; multiplication
by $N$ is \'{e}tale and carries $P$ to the specified branch point.
The further index of $\pi$ is either one or two.  Oddness of $N$ also
shows that the only fixed point of $[-1]$ in $E[N]$ is the origin,
which proves the displayed cardinality of $S_N$.

Normalize $C_N\times_{\mathbf P^1}H$ and choose a connected component
$D$.  Its projection to $C_N$ is \'{e}tale.  Here is the complete local
check.  Where $\sigma$ is unramified, base change is unramified.  At a
ramified point of $\sigma$, choose formal parameters so that its map is
$z=u^2$, after absorbing a unit into $u$.  The other map has the form
$z=v^{2e}$ after absorbing a unit into $v$, since its ramification index
is even and the residue field is algebraically closed of characteristic
zero.  The normalization of $u^2=v^{2e}$ is the disjoint union of the
two branches $u=\pm v^e$, each unramified over the $v$-disc.  This proves
the assertion also at points where several unramified sheets of $\sigma$
lie over the same branch value.  Therefore every connected component is
finite \'{e}tale over $C_N$, with degree at most $d$.  It maps onto $H$:
both maps in the original fiber product are finite and surjective, and
each component of its normalization has dimension one.  Its degree over
$H$ is at most the total base-change degree four.  Composing with
$C_N\to C$ proves both degree bounds. \hfill$\square$

Applied to Schoen's $C$, C15 gives a connected arithmetic ball cover
$X\to Y$ of degree at most $2dN^2$ whose Albanese surjects onto $J(H)$.
For any set $B\subset S_N$ of cardinality $2g+2$, its hyperelliptic
double cover $H_B\to\mathbf P^1$ satisfies the hypotheses, and hence
\[
 \deg(X/Y)\leq4N^2,\qquad
 \operatorname{Alb}(X)\twoheadrightarrow J(H_B).
\]
As $N$ ranges through odd integers, these $H_B$ are dense in the complex
hyperelliptic moduli space of each fixed genus $g\geq1$.  Indeed, odd-order
torsion is dense in the complex torus $E(\C)$, so its image in
$\mathbf P^1(\C)$ is dense.  Any finite ordered configuration of distinct
points can be approximated by distinct points in this image.  The finitely
many odd orders occurring in such an approximation divide a common odd
$N$.  Passing continuously from configurations to the moduli of their
double covers proves density.  The curves are all defined over
$\overline{\Q}$ when $E$ is.  Density by itself does not prove that a
specified countable family avoids the CM locus.  The following arguments
give that additional conclusion explicitly.

\paragraph{An explicit new non-CM isogeny class, with cover degree at most $36$.}
The elliptic curve
\[
 E:\quad y^2=x^3-30x-25
\]
has $j(E)=2048$ and is therefore isomorphic over $\overline{\Q}$ to
Schoen's $F$ with the chosen origin.  Its nonzero three-torsion points have
four distinct $x$-coordinates, the roots of
\[
 \phi(r)=r^4-60r^2-100r-300.
\]
For completeness, the usual duplication formula yields this equation
without requiring a division-polynomial table: the condition
$x(2P)=x(P)$, with $y(P)\neq0$, gives
$3x^4+6Ax^2+12Bx-A^2=0$ on $y^2=x^3+Ax+B$; it is exactly the
nonzero three-torsion condition because $2P=P$ is impossible for such a
point.  Put $A=-30$, $B=-25$ and divide by three.

Fix a root $r$ of $\phi$, omit it from the four finite branch points, and
include infinity.  The resulting elliptic curve $H_r$ is
\[
 y^2=x^3+rx^2+(r^2-60)x+(r^3-60r-100).
\]
Its cubic is $\phi(x)/(x-r)$, so it has three distinct roots.  After
shifting $x$ by $-r/3$, its short Weierstrass coefficients are
\[
 A_r=\frac23r^2-60,\qquad
 B_r=\frac{20}{27}r^3-40r-100.
\]
Substitution in $j=1728\cdot4A_r^3/(4A_r^3+27B_r^2)$ and reduction
modulo $\phi(r)$ give
\[
 j(H_r)=-\frac{512}{405}
       (r^3+24r^2-130r-1200).
\]
Each of the four values is a root of the primitive polynomial
\begin{align*}
 P(J)={}&19683J^4-40310784J^3-5096079360J^2\\
        &+37203006717952J+6412351813189632.
\end{align*}
Both identities are polynomial identities over $\Q$ after reduction by
$\phi$; the accompanying standard-library certificate
\nolinkurl{check_C16_torsion_family.py} verifies them using exact rational
arithmetic, rather than floating-point root approximations.

Here is an explicit irreducibility certificate.  The monic reduction of
$P$ modulo $11$ is
\[
 P_0(J)=J^4-2J^3-5J^2+5J-5.
\]
In $\mathbf F_{11}[J]/(P_0)$ one has
\[
 J^{121}=5J^3-3J^2-5J+4,
 \qquad J^{14641}=J,
\]
and in $\mathbf F_{11}[J]$ one has the exact identity
\begin{align*}
 1={}&(-5J^2-5J-3)P_0(J)\\
 &+(J^3+4J^2-5J+2)(5J^3-3J^2-6J+4).
\end{align*}
The identity $J^{14641}=J$ says that all irreducible factors of $P_0$
have degree dividing four; the Bezout identity excludes factors of degree
one or two.  Thus $P_0$ is irreducible, and so is $P$ over $\Q$.
Since $P$ is primitive with leading coefficient $19683=3^9\neq1$,
its roots are not algebraic integers.  Every $H_r$ is therefore non-CM.
Moreover, it is not isogenous to $E$.  Indeed, a cyclic isogeny gives a
monic integral modular equation in the target $j$-invariant, so integrality
of the source $j$ implies integrality of the target; any isogeny factors
into cyclic isogenies.  Since $j(E)=2048$ is integral, the assertion follows.
The same argument applies place by place.  The standard source statements
are Milne, \emph{Modular Functions and Modular Forms}, version~1.31,
Theorem~6.1 and its proof, pp.~88--90, the cyclic-isogeny interpretation
in Lemma~8.5 and the following discussion, p.~99, and CM integrality in
Corollaries~12.15--12.16, p.~125.

All four branch points of $H_r$ belong to $S_3$, so the torsion-branch
lemma and C15 give the claimed ball cover of degree at most
$4\cdot3^2=36$.  This provides a non-CM elliptic isogeny class different
from Schoen's original one inside the same arithmetic covering tower.
No assertion that the four $H_r$ are mutually nonisogenous is needed.

\paragraph{Infinitely many pairwise nonisogenous non-CM factors.}
There are infinitely many pairwise nonisogenous non-CM elliptic curves
$E_1,E_2,\ldots$ over $\overline{\Q}$, each a quotient of the Albanese of
a connected finite \'{e}tale cover of Schoen's fixed ball surface $Y$.
For every finite initial set, one common connected cover of $Y$ has
$E_1\times\cdots\times E_m$ as an Albanese quotient.

\emph{Proof.}
We first produce a usable ordinary prime outside any prescribed finite
set.  The discriminant of $E:y^2=x^3-30x-25$ is
$2^4 3^6 5^3$.  If $p\equiv1\pmod{20}$, then $p\geq41$, the reduction
is good, and all its two-torsion is rational over $\mathbf F_p$:
\[
 x^3-30x-25=(x+5)(x^2-5x-5),
\]
whose remaining roots are $(5\pm3\sqrt5)/2$.  The field $\mathbf F_p$
contains the fifth roots of unity, and hence $\sqrt5$, because
$p\equiv1\pmod5$.  Consequently $4$ divides $\#E(\mathbf F_p)$ and
\[
 a_p=p+1-\#E(\mathbf F_p)\equiv2\pmod4.
\]
The Hasse bound gives $|a_p|\leq2\sqrt p<p$, so $a_p\not\equiv0\pmod p$.
This implies ordinary reduction.  One can check the last implication
directly: if $\varphi$ is Frobenius, its dual satisfies
$\widehat\varphi=[a_p]-\varphi$ and has degree $p$; its differential is
the nonzero scalar $a_p$.  It is therefore separable of degree $p$.
Since $[p]=\widehat\varphi\varphi$ and Frobenius is bijective on geometric
points, $\overline E[p](\overline{\mathbf F}_p)$ has $p$ points.
The facts used here are explicitly stated and proved in Stoll,
\emph{Elliptic Curves}, summer~2023 notes, Lemma~11.16, p.~42,
Theorem~12.1, p.~44, and Theorem~13.2, pp.~48--49.

There are infinitely many primes $p\equiv1\pmod{20}$, with an elementary
proof sufficient here.  Given a finite set of primes, choose a positive
integer $M$ divisible by $20$ and by every prime in that set.  Any prime
$p$ dividing
\[
 \Phi_{20}(M)=M^8-M^6+M^4-M^2+1>1
\]
is outside the set and is different from $2,5$.  The identity
$(M^2+1)\Phi_{20}(M)=M^{10}+1$ implies $M^{10}=-1$ modulo $p$.
The order of $M$ modulo $p$ is therefore four or twenty.  In the former
case $M^2=-1$, which gives $\Phi_{20}(M)=5$ modulo $p$, a contradiction.
Thus the order is twenty and $p\equiv1\pmod{20}$.

Fix such a prime and an embedding $\overline{\Q}\hookrightarrow
\overline{\Q}_p$.  The reduction map on $p$-torsion has a kernel of
cardinality $p$ and image $\overline E[p](\overline{\mathbf F}_p)$ of
cardinality $p$.  We spell out why no lifting assumption is hidden here.
The finite flat group scheme $E[p]$ over a good-reduction discrete
valuation ring has rank $p^2$.  After a finite unramified extension its
$p$ distinct special-fiber points are rational.  Their idempotent
decomposition lifts over this henselian ring.  Each component has rank
$p$, since translations make all special-fiber local lengths equal.
Its characteristic-zero generic fiber is \'{e}tale with $p$ geometric
points.  Every special point therefore lifts, and exactly $p$ generic
points reduce to the origin.

Choose a nonzero $P\in E[p]$ reducing to the origin, and choose
$Q,R\in E[p]$ whose reductions are nonzero and distinct modulo sign.
There are $(p-1)/2\geq2$ available nonzero sign classes.  The coordinates
satisfy
\[
 v_p(x(P))<0,\qquad
 v_p(x(Q)),v_p(x(R))\geq0,\qquad
 v_p(x(R)-x(Q))=0.
\]
For the first inequality, the integral local parameter $t=-x/y$ at the
origin gives $x=t^{-2}$ times an integral unit, and $v_p(t(P))>0$.
The others follow by reduction at finite distinct $x$-coordinates.

Let $E'_p$ be the double cover of $\mathbf P^1$ branched at
$\infty,x(P),x(Q),x(R)$.  Its branch cross ratio can be chosen as
\[
 \lambda=\frac{x(R)-x(Q)}{x(P)-x(Q)},\qquad v_p(\lambda)>0.
\]
The Legendre formula gives
\[
 j(E'_p)=256\frac{(1-\lambda+\lambda^2)^3}
                         {\lambda^2(1-\lambda)^2},
 \qquad v_p(j(E'_p))=-2v_p(\lambda)<0.
\]
This is a non-CM elliptic curve, by CM integrality.  All its branch points
lie in $S_p$, so the torsion-branch lemma gives a ball cover of degree at
most $4p^2$ whose Albanese surjects onto it.

Construct the sequence inductively.  After finitely many $E_i$ have been
chosen, the set of rational primes at which any conjugate of one of their
$j$-invariants is nonintegral is finite: each $j(E_i)$ is an algebraic
number.  Choose the next prime $p\equiv1\pmod{20}$ outside that set and
apply the construction.  All earlier $j(E_i)$ are integral at every
place above this new $p$, whereas $j(E'_p)$ is nonintegral at the chosen
place.  The local form of the monic modular-equation argument therefore
excludes an isogeny to every earlier $E_i$.  It also excludes an isogeny
to Schoen's original integral-$j$ elliptic curve.  This proves the infinite
assertion.

Finally, any finite list of the resulting connected finite \'{e}tale
ball covers has a common connected finite \'{e}tale cover, obtained by
taking a connected component of their fiber product over $Y$.  Projection
to each connected factor is surjective.  The maps to the $E_i$ persist,
and their distinct simple isogeny classes imply that the direct sum of
their $H^1$ injects into the $H^1$ of this common cover.  By complete
reducibility, the product $E_1\times\cdots\times E_m$ is an Albanese
quotient. \hfill$\square$

\paragraph{Nonzero periods with two of the new factors.}
The two independent curve projections of C12 give a generically finite
map $Y\to C_0^2$ of degree $125$.  For two curves $H_1,H_2$ furnished by
the torsion-branch lemma, choose the corresponding connected covers
$D_i\to C_0$; each $D_i\to H_i$ has degree at most four.  A connected
component $X$ of
\[
 Y\times_{C_0^2}(D_1\times D_2)
\]
is a finite \'{e}tale arithmetic ball cover of $Y$ and maps dominantly to
$H_1\times H_2$, with degree
\[
 1\leq d\leq125\cdot4\cdot4=2000.
\]
Indeed, its degree over $D_1\times D_2$ is at most the total generic
base-change degree $125$, and the remaining product map has degree at
most $16$.  If the $H_i$ are elliptic, choose differentials $\alpha_i$
normalized by $i\int_{H_i}\alpha_i\wedge\overline\alpha_i=1$.
Writing $a_i:X\to H_i$ for the two maps, the projection formula yields
\[
 \int_X a_1^*\alpha_1\wedge a_2^*\alpha_2
       \wedge\overline{a_1^*\alpha_1}
       \wedge\overline{a_2^*\alpha_2}=d>0.
\]
Thus one can obtain fourfold periods with factors from infinitely many
non-CM isogeny classes, with a positive integral degree in each instance.
The map to $H_1\times H_2\times H_1\times H_2$ repeats the two
projections; its image cycle is $d$ times the product of the two diagonals.
These are divisor-generated classes.  The construction substantially
enlarges the non-CM $H^1$ sources, but C14 still prevents these particular
product-of-curves image cycles from detecting the generic exceptional
Weil component.  The geometric factors constructed here are not supplied
by the CM character packets in the congruence $H^1$ statement of C01.

\paragraph{Framework anchor.}
The new Albanese quotients and image surfaces are actual algebraic
correspondences, as required by the realization/polarization viewpoint
of cote~44, scans~7--8.  Their explicit divisor-generated image classes
must be distinguished from the additional algebraic inverse and projector
classes discussed in cote~16, scans~43--46.  The infinite family enlarges
the available non-CM realizations without resolving that higher-degree
existence question.


\subsection{THETA3-C20: a simple non-CM abelian surface in one ball covering tower}\label{C20}

The elliptic factors in C16 can be replaced by a geometrically
simple abelian surface with no extra endomorphisms. This supplies a
non-CM simple weight-one source of dimension two, rather than a
product of the previously constructed elliptic sources.

\paragraph{Statement.}
Let
\[
 E:\ y^2=x^3-30x-25,\qquad
 \phi(x)=x^4-60x^2-100x-300.
\]
For any nonzero point $P\in E[7](\overline{\Q})$, put $a=x(P)$
and let $H_a$ be the smooth projective curve with affine equation
\[
 H_a:\quad y^2=(x-a)\phi(x).
 \tag{C20a}
\]
Then $H_a$ has genus two and
\[
 \operatorname{End}_{\overline{\Q}}J(H_a)=\Z.
 \tag{C20b}
\]
In particular $J(H_a)$ is geometrically simple and is not CM.
For Schoen's fixed arithmetic ball surface $Y$ in C15--C16, there
is a connected finite \'etale cover $X\to Y$ with
\[
 \deg(X/Y)\leq1764,\qquad
 \operatorname{Alb}(X)\twoheadrightarrow J(H_a).
 \tag{C20c}
\]
The surjection is understood over $\overline{\Q}$ and up to the
harmless choice of origins in the Albanese construction; its
associated $H^1$ inclusion is a rational Hodge morphism.

\paragraph{The division polynomial and its conjugate reductions.}
The four roots of $\phi$ are the nonzero three-torsion
$x$-coordinates, by C16. Set
\begin{align*}
 u&=x^3-30x-25,\\
 b&=x^6-150x^4-500x^3-4500x^2-3000x+22000,\\
 \psi_3&=3\phi,\qquad
 \psi_5=32u^2b-\psi_3^3,\qquad
 \psi_7=\psi_5\psi_3^3-128u^2b^3.
\end{align*}
The elementary division-polynomial recursion
$\psi_{2m+1}=\psi_{m+2}\psi_m^3-
\psi_{m-1}\psi_{m+1}^3$, with $\psi_2=2y$ and
$\psi_4=4yb$, gives these formulas. The exact source is
Silverman, \emph{The Arithmetic of Elliptic Curves},
second edition (2009), Chapter~III, Exercise~3.7,
pp.~105--106, especially parts (b), (d) and (f).
Thus $\psi_7$ has degree
24, leading coefficient seven, and its roots are exactly the
$x$-coordinates of the nonzero seven-torsion points.
One may equally obtain the same equation by the addition law
and the equality $x(4P)=x(3P)$, excluding the smaller-order
denominators; the displayed recursion fixes the normalization.

Here is an explicit finite-field irreducibility certificate.
The monic reduction of $\psi_7$ modulo 19 is
\begin{align*}
 P_{19}(x)={}&x^{24}-9x^{22}-4x^{21}-9x^{20}
       +8x^{19}+6x^{18}-4x^{16}+5x^{15}\\
 &+8x^{14}-2x^{13}-4x^{12}+7x^{11}-7x^{10}
       -9x^8-4x^7+7x^6\\
 &+4x^5+2x^4-4x^3+8x^2-7.
\end{align*}
Successive nineteenth powers modulo this polynomial give
\[
 x^{19^{24}}\equiv x,\qquad
 \gcd(P_{19},x^{19^8}-x)=
 \gcd(P_{19},x^{19^{12}}-x)=1.
 \tag{C20d}
\]
For an explicit check of the large exponents, the ascending
coefficient lists of the two intermediate remainders are
\[
\begin{split}
 r_8={}&(15,13,2,16,7,18,2,6,14,16,7,8,\\
       &\hspace{14mm}10,3,17,11,17,17,6,18,18,1,12,16),\\
 r_{12}={}&(8,18,13,9,13,1,6,13,4,16,14,12,\\
       &\hspace{14mm}14,3,11,3,12,17,5,0,12,6,0,3).
\end{split}
\]
Euclid's algorithm applied to $P_{19}$ and $r_j-x$ gives
the two unit gcds in (C20d). The first equality says that all
irreducible factors have degree dividing 24. Every proper
divisor of 24 divides 8 or 12, so the two gcds exclude all
proper factors. Hence $\psi_7$ is irreducible over $\Q$.
The independent integer-polynomial certificate
\path{documentation/check_C20_independent.py} implements this
recursion, modular powering and Euclidean division with exact
integer arithmetic. It does not use a numerical factorization.

Modulo 37, $\psi_7$ has the simple roots $7,20,35$.
In particular $K=\Q(a)$ has embeddings into $\Q_{37}$ taking
$a$ to Hensel lifts of 7 and 20. These give two places of the
same field $K$, with residue field $\F_{37}$, and two good
reductions of the same curve (C20a). No choice of unrelated
characteristic-zero curves is being compared.

\paragraph{The two point counts and absolute simplicity.}
For the two reductions, direct enumeration gives
\[
\begin{array}{c|r|r|l}
 a\bmod37 & \#H(\F_{37}) & \#H(\F_{37^2})
                    &\text{Frobenius polynomial}\\ \hline
 7 &35&1297&T^4-3T^3-32T^2-111T+1369\\
 20&37&1361&T^4-T^3-4T^2-37T+1369 .
\end{array}
 \tag{C20e}
\]
The polynomials $(x-7)\phi(x)$ and $(x-20)\phi(x)$ are
squarefree modulo 37, so these are smooth genus-two reductions.
The counts include the one point at infinity of the odd-degree
hyperelliptic equation. To specify completely the finite
calculation, use
$\F_{37^2}=\F_{37}[w]/(w^2-2)$; 2 is a nonsquare modulo 37.
For $z=r+sw$ evaluate $(z-a)\phi(z)$ by polynomial arithmetic
and count its square roots in this field. The sum of these
counts over the $37^2$ choices of $(r,s)$, plus one, gives the
second count. An independent implementation uses
$\chi_{\F_{37^2}}(v)=
\chi_{\F_{37}}(\operatorname{Norm}v)$ for $v\ne0$.
The two implementations agree and their full outputs are
included with the paper.

If $A_1=37+1-\#H(\F_{37})$, then
\[
 A_2=\frac{\#H(\F_{37^2})-37^2-1+A_1^2}{2},
 \qquad
 f(T)=T^4-A_1T^3+A_2T^2-37A_1T+37^2.
\]
This derives the two polynomials in (C20e) from the counts,
rather than assuming them. Both are irreducible modulo 5:
in each quartic quotient $T^{5^4}=T$, while
$\gcd(f,T^{5^2}-T)=1$. The respective ascending coefficient
lists of $T^{25}\bmod f$ are $(3,1,4,2)$ and $(1,2,3,2)$.
Thus both quartics are irreducible over $\Q$.
Their middle coefficients $-32$ and $-4$ are units at 37,
so the two Jacobians are ordinary and simple.

We use the exact finite-field criterion of Howe--Zhu,
\emph{On the existence of absolutely simple abelian varieties
of a given dimension over an arbitrary field},
arXiv:math/0002205v1, Theorem~6, p.~5, with proof pp.~5--6.
For $T^4+cT^3+dT^2+37cT+37^2$, its four exceptional
conditions are
$c=0$, $c^2=37+d$, $c^2=2d$ and $c^2=3d-111$.
None holds for $(c,d)=(-3,-32)$ or $(-1,-4)$.
Hence both reductions are absolutely simple.
By the same work, \S2 p.~3 and Proposition~3(2) with
Lemma~4 p.~4, their geometric endomorphism algebras are
exactly the quartic CM fields
\[
 K_1=\Q[T]/(T^4-3T^3-32T^2-111T+1369),\quad
 K_2=\Q[T]/(T^4-T^3-4T^2-37T+1369).
 \tag{C20f}
\]
Indeed absolute simplicity in the ordinary case keeps
$\Q(\pi^r)=\Q(\pi)$ over every finite residue extension.

\paragraph{Excluding every extra geometric endomorphism.}
The real subfield of a quartic ordinary Frobenius field is
generated by $s=\pi+37/\pi$. Its equation is
$s^2+c s+(d-74)=0$ for the coefficients just used.
Consequently
\[
 K_1^+=\Q(\sqrt{433}),\qquad
 K_2^+=\Q(\sqrt{313}).
 \tag{C20g}
\]
The polynomial discriminants, computed by the resultant with
the derivative, are
\[
 \operatorname{disc}(f_1)=
       2^4\,3^3\,37^2\,433^2,\qquad
 \operatorname{disc}(f_2)=
       2^4\,3^3\,11\,37^2\,313^2.
 \tag{C20h}
\]
Both are nonsquares. A quartic CM field containing an
imaginary quadratic subfield would be the compositum of
that field and its real quadratic subfield, hence a
biquadratic Galois extension. A quartic field containing
two different quadratic subfields is likewise
biquadratic. In a biquadratic extension the Galois
permutations of four embeddings are all even, so the
discriminant is a square. The discriminant of a defining
polynomial differs from the field discriminant by a
square. Thus each $K_i$ has exactly one quadratic
subfield, namely the one displayed in (C20g).
The two fields have no common nontrivial subfield up to
$\Q$-isomorphism: degree two is excluded by (C20g), and
an isomorphism in degree four would identify their unique
quadratic subfields.

Put $B=J(H_a)$. Every geometric endomorphism is defined
over some finite extension of $K$. At either of the two
places above 37 it extends to the good abelian scheme and
specializes injectively. The exact source is Conrad,
\emph{Gross--Zagier Revisited}, MSRI Publications~49
(2004), Theorem~2.1(2), printed p.~70.
After any such finite extension the geometric targets
remain the fields $K_1,K_2$, as proved above. Passing to
all geometric endomorphisms therefore embeds the finite
$\Q$-algebra
$\operatorname{End}_{\overline\Q}(B)\otimes\Q$
into each field. It is itself a field: a finite
$\Q$-subalgebra of a number field has no zero divisors,
and multiplication by a nonzero element is a bijection
of its finite-dimensional vector space. Its degree
divides four and is therefore one, two or four.
The last two degrees have just been excluded.
We obtain $\operatorname{End}_{\overline\Q}^0(B)=\Q$.
The integral endomorphism ring is an order in this
field containing $\Z$, hence is $\Z$. This proves
(C20b). A nontrivial abelian subvariety would give a
nontrivial rational idempotent by Poincar\'e
reducibility, so $B$ is geometrically simple.
A CM abelian surface would have a commutative
endomorphism algebra of dimension four, which is
impossible here.

\paragraph{Realization on the fixed ball covering tower.}
The branch points of (C20a) are infinity, the four
$x$-coordinates of $E[3]\setminus\{0\}$ modulo sign,
and $x(P)$ for a seven-torsion point. They are distinct:
an equality of the last coordinate with a three-torsion
coordinate would identify $P$ with a point of order
dividing three, up to sign. All six belong to
$S_{21}=\pi(E[21])$ of C16. The torsion-branch lemma
therefore gives a connected finite \'etale cover
$D\to C$ of degree at most $2\cdot21^2$ and a finite
map $D\to H_a$ of degree at most four. Since
$C\to C_0$ is \'etale of degree two, C15 supplies
a connected ball cover of degree at most
$4\cdot21^2=1764$ with Albanese quotient $J(H_a)$.
This proves (C20c). \hfill$\square$

The graph of the map to $H_a$ and the Abel--Jacobi
correspondence give the asserted rational $H^1$
realization. Thus this finding belongs to the
Jacobian-source and algebraic-correspondence programme
of Grothendieck's cote~44, scans~7--8; the
polarization and correspondence interpretation is
that of cote~16, scans~43--46. It supplies a new
simple source in this particular covering tower.
It does not identify a nonzero exceptional Weil
projection of an image cycle: that remains the
separate extraction problem in C13--C14.


\section{Integral models, theta frames and specialization}

\subsection{The precise unramified local normalization}

Let $k_0/\Q_p$ be finite with $p$ odd, residue cardinality $q$, and
uniformizer $\pi$. Let $k/k_0$ be unramified quadratic and
$\mathcal O=\mathcal O_k$. We use Li--Zhang's principally polarized
unitary Rapoport--Zink space $\mathcal N_n$ of signature $(1,n-1)$,
and its nonsplit rank-$n$ Hermitian space $\mathbb V_n$ of special
quasi-homomorphisms. For a full-rank lattice
$L=\mathcal O x_1+\cdots+\mathcal O x_n\subset\mathbb V_n$, define
\[
 \operatorname{Int}(L)=\chi\left(\mathcal N_n,
  \mathcal O_{\mathcal Z(x_1)}\otimes^{\mathbf L}\cdots
                    \otimes^{\mathbf L}\mathcal O_{\mathcal Z(x_n)}\right).
\]
This is a derived Euler characteristic, not a definition by the
ordinary scheme-theoretic intersection length. The source proves
basis independence. Li--Zhang, \emph{Kudla--Rapoport cycles and
derivatives of local densities}, Theorem~1.2.1, equivalently
Theorem~3.4.1, identifies it with
\[
 \operatorname{Int}(L)=-\left.\frac{d}{dX}
                           \operatorname{Den}(X,L)\right|_{X=1}.
\]
Their Theorem~3.5.1 gives the exact polynomial
\begin{align*}
 \operatorname{Den}(X,L)
 &=\sum_{L\subset N\subset N^\vee}
      X^{2\operatorname{length}_{\mathcal O}(N/L)}m(t(N);X),\\
 m(t;X)&=\prod_{i=0}^{t-1}\bigl(1-(-q)^iX\bigr).
\end{align*}
Here $t(N)$ is the number of positive elementary divisors. The
nonsplit space in this convention has odd determinant valuation.
The interpolation normalization is
\[
 \operatorname{Den}((-q)^{-k},L)
 =\frac{\operatorname{Den}(\langle1\rangle^{n+k},L)}
        {\operatorname{Den}(\langle1\rangle^{n+k},\langle1\rangle^n)}.
\]
The identity $\operatorname{Int}=-\operatorname{Den}'(1)$ has no
omitted logarithm. Conversion to a Whittaker derivative introduces
$\log(q^2)$ and explicit unramified factors as in the source's
Remark~3.4.2. The following calculations use this precise
normalization. Author-PDF locators are \S1.2, pp.~3--4;
Theorem~3.4.1, p.~17; Theorem~3.5.1, p.~18; and
Corollary~3.5.3, p.~19.

\subsection{THETA3-C06: an explicit negative derived intersection}\label{C06}

\paragraph{Statement.}
For $n=3$, let $L$ have Hermitian Gram matrix $\pi I_3$.
Then in the specified nonsplit local model,
\[
 \operatorname{Den}(X,L)=(1-X)\bigl(1+(q-q^2)X+X^2\bigr),
 \qquad \operatorname{Int}(L)=q+2-q^2<0.
\]
In particular the value at $q=3$ is $-4$.

\paragraph{Proof.}
The determinant valuation is three, hence odd, so the ambient
Hermitian space is the required nonsplit $\mathbb V_3$. Every
integral overlattice $N\supset L$ has
\[
 \operatorname{val}(N)=3-2\operatorname{length}_{\mathcal O}(N/L)
 \geq0.
\]
Thus its index length is zero or one. The length-zero term is $L$,
of type three. Since $L^\vee=\pi^{-1}L$, a length-one overlattice
has the form $N=L+\mathcal O\pi^{-1}v$, where $v\bmod\pi$ spans
a line in $L/\pi L$. Its pairings with $L$ are integral. Its
remaining integrality condition is
$h(v,v)\in\pi^2\mathcal O$, equivalently that the residual line
is isotropic for the nondegenerate Hermitian form $I_3$ over
$\mathbb F_{q^2}/\mathbb F_q$. Distinct residual lines give
distinct overlattices. Such an $N$ has determinant valuation one,
so its elementary divisors are $(0,0,1)$ and $t(N)=1$.

Here is a direct count of these lines. Let $a_r$ count vectors in
$\mathbb F_{q^2}^r$ whose coordinate norms sum to zero. The norm
map has one preimage of zero and $q+1$ preimages of each nonzero
element, so
\[
 a_1=1,\qquad
 a_r=(q+1)q^{2r-2}-q a_{r-1}.
\]
Consequently $a_2=q^3+q^2-q$ and $a_3=q^5-q^3+q^2$. Dividing the
nonzero isotropic-vector count by $q^2-1$ gives
\[
 \frac{q^5-q^3+q^2-1}{q^2-1}=q^3+1
\]
isotropic lines. The density formula therefore becomes
\begin{align*}
 \operatorname{Den}(X,L)
 &=(1-X)(1+qX)(1-q^2X)+(q^3+1)X^2(1-X)\\
 &=(1-X)\bigl(1+(q-q^2)X+X^2\bigr).
\end{align*}
Differentiating and invoking Li--Zhang gives
$\operatorname{Int}(L)=q+2-q^2$. Since
$q^2-q-2=(q-2)(q+1)>0$ for odd $q\geq3$, it is negative.
$\square$

This is a specialization of the source theorem, not a claim of a
new local-density theorem. It is consistent with the independent
horizontal and vertical calculation in Li--Zhang's proof of
Lemma~6.2.1, pp.~35--36: at $a_3=1$ the two terms are
$(q+1)+(1-q^2)$. A separate finite-field script in this run counted
$28,126,344,1332$ lines for $q=3,5,7,11$, respectively, and obtained
intersection values $-4,-18,-40,-108$. These finite checks do not
replace the general proof.

Thus a positive definite Hermitian norm cannot be identified with
an arbitrary local arithmetic intersection without controlling the
derived signs and the chosen cycle combination. Positivity of
representation densities at their interpolation arguments also
does not imply positivity of the derivative at one.

\paragraph{Framework anchor.}
The precise intersection-theoretic analogue is SGA~4$\frac12$,
\emph{Cycle}, Remark~2.3.9, p.~150, where the multiplicity is an
alternating Tor sum.  The negative example prevents identifying this
derived local number with the positive polarized pairing of cote~44,
scans~7--8, without an additional comparison theorem.

\subsection{THETA3-C07: a positive local sector with an exact bound}\label{C07}

\paragraph{Statement.}
Let $M$ be a self-dual Hermitian $\mathcal O$-lattice of rank $n-1$,
and let $x\perp M$ satisfy $h(x,x)=\pi^{2a+1}$ with $a\geq0$.
For $L=M\perp\mathcal O x$,
\[
 \operatorname{Den}(X,L)=(1-X)\sum_{j=0}^aX^{2j},
 \qquad \operatorname{Int}(L)=a+1\geq1.
\]

\paragraph{Proof.}
An integral overlattice $N\supset L$ splits as
$M\perp(N\cap M_k^\perp)$. Indeed, for $y\in N$, its orthogonal
projection to $M_k$ belongs to $M$ because $h(y,M)\subset\mathcal O$
and the pairing on $M$ is perfect. Subtract this projection.
The complementary lattice is $\pi^{-j}\mathcal O x$. Containment
and integrality give $j\geq0$ and $2a+1-2j\geq0$, hence
$0\leq j\leq a$. Each overlattice has index length $j$ and type
one, since its remaining positive elementary divisor is the odd
integer $2a+1-2j$. The density formula is the displayed sum.
Differentiation at one and the local theorem give $a+1$.
$\square$

The rank-one deformation interpretation after cancellation of the
self-dual part appears in Li--Zhang \S2.11 and Examples~2.4.3 and
3.5.2. The preceding proof does not require that interpretation.
It gives a normalized positive local certificate; comparison with
a different archimedean mixed period still requires a matching
theorem for its test data and factors.

\paragraph{Framework anchor.}
Here the exact local model permits a positive calculation, while the
cycle-class and intersection comparison still has the hypotheses of
SGA~4$\frac12$, \emph{Cycle}, Theorem~2.3.8 and Remark~2.3.9,
pp.~148--150.  Its connection to the correspondence/polarization problem
of cote~44, scans~7--8, requires the geometric identification in C17.

\subsection{THETA3-C08: vertical positivity does not detect a generic class}\label{C08}

\paragraph{Statement and counterexample.}
A positive ordinary local intersection on an integral model need
not detect a nonzero generic-fibre cycle. Let $R$ be a discrete
valuation ring with uniformizer $\pi$, and consider $\mathbb P^1_R$.
Its special fibre $F$ and a section $H$ meet with local intersection
ring $R/(\pi)$ and length one. Nevertheless $F$ has empty generic
fibre, so its restriction to the generic Chow group and cohomology
is zero. $\square$

An actual horizontal flat cycle of the required relative
codimension has an algebraic generic fibre by construction. The
remaining issue is identifying its class with the desired Hodge
tensor. The specialization criterion below makes the available
functoriality precise; it does not construct the missing cycle.

The basic-locus uniformization retains a geometric factor as well.
Vollaard--Wedhorn, \emph{The supersingular locus of the Shimura
variety of $\operatorname{GU}(1,n-1)$ II}, equation~(5.4.2), p.~34,
is
\[
 I(\Q)\backslash
 \bigl(\mathcal N_{\mathrm{red}}\times G(\mathbb A_f^p)/K^p\bigr)
 \simeq \mathcal M^{\mathrm{ss}}.
\]
It does not identify the whole supersingular locus with a finite
Shimura set. In rank three the Rapoport--Zink factor contributes
Fermat curves, as recalled in the source's introduction, p.~4.
Definite mass formulas can describe finite indexing data, while
local intersections depend on the retained geometry.

Li--Zhang's Theorem~4.2.1, p.~21, decomposes the horizontal part of
a rank-$(n-1)$ special cycle into primitive quasi-canonical lifting
cycles indexed by integral overlattices of type at most one.
Theorem~6.1.3, p.~34, compares the horizontal intersection function
with the corresponding horizontal density derivative. These are
more refined existing inputs than an unqualified positivity claim
for the total intersection.

\paragraph{Framework anchor.}
Cote~16, scans~44--46, explicitly discusses specialization in the
standard-conjecture framework.  The example above explains why a
vertical local intersection is not, by itself, a generic realization
class.  C17 states the actual relative-correspondence hypotheses under
which the cited SGA specialization comparison applies.

\subsection{THETA3-C09: finite lifting valuations and the infinite-product obstruction}\label{C09}

\paragraph{Statement.}
For $e_s=q^{s-1}(q+1)$,
\[
 \sum_{s\geq1}\frac1{e_s}=\frac{q}{q^2-1}.
\]
If $t_s\in\C_p$ have normalized valuations $v(t_s)=1/e_s$, the
infinite product $\prod_{s\geq1}t_s$ does not converge in $\C_p$.

\paragraph{Proof.}
The sum is geometric. The valuations of its partial products tend
to the displayed finite value, so the absolute values tend to a
nonzero number. A convergent sequence of these products would have
a nonzero limit, forcing its successive ratios $t_s$ to tend to
one. But $v(t_s)>0$ implies $|t_s|<1$ and $|1-t_s|=1$ for every
$s$. This is impossible. $\square$

The relevant primary source is Gross, \emph{On canonical and
quasi-canonical liftings}, Proposition~5.3(2),(3), manuscript p.~9,
proof pp.~9--11, in the original MPI/SFB~85-30 scan appended to
Gross--Zagier. These are PDF pages~77--78 of the retrieved combined
scan. In the unramified quadratic case it gives the degree
$[W_s:W]=e_s$ and valuation $v_{W_s}(t_s)=1$ of the formal modulus.
The original scan was visually checked. Li--Zhang \S4.1, p.~20,
uses the same extension degree $q^s(1+q^{-1})$.

Normalize the base valuation by $v(\pi)=1$. Then
$v(t_s)=1/e_s$. The canonical modulus $t_0\in\pi W$ satisfies
$v_{W_s}(t_s-t_0)=1$. In $W[[t]]$ the intersection with the canonical
section is therefore
\[
 W_s/(t_s-t_0),
\]
of ordinary length one over $W_s$ and over $W$, since the residue
fields agree. The branch degree $e_s$, its local intersection
length one, and the base-normalized coordinate valuation $1/e_s$
are different invariants. For finite $S$,
\[
 v\left(\prod_{s=1}^St_s\right)
 =\frac{q(1-q^{-S})}{q^2-1},\qquad
 v\left(\prod_{s=1}^S
             \operatorname{Norm}_{W_s/W}(t_s)\right)=S.
\]
The norm uses complete conjugate orbits and has valuation one at
each level. None of these finite identities is itself an identity
for a theta section. A comparison by divisors must specify the
section, its integral frame, and which branches or norms occur.
For the comparison with C11 over $\Q_p$, set this $q=p$; the
Lubin--Tate residue cardinality in C11 is instead $Q=p^2$.

\paragraph{Framework anchor.}
Gross's deformation divisors supply a concrete local intersection problem.
The exact normalization matters before applying the intersection and
cycle-class formalism of SGA~4$\frac12$, \emph{Cycle}, pp.~148--150.
The correspondence viewpoint in cote~44, scans~7--8, does not identify
these local divisors with a chosen analytic theta section.

\subsection{THETA3-C10: non-CM arithmetic theta and its realization}\label{C10}

\paragraph{Elementary observation.}
If an algebraic cycle $Z$ is homologically trivial, every ordinary
cohomological pairing with $Z$ is zero. A nonzero arithmetic height
may detect information about rational equivalence but does not
supply a nonzero ordinary Hodge class.

\paragraph{Proof.}
Homological triviality is precisely
$\operatorname{cl}_{\mathrm B}(Z)=0$. Its cup product with every
Betti class is zero, and the same statement holds in the compatible
\'etale realization. A height does not change this vanishing.
$\square$

There is an existing non-character extension relevant to this
boundary. Li--Liu, \emph{Chow groups and $L$-derivatives of
automorphic motives for unitary groups, II}, Example~1.10,
pp.~6--7, includes the representation attached to
$\operatorname{Sym}^{2r-1}$, $r\geq2$, of a non-CM elliptic curve,
under its modularity and local hypotheses. Good reduction at all
nonsplit finite places gives $S_\pi=\varnothing$; when
$r[F:\Q]$ is odd the stated case applies. Theorem~1.4 produces a
nonzero localized $\CH^r_{\mathrm{hom}}$ from a simple central zero,
with the technical Hypothesis~6.6 removable in this empty-$S_\pi$
case by Remark~1.9. Theorem~1.5's height formula also requires the
stated generating-series modularity hypothesis. These cited source
results are carried as background, not reproved here.

Thus non-CM input is available in arithmetic theta theory, but this
output belongs to homologically invisible cycles and
Abel--Jacobi or motivic extension data. A non-CM Hodge-class problem
still requires an algebraic correspondence with the relevant
nonzero realization and the rank condition C04.

\paragraph{Framework anchor.}
The output in this source is homologically trivial, so its nonvanishing
height does not give a nonzero class in the pure cohomological realization.
This is the precise distinction from the algebraic cohomology classes
and correspondence inverses discussed in cote~16, scans~43--46.

\subsection{THETA3-C11: the exact local derivative slope, with sharpness}\label{C11}

\paragraph{Statement.}
Let $p$ be odd, put $Q=p^2$, and choose $\pi\in\C_p$ with
$v_p(\pi)=1$. Set
\[
 \ell(T)=\sum_{r\geq0}\frac{T^{Q^r}}{\pi^r},\qquad
 E(Z)=\ell^{-1}(Z),\qquad D=\ell'(T)^{-1}\frac{d}{dT}.
\]
For every $f\in\mathcal O_{\C_p}[[T]]$ and every integer $a\geq1$,
\begin{align*}
 v_p(D^af(0))
 &\geq v_p(a!)-\frac{a-1}{p^2-1}\\
 &=\frac{ap}{p^2-1}+\frac1{p^2-1}
                     -\frac{s_p(a)}{p-1},
\end{align*}
where $s_p(a)$ is the base-$p$ digit sum and $v_p(0)=+\infty$.
Moreover $D^p=0$ after reduction modulo $p$.
The bound is attained infinitely often: for $f(T)=T$ and $a=Q^r$,
\[
 v_p(D^{Q^r}T(0))=\frac{(Q^r-1)p}{p^2-1}.
\]

\paragraph{Proof of the bound.}
Choose $\beta$ with $\beta^{Q-1}=\pi$, so
$v_p(\beta)=1/(Q-1)$. The scaled logarithm is
\[
 L(T)=\frac{\ell(\beta T)}\beta
 =\sum_{r\geq0}
     \pi^{(Q^r-1)/(Q-1)-r}T^{Q^r}.
\]
It is integral and has linear coefficient one, since for $r\geq1$
\[
 \frac{Q^r-1}{Q-1}-r=1+Q+\cdots+Q^{r-1}-r\geq0.
\]
Recursive formal inversion over any commutative ring therefore gives
\[
 U(T)=\frac{E(\beta T)}\beta
       \in T\mathcal O_{\C_p}[[T]],\qquad U'(0)=1.
\]
Integrality of $f$ implies
\[
 f(E(\beta T))-f(0)=f(\beta U(T))-f(0)
                         \in\beta T\mathcal O_{\C_p}[[T]].
\]
The chain rule in the logarithmic coordinate gives
\[
 \frac{d}{dZ}f(E(Z))=(Df)(E(Z)),\qquad
 f(E(Z))=\sum_{a\geq0}\frac{D^af(0)}{a!}Z^a.
\]
Its coefficient after substituting $Z=\beta T$ is
$\beta^aD^af(0)/a!$, of valuation at least $v_p(\beta)$ for
$a\geq1$. This proves the bound. Legendre's formula
$v_p(a!)=(a-s_p(a))/(p-1)$ gives the second expression. Since
$s_p(a)\leq(p-1)(1+\lfloor\log_pa\rfloor)$, the error from the
linear term is explicitly logarithmic.

\paragraph{Proof of nilpotence.}
We have
\[
 \ell'(T)=\sum_{r\geq0}(Q/\pi)^rT^{Q^r-1}\equiv1\pmod p.
\]
Thus $D$ reduces to $d/dT$ over $\mathcal O_{\C_p}/p$.
The $p$-th iterate of $d/dT$ is zero in characteristic $p$, because
its coefficient on every monomial contains a product of $p$
consecutive integers. This also holds for this nonreduced
coefficient ring.

\paragraph{Proof of sharpness.}
In $L(T)$ the exponent of $\pi$ is zero for $r=0,1$ and positive
for every $r\geq2$. Modulo the maximal ideal,
$\overline L(T)=T+T^Q$. In characteristic $p$ its inverse is
\[
 V(T)=\sum_{j\geq0}(-1)^jT^{Q^j},\qquad V(T)+V(T)^Q=T.
\]
Formal inversion commutes with reduction. Hence the
$T^{Q^r}$ coefficient of $U$ is a unit. Since
$E(Z)=\beta U(Z/\beta)$, the coefficient of $Z^{Q^r}$ in $E$
has valuation exactly $-(Q^r-1)/(Q-1)$. For $f=T$ and $a=Q^r$,
$D^aT(0)=a![Z^a]E(Z)$ and $s_p(a)=1$, giving
\[
 v_p(D^aT(0))=
 \frac{a-1}{p-1}-\frac{a-1}{p^2-1}
 =\frac{(a-1)p}{p^2-1}.
\]
This proves all assertions. $\square$

\paragraph{Precise relation to the CM local source.}
The coefficient is
\[
 \frac{p}{p^2-1}=\frac1p+\frac1{p^3}+\frac1{p^5}+\cdots,
\]
exactly the linear slope in Alp\"oge's step~2, p.~3.
For this arithmetic application, now require $\pi$ to be a
uniformizer of the unramified quadratic extension of $\Q_p$.
In that local-field setting the displayed logarithm is a permitted
coordinate for the height-two Lubin--Tate formal group.
The purely formal theorem above allows arbitrary $\pi\in\C_p$
of valuation one and does not assert an integral formal group for
that wider class of parameters.
Bannai--Kobayashi, \emph{Integral structures on $p$-adic Fourier
theory}, Proposition~3.5 and proof, pp.~531--532, first use
$\sum T^{q^r}/\pi^r$ and then treat general integral parameters.
An integral change of parameter with unit linear coefficient
preserves integral coefficient functions. Under a formal-group
isomorphism the normalized invariant derivations differ by a
constant unit, so the valuation statement is unchanged.
In their CM elliptic application, \S5, pp.~543--544, $p$ is odd
inert and the curve has good supersingular reduction; its formal
group is the height-two Lubin--Tate group for
$\pi=-\epsilon p$. The preceding proof allows this unit multiple.

The source provides actual integral functions to which the estimate
applies. On p.~546 it sets, for $b$ prime to $p$,
\[
 \wp_b(z)=\wp(z)-b^2\wp(bz),\qquad
 \phi(t)=\wp_b(\lambda(t))\in\mathcal O_K[[t]],
\]
and constructs an integral modified zeta function as well.
Proposition~5.5, pp.~546--547, identifies the invariant derivatives
and Fourier moments with modified Bernoulli--Hurwitz values.
Theorem~5.1, p.~544, states Katz's stronger floor-divisibility
result for the specified modified values, re-proved in that section.
These are existing CM elliptic applications of the same slope; no
new proof of Katz's full theorem is claimed here.

The slope also equals the valuation of a Lubin--Tate period.
Schneider--Teitelbaum, \emph{$p$-adic Fourier theory},
Lemma~3.4(b), gives
\[
 v_p(\Omega_{\mathrm{LT}})=\frac1{p-1}-\frac1{e(q-1)}.
\]
Putting $e=1$, $q=p^2$ gives $p/(p^2-1)$; the source's appendix
gives the period calculation. The $p$-adic period
$\Omega_{\mathrm{LT}}$ is distinct from the complex period
$\Omega_{\C}$ in the analytic normalization of Alp\"oge's theta
values.

\paragraph{Divided powers give a stronger, testable local identity.}
Katz, \emph{Divisibilities, congruences, and Cartier duality},
Identity~2.4 and Congruences~2.5, p.~670, work in the integral
algebra of invariant differential operators of the coordinate model
$[p](T)=pT+\alpha T^{Q}$, with $\alpha$ a unit and $Q=p^2$.
Define $D(n)$ by
\[
 f(X\mathbin\oplus Y)=\sum_{n\geq0}(D(n)f)(X)Y^n,
 \qquad D=D(1).
\]
Then $D(n)=D^n/n!$ for $n<Q$, and
\[
 D(n)^p\equiv
 \begin{cases}
 0,&p\nmid n,\\
 \alpha^{n/p}D(n/p),&p\mid n
 \end{cases}
 \pmod{p\operatorname{Diff}(G)}.
\]
In particular $D^p\equiv0$ and $D(p)^p\equiv\alpha D$ modulo $p$.
For completeness, the operator-valued translation series
$\Phi(U)=\sum D(n)U^n$ satisfies
$\Phi([p](U))=\Phi(U)^p$ by iterated translation.
Reduction modulo $p$ and comparison of the $U^{pn}$ coefficients
give the congruence. Over characteristic zero,
$\Phi(U)=\exp(D\lambda(U))$; since
$\lambda(U)\equiv U\pmod{U^Q}$, comparison below degree $Q$
gives the divided-power identity. This reconstructs the source's
elementary operator argument and exhibits more integral information
than nilpotence alone.

Katz's Main Theorem~3.11, p.~672, gives the stronger bound
\[
 v_p(D^nf(0))\geq\left\lfloor\frac{np}{p^2-1}\right\rfloor
\]
for integral $f$ satisfying the explicit torsion-trace equation
\[
 \sum_{u\in G[p]}f(T\mathbin\oplus u)=0.
\]
This theorem is cited as background, not reproved here.
In \S~IV, equations~(4.8)--(4.10), p.~676, Katz constructs such an
integral function on the supersingular CM elliptic formal group:
$f=\wp_b-[p]^*\wp_b$. Equations~(4.13)--(4.16), p.~677, identify
its derivatives with the modified Bernoulli--Hurwitz values.
Thus the stronger divisibility has a concrete functional equation
and actual integral elliptic examples. Matching a proposed theta
section to this function, or proving its own torsion-trace identity
in a compatible integral frame, remains a specific additional
comparison.

\paragraph{An integral trace-zero operator and the optimal uniform error one.}
In Katz's exact model $[p](T)=pT+\alpha T^Q$, over the ring $R$
of integers of a finite extension of $\Q_p$, the cited theorem
implies the stronger uniform estimate
\[
 v_p(D^af(0))\geq
       \left\lfloor\frac{ap}{p^2-1}\right\rfloor-1
       \qquad(f\in R[[T]],\ a\geq1).
\]
The constant one is optimal uniformly in $f$ and $a$.
Here is a complete reduction to Katz's theorem. Put
$A=R[[S]]$, $B=R[[T]]$, $S=[p](T)$. Weierstrass division for
\[
 T^Q+(p/\alpha)T-S/\alpha=0
\]
gives a free $A$-basis $1,T,\ldots,T^{Q-1}$ of $B$.
Newton identities yield
\[
 \Tr(1)=Q,\qquad \Tr(T^j)=0\ (1\leq j<Q-1),\qquad
 \Tr(T^{Q-1})=-(Q-1)p/\alpha.
\]
Thus $\Tr_{B/A}(B)=pA$. Define the integral functions
\[
 h=\Tr_{B/A}(f)/p,\qquad
 g=pf-[p]^*h.
\]
Their trace satisfies $\Tr(g)=p\Tr(f)-Qh=0$.
Over the split generic kernel this trace is the sum over translations
by $G[p]$, so $g$ satisfies exactly Katz's torsion-trace equation.
The equality holds integrally as well, since the formal rings are
$p$-torsion free. Katz's theorem therefore applies to $g$.
The invariant derivative satisfies $D[p]^*=p[p]^*D$, and $D$
preserves the integral formal functions. Hence
\[
 pD^af(0)=D^ag(0)+p^aD^ah(0)
\]
has valuation at least $\lfloor ap/(p^2-1)\rfloor$; the second
summand has valuation at least $a$, which is no smaller.
Dividing by $p$ proves the estimate. For $a=Q-1$ and
$f=T^{Q-1}$, one has $D^af(0)=a!$ and
\[
 v_p(a!)=p-1,
 \qquad\left\lfloor\frac{ap}{p^2-1}\right\rfloor=p.
\]
A smaller uniform loss is therefore impossible. $\square$

The argument initially has the exact DVR and coordinate hypotheses
of Katz. It includes the CM elliptic groups described in his \S~IV.
After integral base change from this model, it also holds for
$\mathcal O_{\C_p}$ coefficients: a fixed derivative uses only a
finite jet; approximate its integral coefficients by algebraic
integral coefficients in a finite extension and pass to the closed
valuation ball. Integral changes of formal parameter with unit
linear coefficient preserve the valuation of the normalized
invariant derivative. This is a deduction from Katz's theorem,
without a priority claim. It is not an automatic assertion about
an arbitrary analytic theta normalization.

\paragraph{The remaining theta-frame gate.}
A theta function is a section of a line bundle. Extension as an
invertible sheaf over the formal integral neighborhood gives some
integral frame, but differentiating its coefficient depends on the
frame or connection. The comparison still required is an equality
between the normalized analytic derivatives
$D^a\theta(0)/\Omega_{\C}^a$ and $D_G^af(0)$ for an integral
coefficient function in a specified compatible frame. Existence of
an arbitrary integral trivialization does not prove that equality;
a horizontal or cubical analytic frame need not be integral.

This is an actual obstruction to inferring the derivatives from
the divisor. On the formal characteristic-zero germ
$\C_p[[T]]$, take the trivial line bundle, section one, and frame
$e$. Replace the frame by
$e'=(1+p^{-M}T^a)e$. The two frames agree at the origin and the
section has the same empty divisor, but its coefficient becomes
$f'=(1+p^{-M}T^a)^{-1}$ and
\[
 D^af'(0)=-a!p^{-M}.
\]
Its valuation is arbitrarily negative as $M$ grows. The new frame
is a generic formal frame, not an integral one; no assertion about
invertibility on the whole rigid open unit disc is involved.
Even the exact integral Cartier divisor does not determine higher
jets: for $2\leq a<Q$, all
$f_N(T)=T(1+p^NT^{a-1})$ define the divisor $(T)$ and have first
derivative one, whereas
\[
 D^af_N(0)=a!p^N,
\]
because $E(Z)=Z+O(Z^Q)$. The uniform lower bound remains valid,
but the precise derivatives cannot be recovered from that divisor.

\paragraph{A concrete reduced-theta failure of integrality.}
The distinction occurs for a usual normalized reduced theta.
Take $E:y^2=x^3-x$, $\omega=dx/(2y)$, and $t=-x/y$.
At $p=7$ this is a good supersingular CM elliptic curve.
After replacing $y$ by $2y$, it is precisely the example in
Bannai--Kobayashi, arXiv:math/0610163v4, \S4, pp.~49--50.
Their Example~1.9, pp.~10--11, defines
$\theta(z)=\exp(-e_2^*z^2/2)\sigma(z)$ with $\theta'(0)=1$.
For the square lattice $e_2^*=0$, by its $i$-symmetry, also recorded
in Figure~1, p.~49. Thus $\theta=\sigma$ in this normalization.

The equation $t^2=x/(x^2-1)$ gives
\[
 x=\frac{1+\sqrt{1+4t^4}}{2t^2},\qquad y=-x/t,
 \qquad \omega=\frac{dt}{\sqrt{1+4t^4}}.
\]
Consequently
\[
 \lambda(t)=t-\frac25t^5+\frac23t^9+O(t^{13}).
\]
The Weierstrass equation $\wp''=6\wp^2-2$ gives
$\wp(z)=z^{-2}+z^2/5+z^6/75+O(z^{10})$.
Using $(\log\sigma)''=-\wp$ and the sigma normalization yields
\[
 \log(\sigma(z)/z)=-\frac{z^4}{60}-\frac{z^8}{4200}+O(z^{12}),
 \qquad
 \sigma(z)=z-\frac{z^5}{60}-\frac{z^9}{10080}+O(z^{13}).
\]
Exact composition therefore gives
\[
 \theta(\lambda(t))
 =t-\frac5{12}t^5+\frac{1411}{2016}t^9+O(t^{13}).
\]
Since $2016=7\cdot288$ and $7\nmid1411$, the $t^9$ coefficient
has valuation $-1$ at $7$. This actual normalized reduced theta is
not an integral coefficient function in the stated formal
parameter. The finite calculation makes no assertion about an
asymptotic radius of convergence and does not refute Alp\"oge's
accepted asymptotic theorem. It proves that analytic reducedness
alone does not provide the integral-frame hypothesis needed above.
A separate rational symbolic calculation reproduced these
coefficients; it is not a formal proof certificate.

\paragraph{Framework anchor.}
This finding concerns an integral function, invariant derivation, and
specified local frame.  It reaches the correspondence/polarization
problem of cote~44, scans~7--8, only after a section and integral-model
comparison such as C17 is supplied.  The explicit reduced-theta example
shows why that normalization datum cannot be suppressed.

\subsection{THETA3-C17: the cycle-moduli comparison and its specialization criterion}
\label{C17}

The specialization statement below is a standard consequence of the cited
EGA/SGA functorialities.  The contribution here is the explicit comparison
criterion and identification of the geometric data that an application
to the proposed bridge would have to supply.

\paragraph{The two cycles have different definitions and targets.}
The 032 construction starts with maps from a ball surface to fixed
CM abelian varieties, induced through the Albanese. For
$f:S\to A=\prod_iB_i$, it produces the image-surface cycle
$f_*[S]\in\CH_2(A)$. It is a cycle on the fixed abelian variety
$A$.

By contrast, Li--Zhang \S2.3, p.~11, defines a local special cycle
$\mathcal Z(L)$ by an extension condition on special
quasi-homomorphisms: after applying the framing quasi-isogenies,
each $x\in L$ must extend from the special fibre to an actual
$\mathcal O_k$-linear homomorphism of the indicated
$p$-divisible groups. Globally, \S\S13.2--13.3, p.~82, starts
with the unitary moduli object $(A_0,A)$, its actions,
polarizations, and level, and considers tuples of special
homomorphisms in the allowed prime-to-$p$ isogeny category with
prescribed Hermitian Gram matrix $T$ and prescribed finite test
function. This is a moduli condition on a unitary Shimura model.
It is not supplied by an arbitrary map from its surface to fixed
abelian varieties. In particular the two definitions neither have
the same ambient variety nor automatically the same codimension.

A literal comparison therefore requires a specified algebraic
correspondence from the relevant Shimura model to the fixed target
$A$, a theorem identifying its action on the selected special cycle
with the projected image-surface cycle, and integral extensions
compatible with the desired local model. This is an identification
gate, not an assertion that such a correspondence is impossible.
C04 still supplies the separate rank requirement on the target
Hodge summand.

\paragraph{Exact primary EGA and SGA locators.}
EGA~IV$_3$, Theorem~8.8.2, pp.~28--29, and Scholie~8.8.3,
pp.~33--34, descend finite-presentation schemes, morphisms, and
finite collections of relations through filtered limits.
EGA~IV$_2$, Theorem~6.9.1, p.~153, is generic flatness. These give
spreading of finite specified geometric data and flatness after
shrinking a suitable integral base. They do not produce a desired
cycle, and shrinking does not ensure good behavior at an arbitrarily
specified bad prime.

For specialization use the original SGA~4, tome~III,
Expos\'e~XVI, Theorem~2.1, p.~213, and Corollary~2.3, p.~215,
with locally constant prime-to-residue-characteristic coefficients.
The locally constant hypothesis is explicit in those statements;
one need not rely on the misprinted formulation of Corollary~2.2.
For cycle classes use SGA~4$\frac12$, \emph{La classe de cohomologie
associ\'ee \`a un cycle}: Definition~2.3.2, p.~145;
the perfect-weight construction~2.3.3, p.~146;
Lemma~2.3.7 and Theorem~2.3.8, p.~148; and
Remark~2.3.9, p.~150. Lemma~2.3.7 gives base-change compatibility.
Theorem~2.3.8 gives pullback and proper-intersection compatibility
under its fibrewise codimension and perfect-weight hypotheses.
Remark~2.3.9 explicitly uses alternating sums of Tor for proper
intersection multiplicities. These original scans were inspected;
none of these source theorems is claimed to have been reproved here.

\paragraph{A precise cohomological criterion.}
Let $R$ be a strictly henselian discrete valuation ring,
$\ell$ invertible in $R$, and $X,Y$ smooth proper over $R$.
Let $\Gamma\subset X\times_RY$ be a relative flat cycle of pure
codimension $c$ with the preceding cycle-class hypotheses; for
example its structure sheaf is a perfect weight on the regular
ambient product. Fix geometric generic and special fibres and the
specialization isomorphisms. Then the action of $\Gamma$ commutes
with specialization:
\[
 \operatorname{sp}_Y\circ(\Gamma_{\bar\eta})_*
       =(\Gamma_{\bar s})_*\circ\operatorname{sp}_X.
\]
The correspondence action has the usual degree shift
$2c-2\dim(X/R)$ and Tate twist $c-\dim(X/R)$.

\paragraph{Proof from the cited functorialities.}
Write the action as pullback to the product, cup product with
$\operatorname{cl}(\Gamma)$, and the proper smooth projection's
trace. The cited specialization theorem identifies the fibre
cohomologies. Pullback and cup product commute with base change;
Lemma~2.3.7 identifies the two restrictions of the cycle class;
and the trace functoriality in the construction~2.3.1--2.3.4
commutes with that base change. Combining the three operations
proves the identity. Flatness makes the fibre restriction of
$\mathcal O_\Gamma$ agree with its derived base change, so it is
the stipulated fibre cycle class. $\square$

Consequently, once actual correspondences and their relative cycle
classes are provided, their cohomological action and any certified
rank calculation can be transported through this comparison. The
criterion neither lifts arbitrary special-fibre cycles nor proves
that an arbitrary Hodge projector is algebraic. The same
cycle-class/intersection compatibility shows that homological
triviality implies numerical triviality: pair a zero cycle class
with each complementary algebraic class and apply the trace.
The converse is not supplied by EGA/SGA or by the local
intersection formula. Passing among a geometric correspondence,
its homological realization, and a desired Hodge tensor therefore
requires the explicit maps and rank statements described above.

\paragraph{Archival anchor.}
Cote~16, scan~43, states the complementary-dimension issue for algebraic
classes; scans~44--45 formulate cohomological inverse maps induced by
algebraic correspondences; scan~46 continues the discussion of
specialization and hyperplane reductions.  C17 concerns the functoriality
of already specified relative correspondences.  It does not establish the
existence of the conjectural correspondences in that letter.  Cote~44,
scans~7--8, supplies the related Albanese and polarization interpretation.


\subsection{THETA3-C18: tame theta ratios have an integral common frame}\label{C18}

Let $p\geq5$, let $R$ be the ring of integers of a finite extension of
$\Q_p$, and let $\mathcal E/R$ be the smooth elliptic curve
$y^2=4x^3-g_2x-g_3$ with $g_2,g_3\in R$ and unit discriminant.
Use its integral parameter $t=-2x/y$ and its formal logarithm
$\lambda(t)$ normalized by the invariant differential $dx/y$.
The letter $\mathcal E$ here denotes a curve, not a CM field.
Enlarge $R$ by a finite extension whenever the torsion points below
must be made rational.

For a compatible system of division points $\widetilde T=(T_n)$,
suppose $T_1$ is a nonzero torsion point of order $N$ prime to $p$.
Write $s_T(t)$ for its algebraic sigma series in
Katagiri--Kobayashi, \emph{p-adic properties of division polynomials
and algebraic sigma functions}, arXiv:2609.07024v1,
Definition~3.1, p.~9. Their Remark~3.2 on that page identifies it over
$\C$ with
\[
 s_T(t)=\sigma_{z_T}(\lambda(t)),\qquad
 \sigma_{z_T}(z)=
 e^{-\eta(z_T)(z+z_T/2)}\sigma(z+z_T).
\]
The torsion lift, differential, and normalization are part of this
notation. Theorem~2.7, pp.~7--8, supplies the algebraic value and
Taylor-coefficient statement. For an arbitrary compatible system,
one can choose a rational lattice lift realizing the finite datum
$T_{4N}$ on which these constructions depend.
The following is a precise integral consequence of their
construction, rather than an integrality assertion for the raw sigma
series.

\paragraph{Common-factor statement.}
There are units $u_T(t)\in R[[t]]^\times$ such that
\begin{equation}\label{C18-factor}
 s_T(t)=q(t)u_T(t),\qquad
 q(t)=\frac{\sigma(\lambda(t))}{t},\qquad
 u_T(0)=\sigma_{\rm alg}(\widetilde T)\in R^\times.
\end{equation}
The same $q(t)\in\operatorname{Frac}(R)[[t]]^\times$ works for every
such tame torsion point. Consequently, for two lists of the same
length,
\begin{equation}\label{C18-ratio}
 R_{\mathbf T,\mathbf U}(t):=
 \frac{\prod_{i=1}^r s_{T_i}(t)}
      {\prod_{i=1}^r s_{U_i}(t)}
       =\frac{\prod_i u_{T_i}(t)}{\prod_i u_{U_i}(t)}
       \in R[[t]]^\times.
\end{equation}
After pullback by a common prime-to-$p$ multiplication isogeny,
this is the completion of a rational function on
$\mathcal E$ that is a unit along the identity residue disc.
On that cover it is the ratio of sections of equal tensor powers
of the same line bundle; before pullback a torsion flat twist can
remain. This integral-unit conclusion is specific to equal tensor
degrees and prime-to-$p$ torsion. A weaker derivative estimate for
arbitrary integer tensor degree is proved below.

\paragraph{Proof, including the integral model.}
The rational function $g_T$ used in Definition~3.1 has
\[
 \operatorname{div}(g_T)=[N]^*(\tau_{T_1}^*(O)-(O)),\qquad
 (t g_T)(O)=N^{-1}\sigma_{\rm alg}(\widetilde T).
\]
By their Lemma~3.3(1), pp.~9--10, the latter sigma value is a
unit. Here is the resulting local integral calculation.
Because $N$ is prime to $p$, all the torsion sections in this
divisor are disjoint over $R$ after the finite extension. The
only one through the special-fiber origin is the simple pole
along $O$. Any vertical contribution to the divisor of $g_T$
is an integer multiple of the smooth irreducible special fiber.
In the regular local ring at its origin, cancellation of the
horizontal pole gives $t g_T=\varpi^a v$ with $v$ a unit and
$a\in\Z$. Its restriction to $O$ is a unit by the displayed
normalization. Hence $a=0$, and
\[
 h_T(t):=t g_T(t)\in R[[t]]^\times.
\]
This argument can equivalently be made directly with the
horizontal rational function in the proof of Lemma~3.3.

Multiplication by $N$ on the integral formal group has expansion
$[N](t)=Nt+\cdots$ with $N\in R^\times$. Its compositional inverse
$a_N(t)=[N^{-1}](t)$ belongs to $tR[[t]]$, with linear coefficient
$N^{-1}$. Therefore
\[
 u_T(t):=t g_T(a_N(t))
       =\frac{t}{a_N(t)}h_T(a_N(t))\in R[[t]]^\times,
 \qquad u_T(0)=N\,h_T(0)=\sigma_{\rm alg}(\widetilde T).
\]
Definition~3.1 now gives \eqref{C18-factor}, and cancellation
gives \eqref{C18-ratio}.

To verify the rational-function assertion, take a common multiple
$M$ of all the orders $N_i,N_i'$ in the two lists, still prime
to $p$. The formal group identity
$a_{N_i}([M](t))=[M/N_i](t)$ gives directly
\[
 R_{\mathbf T,\mathbf U}([M](t))=
 \frac{\prod_i[M/N_i]^*g_{T_i}(t)}
      {\prod_i[M/N_i']^*g_{U_i}(t)}.
\]
The common sigma factors cancel before this substitution.
The right side is an explicit rational function on $\mathcal E$
over the same finite torsion extension. Its completion is an
integral unit by \eqref{C18-ratio} and the integral multiplication
series. This proves the asserted rationality on the finite \'etale
isogeny cover without asserting descent to the original curve.
\hfill$\square$

\paragraph{Derivative consequence.}
Suppose the integral formal parameter has been identified by a
unit-linear integral change of coordinates with the exact height-two
model of C11, $[p](T)=pT+\alpha T^{p^2}$ with $\alpha\in R^\times$.
For its normalized invariant derivation $D$, C11 applies to
\eqref{C18-ratio} and gives, for every $a\geq1$,
\[
 v_p\bigl(D^a R_{\mathbf T,\mathbf U}(0)\bigr)
       \geq\left\lfloor\frac{ap}{p^2-1}\right\rfloor-1.
\]
The assertion survives multiplication by an integral unit and
unit-linear integral changes of parameter, with the invariant
differential normalization retained. This is an actual theta
quotient to which the integral-function estimate applies.

For a single translated theta section the same calculation gives
the exact connection conversion
\[
 q^{-1}D(qf)=Df+\frac{Dq}{q}f.
\]
It separates the integral function $u_T$ from the common connection
term $Dq/q$. The coefficient example in C11 shows why $q$ cannot
simply be declared integral. For a tensor quotient of total degree
zero this common factor cancels as in \eqref{C18-ratio}; for a
tensor of degree $d$ it remains as $q^d$. Thus the remaining
comparison with Alp\"oge's harmonic covariant theta derivatives
is the specified connection and its degree, rather than an
unspecified raw-theta integrality hypothesis. This statement does
not identify those archive derivatives with the degree-zero
quotients without their actual connection data.

\paragraph{The derivative of the common connection is integral and rational.}
Put $A=Dq/q$, with $D$ normalized by $dx/y$.  The classical identity
$(\log\sigma)''=-\wp$ and $x=\wp(z)$ give the exact identity
\begin{equation}\label{C18-connection-rational}
 B:=DA=D^2\log q=-x-D^2\log t
   =\frac{(Dt)^2-tD^2t-t^2x}{t^2}.
\end{equation}
This is a rational function on the elliptic curve, and
$B\in t^2R[[t]]$.  Here is a direct integral verification, which does
not assume that $q$ or $A$ is integral.  Put $u=t^2x$.  The Weierstrass
equation and $t=-2x/y$ give
\[
 4u^2(u-1)=g_2u t^4+g_3t^6.
\]
The implicit-function derivative at $(u,t)=(1,0)$ is $4\in R^\times$.
Thus $u\in1+t^4R[[t]]$, so $x=t^{-2}+t^2R[[t]]$.
Since $Dx=y$ and $Dy=6x^2-g_2/2$, direct differentiation gives
\[
 Dt=\frac{4x^3+g_2x+2g_3}{4x^3-g_2x-g_3}
    =1+\frac{g_2t^4}{2u}+\frac{3g_3t^6}{4u^2}
       \in1+t^4R[[t]].
\]
Consequently $D=(Dt)\,d/dt$ preserves $R[[t]]$ and
$D^2t\in t^3R[[t]]$.  The numerator in
\eqref{C18-connection-rational} belongs to $t^4R[[t]]$, proving the
claim.  An explicit expression, useful independently of the parameter,
is
\begin{align*}
 B&=9x-\frac{g_2}{2x}-\frac{g_3}{x^2}
       -\frac{(6x^2-g_2/2)^2}{y^2}\\
  &=\frac{-5g_2x^4-13g_3x^3+(g_2^2/4)x^2
                +(3g_2g_3/2)x+g_3^2}{x^2y^2}.
\end{align*}
Since $q\in1+t^4\operatorname{Frac}(R)[[t]]$, both
$\log q(0)$ and $D\log q(0)$ are zero.  Under the exact height-two
coordinate hypothesis already stated, C11 now gives
\[
 v_p\bigl(D^a\log q(0)\bigr)
   \geq\left\lfloor\frac{(a-2)p}{p^2-1}\right\rfloor-1
       \qquad(a\geq3),
\]
while $D^2\log q(0)=B(0)=0$.
The integral rational derivative need not have an integral primitive.
For example, C11's curve with $g_2=4,g_3=0$ gives
\[
 A=-\frac53t^3+\frac{11}{7}t^7+O(t^{11}),\qquad
 B=-5t^2+t^6+O(t^{10}).
\]
Thus at $p=7$ the connection term itself is not integral in this frame.

\paragraph{A weaker positive slope for every tame tensor degree.}
Retain the exact height-two coordinate hypothesis, so that the
operator congruence from C11 holds:
\[
 D^p\bigl(R[[t]]\bigr)\subset pR[[t]].
\]
For every integer $d$, every $v(t)\in R[[t]]$, and every $n\geq1$,
the series $f=q^dv$ satisfies
\begin{equation}\label{C18-weaker-slope}
 v_p\bigl(D^nf(0)\bigr)
 \geq\max\left\{0,\ v_p(n!)-\frac{np}{p^2-1}\right\}
 =\max\left\{0,\frac{n}{p^2-1}
                         -\frac{s_p(n)}{p-1}\right\}.
\end{equation}
Here $s_p(n)$ is the sum of the base-$p$ digits of $n$.
This applies in particular to every finite product or quotient of the
tame translated sigma series: its integer tensor degree is $d$, and
the remaining factor $v$ is an integral unit.  Degree zero still has
the stronger C11 estimate with slope $p/(p^2-1)$.
The slope in \eqref{C18-weaker-slope} is $1/(p^2-1)$; it does not assert
the stronger slope for a single theta section.

\emph{Proof.}
The operator congruence implies
$D^k(R[[t]])\subset p^{\lfloor k/p\rfloor}R[[t]]$ by iteration.
Let $l_j=D^j\log q(0)$.  In the additive coordinate $z=\lambda(t)$,
both $t(z)$ and $\sigma(z)$ are odd.  Hence $q(t(z))$ and its formal
logarithm are even.  As the logarithm starts in degree four,
$l_j=0$ for odd $j$ and for $j<4$.  For even $j\geq4$,
\[
 l_j=D^{j-2}B(0),\qquad
 v_p(l_j)\geq\left\lfloor\frac{j-2}{p}\right\rfloor.
\]
For every such $j$ one has the elementary inequality
\begin{equation}\label{C18-Bell-weight}
 v_p(j!)-\left\lfloor\frac{j-2}{p}\right\rfloor
          +\frac1{p-1}\leq\frac{jp}{p^2-1}.
\end{equation}
For $j\leq p-1$ this is immediate; for $j=p+1$ it is an equality.
The remaining even integers satisfy $j\geq p+3$.  Using
$v_p(j!)\leq(j-1)/(p-1)$ and
$\lfloor(j-2)/p\rfloor\geq(j-p-1)/p$, the left side is at most
$j/(p(p-1))+(p+1)/p$.  Its difference from the right side is
nonpositive because already at $j=p+3$ the right side minus this
upper bound is
$(p^2-3p-2)/(p(p^2-1))>0$; the difference then increases with $j$.

Bell's formula for the exponential gives
\[
 D^n(q^d)(0)=
 n!\!\sum_{\substack{m_j\geq0,\ \sum jm_j=n\\j\geq4\ {\rm even}}}
       \prod_j\frac{(d l_j)^{m_j}}{m_j!(j!)^{m_j}}.
\]
It is also a polynomial with integer coefficients in the integral
quantities $d l_j$, so these derivatives are integral.  In each term,
use $v_p(m_j!)\leq m_j/(p-1)$ and
\eqref{C18-Bell-weight}.  The resulting valuation is at least
$v_p(n!)-np/(p^2-1)$.  This proves the stated bound for $q^d$,
including negative $d$; $d=0$ is immediate.

Finally write the additive-coordinate series for $q^d$ and $v$ in
exponential Taylor form.  The coefficient of $z^k$ in $q^d$ has
valuation at least $-kp/(p^2-1)$.  The coefficient of $z^a$ in $v$
has valuation at least
\[
 \left\lfloor\frac ap\right\rfloor-v_p(a!)
  =-\sum_{r\geq2}\left\lfloor\frac a{p^r}\right\rfloor
  \geq-\frac{a}{p(p-1)}
  \geq-\frac{ap}{p^2-1}.
\]
Convolution gives the same coefficient bound for $q^dv$, and
multiplication by $n!$ gives the factorial lower bound in
\eqref{C18-weaker-slope}.  Integrality of all its derivatives follows
separately from Leibniz's formula, because the derivatives of $q^d$
and $v$ are integral.  Legendre's factorial formula gives the final
digit-sum expression. \hfill$\square$

\framework{This is an integral trivialization and specialization
calculation for a line-bundle ratio on a smooth proper model.
Its geometric setting is EGA~IV$_2$, Theorem~6.9.1, p.~153
(flat spreading after localization), and SGA~4$\frac12$,
\emph{Cycle}, Definition~2.3.2, p.~145, and Theorem~2.3.8,
p.~148 (cycle classes and pullback). The connection normalization
is additional arithmetic data, not a change to Grothendieck's
cycle-class or polarization statements.}


\subsection{THETA3-C19: intrinsic CM derivative bounds and the reduced theta frame}\label{C19}

This subsection removes the exact Lubin--Tate coordinate hypothesis
from the derivative conclusions of C18 for a good-reduction CM
elliptic curve at an inert prime.  The argument uses integral
translation operators and the CM action, rather than an asserted
integral isomorphism with a specified polynomial formal group.

\paragraph{Hypotheses and normalization.}
Let $p\geq5$ be inert in an imaginary quadratic field $K$, put
$q=p^2$ and $N=q-1$, and fix the embedding $K\hookrightarrow\C_p$
through which the CM action acts on invariant differentials.
Let $E$ be an elliptic curve over a number field with CM by an order
$\mathcal O\subset K$ whose conductor is prime to $p$.
After a finite extension, suppose that $E$ has good reduction over
$R=\mathcal O_L$, where $L/\Q_p$ is finite and contains $K_p$,
that all its CM endomorphisms are defined, and that its chosen
invariant differential $\omega$ generates the integral invariant
differentials of the smooth model $\mathcal E/R$.
CM endomorphisms extend to this abelian scheme by the N\'eron
mapping property: see Conrad, \emph{Minimal models for elliptic
curves}, author version dated 21 November 2015,
Definition~5.1 and Example~5.3, p.~13.
Uniqueness of the extension preserves the group-homomorphism
identities, since they already hold on the generic fiber.  Thus
$\mathcal O\otimes\Z_p=\mathcal O_{K_p}$, with residue field
$\mathbb F_q$, acts on its formal completion.
The completion of this action is proved below, so no continuity
assumption about an unspecified action is needed.

Choose an integral formal parameter $t$ with
$\omega=(1+O(t))\,dt$, and let $D$ be the invariant derivation dual
to $\omega$.  Thus $Dt(0)=1$.  All valuations below satisfy
$v_p(p)=1$.  The ring $R$ may be ramified over $\Z_p$; in
particular, $R/pR$ is not assumed reduced.

\paragraph{Statements.}
For every $f\in R[[t]]$ and every $n\geq1$,
\begin{equation}\label{C19-strong}
 v_p\bigl(D^nf(0)\bigr)
 \geq\max\left\{0,
       \left\lfloor\frac{np}{p^2-1}\right\rfloor-1\right\}.
\end{equation}
The loss of one cannot be replaced uniformly by zero: at
$n=p^2-1$, the function $f=t^{p^2-1}$ has
$v_p(D^nf(0))=p-1$, whereas
$\lfloor np/(p^2-1)\rfloor=p$.
If the torsion-translation trace of $f$ is zero, the stronger bound
\begin{equation}\label{C19-trace-zero}
 v_p\bigl(D^nf(0)\bigr)
       \geq\left\lfloor\frac{np}{p^2-1}\right\rfloor
\end{equation}
holds.  Independently, the full integral operator satisfies
\begin{equation}\label{C19-Dp}
 D^p(R[[t]])\subset pR[[t]].
\end{equation}

For the theta assertion use a short Weierstrass equation and
differential defined over a number field and integral at the
chosen place.  Write the equation as
$y^2=4x^3-g_2x-g_3$ with unit discriminant,
$\omega=dx/y$, and $t=-2x/y$.  Let $\Lambda$ be the complex period
lattice for this algebraic differential, and let $e_2^*(\Lambda)$
be the coefficient in the classical sigma quasiperiod formula.
Then this coefficient is algebraic and
\begin{equation}\label{C19-e2-integral}
 e_2^*(\Lambda)\in R.
\end{equation}
For the tame algebraic sigma translates in C18, put
\[
 s_T^{\rm red}(t)=
 \exp\bigl(-e_2^*(\Lambda)\lambda(t)^2/2\bigr)s_T(t).
\]
These are the reduced, or Fock-translated, theta normalizations
with the same torsion lifts as C18.  Every finite product and
quotient of these series of total integer tensor degree $d$ has
the form $Q(t)^d v(t)$ with $v\in R[[t]]^\times$ and
\[
 Q(t)=\exp\bigl(-e_2^*(\Lambda)\lambda(t)^2/2\bigr)
            \frac{\sigma(\lambda(t))}{t}.
\]
For every $v\in R[[t]]$, every integer $d$, and $n\geq1$,
\begin{equation}\label{C19-theta-weak}
 v_p\bigl(D^n(Q^dv)(0)\bigr)
 \geq\max\left\{0,
       \frac{n}{p^2-1}-\frac{s_p(n)}{p-1}\right\}.
\end{equation}
Here $s_p(n)$ is the base-$p$ digit sum.  In degree zero the common
factor cancels, and \eqref{C19-strong} applies directly.
These are respectively slopes $1/(p^2-1)$ and $p/(p^2-1)$;
no equality of the two estimates is asserted.

\paragraph{Proof of the operator congruence over $R/pR$.}
Let $S=R/pR$, and reduce the formal group and $D$ to $S$.
A continuous $S$-linear derivation of $S[[t]]$ is of the form
$b(t)\,d/dt$.  Translation invariance identifies such a derivation
with its tangent value: if $D=w(t)\,d/dt$ and $w(0)=1$, the
invariance identity for the group law gives
$b(t)=b(0)w(t)$.  Thus invariant derivations form the free
rank-one $S$-module generated by $D$.
In characteristic $p$, the iterated Leibniz formula shows that
$D^p$ is a derivation, since its middle binomial coefficients
vanish.  It is still translation invariant.  Consequently
$D^p=H D$ for some $H\in S$.

Choose $a\in\mathcal O$ whose image in $\mathbb F_q$ lies outside
$\mathbb F_p$.  The formal endomorphism $[a]$ has unit tangent
coefficient $a$, and so is a formal automorphism.  The invariant
differential identity $[a]^*\omega=a\omega$ gives
$D[a]^*=a[a]^*D$.  Comparing its $p$-fold iterate with
$D^p=HD$ gives $(a^p-a)H=0$ in $S$.
The element $a^p-a$ has nonzero residue and is a unit in $S$,
including when $S$ is nonreduced.  Hence $H=0$, proving
\eqref{C19-Dp}.  Iteration also gives
\begin{equation}\label{C19-iterated-Dp}
 D^k(R[[t]])\subset p^{\lfloor k/p\rfloor}R[[t]]
       \qquad(k\geq0).
\end{equation}

\paragraph{Completion and linearization of the CM action.}
For completeness, let $I=(p,t)\subset R[[t]]$.
For every integral one-dimensional formal group,
$[p](t)=pt+O(t^2)\in I^2$ and
$[p^r](t)\in I^{r+1}$, by induction: substitution of
$u\in I^m$ into $[p](u)$ gives an element of $I^{m+1}$.
If $a-b\in p^r\mathcal O$, then the formal group identity
$[a]=F([b],[a-b])$ shows that
$[a](t)-[b](t)\in I^{r+1}$.
Approximating an element of $\mathcal O_{K_p}$ by elements of
$\mathcal O$ therefore gives a convergent sequence of formal
endomorphisms.  The ring $R[[t]]$ is $I$-adically complete;
for each fixed coefficient this convergence is $p$-adic.
The limit respects the group law, addition, composition, and
tangent coefficients.  This constructs the asserted continuous
$\mathcal O_{K_p}$ action.

Let $\mu_N\subset\mathcal O_{K_p}^\times$ be the Teichm\"uller
roots of unity.  The integral parameter
\[
 X(t)=\frac1N\sum_{\zeta\in\mu_N}\zeta^{-1}[\zeta](t)
\]
has linear coefficient one and satisfies
$X([\eta](t))=\eta X(t)$ for every $\eta\in\mu_N$.
It is thus an invertible integral change of parameter, in which
$[\zeta](X)=\zeta X$.
Let $L(X)$ be the formal logarithm normalized by $\omega$.
Its identity $L(\zeta X)=\zeta L(X)$ implies that only powers
congruent to one modulo $N$ occur.  In particular,
\begin{equation}\label{C19-log-gap}
 L(X)=X+O(X^q).
\end{equation}

Define integral translation operators $D^{[j]}$ by
\[
 f(F(X,Y))=\sum_{j\geq0}(D^{[j]}f)(X)Y^j.
\]
They preserve $R[[X]]$, are $R$-linear and $p$-adically continuous,
and commute with one another.  The last assertion follows by
comparing the coefficients in
$f(F(F(X,Y),Z))$, which is symmetric in $Y,Z$.
Over $L$, the invariant differential Taylor formula is
\[
 f(F(X,Y))=\sum_{k\geq0}\frac{D^kf(X)}{k!}L(Y)^k.
\]
Therefore \eqref{C19-log-gap} gives
\begin{equation}\label{C19-small-Hasse}
 D^{[j]}=D^j/j!\qquad(0\leq j<q).
\end{equation}
No such identification is used for the higher-index operators.

\paragraph{The torsion power sums.}
The generic $p$-torsion is a one-dimensional vector space over
$\mathcal O/p\mathcal O=\mathbb F_q$.  The special fiber's
geometric $p$-torsion is also an $\mathbb F_q$-module, but it has
at most $p$ points: multiplication by $p$ has degree $p^2$ and
inseparable degree at least $p$, since its differential is zero.
It must consequently have only the identity point.
Thus all the generic $p$-torsion reduces to the identity, and
the formal multiplication map has degree $q$.
For any nonzero torsion point with coordinate $\xi$, all its
nonzero torsion coordinates are $\zeta\xi$, $\zeta\in\mu_N$.

Weierstrass preparation, followed by this description of the
roots over characteristic zero, gives
\[
 [p](X)=U(X)(X^q-\beta X),\qquad
 U\in R[[X]]^\times,\quad
 \beta=\xi^N=-p/\alpha,\quad \alpha=U(0)\in R^\times.
\]
Indeed, the prepared polynomial is monic of degree $q$, and its
simple roots are exactly the $q$ torsion coordinates; comparison
of the derivative at zero gives $-U(0)\beta=p$.
Consequently the torsion power sums are exactly
\begin{equation}\label{C19-power-sums}
 \sum_{u\in E[p]}X(u)^j=
 \begin{cases}
 q,&j=0,\\
 0,&j>0\text{ and }N\nmid j,\\
 N(-p/\alpha)^r,&j=rN, r\geq1.
 \end{cases}
\end{equation}
Only the torsion polynomial has this special form; no assertion
that $[p](X)$ itself equals $pX+\alpha X^q$ has been made.

\paragraph{The trace-zero estimate without a polynomial coordinate.}
Define the torsion-translation trace after a finite extension
containing the torsion by
\[
 \mathcal T f(X)=\sum_{u\in E[p]}f(F(X,X(u))).
\]
All $X(u)$ lie in the maximal ideal.  Taylor substitution and
summation thus converge coefficientwise $p$-adically for every
integral $f$.  If $\mathcal T f=0$, equations
\eqref{C19-power-sums} and \eqref{C19-small-Hasse} give
\[
 0=p^2f+N\sum_{r\geq1}(-p/\alpha)^rD^{[rN]}f,
 \qquad
 D^{[N]}f=\frac{\alpha p}{N}(1+H)f,
\]
where
\begin{equation}\label{C19-H}
 H=N\sum_{r\geq2}(-1)^r p^{r-2}\alpha^{-r}D^{[rN]}.
\end{equation}
This series is an integral continuous operator on $R[[X]]$:
each summand preserves the integral ring and its scalar has
valuation tending to infinity.  The ring is complete for this
$p$-adic topology.  Every translation operator commutes with
$D$ and with $\mathcal T$; hence $H$ does also, and all these
operators preserve the trace-zero submodule.

Since $v_p(N!)=p-1$, multiplication by $N!$ gives on that
submodule
\[
 D^N=p^p u(1+H),\qquad
 u=\frac{\alpha N!}{Np^{p-1}}\in R^\times.
\]
Write $n=mN+r$ with $0\leq r<N$.  The commutation and
trace-zero stability just proved allow iteration, giving
\[
 D^nf=p^{pm}u^m r!D^{[r]}(1+H)^m f.
\]
The last factor is integral, so its value at zero has valuation
at least zero.  Furthermore
\[
 v_p(r!)=\left\lfloor\frac rp\right\rfloor
        =\left\lfloor\frac{rp}{p^2-1}\right\rfloor
          \qquad(0\leq r<p^2-1).
\]
For the second equality, write $r=ap+b$, $0\leq b<p$;
then $rp=a(p^2-1)+(bp+a)$, and
$0\leq bp+a<p^2-1$ in the stated range.
This proves \eqref{C19-trace-zero}.

This is the divisibility argument of Katz,
\emph{Divisibilities, congruences, and Cartier duality},
\S III, equations (3.1)--(3.8), pp.~671--672, and Part I of
the proof of Theorem~3.11, pp.~672--673, applied using
\eqref{C19-power-sums} rather than an exact polynomial
coordinate.  The stronger Cartier-operator congruences used
for the other conclusions of that theorem are not claimed here.

\paragraph{The trace ideal and removal of the trace-zero condition.}
Put $A=R[[S]]$, $B=R[[X]]$, with $S=[p](X)$.
This is a finite free extension of degree $q$.
Weierstrass preparation gives
\[
 [p](X)-S=U_S(X)P_S(X),
\]
where $P_S$ is monic of degree $q$ and $U_S$ is a unit.
Write $\omega=a(X)dX$ with $a(X)$ a unit.
The invariant differential identity implies
\[
 [p]'(X)=p\,\frac{a(X)}{a([p](X))}.
\]
In $B=A[X]/(P_S)$ it follows that
$P_S'(X)=p\,v(X)$ with $v\in B^\times$.

Here is an elementary proof of the exact trace ideal, avoiding
an uncomputed different.  If $\ell_{
\rm top}:B\to A$ takes the coefficient of $X^{q-1}$ in the
unique representative of degree less than $q$, then
\[
 \operatorname{Tr}_{B/A}\bigl(b/P_S'(X)\bigr)
      =\ell_{\rm top}(b).
\]
This identity follows over the fraction field by Lagrange
interpolation at the distinct roots of $P_S$; it then holds
in the fraction algebra for every $b\in B$.
Substituting $vb$ gives
\[
 \operatorname{Tr}_{B/A}(b)/p=\ell_{\rm top}(vb)\in A.
\]
Taking $b=v^{-1}X^{q-1}$ makes the right side one.
We have therefore proved the exact equality
\begin{equation}\label{C19-trace-ideal}
 \operatorname{Tr}_{B/A}(B)=pA.
\end{equation}

For any integral $f$, set
$h=\operatorname{Tr}_{B/A}(f)/p\in A$ and
$g=pf-[p]^*h$.  Since the rank is $q=p^2$,
$\operatorname{Tr}(g)=p\operatorname{Tr}(f)-qh=0$.
The algebra trace, pulled back to $B$, equals $\mathcal T$:
after adjoining the torsion, the roots of
$[p](Y)=[p](X)$ are precisely $F(X,X(u))$.
The equality can first be checked in the separable fraction
algebra and then in the integral power-series ring.
Thus \eqref{C19-trace-zero} applies to $g$.
On the other hand, invariant differentiation gives
\[
 D^n[p]^*h(0)=p^nD^nh(0),
\]
where on the right $D$ is the same invariant derivation in the
target parameter.  Since $D$ preserves the integral ring and
$\lfloor np/(p^2-1)\rfloor\leq n$, the identity
$pf=g+[p]^*h$ proves \eqref{C19-strong}.
The additional lower bound zero follows from integrality of
$D^nf$.  Finally, $D=w(t)d/dt$ with $w(0)=1$ gives
$D^n(t^n)(0)=n!$: terms with fewer than $n$ ordinary derivatives
still vanish at zero.  For $n=p^2-1$, Legendre's formula gives
$v_p(n!)=p-1$, proving the claimed sharpness.
The estimates also extend to $f\in\mathcal O_{\C_p}[[t]]$
on the base-changed fixed formal group.  For fixed $n$, the
value $D^nf(0)$ depends only on the first $n+1$ coefficients
of $f$.  Approximate these by integral algebraic coefficients,
apply the finite-extension assertion, and use continuity and
closedness of each stated valuation inequality.

\paragraph{Integral de Rham CM splitting and $e_2^*$.}
Use the short Weierstrass model specified above and set
$\eta=x\omega$.  C18's integral expansion gives
\[
 x=t^{-2}+O(t^2),\qquad
 \omega=(1+O(t^4))dt,\qquad
 \eta=t^{-2}dt+O(t^2)dt.
\]
On $\mathcal E\setminus O$ the differential $\eta$ is regular.
Near $O$ the differential $\eta-d(-1/t)$ is regular.  These
two differentials and the transition function $-1/t$ define
an integral \v Cech--de Rham class.  The image of this class
in $H^1(\mathcal E,\mathcal O_{\mathcal E})$ pairs with
$\omega$ by the residue $\pm\operatorname{Res}_O(\omega/t)
=\pm1$.  Relative Serre duality and the Hodge exact sequence
therefore show that $(\omega,\eta)$ is an integral basis of
$H^1_{\rm dR}(\mathcal E/R)$.  Equivalently, its de Rham cup
pairing is a unit; the sign of this pairing is immaterial here.
For a curve, the Hodge exact sequence used here follows also
directly from the two-term \v Cech--de Rham complex: the map
$H^1(\mathcal O)\to H^1(\Omega^1)$ induced by $d$ is zero,
as the residue of an exact differential is zero.

Choose the global CM element $a$ above.  Both $a$ and
$a-\bar a$ are units in $R$, since reduction of complex
conjugation on $K_p$ is the $p$-power automorphism of
$\mathbb F_q$.  The integral endomorphism acts as
\[
 [a]^*\omega=a\omega,
 \qquad [a]^*\eta=b\eta+c\omega,
 \qquad b,c\in R.
\]
The cup product is multiplied by $\deg[a]=a\bar a$;
comparison in the integral basis gives $ab=a\bar a$ and
hence $b=\bar a$.

Here is the comparison with the specified analytic coefficient.
Under complex uniformization, $x=\wp(z)$ and $\omega=dz$.
Writing $\zeta=\sigma'/\sigma$, one has
$\zeta'=-\wp$.  Thus the periods of $\eta$ are the negatives
of the sigma quasiperiods.  The classical formula is
\[
 \eta_{\rm per}(\gamma)
    =e_2^*(\Lambda)\gamma+C\bar\gamma
       \qquad(\gamma\in\Lambda),
\]
for the nonzero area constant $C$.
Consequently the periods of
$\eta+e_2^*(\Lambda)\omega$ are $-C\bar\gamma$;
this class is a $\bar a$-eigenvector for $[a]^*$.
Comparison with the integral matrix above gives exactly
\[
 c=(\bar a-a)e_2^*(\Lambda).
\]
It proves algebraicity and \eqref{C19-e2-integral}, since
$\bar a-a$ is a unit.  The source for the quasiperiod formula
and its normalization is Bannai--Kobayashi,
\emph{Algebraic theta functions and the $p$-adic interpolation
of Eisenstein--Kronecker numbers}, arXiv:math/0610163v4,
Example~1.9, pp.~10--11.  That example explicitly gives
$\theta_{\rm red}(z)=e^{-e_2^*z^2/2}\sigma(z)$ and
$\theta_{\rm red}'(0)=1$.

\paragraph{Completion of the reduced theta estimate.}
Let $q_0(t)=\sigma(\lambda(t))/t$; this is the common factor
denoted $q(t)$ in C18 and is distinct from the number $q=p^2$.
C18 proves $s_T=q_0u_T$, with $u_T$ an integral unit, and
\[
 B=D^2\log q_0=-x-D^2\log t\in t^2R[[t]].
\]
The sigma quasiperiod identity also gives, by direct substitution
in the translation formula of C18,
\[
 s_T^{\rm red}(t)=Q(t)u_T(t),\qquad
 Q(t)=e^{-e_2^*\lambda(t)^2/2}q_0(t).
\]
In particular every stated product or quotient has the announced
factorization.  No integrality of the raw series $Q$ is assumed.

Put $l_j=D^j\log Q(0)$.  In the additive coordinate
$z=\lambda(t)$, the series $Q$ and its logarithm are even,
and $Q(0)=1$.  Hence $l_0=l_1=0$, all odd $l_j$ vanish,
$l_2=-e_2^*\in R$, and for even $j\geq4$,
\[
 l_j=D^{j-2}B(0),\qquad
 v_p(l_j)\geq\left\lfloor\frac{j-2}{p}\right\rfloor
\]
by \eqref{C19-iterated-Dp}.
The Bell-polynomial proof in C18 now applies verbatim, allowing
the additional even index $j=2$.  Its required numerical
inequality for that index is
\[
 v_p(2!)-0+\frac1{p-1}
       \leq\frac{2p}{p^2-1},
\]
which holds for $p\geq5$.  For clarity, in each exponential
Bell term the inequality
\[
 v_p(j!)-\left\lfloor\frac{j-2}{p}\right\rfloor
       +\frac1{p-1}\leq\frac{jp}{p^2-1}
       \qquad(j\geq2\text{ even})
\]
and $v_p(m!)\leq m/(p-1)$ bound the total denominator loss
by $np/(p^2-1)$.  All $D^n(Q^d)(0)$ are also integral,
since exponential Bell polynomials have integer coefficients
in the integral variables $d l_j$, including for $d<0$.
Thus
\[
 v_p(D^n(Q^d)(0))
       \geq\max\{0,v_p(n!)-np/(p^2-1)\}.
\]
For $v\in R[[t]]$, \eqref{C19-iterated-Dp} bounds the
coefficient of $z^a$ in $v(t(z))$ below by
$\lfloor a/p\rfloor-v_p(a!)\geq-a/(p(p-1))
\geq-ap/(p^2-1)$.
Convolution gives the same factorial bound for $Q^dv$,
and Leibniz's rule gives integrality of its derivatives.
Legendre's formula yields \eqref{C19-theta-weak}.
The degree-zero assertion follows from cancellation of $Q$
and the already proved \eqref{C19-strong}.
\hfill$\square$

\paragraph{Formal-module generalization and several independent factors.}
The formal-group assertions have a slightly wider scope than
the elliptic application.  Let $R$ be any complete discrete
valuation ring of mixed characteristic $(0,p)$ containing
the integers of the unramified quadratic extension of $\Q_p$.
Let $G/R$ be a one-dimensional smooth commutative formal group
with a strict action of those integers, meaning that the action
of $a$ on the tangent space is multiplication by $a$.
Assume that $[p]:G\to G$ is finite flat of degree $p^2$ and
that its reduction has its entire kernel supported at the
identity.  Equivalently in this setting one may assume that
the closed-fiber formal group has height two.
Then all of \eqref{C19-Dp}, \eqref{C19-trace-zero},
\eqref{C19-trace-ideal}, and \eqref{C19-strong} hold.
For these formal assertions alone $p\geq2$ suffices.

Indeed, the integral action is now supplied as a hypothesis,
and averaging over its Teichm\"uller units gives the same
coordinate $X$.  The prepared $p$-torsion polynomial has
$p^2$ simple roots in characteristic zero, all in the maximal
ideal.  Its roots form a one-dimensional vector space over
$\mathbb F_{p^2}$, so the same orbit and power-sum calculation
applies.  Every subsequent translation-operator, convergence,
invariant-differential and trace argument was purely formal.
The elementary factorial identities remain valid at $p=2,3$.
No elliptic sigma or de Rham assertion at those two primes
is included in this generalization.

The proof in fact gives operator divisibility, not only a
bound at the origin.  With
\[
 b(n)=\max\{0,\lfloor np/(p^2-1)\rfloor-1\},
\]
one has $D^nR[[t]]\subset p^{b(n)}R[[t]]$.
For a product $G_1\times\cdots\times G_r$ of such formal
modules and its independently normalized derivations $D_i$,
it follows for every integral series $f$ that
\[
 v_p\bigl(D_1^{n_1}\cdots D_r^{n_r}f(0)\bigr)
       \geq\sum_{i=1}^r b(n_i).
\]
To justify coefficient rings in this last assertion, apply the
one-variable proof coefficientwise in the other variables:
all translation operators preserve their integral coefficients,
the same $p$-adically convergent sums still converge, and the
finite-free trace identity is preserved by this power-series
base change.  The $D_i$ commute, so their integral operator
divisibilities multiply.

Thus these bounds apply whenever the formal completion of a
good-reduction CM abelian scheme has been decomposed into the
specified one-dimensional height-two factors by its integral
local CM idempotents.  The precise decomposition, its integral
differentials and the theta-section connection are still part
of the input when using this multi-factor consequence.

\paragraph{Scope of the arithmetic bridge.}
The only formal-group hypothesis used here beyond good reduction
is the maximal local unramified quadratic CM action; a CM order
with conductor divisible by $p$ is not covered by this proof.
The arguments apply after every finite extension carrying the
stated good model and normalized differential.  In particular,
they never replace $R/pR$ by its reduced quotient.
The result supplies intrinsic derivative estimates for the actual
tame theta functions with the named differential and torsion
normalization.  A comparison with a differently normalized
harmonic or covariant derivative still uses its explicit analytic
normalization; the bounds do not make that identification by
the existence of an arbitrary line-bundle trivialization.

The de Rham part is a concrete realization of the integral
correspondence and polarization setting, and its algebraic
splitting is computed rather than assumed.  It is compatible
with the Albanese and polarization discussion in Grothendieck's
cote~44, scans~7--8, and with C17's precisely stated cycle-class
specialization formalism.  Those archive passages are not cited
as proofs of the local valuation estimates established here.


\subsection{THETA3-C21: split CM formal factors and higher-dimensional theta ratios}\label{C21}

\paragraph{Statement.}
Let $A$ be a $g$-dimensional abelian variety over a number field,
with CM by an order $\mathcal O$ in a CM field $E$ of degree
$2g$.  Write $F$ for the maximal totally real subfield of $E$.
Let $p\geq5$ split completely in $F$, and suppose every prime
$\nu$ of $F$ above $p$ is unramified and inert in $E$.
Assume that $\mathcal O$ is maximal at $p$.
After a finite extension choose a good abelian scheme
$\mathcal A/R$, with $R$ the integers of a finite extension
of $\Q_p$, and define all CM endomorphisms over it.
Enlarge the field further to contain the CM embeddings whenever
needed.  Then
\begin{equation}\label{C21-formal-product}
 \widehat{\mathcal A}\simeq
       \prod_{\nu\mid p}G_\nu,
\end{equation}
where every $G_\nu$ is an integral one-dimensional formal
module for the unramified quadratic extension of $\Q_p$,
with height-two special fiber.

Choose integral invariant differentials on the factors and
write $D_\nu$ for their dual invariant derivations, normalized
by integral parameters with $D_\nu t_\nu(0)=1$.
For $f\in R[[t_\nu:\nu\mid p]]$ and a multi-index
$\mathbf n=(n_\nu)$, put
\[
 b_p(n)=\max\left\{0,
            \left\lfloor\frac{np}{p^2-1}\right\rfloor-1\right\}.
\]
Then the intrinsic mixed-derivative estimate is
\begin{equation}\label{C21-mixed-bound}
 v_p\left(\left(\prod_{\nu\mid p}D_\nu^{n_\nu}\right)f(0)\right)
       \geq\sum_{\nu\mid p}b_p(n_\nu).
\end{equation}
In particular the right side is at least
$p|\mathbf n|/(p^2-1)-2r$, where $r$ is the number of
nonzero entries in $\mathbf n$.

There is a direct theta application with a concrete integral
hypothesis.  Suppose $s_1,\ldots,s_m$ and
$u_1,\ldots,u_m$ are integral sections of the same power
$\mathscr L^{\otimes k}$ of a line bundle on $\mathcal A$,
and every $u_i$ is a unit at the identity in an integral local
frame.  Their equal-degree ratio
\begin{equation}\label{C21-section-ratio}
 f=\frac{\prod_i s_i}{\prod_i u_i}
\end{equation}
belongs to the integral formal ring and satisfies
\eqref{C21-mixed-bound}.  The same conclusion holds after
an explicitly chosen common prime-to-$p$ isogeny pullback
has trivialized the flat torsion twists of translated theta
sections, provided the resulting sections have the stated
integral normalization and the denominator is a unit at the
identity.  No assertion that all higher-dimensional tame
theta values are automatically units is used.

\paragraph{Integral CM idempotents and their heights.}
The local splitting hypotheses give the exact algebra
\[
 \mathcal O\otimes\Z_p
       =\mathcal O_E\otimes\Z_p
       \simeq\prod_{\nu\mid p}\mathcal O_{E_\nu},
 \qquad E_\nu/\Q_p\text{ unramified quadratic}.
\]
Let $e_\nu$ be its mutually orthogonal idempotents.
These act on $\mathcal A[p^n]$ for every $n$: the action
factors through $\mathcal O/p^n\mathcal O$, and the
idempotents are compatible as $n$ varies.
An idempotent on a finite flat commutative group splits it
into its image and its complementary image, with inverse
maps given by the idempotent and its complement.
Thus
\[
 \mathcal A[p^\infty]=\prod_{\nu\mid p}\mathcal G_\nu,
       \qquad \mathcal G_\nu=e_\nu\mathcal A[p^\infty].
\]
The factors are finite flat at every level; equivalently,
their coordinate algebras are direct retracts of the finite
flat coordinate algebra before splitting.

The rational Betti homology of $A$ is a rank-one $E$-module:
it is a faithful $E$-vector space of rational dimension $2g$.
Using complex uniformization, its $p$-adic completion identifies
the rational Tate module as a rank-one
$E\otimes\Q_p$-module.  Hence every $\mathcal G_\nu$
has height two.
The CM type selects one embedding out of every conjugate
pair of embeddings of $E$.  Because $p$ splits completely
in $F$, the two embeddings extending a fixed $\nu$ form
exactly such a pair.  Therefore
$e_\nu\operatorname{Lie}(\mathcal A)$ is an integral direct
summand of rank one.  Each $\mathcal G_\nu$ has dimension
one, and its $\mathcal O_{E_\nu}$ action on the tangent
line is through one specified embedding into $R$.

Here is why no \'etale part remains in its special fiber.
The Tate module of that \'etale part is a free $\Z_p$-module
of some rank $h_{\rm et}\leq2$, stable under
$\mathcal O_{E_\nu}$.  Its rationalization is an
$E_\nu$-vector space, so $h_{\rm et}$ is even.
It cannot be two: the remaining connected part would have
height zero, making the whole factor \'etale of dimension
zero, whereas its dimension is one.  Thus $h_{\rm et}=0$.
The special fiber is connected.  Equivalently its geometric
$p$-torsion, an $\mathbb F_{p^2}$-module, has only the
identity point.  All the finite kernels are therefore supported
at the identity in the special fiber.

\paragraph{A direct construction of the formal product.}
The completed local CM action exists on
$\widehat{\mathcal A}$ without an algebraization claim for
the idempotents.  In arbitrary integral coordinates
$t_1,\ldots,t_g$, put $I=(p,t_1,\ldots,t_g)$.
Each component of $[p](\mathbf t)$ equals $pt_i$ plus
terms of total degree at least two.  The same induction as
in C19 gives $[p^n](\mathbf t)\in I^{n+1}$ componentwise.
Approximating local CM elements by global ones consequently
converges $I$-adically to formal endomorphisms, respecting
addition, composition and tangent maps.

Choose the initial coordinates so that their tangent
covectors respect the rank-one direct-sum decomposition
defined by the $e_\nu$.  Set
\[
 X_\nu=t_\nu\circ e_\nu.
\]
The linear part of the tuple $(X_\nu)$ is the identity.
The integral formal inverse-function theorem therefore makes
it a system of integral coordinates.  Orthogonality gives
$X_\mu(e_\nu P)=0$ for $\mu\ne\nu$, and
$X_\nu(e_\nu P)=X_\nu(P)$.
Thus the image of $e_\nu$ is precisely the coordinate axis
on which all other $X_\mu$ vanish.  Because the $e_\nu$ are
group homomorphisms and their sum is the identity, the group
law is the product of its restrictions to these axes.
This proves \eqref{C21-formal-product} with each factor
smooth of dimension one.

The connectedness proved above identifies the $p^n$-kernels
of these formal factors with the entire finite flat groups
$e_\nu\mathcal A[p^n]$.  In particular $[p]$ on each
$G_\nu$ is finite flat of degree $p^2$, with height-two
special fiber.  Identifying $\mathcal O_{E_\nu}$ through
its tangent embedding makes its action strict, exactly as
required by the formal-module generalization of C19.
The exact basic references are Tate, \emph{$p$-Divisible
Groups}, \S1.4, p.~160 (connected--\'etale sequence),
\S2.2, Proposition~1, p.~162 with proof pp.~162--164
(connected groups and divisible formal groups), the
dimension definition following (4), p.~164, and
\S2.3, Proposition~3 and Example~(b), p.~166
(invariance of height and dimension under reduction and
the formal completion of an abelian scheme).
The idempotent splitting and coordinate construction have
been given explicitly above.

\paragraph{Mixed derivatives.}
C19 gives the full one-variable operator divisibility
\[
 D_\nu^nR[[X_\nu]]\subset
              p^{b_p(n)}R[[X_\nu]].
\]
Apply it coefficientwise in the other formal variables.
The factors are independent, so their derivations commute,
and successive application multiplies these integral powers
of $p$.  This proves \eqref{C21-mixed-bound}.
For $n>0$, $b_p(n)\geq np/(p^2-1)-2$; summing over the
nonzero entries gives the displayed total-degree consequence.
As in C19, finite-jet approximation extends the assertion to
integral $\C_p$ coefficients on the fixed base-changed group.
Rescaling any differential by an integral unit preserves the
valuation bound.  A nonunit analytic period rescaling must
instead be retained explicitly in the derivative formula.

\paragraph{Integral sections, common frames, and torsion twists.}
Trivialize $\mathscr L$ on the formal neighborhood of the
identity by an integral frame.  The numerator sections in
\eqref{C21-section-ratio} then have integral coefficients,
and every denominator series has a unit constant term.
The quotient is consequently an integral formal function.
Changing the common frame multiplies each section by the
same factor to the same tensor degree, so that factor cancels
exactly.  In the canonical complex theta trivialization this
includes the common quadratic exponential potential.
Thus the canonical theta ratio represents the same formal
function in the exponential coordinates normalized by the
chosen invariant differentials.  The analytic matching is
the degree-zero case of A8.  If the supplied theta sections
are rescaled by constants, their ratio is correspondingly
rescaled by their product quotient $c$; the resulting bound
is \eqref{C21-mixed-bound} plus $v_p(c)$.

For translated sections, $\tau_T^*\mathscr L$ differs from
$\mathscr L$ by the flat line bundle attached to
$\phi_{\mathscr L}(T)\in\mathcal A^\vee$.
If all the translations have prime-to-$p$ torsion order,
choose a common prime-to-$p$ integer $M$ killing these
dual torsion points.  Pullback by $[M]$ trivializes their
flat twists: on rigidified degree-zero line bundles it
corresponds to multiplication by $M$ on the dual variety.
Choose these integral identifications and retain the
integrality and unit-denominator conditions in the statement.
The pulled-back equal-degree ratio is then a rational function
on the isogeny cover whose completion is integral.
This assertion does not require it to descend to the
unpulled-back abelian variety.

On the formal neighborhood $[M]$ is an integral isomorphism,
with inverse having integral coefficients because its tangent
matrix is $M$ times the identity and $M\in R^\times$.
Also
\[
 D_\nu[M]^*=M[M]^*D_\nu.
\]
For a mixed derivative at the identity, passage to the
pullback therefore multiplies its value by
$M^{|\mathbf n|}$, a unit.  The same valuation estimate
holds in the original local translated-theta coordinates
with the specified identifications. \hfill$\square$

\paragraph{Remaining scope.}
This supplies the strong local estimate on an actual
higher-dimensional CM formal group at the stated split-in-$F$,
inert-in-$E$ primes, and on canonical theta ratios meeting
the explicit integral-section hypothesis.  Positive tensor
degree retains the canonical connection; C22 shows why
that issue cannot be removed by declaring every theta
section to have the ratio estimate.  Neither the formal
idempotents nor this analytic estimate construct the
special-cycle correspondence required in C17.


\subsection{THETA3-C22: the weaker slope is sharp for an actual tame theta section}\label{C22}

The stronger integral-function estimate in C19 cannot be extended
with a bounded loss to every individual tame reduced theta section.
The following explicit example determines the optimal possible
linear slope for that proposed extension.

\paragraph{Statement with the actual theta normalization.}
Take
\[
 E:\quad y^2=4x^3-4x,\qquad p=7,\qquad
 \omega=dx/y,\qquad t=-2x/y,
 \qquad T=(0,0).
\]
This curve has good reduction at 7 and CM by $\Z[i]$, and 7 is
inert in $\Q(i)$.  The point $T$ is a nonzero point of order
two.  Let $z_T$ be a half-period representing $T$, with a
compatible torsion lift as in C18, and use the actual Fock
translation
\[
 \sigma_{z_T}(z)=
 e^{-\eta_{\rm per}(z_T)(z+z_T/2)}\sigma(z+z_T).
\]
Here the real-linear quasiperiod map is extended to the
half-period in the indicated formula.  The period lattice is
$i$-stable, so $e_2^*=0$ and this is already the reduced theta
translation.  Normalize its nonzero constant by putting
\[
 F(z)=\frac{\sigma_{z_T}(z)}{\sigma_{z_T}(0)},\qquad
 f_T(t)=F(\lambda(t)),\qquad F(0)=1.
\]
The algebraic value $\sigma_{z_T}(0)$ is a 7-adic unit by
Katagiri--Kobayashi, arXiv:2609.07024v1,
Theorem~2.7, pp.~7--8, and Lemma~3.3(1), pp.~9--10.
Thus this normalization does not change any derivative valuation.
If $D$ is the invariant derivation dual to $\omega$, then for
every integer $r\geq0$, with $n_r=8\cdot7^r$,
\begin{equation}\label{C22-exact}
 v_7\bigl(D^{n_r}f_T(0)\bigr)
      =\frac{7^r-1}{6}
      =\frac{n_r}{48}-\frac16.
\end{equation}
Consequently no bound
$v_7(D^nf_T(0))\geq c n-C$ with $C$ independent of $n$
can hold for any $c>1/48$.
In particular the coefficient $7/48$ from the degree-zero
integral-function theorem cannot hold for this degree-one
theta section, even with an arbitrarily large fixed loss.

\paragraph{Evenness and the logarithmic derivatives.}
The sigma normalization and quasiperiod identities are those of
Bannai--Kobayashi, arXiv:math/0610163v4,
Example~1.9, pp.~10--11.
The automorphism $(x,y)\mapsto(-x,iy)$ pulls $\omega$ back to
$i\omega$, so the lattice is $i$-stable.  The weight-two
transformation of $e_2^*$, or comparison in
$\eta_{\rm per}(\gamma)=e_2^*\gamma+C\bar\gamma$, gives
$e_2^*=-e_2^*$ and hence $e_2^*=0$.

Since $2z_T$ is a lattice element not divisible by two in the
lattice, the sigma multiplier at $2z_T$ is $-1$.
The translation formula and oddness of sigma therefore give
\[
 \sigma(z+z_T)
     =-e^{2\eta_{\rm per}(z_T)z}\sigma(z-z_T),
 \qquad
 \sigma(z_T-z)
     =e^{-2\eta_{\rm per}(z_T)z}\sigma(z+z_T).
\]
After multiplication by the defining exponential factor these
identities say $F(-z)=F(z)$.  Thus $\log F$ is even, with
constant and linear terms zero.
The second derivative has the particularly simple algebraic form
\begin{equation}\label{C22-B}
 \frac{d^2}{dz^2}\log F(z)=-\wp(z+z_T).
\end{equation}
The right side is $-x(T+P(t))$ in the formal parameter at the
identity.  It belongs to $R[[t]]$ after a finite extension
carrying the torsion data: translation by $T$ is integral,
and its residue disc is disjoint from the sole pole of $x$,
because $T$ is nonzero prime-to-7 torsion.

Put $l_j=(d/dz)^j\log F(0)$.  Equation \eqref{C22-B}
and the intrinsic congruence $D^7R[[t]]\subset7R[[t]]$ of
C19 give
\begin{equation}\label{C22-log-bound}
 l_j=0\quad(j\text{ odd}),\qquad
 v_7(l_j)\geq\left\lfloor\frac{j-2}{7}\right\rfloor
       \quad(j\geq2\text{ even}).
\end{equation}
The decisive logarithmic derivative is a unit.  Indeed,
$Dx=y$ and $Dy=6x^2-2$ imply by direct differentiation
\[
 D^2x=6x^2-2,\qquad
 D^4x=120x^3-72x,\qquad
 D^6x=5040x^4-4032x^2+144.
\]
Evaluating at $T$ gives
\begin{equation}\label{C22-unit}
 l_8=-D^6x(T)=-144\in R^\times.
\end{equation}
The same computation gives $l_2=0$ and $l_4=2$, although
only \eqref{C22-unit} is needed for sharpness.

\paragraph{The unique lowest-valuation Bell term.}
Set $s=7/48$.  For a nonzero even logarithmic derivative
define the surplus
\[
 \epsilon_j=v_7(l_j)-v_7(j!)+js-\frac16.
\]
We have $\epsilon_8=0$, and
\begin{equation}\label{C22-strict}
 \epsilon_j>0\qquad(j\ne8\text{ even},\ l_j\ne0).
\end{equation}
Here is the strict inequality in full, rather than an appeal
only to the non-strict Bell bound of C18.
For $j=2,4,6$, \eqref{C22-log-bound} gives
$\epsilon_j\geq(7j-8)/48>0$.
For every even $j\geq10$, use
\[
 v_7(j!)\leq\frac{j-1}{6},\qquad
 \left\lfloor\frac{j-2}{7}\right\rfloor
       \geq\frac{j-8}{7}.
\]
It follows that
\[
 \epsilon_j\geq
   \frac{j-8}{7}-\frac{j-1}{6}
          +\frac{7j}{48}-\frac16
      =\frac{41j-384}{336}>0.
\]
At $j=8$, $v_7(8!)=1$ and \eqref{C22-unit} gives the
claimed equality.  Vanishing logarithmic derivatives simply
contribute no Bell monomials.

Write $F(z)=\sum_n a_nz^n$.  Its coefficient formula is
\[
 a_n=
 \sum_{\substack{m_j\geq0,\ \sum jm_j=n\\j\geq2\ {\rm even}}}
       \prod_j\frac{l_j^{m_j}}{m_j!(j!)^{m_j}}.
\]
For any nonzero summand, add $ns$ to its valuation and use
Legendre's formula.  The result is exactly
\begin{equation}\label{C22-partition}
 \sum_j m_j\epsilon_j
       +\frac16\sum_j s_7(m_j).
\end{equation}
Take $n=8\cdot7^r$.
The partition with $m_8=7^r$ and all other multiplicities
zero has value $1/6$ in \eqref{C22-partition}.
It is the only such partition: any other partition uses an
index $j\ne8$, so the first sum is strictly positive by
\eqref{C22-strict}; the second sum is at least $1/6$,
since some multiplicity is positive.
Thus the coefficient sum has a unique term of least valuation,
and the nonarchimedean triangle rule gives
\[
 v_7(a_n)=-\frac{7n}{48}+\frac16.
\]
Since the base-7 digits of $8\cdot7^r$ sum to two,
$v_7(n!)=(n-2)/6$.  Multiplying the coefficient by $n!$
therefore gives
\[
 v_7(F^{(n)}(0))=
    \frac{n-2}{6}-\frac{7n}{48}+\frac16
    =\frac{n}{48}-\frac16.
\]
The identity $D=d/dz$ in the additive coordinate proves
\eqref{C22-exact}.  Finally, if $c>1/48$, subtracting the
right side of \eqref{C22-exact} from $c n_r$ tends to
infinity.  This proves the asserted failure of every bounded-loss
estimate with the larger slope. \hfill$\square$

\paragraph{A second instance over $\Q(\sqrt{-3})$.}
The same phenomenon occurs on
$E:y^2=4x^3-4$ at $p=5$, with the nonzero point of
order two $T=(1,0)$ and $\omega=dx/y$.
This curve has good reduction at 5, CM by
$\Z[\zeta_3]$, and inert reduction of the CM prime.
The $\zeta_3$-symmetry forces $e_2^*=0$.
The same half-period argument makes the unit-normalized
actual Fock translate $F$ even, and again
$D^2\log F=-x(T+P(t))$ is integral.
Here $Dx=y$, $Dy=6x^2$ and
$D^4x=120x^3-48$, so
$l_6=-D^4x(T)=-72$ is a 5-adic unit.

For completeness the strict surplus calculation is also
explicit in this case.  Put
$\epsilon_j=v_5(l_j)-v_5(j!)+5j/24-1/4$.
The bound $v_5(l_j)\geq\lfloor(j-2)/5\rfloor$
gives $\epsilon_j>0$ for $j=2,4$;
$\epsilon_6=0$; and for all even $j\geq8$ it gives
\[
 \epsilon_j\geq
    \frac{j-6}{5}-\frac{j-1}{4}
          +\frac{5j}{24}-\frac14
       =\frac{19j-144}{120}>0.
\]
The identical Bell-partition argument, now with the
unique multiplicity $m_6=5^r$, yields
\[
 v_5\bigl(D^{6\cdot5^r}f_T(0)\bigr)
       =\frac{5^r-1}{4}
       =\frac{6\cdot5^r}{24}-\frac14.
\]
Thus the obstruction is present for this second CM field
and its stated actual theta normalization as well.

\paragraph{Positive tensor powers.}
In either example the same exact valuation formula holds
with $F$ replaced by $F^d$ for every positive integer $d$
prime to the displayed prime.  Indeed, the normalized
tensor section is still even and has constant term one,
and every logarithmic derivative $l_j$ is replaced by
$d l_j$.  Its valuation, and hence every surplus
$\epsilon_j$, is unchanged.  The same unique least
Bell monomial proves the same equality at every stated
index.  The obstruction therefore persists for these
actual positive tensor powers, with no adjustable
nonintegral scalar normalization.

\paragraph{What this decides.}
C19's weaker estimate for actual reduced theta tensors has
linear coefficient $1/(p^2-1)$, and this example shows that
coefficient is optimal for the class under consideration.
Equal-degree ratios still have the stronger intrinsic bound
because their common canonical frame cancels.
The example rules out transferring that stronger bound to
all individual sections by an unspecified torsion trace or
frame choice.  It makes no assertion that a differently
normalized archive theorem has these same degree, period,
or jet hypotheses; those data must be matched explicitly.


\part{Unresolved comparisons and current gates}
\section{Statements not asserted as proved}

There is no unchecked theorem silently used in the preceding part.
The implications with hypotheses, including B8 and C15's universality
reduction, are proved implications; their hypotheses are not supplied
by this run. The following are the remaining identifications.

\paragraph{G1: the exact iterated character-independent functional.}
The common oscillator kernel, A1's normalized zero orbit, B11's
restriction seesaw, A5's finite-frame identity and A7's actual single
theta lift have been established. B17 now makes the actual wedge
a literal double lift $V\to P\to V$ by the genuine adjoint
kernel and a common positive inverse $Q^{-1}$.
It specializes the analytic last-step Bessel seesaw without a
new occurrence hypothesis. The operator is common to every
character for fixed spaces, kernel and levels.
The remaining \texttt{MainX} question is its stronger arithmetic
uniform relation-preservation statement. To apply it, one must
construct its actual common geometry,
tuples, archimedean maps, definite class kernels and arithmetic links,
and identify its normalized functional on the selected 032 vectors.
The concrete all-place tower fails the canonical node-gap and
signature requirements. The abstract two-step case is more flexible:
with global Witt index zero, one can take $m=2$ and the Bessel
unipotent rank zero. Its middle space must be totally definite.
A7's natural multiplicity space is totally definite precisely when
the selected first-list CM types agree; that direct realization then
gives a divisor-generated balanced line. B16 supplies a conditional
pole criterion for a different coherent definite source.
The arithmetic kernel/covolume and terminal-mass equations,
and an arithmetic interpretation of the analytic correction,
are the remaining matching problems. B17's upper space is
$A\perp\mathbb H$, of signature $(3,1)$ at every real place;
it is not a hidden definite inner form.

\paragraph{G2: turning the explicit certificates into a uniform
arithmetic lower bound for the prescribed iterated family.}
Q1--Q11 and B5--B7 give actual formulas under specified inputs.
The finite unipotent search yields $|J|\geq1/(4S_{\rm ext})$ with
an explicit grid size after its Gaussian envelope and lattice basis
are supplied. An arithmetic Gram bound needs a coefficient field,
denominator and bounds at all embeddings. Q9 provides resultant
denominators for specified integral Hecke matrices. Q11 transfers
the definite period-square certificate directly to the geometric
ball norm using explicit injective-chart constants; no unknown
ball Gram upper bound is needed. The remaining
step is the normalization and overlap data of the particular
iterated family in G1. Positivity alone cannot supply it, as B4
and B6 show. No universal cycle-degree estimate is asserted.

\paragraph{G3: an exceptional non-CM Hodge summand.}
C16 supplies infinitely many non-CM elliptic source factors, and
C20 adds a geometrically simple genus-two Jacobian with
endomorphism ring $\Z$, on a ball cover of degree at most1764.
C13 gives
an explicit algebraic projector onto a Weil space. One nonzero
projected cycle would generate that two-dimensional rational space.
The known curve-product maps in C12, C16 and C20 are subject to the
same C14 vanishing whenever the target has the generic Weil-type
divisor algebra. The remaining task is a specified map or correspondence
with nonzero Weil projection, or the full-rank pairing data in C04.
C15's characteristic-zero universality implication is a source
reduction, not a proof of that conjecture or of the extraction gate.
The countability result restricts fixed arithmetic towers over
$\C$; it does not preclude the dense algebraic families constructed
here or a further deformation argument.

\paragraph{G4: the chosen integral correspondence and connection.}
C17 identifies the exact horizontal correspondence and local-model
data needed to compare the 032 image cycle with a Kudla--Rapoport
special cycle. C06--C08 prevent replacing those data by an
undifferentiated positivity claim. C11 proves the sharp integral
height-two derivative estimate. C19 removes its artificial
polynomial-coordinate hypothesis and proves integrality of the
CM de Rham normalization coefficient $e_2^*$.
A8 identifies the actual completed 032 tensor with the matched
canonical harmonic jet, including its embedding labels and
factorials. On a matched CM elliptic factor, A8/C19 give the
actual reduced-theta derivative bounds and their prime-to-$p$
level normalization. C21 proves the actual product decomposition
at primes split in $F$ and inert in $E$, and the strong
mixed-derivative bound for canonically matched equal-degree
ratios with integral sections and unit denominator.
For arbitrary integer tensor degree the elliptic bound has
the weaker positive slope, which C22 proves optimal through
exact infinite subsequences on two actual tame theta sections.
Thus transferring the ratio slope to all individual sections
is decided negatively in this normalization.
The general positive-degree higher-dimensional theta estimate
and the particular horizontal Kudla--Rapoport correspondence
remain to be constructed. Neither a raw analytic
frame nor a vertical local intersection is silently used in
that role.

These gates are recorded with dated changes in \texttt{LEDGER.md}.
The mandated three-hour progress review has not been replaced by
this mathematical inventory. Its timing is controlled by the actual
run clock and recorded separately.


\appendix
\section{Realization and archival crosswalk for every finding}

The following anchors locate the mathematical framework of a
finding, not its theta formula in Grothendieck's handwriting.
Write G44 for cote~44, scans~7--8 (12 August 1964), and G16 for
cote~16, scans~43--46 (4 January 1967). The original scans are
\url{https://grothendieck.umontpellier.fr/44.pdf} and
\url{https://grothendieck.umontpellier.fr/16.pdf}.
The published \emph{Standard conjectures on algebraic cycles},
pp.~196--198, is an additional polarization/realization anchor.

\begin{longtable}{@{}p{.18\linewidth}p{.48\linewidth}p{.27\linewidth}@{}}
\toprule Finding IDs & Geometric or motivic meaning & Exact framework anchor\\\midrule\endhead
A1--A3 & Oscillator period functionals on a Betti realization; extraction before an algebraic image cycle. & G44, scans 7--8.\\
A4 & Restriction to ample complete intersections; compatibility of the cup pairing. & G16, scans 43--45.\\
A5 & Positive analytic recovery operator on canonical forms; algebraicity is a separate question. & G44, scans 7--8; G16, scans 44--45.\\
A6 & Actual divisor contractions and image-cycle pushforwards, hence algebraic correspondences modulo homological equivalence. & G44, scans 7--8; G16, scans 43--46.\\
A7 & Exact analytic theta realization; its source vector is not asserted to be a motivic correspondence. & G44, scans 7--8.\\
A8 & Rigidified theta line bundle and canonical harmonic jet; actual elliptic integral normalization. & G44, scans 7--8; G16, scans 43--46.\\
B0--B3 & Hermitian realization data and the local representation needed for a period comparison. & G44, scans 7--8, as the eventual correspondence setting.\\
B4--B8 & Positivity, determinant and overlap estimates after fixing realization norms. & G44, scans 7--8; published Standard conjectures, pp.196--198.\\
B9--B12 & Local and global definite-subgroup coefficients of the canonical line; evaluation is not pullback to degree-two cohomology of a point. & G44, scans 7--8.\\
B13 & Uniform functional on an iterated theta family, with arithmetic links retained. & G44, scans 7--8; G16, scans 44--45.\\
B15 & Explicit local theta representations, global splitting characters and exact normalized two-step kernels. & G44, scans 7--8; G16, scans 44--45.\\
B16--B18 & Conditional coherent descent and an unconditional analytic adjoint seesaw on the actual theta image. & G44, scans 7--8; G16, scans 44--45.\\
B14 & Algebraic diagonal contraction of specified tensor realizations, contrasted with a different coherent line-bundle power. & G16, scans 43--46.\\
Q1--Q3 & Quantitative selection of forms and rational multiplicity data preceding cycle construction. & G44, scans 7--8.\\
Q4--Q8 & Explicit normalized period and polarization certificates; Q4,Q5 are aliases of A2,A3. & G44, scans 7--8; published Standard conjectures, pp.196--198.\\
Q9 & Algebraic endomorphism polynomials; idempotence modulo homological equivalence, with denominators. & G16, scans 43--45.\\
Q11 & Definite finite periods control the geometric ball norm through explicit holomorphic chart constants. & G44, scans 7--8; G16, scans 43--46.\\
Q10 & All-degree determinant and period certificates for the A6 cycle construction. & G44, scans 7--8; G16, scans 43--46.\\
C01--C03 & Albanese realization, trace projector and genuine non-CM factors on noncongruence covers. & G44, scans 7--8.\\
C04--C05 & Correspondence-module extraction and its Hodge-group scope. & G44, scans 7--8; G16, scans 43--45.\\
C06--C08 & Derived local intersections and the horizontal/special-fiber distinction. & SGA4$\frac12$, Cycle, Rem.2.3.9 p.150; C17's exact specialization locators.\\
C09--C11 & Local deformation and derivative valuations; an integral frame is additional input to a cycle comparison. & EGA IV$_2$, Thm.6.9.1 p.153; G44 scans 7--8 as the eventual realization setting.\\
C12 & Image cycle with degree500; product of diagonal divisors. & G44, scans 7--8; G16, scans 43--45.\\
C13--C14 & Explicit Weil projector and its action on the Lefschetz algebra. & G16, scans 43--46; G44, scans 7--8.\\
C15--C16 & Albanese quotients from finite \'etale covers and exact image-cycle degrees. & G44, scans 7--8; G16, scans 43--46.\\
C17 & Relative algebraic correspondences and specialization. & EGA IV$_3$, 8.8.2--8.8.3 pp.28--34; SGA4 III, XVI 2.1 p.213, 2.3 p.215; SGA4$\frac12$, Cycle 2.3.2--2.3.9 pp.145--150.\\
C18 & Integral line-bundle ratios on an explicit tame isogeny cover; connection data retained. & EGA IV$_2$, 6.9.1 p.153; SGA4$\frac12$, Cycle 2.3.2 p.145, 2.3.8 p.148.\\
C19 & Integral de Rham CM eigensplitting and intrinsic height-two derivative estimates, with the actual reduced theta frame. & G44, scans 7--8; C17's exact specialization framework.\\
C21--C22 & Local CM product factors and integral canonical theta ratios; an exact sharpness example for individual sections. & G44, scans 7--8; C17 specialization framework; Tate, p-Divisible Groups, connected formal factors.\\
C20 & A simple non-CM abelian-surface quotient of a ball-cover Albanese; injective specialization of endomorphisms. & G44, scans 7--8; G16, scans 43--46; EGA IV$_3$, 11.10 as cited by Conrad, Thm.2.1(2) p.70.\\
\bottomrule
\end{longtable}

\section{Source register and exact version conventions}

All cited supplied works remain in the package under their original
filenames. In particular, 032 means
\emph{The rational Hodge conjecture for CM abelian varieties}; its
printed theorem/page locators are used throughout. The companion
van Hoften--Sempliner paper is the supplied Hexagon 2610.00259v2.
Milne's supplied Hexagon 2611.00244v1, Theorems~2.1, 3.1 and~3.5,
is carried as the stated CM-Hodge, finite-field Tate and Hodge-standard
background. It is distinct from his 1999 article and modular-forms notes.

The supplied Markman source is \emph{Cycles on abelian $2n$-folds of
Weil type from secant sheaves on abelian $n$-folds},
\path{LIT_Markman_Hodge_abelian_fourfolds_2502.03415/250531.tex}.
Its theorem labelled \texttt{thm-algebraicity} in Section~1.5 and
corollary in Section~1.6 cover the stated Weil-sixfold and
abelian-fourfold results. These are background when determining
whether a proposed non-CM target is already known; this run does
not claim to improve those general results.

The external primary-source URLs and exact reading locators are
listed in \path{documentation/external_sources_C.tsv} and the
dated source-check records. Important additional versions are:
\begin{itemize}[leftmargin=*]
\item BMM: Bergeron--Millson--M\oe glin, published 2016 version,
125 PDF pages; Lemmas~5.5, 5.18, Propositions~5.17, 5.19,
Theorem~9.1. The 97-page preprint has different page numbers.
\item Liu: \emph{Fourier--Jacobi cycles and arithmetic relative
trace formula}, \url{https://yhzhumath.github.io/FJ.pdf};
Definitions~4.3, 4.5, 4.12; Theorems~4.15, 4.18;
Corollary~4.20; Appendix~D.
\item Haan: arXiv:1403.5061v7, Theorem~3.4, p.8,
\url{https://arxiv.org/pdf/1403.5061v7}.
\item Gan--Qiu--Takeda: arXiv:1207.4709v3,
Sections~11.3--11.7; its stated dimension range is retained.
\item Gross: \emph{Incoherent definite spaces and Shimura varieties},
25 February 2018, Theorem~3.1, pp.6--7.
\item Milne: \emph{Lefschetz classes on abelian varieties},
\url{https://jmilne.org/math/articles/1999aP.pdf},
Proposition~5.4 and Corollary~5.5, p.25, used in C14.
\item Katagiri--Kobayashi: arXiv:2609.07024v1,
Definition~3.1, Remark~3.2 and Lemma~3.3, pp.9--11,
\url{https://arxiv.org/pdf/2609.07024}.
\item Howe--Zhu: arXiv:math/0002205v1,
\S2 p.3, Proposition~3 and Lemma~4 p.4,
Theorem~6 pp.5--6; \url{https://arxiv.org/pdf/math/0002205}.
\item Conrad: \emph{Gross--Zagier Revisited}, MSRI
Publications~49 (2004), Theorem~2.1(2), p.70,
\url{https://library.slmath.org/books/Book49/files/05conrad.pdf}.
\item Silverman: \emph{The Arithmetic of Elliptic Curves},
second edition (2009), Chapter~III, Exercise~3.7,
pp.105--106.
\end{itemize}

Alp\"oge locators use the pinned mirror commit given at the start
of the paper. The full relative file path, declaration, and where
available line numbers are recorded with the relevant finding.
\path{documentation/alpoge_mirror/pinned_manifest.json} and
\path{documentation/step2_archive_manifest.json} preserve recovery
URLs and hashes. A failed retrieval is explicitly recorded rather
than represented as a source reading. No source-module text was
executed, and no Lean build certificate is claimed.

The mathematical deliverable also contains:
\begin{itemize}[leftmargin=*]
\item \path{CHECK_EVENTS.tsv} and \path{QUESTIONS.tsv};
\item \path{LEDGER.md} and \path{IMPLICATIONS.md};
\item \path{TRANSCRIPTION_NOTES.tsv}.
\end{itemize}
An empty transcription register means that no
supplied text-of-record reading was available for a justified
Grothendieck correction; it is not a claim that all possible
transcriptions are correct.


\end{document}
